% Generated by build_chapter.py from the GY1981 files. Do not edit by hand; edit the files and rebuild.
\chapter{The academic year 1981--82}\label{ch:1981}
\noindent\textbf{Scope and editorial standard.}
This chapter covers all 30 catalogue records, 530--559, and all 45 manuscript
pages P000981--P001025. There are 26 talk records, two headings, one plan and
one intervention recorded separately. The occurrence index contains 210
items; the problem index contains 17 definite questions, including the reused
GP-1979-41. The manuscript order is retained, including 549 before 548 and
the Fuchs intervention 551 beside the Haefliger material 553. Catalogue page
windows are preserved in the talk metadata, while summaries and item pages
identify actual boundaries. All current-status statements are assessed as of
11 November 2026. Notebook wording, verified literature and editorial
reconstruction are separated throughout.

\begin{center}
\begin{tabular}{clrr}
\toprule
Class&Historical assessment&Items&Share\\
\midrule
A&Speaker/Gelfand publication soon afterwards&31&14.8\%\\
B&Established subsequent development by others&9&4.3\%\\
C&Partial development; a precise remainder&5&2.4\%\\
D&No development identified&0&0.0\%\\
E&Known before the talk&149&71.0\%\\
F&Exact assertion or fate not identifiable&16&7.6\%\\
\bottomrule
\end{tabular}
\end{center}
\noindent Percentages are rounded. A class records the fate of the identifiable
mathematical thread, not the correctness of every notebook formula. Thirty-nine
occurrences are false as written; most require a sign, normalization, domain
or quantifier correction rather than a new research problem.

\paragraph{The year's main lines.}
The autumn concentrates on Lie algebras, quivers, multiplicities and
geometric representation theory. Serre relations, classical incidence
geometry, Gelfand--Tsetlin patterns and Verma modules are interwoven with
recent results. Abeasis--Del Fra--Kraft supply a close quiver-publication
match; the Kazhdan--Lusztig multiplicity discussion continues earlier
seminar questions. Elashvili's short classification announcement was received
only a week after his talk. The Jantzen discussion requires a particularly
important distinction: dominant deformation directions have the established
invariance property, whereas an unrestricted regular direction can change
the layer multiplicities. A complete exact computation reproduces the
published type-$A_3$ counterexample.

The winter and spring move through nonstandard perturbation theory,
integrable systems, discrete curvature, microlocal analysis, diffeomorphism
groups and foliations. M\"uller's curvature talk is matched to
Cheeger--M\"uller--Schrader, with its early announcement distinguished from
the 1984 full paper. Convergence of curvature measures does not prove
convergence of solutions of the Regge equations. McDuff's parity statements
and Milnor's discrete-Lie-group results have close subsequent publication
matches. The finite-coefficient Friedlander--Milnor comparison remains a
precise conjectural thread; broad proof claims are not treated as established
merely because a preprint exists. Kirillov and Givental lead into the final
two Wodzicki talks on index theory, eta invariants and spectral cuts.

\paragraph{What is developed here.}
Two C questions permit bounded proofs, without claims of historical priority.
For Haefliger's action of \(T(x)=x+1\) and \(S(x)=-1/x\) on
\(\mathbb {RP}^1\), every continuous function is proved to be
\(u-u\circ T+v-v\circ S\) with continuous \(u,v\). The proof gives a
uniform-norm existence estimate. It does not establish the notebook's
smooth surjectivity or smooth density; the missing distributional and
derivative estimates are stated explicitly.

For Goldsheid's coordinate-product approximation of a real Lagrangian plane,
an explicit Fourier construction disproves the proposed orthogonal-projector
gap bound \(1/\sqrt2\) already for a six-plane in \(\mathbb R^{12}\).
An odd Fourier family forces every dimension-independent constant to be at
least \(\cos(\pi/8)\). A separate argument gives the fixed-dimension upper
bound \(\cos(\pi/(2n))\), sharp for \(n=2\). The existence of any universal
constant in \([\cos(\pi/8),1)\) is not decided. Independent derivations and
the supplied computational checks corroborate these proofs; numerical
minimization is never used as a certificate of a global optimum.

\paragraph{Earlier connections and critical findings.}
The chapter links the KL thread to GP-1979-41, primitive enveloping algebras
to the 1964 occurrences, torus convexity to 1980-508.15, and eta regularity
to 1972-296.2--3 and 1972-296.5. Gilkey's no-filling regularity theorem and
Wodzicki's February 1982 publication precede the April talks; publication
dates therefore prevent a simple narrative of new seminar discoveries.
The final notes also omit a symplectic quotient in a multiplicity formula,
a gamma normalization in a fractional integral, signs in a vector-calculus
matrix, and powers of the shift parameter in a zeta identity. The correction
log separates these faithful notebook errors from doubtful readings.
Spectral-cut independence at zero does not justify independence at every
negative integer. These distinctions, and the 16 genuinely underdetermined
items, are part of the critical record rather than gaps filled by conjecture.

\GStalk{530}{1981-09-14}{I. M. Gelfand}{P000981-P000982}{Heading: on contragredient Lie algebras}
Only a dated heading and the title \ru{О контраградиентных алгебрах Ли} survive before the next dated heading, 21 September, on P000981. The catalog range extends into P000982, but those mathematical notes belong to the following meeting. No mathematical claim is assigned to this heading-only record.


\GStalk{531}{1981-09-21}{S. I. Gelfand}{P000981-P000985}{Reading list and finite-dimensional Lie theory}
The meeting begins with a reading list on highest-weight modules and Kac--Moody algebras. Macdonald's June 1981 Bourbaki lecture can be identified exactly as expos\'e 577, \emph{Affine Lie algebras and modular forms}; Humphreys' contribution is the 1980 LNM 831 chapter, pp.~72--103 \citep{macdonald1981affine,humphreys1980highest}. The intervening exposition, explicitly marked I.~M.~G., reviews roots, the Weyl group, Chevalley--Serre relations, highest weights and characters. A speculative Lorentz/neutrino aside has insufficient mathematical specification. The content ends at Rudakov's heading on P000984, not at the catalog endpoint. Several mistakes are in the notebook itself: the order of the symmetric group, the root/weight lattice distinction and omitted off-diagonal restrictions.

\begin{GSitem}{1981-531.1}{E}{Adjoint representation and the Killing form}
\GSfield{Kind; source}{definition; \GSpage{P000981}; I. M. Gelfand}
\GSfield{The notes}{The notes introduce \(\operatorname{ad}x(y)=[x,y]\) and a \ru{Картаново скал. пр-ние}; the displayed trace contains a short stroke between \(\operatorname{ad}x\) and \(\operatorname{ad}y\).}
\GSfield{Statement}{For a finite-dimensional complex semisimple Lie algebra \(\mathfrak g\), the Killing form is \(B(x,y)=\operatorname{tr}(\operatorname{ad}x\circ\operatorname{ad}y)\). It is a nondegenerate invariant symmetric complex bilinear form. It is not a positive-definite Hermitian scalar product on \(\mathfrak g\). Its restriction to a Cartan subalgebra is nondegenerate, allowing the induced bilinear form on \(\mathfrak h^*\).}
\GSfield{Reconstruction}{Composition is the mathematical emendation. The scan does not warrant silently treating the short stroke as an unambiguous multiplication dot; the ambiguity is logged separately.}
\GSfield{Fate}{This is classical background explicitly reviewed in \citet[\S1, p.~259]{macdonald1981affine}; \citet[II.2 and III.5]{serre1987complex} supplies the Cartan--Killing criterion and Cartan restriction. The former was delivered in June 1981, before this meeting.}
\GSfield{Present status}{As of November 2026, this is an established result with the hypotheses specified above; the notebook occurrence is expository.}
\GSfield{Assessment}{literal statement: not determinable; publication match: exact; verification: read the result; development: annotated; confidence: high}
\end{GSitem}
\begin{GSitem}{1981-531.2}{F}{Approximate Lorentz representations and a neutrino-mass aside}
\GSfield{Kind; source}{remark; \GSpage{P000981}; I. M. Gelfand}
\GSfield{The notes}{Gelfand mentions infinite-dimensional representations close to the standard Lorentz-group representation, an approximate \ru{смещение}, and a small displacement allegedly giving the neutrino a mass; the attribution near Novikov is uncertain.}
\GSfield{Statement}{No unique mathematical problem is recoverable. To formulate one, one would have to specify the Lorentz group or covering group, a class of representations, a topology or deformation parameter, and a Lorentz-covariant dynamical operator whose mass is being computed. None of these data is supplied.}
\GSfield{Reconstruction}{No technical reconstruction is imposed on this aside.}
\GSfield{Fate}{Searches combining Gelfand, Novikov, Lorentz representations, deformation and neutrino mass did not identify a result with these quantifiers. The juxtaposition with highest-weight theory does not identify a mass-generation mechanism. No later neutrino result is claimed as a solution.}
\GSfield{Present status}{As of November 2026, the historical allusion remains unidentified in this bounded search; this is an identification gap, not evidence that a well-posed problem is open.}
\GSfield{Assessment}{literal statement: not determinable; publication match: none found; verification: unresolved lead; development: annotated; confidence: low}
\end{GSitem}
\begin{GSitem}{1981-531.3}{E}{Cartan roots, root spaces and the real root span}
\GSfield{Kind; source}{definition; \GSpage{P000981-P000982}; I. M. Gelfand}
\GSfield{The notes}{A root vector is written \(g_\alpha\), and the equation is displayed in the order \([g_\alpha,h]=\alpha(h)g_\alpha\). The exposition then treats the roots as vectors with a positive scalar product.}
\GSfield{Statement}{Fix a Cartan subalgebra \(\mathfrak h\subset\mathfrak g\). With the usual positive-root convention, \(\mathfrak g_\alpha=\{x:[h,x]=\alpha(h)x\text{ for all }h\in\mathfrak h\}\), \(\alpha\ne0\), and \(\mathfrak g=\mathfrak h\oplus\bigoplus_{\alpha\in\Phi}\mathfrak g_\alpha\). Each root space has dimension one. The real span \(E=\operatorname{span}_{\mathbb R}\Phi\subset\mathfrak h^*\) carries a positive-definite Weyl-invariant form; positivity is a statement on this real space, not on the whole complex dual.}
\GSfield{Reconstruction}{The separate symbols \(x_\alpha\) and \(\mathfrak g_\alpha\), and the distinction between \(E\) and \(\mathfrak h^*\), are editorial clarifications. The adopted convention agrees with the subsequent raising-operator and highest-weight formulas.}
\GSfield{Fate}{The exact finite-dimensional statement is in \citet[\S1, p.~259]{macdonald1981affine} and \citet[VI.1--2]{serre1987complex}. Changing the order of the bracket changes the root sign and must be propagated consistently.}
\GSfield{Present status}{As of November 2026, this is an established result with the hypotheses specified above; the notebook occurrence is expository.}
\GSfield{Assessment}{literal statement: correct with missing hypotheses; publication match: exact; verification: read the result; development: annotated; confidence: high}
\end{GSitem}
\begin{GSitem}{1981-531.4}{E}{Simple roots, positivity and the highest root}
\GSfield{Kind; source}{theorem; \GSpage{P000982}; I. M. Gelfand}
\GSfield{The notes}{Every root is an integral combination of simple roots with coefficients of one sign. A highest root is singled out.}
\GSfield{Statement}{For a finite reduced crystallographic root system \(\Phi\), a base \(\Delta=\{\alpha_1,\ldots,\alpha_r\}\) is a real basis such that each \(\alpha\in\Phi\) is \(\sum n_i\alpha_i\), where all \(n_i\) are integers and either all are nonnegative or all are nonpositive. For an irreducible system there is a unique highest positive root in the root order; for a reducible system this assertion is made separately in each irreducible component.}
\GSfield{Reconstruction}{Irreducibility, needed for a unique highest root, is supplied. No highest-root assertion is carried over to a general infinite Kac--Moody root system.}
\GSfield{Fate}{This is already stated in the June 1981 reading-list source, \citet[\S1, p.~259]{macdonald1981affine}; the base axioms and highest-root assertion belong to classical finite root-system theory, cf. \citet[V.8--13]{serre1987complex}.}
\GSfield{Present status}{As of November 2026, this is an established result with the hypotheses specified above; the notebook occurrence is expository.}
\GSfield{Assessment}{literal statement: correct with missing hypotheses; publication match: exact; verification: read the result; development: annotated; confidence: high}
\end{GSitem}
\begin{GSitem}{1981-531.5}{E}{Coroots and the dual root system}
\GSfield{Kind; source}{construction; \GSpage{P000982}; I. M. Gelfand}
\GSfield{The notes}{The roots \(2\alpha/(\alpha,\alpha)\) are declared to form a root system corresponding to a dual group.}
\GSfield{Statement}{For a reduced crystallographic root system in a Euclidean space, \(\Phi^\vee=\{\alpha^\vee=2\alpha/(\alpha,\alpha):\alpha\in\Phi\}\) is again a reduced crystallographic root system, with transposed Cartan matrix. Passing from a connected reductive group to a dual group requires exchanging the complete root datum \((X^*,\Phi,X_*,\Phi^\vee)\), including its character and cocharacter lattices. The root system alone determines the dual semisimple Lie algebra but not the isogeny class of a group.}
\GSfield{Reconstruction}{The distinction between a dual root system and a dual group with specified root datum is an editorial qualification, not an additional theorem attributed to the speaker.}
\GSfield{Fate}{The inverse-root construction predates the meeting; see \citet[V.6]{serre1987complex}. The group-language shorthand is compatible with root-datum duality only after the lattice choice is supplied; the notes do not make that choice.}
\GSfield{Present status}{As of November 2026, this is an established result with the hypotheses specified above; the notebook occurrence is expository.}
\GSfield{Assessment}{literal statement: correct with missing hypotheses; publication match: exact; verification: read the result; development: annotated; confidence: high}
\end{GSitem}
\begin{GSitem}{1981-531.6}{E}{Weyl reflections and the examples A2 and sl(n)}
\GSfield{Kind; source}{theorem; \GSpage{P000982}; I. M. Gelfand}
\GSfield{The notes}{Reflections in the roots generate a finite group. For \(\mathfrak{sl}(n)\) the notes say the symmetric group is \ru{n-го порядка}. The \(\mathfrak{sl}_3\) sketch has three vectors at \(120^\circ\) whose sum is zero.}
\GSfield{Statement}{The reflection is \(s_\alpha(v)=v-\langle v,\alpha^\vee\rangle\alpha\), and the finite Weyl group of \(\mathfrak{sl}_n\) is \(S_n\), acting by coordinate permutation on \(\sum x_i=0\), of order \(n!\). Its roots are \(e_i-e_j\), \(i\ne j\). In type \(A_2\) there are six roots. Three cyclic roots \(e_1-e_2,e_2-e_3,e_3-e_1\) sum to zero and have pairwise angle \(120^\circ\); they are neither the complete root set nor a simple-root base, which has two members.}
\GSfield{Fate}{The correct statements are classical; see \citet[V.3--4 and VII.6]{serre1987complex}. The sketch can be understood accurately as a three-vector subsystem of the drawing after adjoining their negatives. The scan confirms that the order error is in the notes, not generated by translation.}
\GSfield{Present status}{As of November 2026, this is an established result with the hypotheses specified above; the notebook occurrence is expository.}
\GSfield{Assessment}{literal statement: false as written; publication match: exact; verification: read the result; development: annotated; confidence: high}
\end{GSitem}
\begin{GSitem}{1981-531.7}{E}{Finite Cartan matrices and their symmetrization}
\GSfield{Kind; source}{definition; \GSpage{P000982-P000983}; I. M. Gelfand}
\GSfield{The notes}{The entries \(a_{ij}=\langle\alpha_j,\alpha_i^\vee\rangle\) are described as nonpositive integers, followed by positive definiteness.}
\GSfield{Statement}{A finite-type Cartan matrix has \(a_{ii}=2\), \(a_{ij}\in\mathbb Z_{\le0}\) for \(i\ne j\), and \(a_{ij}=0\iff a_{ji}=0\). It is symmetrizable: a positive diagonal matrix \(D\) makes \(DA\) symmetric positive definite. Equivalently all principal minors of \(A\) are positive. Unless all simple roots have the same length, \(A\) itself need not be symmetric; for example \(\bigl(\begin{smallmatrix}2&-2\\-1&2\end{smallmatrix}\bigr)\) is a finite-type Cartan matrix.}
\GSfield{Fate}{\citet[\S1, condition (C2), p.~260]{macdonald1981affine} uses the principal-minor formulation. The notebook's wording compresses diagonal normalization and symmetrization, so it must not be read as \(a_{ii}\le0\) or as symmetry of every Cartan matrix.}
\GSfield{Present status}{As of November 2026, this is an established result with the hypotheses specified above; the notebook occurrence is expository.}
\GSfield{Assessment}{literal statement: false as written; publication match: exact; verification: read the result; development: annotated; confidence: high}
\end{GSitem}
\begin{GSitem}{1981-531.8}{E}{Chevalley–Serre presentation}
\GSfield{Kind; source}{theorem; \GSpage{P000983}; I. M. Gelfand}
\GSfield{The notes}{The generators \(e_i,f_i,h_i\) obey the Cartan relations and the iterated-adjoint relations with exponent \(1-a_{ij}\); the off-diagonal restriction is not written.}
\GSfield{Statement}{For a finite-type Cartan matrix \(A\), the corresponding complex semisimple Lie algebra is generated by \(e_i,f_i,h_i\), with \([h_i,h_j]=0\), \([h_i,e_j]=a_{ij}e_j\), \([h_i,f_j]=-a_{ij}f_j\), \([e_i,f_j]=\delta_{ij}h_i\), and \((\operatorname{ad}e_i)^{1-a_{ij}}e_j=(\operatorname{ad}f_i)^{1-a_{ij}}f_j=0\) only for \(i\ne j\). Its positive nilpotent subalgebra is presented by the \(e_i\) and the positive Serre relations alone.}
\GSfield{Fate}{This is Serre's presentation, already part of the 1965 lecture course underlying \citet[VI.4, Theorems 6--7, and Appendix]{serre1987complex}; it is also reviewed in \citet[\S1, (1.2)]{macdonald1981affine}. Inserting \(i=j\) would ask for an exponent \(-1\), so the missing qualification is substantive.}
\GSfield{Present status}{As of November 2026, this is an established result with the hypotheses specified above; the notebook occurrence is expository.}
\GSfield{Assessment}{literal statement: correct with missing hypotheses; publication match: exact; verification: read the result; development: annotated; confidence: high}
\end{GSitem}
\begin{GSitem}{1981-531.9}{E}{Weight-space decomposition and the weight lattice}
\GSfield{Kind; source}{theorem; \GSpage{P000983}; I. M. Gelfand}
\GSfield{The notes}{Finite-dimensional modules split into common eigenspaces of the Cartan subalgebra. The weights are then said to lie in the lattice generated by roots (or coroots).}
\GSfield{Statement}{For a finite-dimensional module over a complex semisimple Lie algebra, \[V=\bigoplus_{\lambda\in P}V_\lambda,\qquad P=\{\lambda\in\mathfrak h^*: \lambda(h_i)\in\mathbb Z\},\] where \(P\) is the weight lattice and \(h_i\) are simple coroots. The root lattice \(Q=\sum_i\mathbb Z\alpha_i\) is a subgroup of \(P\), usually proper. Weights of one irreducible module lie in one coset of \(Q\). For \(\mathfrak{sl}_2\), the defining representation has weights \(\pm\alpha/2\), which are not in \(Q=\mathbb Z\alpha\).}
\GSfield{Fate}{\citet[\S1, pp.~260--261]{macdonald1981affine} explicitly distinguishes \(Q\subset P\) and places representation weights in \(P\). Thus the notebook's lattice statement is false as written, even though its weight decomposition is correct. It cannot be repaired by replacing roots with coroots: these naturally belong to the dual Cartan space.}
\GSfield{Present status}{As of November 2026, this is an established result with the hypotheses specified above; the notebook occurrence is expository.}
\GSfield{Assessment}{literal statement: false as written; publication match: exact; verification: read the result; development: annotated; confidence: high}
\end{GSitem}
\begin{GSitem}{1981-531.10}{E}{Root vectors shift weights}
\GSfield{Kind; source}{theorem; \GSpage{P000983}; I. M. Gelfand}
\GSfield{The notes}{A root vector sends a vector of weight \(\lambda\) to weight \(\lambda+\alpha\); one occurrence writes \(e_\alpha v_\alpha\) where the accompanying calculation uses \(v_\lambda\).}
\GSfield{Statement}{For every weight module and \(x_\alpha\in\mathfrak g_\alpha\), \(x_\alpha V_\lambda\subset V_{\lambda+\alpha}\). Indeed \(h(x_\alpha v)=[h,x_\alpha]v+x_\alpha hv=(\alpha(h)+\lambda(h))x_\alpha v\). The resulting vector may be zero; the assertion does not guarantee that \(\lambda+\alpha\) occurs as a weight.}
\GSfield{Reconstruction}{The nonzero-vector qualification and uniform weight subscript clarify the calculation; the one-line proof is included to show why the convention matters.}
\GSfield{Fate}{This elementary relation is \citet[VII.1, Proposition 1(a)]{serre1987complex}. It fixes the bracket convention adopted earlier in the notebook and supplies a direct internal check on the sign of \(\alpha\).}
\GSfield{Present status}{As of November 2026, this is an established result with the hypotheses specified above; the notebook occurrence is expository.}
\GSfield{Assessment}{literal statement: correct with missing hypotheses; publication match: exact; verification: read the result; development: annotated; confidence: high}
\end{GSitem}
\begin{GSitem}{1981-531.11}{E}{Classification by dominant integral highest weights}
\GSfield{Kind; source}{theorem; \GSpage{P000983}; I. M. Gelfand}
\GSfield{The notes}{A vector annihilated by the positive root operators is a highest-weight vector. Cartan's theorem is paraphrased as uniqueness of a representation determined by its highest weight.}
\GSfield{Statement}{For a fixed positive system in a complex semisimple \(\mathfrak g\), finite-dimensional irreducible modules are classified up to isomorphism by \(P_+=\{\lambda:\lambda(h_i)\in\mathbb Z_{\ge0}\}\). The irreducible \(L(\lambda)\) has a one-dimensional highest-weight space annihilated by \(\mathfrak n_+\). For any \(\lambda\in\mathfrak h^*\) there is a unique irreducible highest-weight module, but it is finite-dimensional exactly for \(\lambda\in P_+\). A reducible module is not determined merely by one highest-weight vector.}
\GSfield{Fate}{The finite-dimensional statement is stated in \citet[\S1, p.~261]{macdonald1981affine}; the universal highest-weight construction and finite-dimensional criterion are \citet[VII.3--4, Theorems 1--3]{serre1987complex}. The qualifications separate this theorem from the much larger category of Verma modules treated in record 533.}
\GSfield{Present status}{As of November 2026, this is an established result with the hypotheses specified above; the notebook occurrence is expository.}
\GSfield{Assessment}{literal statement: correct with missing hypotheses; publication match: exact; verification: read the result; development: annotated; confidence: high}
\end{GSitem}
\begin{GSitem}{1981-531.12}{E}{Formal characters and the Weyl vector}
\GSfield{Kind; source}{definition; \GSpage{P000983-P000984}; I. M. Gelfand}
\GSfield{The notes}{The group ring of the weight lattice is introduced using formal exponentials, and \(\rho\) is characterized by \(\rho(h_i)=1\).}
\GSfield{Statement}{In \(\mathbb Z[P]\) use symbols \(e^\lambda\) with \(e^\lambda e^\mu=e^{\lambda+\mu}\). A finite-dimensional weight module has character \(\operatorname{ch}V=\sum_\lambda (\dim V_\lambda)e^\lambda\). For a finite root system \(\rho=\frac12\sum_{\alpha>0}\alpha\) satisfies \(\rho(h_i)=1\). In a Kac--Moody setting the latter equations remain meaningful in a realization, whereas a literal half-sum of all positive roots is generally not a finite vector.}
\GSfield{Fate}{The finite-dimensional definitions are exactly the setup of \citet[\S1, p.~261]{macdonald1981affine}. The last warning matters because the later Verma-module discussion shifts to \(G(A)\); one must not carry an infinite half-sum there without a completion or replacement definition.}
\GSfield{Present status}{As of November 2026, this is an established result with the hypotheses specified above; the notebook occurrence is expository.}
\GSfield{Assessment}{literal statement: correct; publication match: exact; verification: read the result; development: annotated; confidence: high}
\end{GSitem}
\begin{GSitem}{1981-531.13}{E}{Weyl character formula}
\GSfield{Kind; source}{theorem; \GSpage{P000984}; I. M. Gelfand}
\GSfield{The notes}{The surviving, uncancelled formula is the alternating Weyl numerator divided by a product over positive roots. A preceding attempt is crossed out in the source.}
\GSfield{Statement}{For \(\lambda\in P_+\), \[\operatorname{ch}L(\lambda)=\frac{\sum_{w\in W}(-1)^{\ell(w)}e^{w(\lambda+\rho)-\rho}}{\prod_{\alpha\in\Phi_+}(1-e^{-\alpha})}.\] The displayed quotient can be taken in the fraction field of \(\mathbb Z[P]\); the theorem says it equals the finite integral character. Equivalently divide \(\sum_w(-1)^{\ell(w)}e^{w(\lambda+\rho)}\) by the Weyl denominator.}
\GSfield{Fate}{This is the classical Weyl formula, reviewed in \citet[\S1, (1.3)]{macdonald1981affine} and \citet[VII.8]{serre1987complex}. The crossed-out formula is evidence of a correction during note-taking and is not treated as a second asserted theorem.}
\GSfield{Present status}{As of November 2026, this is an established result with the hypotheses specified above; the notebook occurrence is expository.}
\GSfield{Assessment}{literal statement: correct with missing hypotheses; publication match: exact; verification: read the result; development: annotated; confidence: high}
\end{GSitem}
\begin{GSitem}{1981-531.14}{E}{Weyl denominator and the Vandermonde model}
\GSfield{Kind; source}{theorem; \GSpage{P000984}; I. M. Gelfand}
\GSfield{The notes}{The character formula at the trivial highest weight becomes an alternating-sum/product identity. A Vandermonde determinant is displayed as its elementary model.}
\GSfield{Statement}{For a finite root system, \[\sum_{w\in W}(-1)^{\ell(w)}e^{w\rho}=e^\rho\prod_{\alpha>0}(1-e^{-\alpha}).\] For variables \(x_1,\ldots,x_n\), the column convention \(V_{ij}=x_j^{i-1}\) gives \(\det V=\prod_{i<j}(x_j-x_i)\). Type \(A\) specializes the alternating numerator to a determinant after an appropriate monomial factor is removed.}
\GSfield{Reconstruction}{The indices on the Vandermonde product and the determinant convention are supplied explicitly.}
\GSfield{Fate}{The denominator identity is \citet[\S1, (1.4)]{macdonald1981affine}. The Vandermonde analogy is exact at the determinant level; the source does not specify an ordered product convention, so no independent sign error is inferred from the order of its two variables.}
\GSfield{Present status}{As of November 2026, this is an established result with the hypotheses specified above; the notebook occurrence is expository.}
\GSfield{Assessment}{literal statement: correct; publication match: exact; verification: read the result; development: annotated; confidence: high}
\end{GSitem}

\GStalk{532}{1981-09-21}{A. N. Rudakov}{P000983-P000987}{Contragredient Lie algebras and Cartan-matrix types}
Rudakov starts on P000984 with generators, relations, an enlarged Cartan space and a maximal ideal avoiding that space. He then removes redundant central directions and states the finite/affine/indefinite trichotomy for indecomposable generalized Cartan matrices. The literal affine alternative is internally contradictory. At the top of P000986 Gelfand adds the two Rogers--Ramanujan identities. The arrow returning to S.~I.~Gelfand on that page marks the next record. The historical construction should be compared with \citet[Chapter II]{kac1968graded} and \citet[\S2]{macdonald1981affine}; the latter's Serre-presentation convention must not silently replace the maximal-ideal definition for nonsymmetrizable matrices.

\begin{GSitem}{1981-532.1}{E}{The maximal-ideal construction of a contragredient Lie algebra}
\GSfield{Kind; source}{construction; \GSpage{P000984-P000985}}
\GSfield{The notes}{For a matrix \(A=(a_{ij})\), Rudakov starts with generators \(e_i,f_i,h_i,\epsilon_i\), all Cartan generators commuting, with \([e_i,f_j]=\delta_{ij}h_i\), \([h_i,e_j]=a_{ij}e_j\), \([h_i,f_j]=-a_{ij}f_j\), \([\epsilon_i,e_j]=\delta_{ij}e_j\), and \([\epsilon_i,f_j]=-\delta_{ij}f_j\). He factors by a maximal \ru{неповреждающий} ideal.}
\GSfield{Statement}{Over \(\mathbb C\), let \(\widetilde{\mathfrak h}\) have independent basis \(h_i,\epsilon_i\), and put \(\alpha_j(h_i)=a_{ij}\), \(\alpha_j(\epsilon_i)=\delta_{ij}\). Form the Lie algebra \(\widetilde{\mathfrak g}\) with these Cartan and cross relations before imposing Serre relations. It has free positive and negative parts and a triangular decomposition. There is a unique largest ideal \(\mathfrak r\) satisfying \(\mathfrak r\cap\widetilde{\mathfrak h}=0\); the quotient is the contragredient algebra for this enlarged realization. The word `largest' means the sum of all such ideals, not an arbitrary maximal ideal of the whole algebra.}
\GSfield{Reconstruction}{The source's informal requirement that generators survive is expressed as disjointness from the explicitly specified Cartan space. The distinction between the initial presentation and the maximal quotient is retained; no nonsymmetrizable Gabber--Kac theorem is assumed.}
\GSfield{Fate}{The maximal contragredient construction is already in \citet[Chapter II, \S1]{kac1968graded}; the independent derivations are a way to distinguish simple-root coordinates when \(A\) is singular. Modern minimal realizations use Cartan dimension \(2n-\operatorname{rank}A\); \citet[Definition 1.1 and Proposition 1.3]{polo2026realizations} makes that dimension and uniqueness precise.}
\GSfield{Present status}{As of November 2026, this is an established result with the hypotheses specified above; the notebook occurrence is expository.}
\GSfield{Assessment}{literal statement: correct with missing hypotheses; publication match: exact; verification: read the result; development: annotated; confidence: high}
\end{GSitem}
\begin{GSitem}{1981-532.2}{E}{Removing redundant central Cartan directions}
\GSfield{Kind; source}{construction; \GSpage{P000985}}
\GSfield{The notes}{A central ideal \(J\) is chosen disjoint from the span of the \(h_i\), and the algebra is divided by \(J\). A sentence says the operation is no longer needed when \(A\) is nondegenerate. The example \((2)\) is followed by `matrices of order 2 with zero trace'.}
\GSfield{Statement}{In the preceding enlarged Cartan realization, \(z_i=h_i-\sum_j a_{ij}\epsilon_j\) are central and independent. Their span \(Z\) has dimension \(n\), and \(\dim(Z\cap\operatorname{span}\{h_i\})=n-\operatorname{rank}A\). Quotient by a complement \(J\) of this intersection in \(Z\); then \(\dim J=\operatorname{rank}A\), the \(h_i\) survive, and the Cartan dimension becomes \(2n-\operatorname{rank}A\). For \(A=(2)\), the minimal algebra is \(\mathfrak{sl}_2\), the trace-zero \(2\times2\) matrices.}
\GSfield{Reconstruction}{The central basis and dimension calculation are explicit editorial checks. The words about trace-zero matrices describe the resulting sl2 algebra, not the Cartan matrix (2).}
\GSfield{Fate}{This agrees with the realization dimension in \citet[Definition 1.1]{polo2026realizations} and with the complementary-derivation construction in \citet[\S2, p.~262]{macdonald1981affine}. In the enlarged presentation, the quotient is still needed when \(A\) is invertible: for \(A=(2)\), \(h-2\epsilon\) is already a nonzero central element. The source sentence is correct only if it means that one can omit all added derivations when starting with an invertible matrix.}
\GSfield{Present status}{As of November 2026, this is an established result with the hypotheses specified above; the notebook occurrence is expository.}
\GSfield{Assessment}{literal statement: false as written; publication match: exact; verification: read the result; development: annotated; confidence: high}
\end{GSitem}
\begin{GSitem}{1981-532.3}{E}{Roots and root multiplicities of the contragredient quotient}
\GSfield{Kind; source}{definition; \GSpage{P000985}}
\GSfield{The notes}{The notebook defines \(G_\alpha=\{x:[h,x]=\alpha(h)x\}\) and calls \(\alpha\) a root if this space is nonzero.}
\GSfield{Statement}{For a realization with independent simple roots, the algebra has a root-lattice grading and a Cartan weight decomposition \(\mathfrak g=\mathfrak h\oplus\bigoplus_{\alpha\in\Delta}\mathfrak g_\alpha\), with \(\Delta\subset Q\setminus\{0\}\). Every root space of a finite-rank contragredient algebra is finite-dimensional, because the homogeneous components of the free positive and negative Lie algebras are finite-dimensional. The root multiplicity is \(\operatorname{mult}\alpha=\dim\mathfrak g_\alpha\); it need not be one.}
\GSfield{Reconstruction}{The modern convention excludes \(0\) from the root set and names its weight space \(\mathfrak h\). Kac's 1968 paper deliberately uses a `zero root'; the notebook's unqualified wording is therefore a convention ambiguity, not proof of a transcription mistake.}
\GSfield{Fate}{\citet[Chapter II, \S2]{kac1968graded} explicitly distinguishes the formal root-lattice degree from evaluation on a possibly too-small Cartan subalgebra. \citet[\S2, pp.~262--263]{macdonald1981affine} describes multiplicities in the symmetrizable realization. This is the setup for Gelfand's question in 1981-533.4.}
\GSfield{Present status}{As of November 2026, this is an established result with the hypotheses specified above; the notebook occurrence is expository.}
\GSfield{Assessment}{literal statement: correct with missing hypotheses; publication match: exact; verification: read the result; development: annotated; confidence: high}
\end{GSitem}
\begin{GSitem}{1981-532.4}{E}{Generalized Cartan-matrix axioms}
\GSfield{Kind; source}{definition; \GSpage{P000985}}
\GSfield{The notes}{Two numbered conditions normalize the diagonal at 2, require nonpositive integral off-diagonal entries, and impose \(a_{ij}=0\Longleftrightarrow a_{ji}=0\). A separate indecomposability restriction precedes the classification.}
\GSfield{Statement}{A generalized Cartan matrix is an integral matrix \(A\) satisfying \(a_{ii}=2\), \(a_{ij}\le0\) for \(i\ne j\), and \(a_{ij}=0\iff a_{ji}=0\). Indecomposable means that simultaneous permutation of rows and columns cannot make it a nontrivial block diagonal matrix. These axioms do not imply symmetrizability: the latter is the additional existence of a positive diagonal \(D\) such that \(DA\) is symmetric.}
\GSfield{Fate}{These are the conditions recorded as (M2), (M3), (M5) in \citet[Chapter II, \S1]{kac1968graded}, distinct from its extra invariant-form condition. The distinction is decisive for the scope of Berman--Moody's multiplicity formula, which is not a theorem for all matrices allowed in this notebook.}
\GSfield{Present status}{As of November 2026, this is an established result with the hypotheses specified above; the notebook occurrence is expository.}
\GSfield{Assessment}{literal statement: correct with missing hypotheses; publication match: exact; verification: read the result; development: annotated; confidence: high}
\end{GSitem}
\begin{GSitem}{1981-532.5}{E}{Finite, affine and indefinite trichotomy}
\GSfield{Kind; source}{theorem; \GSpage{P000985-P000986}}
\GSfield{The notes}{The three alternatives are labelled finite-dimensional, extended and all remaining algebras. In alternative II the source simultaneously writes \(Au\ge0\Rightarrow u=0\) and the existence of \(u>0\) with \(Au=0\).}
\GSfield{Statement}{For an indecomposable generalized Cartan matrix exactly one case occurs. Finite type: \(A\) is invertible and \(Au\ge0\) implies \(u=0\) or \(u>0\). Affine type: \(\operatorname{rank}A=n-1\), \(\ker A=\mathbb Rv\) for some \(v>0\), and \(Au\ge0\) implies \(Au=0\). Indefinite type: some \(v>0\) satisfies \(Av<0\), and \(u\ge0,Au\ge0\) implies \(u=0\). Inequalities are componentwise and the unrestricted vectors are real. Equivalently, for the nonnegative irreducible matrix \(B=2I-A\), compare its Perron root with 2.}
\GSfield{Reconstruction}{The spectral formulation supplies an independently checkable precise statement and resolves which zero was intended. Indecomposability is essential; a disconnected matrix can mix types.}
\GSfield{Fate}{The finite and affine classes belong to the pre-1981 Kac--Moody classification; see \citet[Chapter II, \S\S4--7]{kac1968graded} and the finite/affine treatment in \citet[\S\S1--3]{macdonald1981affine}. The stated matrix test follows directly from Perron--Frobenius: \(Av=(2-\rho(B))v\), with a positive left Perron vector giving the affine and indefinite implications. Thus the literal affine line is a notebook error, not a new exceptional alternative.}
\GSfield{Present status}{As of November 2026, this is an established result with the hypotheses specified above; the notebook occurrence is expository.}
\GSfield{Assessment}{literal statement: false as written; publication match: exact; verification: read the result; development: annotated; confidence: high}
\end{GSitem}
\begin{GSitem}{1981-532.6}{E}{Second Rogers–Ramanujan identity}
\GSfield{Kind; source}{theorem; \GSpage{P000986}; I. M. Gelfand}
\GSfield{The notes}{Gelfand first writes the product with residues 2 and 3 modulo 5 and the sum with exponent \(m(m+1)\).}
\GSfield{Statement}{With \((q;q)_m=\prod_{j=1}^{m}(1-q^j)\) and \((q;q)_0=1\), \[\sum_{m\ge0}\frac{q^{m(m+1)}}{(q;q)_m}=\prod_{n\ge1}\frac1{(1-q^{5n-2})(1-q^{5n-3})}.\] This holds as an identity of formal power series, and analytically for \(|q|<1\). The empty denominator at \(m=0\) is one.}
\GSfield{Fate}{The identity is classical. A directly relevant Lie-theoretic interpretation was already announced in February 1981 by \citet[pp.~699--701]{lepowskywilson1981rogers}; it interprets the Rogers--Ramanujan sides through vacuum-state characters of standard \(A_1^{(1)}\)-modules. That publication precedes this September remark and is not its later solution.}
\GSfield{Present status}{As of November 2026, this is an established result with the hypotheses specified above; the notebook occurrence is expository.}
\GSfield{Assessment}{literal statement: correct; publication match: exact; verification: publisher metadata checked; development: annotated; confidence: high}
\end{GSitem}
\begin{GSitem}{1981-532.7}{E}{First Rogers–Ramanujan identity}
\GSfield{Kind; source}{theorem; \GSpage{P000986}; I. M. Gelfand}
\GSfield{The notes}{The next product has residues 1 and 4 modulo 5, while the numerator on the sum side is \(q^{m^2}\).}
\GSfield{Statement}{In the same formal or analytic range, \[\sum_{m\ge0}\frac{q^{m^2}}{(q;q)_m}=\prod_{n\ge1}\frac1{(1-q^{5n-1})(1-q^{5n-4})}.\] Combinatorially, the right side counts partitions whose parts are congruent to 1 or 4 modulo 5; the sum also enumerates partitions with adjacent parts differing by at least two. This explanatory interpretation is additional to the notebook's displayed identity.}
\GSfield{Fate}{\citet[abstract and pp.~699--701]{lepowskywilson1981rogers} expressly treats the two classical identities and their distinct standard affine modules. The notebook's order is second identity followed by first; changing that order is unnecessary.}
\GSfield{Present status}{As of November 2026, this is an established result with the hypotheses specified above; the notebook occurrence is expository.}
\GSfield{Assessment}{literal statement: correct; publication match: exact; verification: publisher metadata checked; development: annotated; confidence: high}
\end{GSitem}

\GStalk{533}{1981-09-21}{S. I. Gelfand}{P000985-P000987}{Verma embeddings and a root-multiplicity question}
The return arrow on P000986 begins a review of Verma modules and the Bernstein--Gelfand--Gelfand criterion for embeddings. The formulas mix the notebook's unshifted highest weight with a negative rho shift and do not delimit the finite-dimensional semisimple scope of the quoted theorem. At the top of P000987, Gelfand asks for root-space dimensions and independent words for general Cartan matrices. The next dated heading starts Seifert's talk. The multiplicity question is indexed, but neither efficiency nor a uniform positive combinatorial formula is demanded in the surviving wording; no artificial open conjecture is extracted.

\begin{GSitem}{1981-533.1}{E}{Verma modules and their PBW weight structure}
\GSfield{Kind; source}{construction; \GSpage{P000986}}
\GSfield{The notes}{For \(\lambda\in H^*\), impose \(e_iv_0=0\) and \(hv_0=\lambda(h)v_0\), and apply words in the \(f_i\). The notes calculate the weight \(\lambda-\alpha_{i_1}-\cdots-\alpha_{i_k}\).}
\GSfield{Statement}{For a complex Lie algebra with triangular decomposition \(\mathfrak g=\mathfrak n_-\oplus\mathfrak h\oplus\mathfrak n_+\) of the semisimple or contragredient type under discussion, set \(\mathfrak b_+=\mathfrak h\oplus\mathfrak n_+\), and let \(\mathbb C_\lambda\) be its one-dimensional module with \(\mathfrak n_+\) acting by zero. Then \(M(\lambda)=U(\mathfrak g)\otimes_{U(\mathfrak b_+)}\mathbb C_\lambda\cong U(\mathfrak n_-)\) as a vector space. It is the universal highest-weight module of highest weight \(\lambda\); its weight set is contained in \(\lambda-Q_+\). Arbitrary \(f_i\)-words span, while a PBW basis uses ordered monomials in a basis of \(\mathfrak n_-\).}
\GSfield{Reconstruction}{The induction notation specifies precisely the universal quotient. No additional freeness of arbitrary simple-generator words is inferred.}
\GSfield{Fate}{This construction predates the seminar. \citet[\S3, Definition 1 and Lemma 1]{bernsteingelfandgelfand1971structure} uses a shifted parameter: its module \(M_\xi\) has highest weight \(\xi-\rho\). Keeping that convention separate from the notebook's \(M(\lambda)\) is necessary for the subsequent embedding test. The prior index's 1980-520.1 concerns the same construction but not a previously indexed open question.}
\GSfield{Present status}{As of November 2026, this is an established result with the hypotheses specified above; the notebook occurrence is expository.}
\GSfield{Assessment}{literal statement: correct with missing hypotheses; publication match: exact; verification: read the result; development: annotated; confidence: high}
\end{GSitem}
\begin{GSitem}{1981-533.2}{E}{One-dimensional Hom spaces and injectivity of Verma homomorphisms}
\GSfield{Kind; source}{theorem; \GSpage{P000986}}
\GSfield{The notes}{The notes assert \(\operatorname{Hom}(M(\lambda),M(\mu))=0\) or \(\mathbb C\), and every nonzero map is injective, immediately after introducing \(G(A)\).}
\GSfield{Statement}{For a finite-dimensional complex semisimple Lie algebra and any \(\lambda,\mu\in\mathfrak h^*\), \(\dim\operatorname{Hom}_{\mathfrak g}(M(\lambda),M(\mu))\le1\), and every nonzero homomorphism is an embedding. The image of the highest vector is a singular vector of weight \(\lambda\). Injectivity alone also follows from the freeness over \(U(\mathfrak n_-)\) and the domain property of that enveloping algebra. The one-dimensional bound requires its own theorem and is not inferred merely from PBW.}
\GSfield{Reconstruction}{The verified theorem is deliberately stated in its finite-dimensional semisimple scope. This is not a claim that every assertion fails in every larger class.}
\GSfield{Fate}{\citet[\S3, Theorem 1, attributed there to Verma]{bernsteingelfandgelfand1971structure} gives exactly this finite-dimensional semisimple statement. The notebook has not supplied the restrictions needed for an arbitrary Kac--Moody generalization; imaginary-root and critical-weight phenomena cannot be discarded by calling the theorem formal.}
\GSfield{Present status}{As of November 2026, this is an established result with the hypotheses specified above; the notebook occurrence is expository.}
\GSfield{Assessment}{literal statement: correct with missing hypotheses; publication match: partial; verification: read the result; development: annotated; confidence: high}
\end{GSitem}
\begin{GSitem}{1981-533.3}{E}{The BGG strong-linkage criterion and the rho-sign error}
\GSfield{Kind; source}{theorem; \GSpage{P000986}}
\GSfield{The notes}{A nonzero homomorphism is said to imply \(\lambda\le\mu\), \(\lambda-\rho=w(\mu-\rho)\), and an ordered sequence of integral reflections. The same three conditions are declared sufficient. The scanned minus signs before both rho terms are clear.}
\GSfield{Statement}{In the preceding finite-dimensional semisimple setting, an embedding \(M(\lambda)\hookrightarrow M(\mu)\) exists exactly when there is a sequence \(\nu_0=\mu,\nu_r=\lambda\), positive roots \(\beta_j\), and integers \(m_j=\langle\nu_{j-1}+\rho,\beta_j^\vee\rangle>0\) with \(\nu_j=s_{\beta_j}(\nu_{j-1}+\rho)-\rho=\nu_{j-1}-m_j\beta_j\), allowing the empty sequence. In particular \(\lambda+\rho=w(\mu+\rho)\); Weyl-orbit membership and the root order alone are weaker than the full ordered condition.}
\GSfield{Reconstruction}{The chain is rewritten with a single forward index to remove the notebook’s dangling delta subscript i-1. Positivity of the roots and integrality at each reflection are explicit. The sl2 test is an editorial verification, not an additional claim from the notes.}
\GSfield{Fate}{\citet[\S4, condition (A), Theorem 2, and \S5, Theorem 3]{bernsteingelfandgelfand1971structure} proves necessity and sufficiency, using shifted parameters from the outset; the paper was received 5 October 1970 and published in 1971. Under the notebook's explicit unshifted definition, the minus sign is false: for \(\mathfrak{sl}_2\), the singular vector \(fv_0\) gives \(M(-2)\hookrightarrow M(0)\), and \(-2+1=-(0+1)\), whereas \(-2-1\) is not \(\pm(0-1)\).}
\GSfield{Present status}{As of November 2026, this is an established result with the hypotheses specified above; the notebook occurrence is expository.}
\GSfield{Assessment}{literal statement: false as written; publication match: exact; verification: read the result; development: annotated; confidence: high}
\end{GSitem}
\begin{GSitem}{1981-533.4}{E}{Computing root-space dimensions and independent root words}
\GSfield{Kind; source}{problem; \GSpage{P000987}; I. M. Gelfand}
\GSfield{The notes}{Gelfand labels the question \ru{Задача (не банальная)}: for a Cartan matrix satisfying conditions 1 and 2, \ru{Какова размерность пр-ва корней?} The next line mentions constructing linearly independent words for every root.}
\GSfield{Statement}{Given a generalized Cartan matrix \(A\), the maximal contragredient quotient defined in record 532, and a nonzero degree \(\alpha=\sum_i n_i\alpha_i\), determine \(\dim\mathfrak g_\alpha\), and obtain representatives by Lie words when that component is nonzero. The wording does not require a positive combinatorial formula, a canonical basis or a complexity bound. For symmetrizable \(A\), let \(D=\sum_{w\in W}(-1)^{\ell(w)}e^{w\rho-\rho}\) and \(-\log D=\sum_{\gamma>0}c_\gamma e^{-\gamma}\). Denominator inversion gives \(\operatorname{mult}\alpha=\sum_{d\mid\alpha}\mu(d)c_{\alpha/d}/d\), with \(d\mid\alpha\) meaning divisibility of every simple-root coefficient.}
\GSfield{Reconstruction}{The explicit computational interpretation and the recursive lowering proof are editorial. No request for a uniform cancellation-free or efficient formula is added. The notes may have been seeking something more illuminating than elimination, but that stronger demand is not specified.}
\GSfield{Fate}{\citet[\S2 and Theorem 2]{bermanmoody1979multiplicities}, received 28 September 1978 and published September 1979, already gives an explicit inversion formula for \emph{symmetrizable} Cartan matrices. That hypothesis is not in Gelfand's question. In the wider maximal-quotient definition, each specified degree is nevertheless a finite linear-algebra computation: for height greater than one, \(x\mapsto([f_i,x])_i\) embeds \(\mathfrak g_\alpha\) into the sum of lower positive components. Its rank on a finite free-Lie spanning family computes the multiplicity, and pivot columns select independent words. The annotation proves this elementary recursive observation, rather than attributing a nonsymmetrizable formula to Berman--Moody.}
\GSfield{Present status}{As of November 2026, the finite computational problem as stated is answerable; symmetrizable multiplicity formulas were available before the meeting. Difficult current questions about positive formulas, growth or complexity are not declared identical to this broad notebook question.}
\GSfield{Assessment}{literal statement: correct; publication match: partial; verification: read the result; development: annotated; confidence: high; problem \hyperref[prob:GP-1981-01]{GP-1981-01}}
\end{GSitem}

\GStalk{534}{1981-09-28}{Benedict Seifert}{P000985-P000990}{Transformation-group geometry and a Bruhat-type decomposition}
The actual heading is on P000987 and includes a question mark after the Russian surname and the affiliation Paris. Fixed sets of pointwise stabilizers define a closure geometry; the speaker builds flags, proposes four axioms and asserts a double-coset decomposition. The closest publication is Seifert's \emph{Transitivity Properties, Bruhat Decomposition and Tits Systems for Transformation Groups}, received 5 March 1983, published June 1984. Its exact axiom list could not be read in this run, so the identification is partial. Seifert's later paper explicitly distinguishes its new tree transitivity and proof from the earlier work \citep[Introduction and reference S.1]{seifert1986highly}. Gelfand's three concrete double-coset exercises occur on P000989 before the Zelevinsky heading.

\begin{GSitem}{1981-534.1}{E}{Closed sets from pointwise stabilizers}
\GSfield{Kind; source}{definition; \GSpage{P000987}}
\GSfield{The notes}{For an action \(G\curvearrowright X\), \(S\) belongs to \(\mathcal F\) when every \(x\notin S\) can be moved by some group element fixing every point of \(S\). Arbitrary intersections of members of \(\mathcal F\) are closed, and the intersection of all closed supersets defines \(S^{\mathrm{cl}}\).}
\GSfield{Statement}{Write \(G_{(S)}=\{g:gs=s\text{ for every }s\in S\}\), and \(\operatorname{Fix}(H)=\{x:hx=x\text{ for every }h\in H\}\). Then \(\mathcal F=\{S:S=\operatorname{Fix}(G_{(S)})\}\), and \(\operatorname{cl}(S)=\operatorname{Fix}(G_{(S)})\). This operation is extensive, monotone, idempotent and equivariant under \(G\); intersections of closed sets are closed. The bottom is \(\operatorname{cl}(\varnothing)=\operatorname{Fix}(G)\), which need not be empty, and the top is \(X\).}
\GSfield{Reconstruction}{The closure formula and bottom element make the source’s terminal (0) notation precise. They are elementary reconstruction of the displayed definition, not an inferred extra axiom.}
\GSfield{Fate}{These assertions are elementary consequences of the order-reversing stabilizer/fixed-set correspondence: \(S\subseteq\operatorname{Fix}(H)\) iff \(H\subseteq G_{(S)}\). Their particular use in Seifert's Bruhat program is linked to \citet{seifert1984transitivity} by \citet[Introduction and reference S.1]{seifert1986highly}, but the formal closure argument needs no transitivity or finiteness assumption.}
\GSfield{Present status}{As of November 2026, this is an established result with the hypotheses specified above; the notebook occurrence is expository.}
\GSfield{Assessment}{literal statement: correct; publication match: partial; verification: indirect; development: annotated; confidence: high}
\end{GSitem}
\begin{GSitem}{1981-534.2}{E}{Setwise stabilizers and the restriction exact sequence}
\GSfield{Kind; source}{construction; \GSpage{P000987}}
\GSfield{The notes}{The notebook distinguishes \(G_S=\bigcap_{s\in S}G_s\) from \(G^S=\{g:gS=S\}\), and gives an extension by a barred permutation group on \(S\).}
\GSfield{Statement}{For every subset \(S\subseteq X\), restriction induces an exact sequence \[1\longrightarrow G_{(S)}\longrightarrow G_{\{S\}}\longrightarrow G_{\{S\}}|_S\longrightarrow1,\] where \(G_{\{S\}}\) is the setwise stabilizer and the right-hand group is the actual image in \(\operatorname{Sym}(S)\). The image need not be every automorphism of an independently defined geometry on \(S\).}
\GSfield{Fate}{This is the first-isomorphism theorem for the restriction homomorphism. It is the exact algebraic content of the notebook's extension diagram; the later Seifert citation provides context, not an additional surjectivity onto a larger group \citep[Introduction]{seifert1986highly}.}
\GSfield{Present status}{As of November 2026, this is an established result with the hypotheses specified above; the notebook occurrence is expository.}
\GSfield{Assessment}{literal statement: correct; publication match: partial; verification: indirect; development: annotated; confidence: high}
\end{GSitem}
\begin{GSitem}{1981-534.3}{E}{Ordered and saturated flags in the closure geometry}
\GSfield{Kind; source}{construction; \GSpage{P000987}}
\GSfield{The notes}{For a partially ordered set \(I\), the speaker considers maps \(f:I\to\mathcal F\) satisfying \(f(i)\subseteq f(j)\iff i\le j\), and then requires adjacent elements of \(I\) to map to adjacent closed sets.}
\GSfield{Statement}{The first condition says that \(f\) is an order embedding. The second is preservation of covering relations: if \(i<j\) with no intermediate element, there is no closed set strictly between \(f(i)\) and \(f(j)\). Equivariance of closure gives \((g\cdot f)(i)=g f(i)\), an action on both families. When \(I\) is a chain and the image is maximal, its stabilizer is \(B_f=\bigcap_iG_{\{f(i)\}}\). These constructions alone assert neither flag transitivity nor a Bruhat decomposition.}
\GSfield{Fate}{This is a precise formalization of the displayed definitions. Seifert's subsequently published transformation-group program concerns additional transitivity hypotheses; his 1986 Introduction explicitly says that its tree formulation is different from the 1984 approach \citep{seifert1986highly}.}
\GSfield{Present status}{As of November 2026, this is an established result with the hypotheses specified above; the notebook occurrence is expository.}
\GSfield{Assessment}{literal statement: correct; publication match: partial; verification: indirect; development: annotated; confidence: high}
\end{GSitem}
\begin{GSitem}{1981-534.4}{E}{Singleton closures and an incorrect stabilizer-equality assertion}
\GSfield{Kind; source}{claim; \GSpage{P000988}}
\GSfield{The notes}{For a transitive action, the closure of \(\{x\}\) is called a line and is said to be the set of points with \ru{тем же стабилизатором}, the same stabilizer as \(x\).}
\GSfield{Statement}{Always \(\operatorname{cl}\{x\}=\{y:G_x\subseteq G_y\}\). Equality of stabilizers follows, for example, when \(G\) is finite and transitive, because conjugate stabilizers then have the same finite order. It is false for arbitrary infinite transitive actions. Let \(G=\mathbb Z[1/2]\rtimes\langle t\rangle\), with \(tbt^{-1}=2b\), and let \(H=\mathbb Z\) be the translation subgroup. In the transitive action on left cosets \(G/H\), the points \(x=H\) and \(y=t^{-1}H\) have stabilizers \(H\) and \(t^{-1}Ht=\tfrac12\mathbb Z\), respectively. Thus \(G_x\subsetneq G_y\), so \(y\in\operatorname{cl}\{x\}\) but their stabilizers differ.}
\GSfield{Reconstruction}{The infinite affine-group example and the finite-action repair are editorial checks. No exact published version of this sentence has been verified.}
\GSfield{Fate}{The corrected inclusion formula follows directly from the definition in 1981-534.1. The displayed affine-group counterexample checks all hypotheses of the literal claim and shows that transitivity alone cannot justify equality. This audit concerns the notebook; it does not allege that the inaccessible 1984 paper contains the same error \citep{seifert1984transitivity}.}
\GSfield{Present status}{As of November 2026, this is an established result with the hypotheses specified above; the notebook occurrence is expository.}
\GSfield{Assessment}{literal statement: false as written; publication match: none found; verification: indirect; development: annotated; confidence: high}
\end{GSitem}
\begin{GSitem}{1981-534.5}{E}{Double transitivity and its singleton-closure obstruction}
\GSfield{Kind; source}{theorem; \GSpage{P000988}}
\GSfield{The notes}{Double transitivity is identified with transitivity on \((F\times F)\setminus\Delta\) and with transitivity of a point stabilizer on the remaining points. A non-singleton singleton closure is said to prevent double transitivity.}
\GSfield{Statement}{For a transitive action on a set \(F\) with at least two points, 2-transitivity is equivalent to transitivity of \(G_x\) on \(F\setminus\{x\}\). If \(|F|\ge3\), a 2-transitive action has \(\operatorname{Fix}(G_x)=\{x\}\), so a larger singleton closure obstructs it. The condition \(|F|\ge3\) matters: the regular action of the group of order two on two points is 2-transitive, but \(G_x=1\) and \(\operatorname{cl}\{x\}=F\). Alternatively, the source's pictured obstruction is valid whenever \(\{x\}\subsetneq\operatorname{cl}\{x\}\subsetneq F\).}
\GSfield{Fate}{These are immediate orbit-stabilizer criteria; they explain precisely the argument in the notebook's three-point drawing. Seifert's later paper makes multiple transitivity a structural input rather than a consequence of single transitivity \citep[Introduction and Definition 2.1]{seifert1986highly}. The two-point counterexample is an editorial boundary check.}
\GSfield{Present status}{As of November 2026, this is an established result with the hypotheses specified above; the notebook occurrence is expository.}
\GSfield{Assessment}{literal statement: correct with missing hypotheses; publication match: partial; verification: indirect; development: annotated; confidence: high}
\end{GSitem}
\begin{GSitem}{1981-534.6}{A}{Four geometric axioms and the asserted Bruhat decomposition}
\GSfield{Kind; source}{claim; \GSpage{P000988-P000989}}
\GSfield{The notes}{The axioms are transitivity on maximal flags; finite flag dimension; a dimension-intersection condition; and surjectivity of restriction to \(S\) for elements also preserving a disjoint singleton closure. They are followed by \(G=\coprod_WG^fwG^f\).}
\GSfield{Statement}{Here is the strongest literal formulation, with normalizations made explicit. Put \(0=\operatorname{Fix}(G)\). Assume maximal flags are a single \(G\)-orbit; the rank \(r(S)\) is finite and given by the length of maximal chains from \(0\) to \(S\); if \(S,S'\subseteq T\) and \(r(S)+r(S')>r(T)\), then \(S\cap S'\ne0\); and whenever \(C=\operatorname{cl}\{x\}\) has \(C\cap S=0\), restriction \(G_{\{S\}}\cap G_{\{C\}}\to G_{\{S\}}|_S\) is onto. The notebook claims that for a maximal flag with stabilizer \(B\) and its frame group \(N\), representatives of \(W=N/T\) give distinct double cosets \(G=\coprod_{w\in W}B\dot wB\).}
\GSfield{Reconstruction}{Rank normalization, the meaning of the restriction image, and representatives dot-w are editorial. The displayed assertion remains an attributed notebook claim; no proof is certified here solely from the abridged axioms.}
\GSfield{Fate}{The author, Paris location and subject fit \citet{seifert1984transitivity}, received 5 March 1983 and published June 1984, and Seifert's own \citet[Introduction, reference S.1]{seifert1986highly} confirms the predecessor. The 1984 full text was not accessible: its hypotheses and proof have not been checked against these four lines. The 1986 tree theorem is a different theorem and is not substituted for them. Classification A records the credible publication lineage; the publication match remains partial.}
\GSfield{Present status}{As of November 2026, the publication lineage is identified, but exact theorem-level verification of the 1984 axiom system remains an access/identification limitation. It is not classified as an unresolved mathematical conjecture.}
\GSfield{Assessment}{literal statement: not determinable; publication match: partial; verification: indirect; development: annotated; confidence: medium}
\end{GSitem}
\begin{GSitem}{1981-534.7}{A}{Frames, their normalizer and the Weyl quotient}
\GSfield{Kind; source}{construction; \GSpage{P000989}}
\GSfield{The notes}{Every maximal flag is said to have \(f(i)=\operatorname{cl}\{e_j:j\le i\}\). Put \(C_i=\operatorname{cl}\{e_i\}\); \(N\) permutes these sets and \(T\) preserves each one. An exact sequence ends in \(W=N/T\).}
\GSfield{Statement}{For a chosen finite family of distinct closed sets \(\mathcal C=\{C_1,\ldots,C_r\}\), let \(N\) be its setwise stabilizer and \(T=\bigcap_iG_{\{C_i\}}\). The permutation action gives \(1\to T\to N\to W\to1\), with \(W\) its image in \(S_r\). This exact sequence is automatic. The additional statement that every maximal flag admits such a point-generated frame, and that this particular \(W\) parametrizes all \(B\)-double cosets, belongs to the preceding claimed geometric theorem; it is not true for an arbitrary closure operator without extra hypotheses.}
\GSfield{Fate}{The quotient construction is elementary, while its Bruhat role matches Seifert's 1984 program only partially verified through \citet[Introduction and reference S.1]{seifert1986highly}. The notes do not justify importing all Coxeter or BN-pair axioms from the word `Weyl'; the 1986 paper itself exhibits a modified Tits-system setting.}
\GSfield{Present status}{As of November 2026, the exact sequence is fully determined by its definition. The asserted frame theorem remains tied to the incompletely verified 1984 theorem, rather than a new open problem.}
\GSfield{Assessment}{literal statement: correct with missing hypotheses; publication match: partial; verification: indirect; development: annotated; confidence: medium}
\end{GSitem}
\begin{GSitem}{1981-534.8}{F}{Unspecified examples outside the axioms}
\GSfield{Kind; source}{remark; \GSpage{P000989}}
\GSfield{The notes}{The speaker knows no examples other than classical ones satisfying all axioms, but says that some groups have a Bruhat decomposition while failing the axioms.}
\GSfield{Statement}{The first sentence records the speaker's knowledge at that date. The second asserts the existence of a group/action with a double-coset decomposition of the indicated kind that fails at least one displayed axiom, but it supplies no group, action or failed axiom. Consequently no exact example can be identified from this line.}
\GSfield{Fate}{\citet[Introduction and Theorem 2.2]{seifert1986highly} later supplies concrete tree-action examples with a modified Tits system and Bruhat decomposition. That is relevant thematic development, but its deliberately different transitivity framework does not identify the unnamed 1981 examples.}
\GSfield{Present status}{As of November 2026, the unnamed examples remain unidentified; neither the historical knowledge claim nor failure of the particular 1981 axioms is inferred from later thematic examples.}
\GSfield{Assessment}{literal statement: not determinable; publication match: thematic only; verification: read the result; development: annotated; confidence: low}
\end{GSitem}
\begin{GSitem}{1981-534.9}{E}{Unitary double cosets and singular values}
\GSfield{Kind; source}{problem; \GSpage{P000989}; I. M. Gelfand}
\GSfield{The notes}{Gelfand's first exercise is to describe \(U(n)\backslash GL(n)/U(n)\); neither the field nor \(U(n)\) is expanded in the surviving line.}
\GSfield{Statement}{Under the natural unitary interpretation \(G=GL_n(\mathbb C)\), \(U(n)\) the unitary group, each double coset has exactly one representative \(\operatorname{diag}(s_1,\ldots,s_n)\) with \(s_1\ge\cdots\ge s_n>0\). These are the singular values: diagonalize the positive Hermitian matrix \(g^*g\), take its positive square root, and absorb the two unitary factors. Thus the quotient is the closed positive Weyl chamber in \(\mathbb R_{>0}^n\).}
\GSfield{Reconstruction}{The interpretation of U as unitary and the complex ground field are explicit assumptions. If U meant an upper unitriangular subgroup, the answer would be different; the ambiguous symbol is logged.}
\GSfield{Fate}{This is the classical singular-value/polar-decomposition calculation, not an open problem. The proof above supplies the precise quotient directly. Compact/complex group comparison is reviewed in \citet[VIII.7]{serre1987complex}; that citation is contextual, not a claim that this exact exercise originates there.}
\GSfield{Present status}{As of November 2026, this is an established result with the hypotheses specified above; the notebook occurrence is expository.}
\GSfield{Assessment}{literal statement: correct with missing hypotheses; publication match: partial; verification: indirect; development: annotated; confidence: high; problem \hyperref[prob:GP-1981-02]{GP-1981-02}}
\end{GSitem}
\begin{GSitem}{1981-534.10}{E}{Borel double cosets and permutation matrices}
\GSfield{Kind; source}{problem; \GSpage{P000989}; I. M. Gelfand}
\GSfield{The notes}{The second exercise is \(B(n)\backslash GL(n)/B(n)\), with `triangular' written under \(B(n)\).}
\GSfield{Statement}{For \(G=GL_n(k)\) over a field and \(B\) its invertible upper triangular subgroup, \(G=\coprod_{w\in S_n}B\dot wB\), where \(\dot w\) is the permutation matrix. Hence the double-coset set is \(S_n\). Equivalently, two complete flags have a relative-position permutation, determined by the dimensions of intersections of their steps. This should be distinguished from merely counting \(n!\) possible representatives without establishing disjointness.}
\GSfield{Fate}{This is the classical Bruhat decomposition. \citet[VIII.4, Theorems 7--8 and Corollary]{serre1987complex} states the bijection between the Weyl group and Borel double cosets in the complex group setting. It explains the concrete model for Seifert's abstract decomposition.}
\GSfield{Present status}{As of November 2026, this is an established result with the hypotheses specified above; the notebook occurrence is expository.}
\GSfield{Assessment}{literal statement: correct with missing hypotheses; publication match: exact; verification: read the result; development: annotated; confidence: high; problem \hyperref[prob:GP-1981-02]{GP-1981-02}}
\end{GSitem}
\begin{GSitem}{1981-534.11}{E}{Mixed unitary–Borel double cosets}
\GSfield{Kind; source}{problem; \GSpage{P000989}; I. M. Gelfand}
\GSfield{The notes}{The third exercise is \(U(n)\backslash GL(n)/B(n)\).}
\GSfield{Statement}{For \(GL_n(\mathbb C)\), with \(U(n)\) unitary and \(B\) invertible upper triangular, Gram--Schmidt gives \(g=ur\) with \(u\in U(n)\), \(r\in B\). Thus there is exactly one \(U(n)\)-\(B\) double coset. Equivalently, the unitary group acts transitively on the complex complete flag manifold; the stabilizer of the standard flag in \(U(n)\) is the diagonal compact torus.}
\GSfield{Reconstruction}{As in the first exercise, unitary U and the complex field are explicit reconstructed conventions.}
\GSfield{Fate}{This is a pre-existing classical matrix-factorization exercise, with a direct proof by orthonormalizing the columns. The compact-group/complex-flag framework is summarized in \citet[VIII.4 and VIII.7]{serre1987complex}. No additional double-coset invariant survives in this mixed quotient.}
\GSfield{Present status}{As of November 2026, this is an established result with the hypotheses specified above; the notebook occurrence is expository.}
\GSfield{Assessment}{literal statement: correct with missing hypotheses; publication match: partial; verification: indirect; development: annotated; confidence: high; problem \hyperref[prob:GP-1981-02]{GP-1981-02}}
\end{GSitem}

\GStalk{535}{1981-09-28}{A. V. Zelevinsky}{P000988-P000991}{Equioriented type-A quiver orbits and their geometry}
The actual exposition begins in the lower part of P000989 and explicitly reports \citet{abeasisdelfrakraft1981geometry}, whose September 1981 issue and receipt date of 1 April 1981 precede the meeting. It explains the finite orbit classification, rank invariants, good singularities of orbit closures and the determinantal special case. The closing analogy with Schubert geometry has a precise later realization in Zelevinsky's 1985 paper; the accompanying p-adic multiplicity connection must be distinguished from that purely geometric identification. The next dated heading on P000990 begins Kirillov's talk.

\begin{GSitem}{1981-535.1}{E}{The representation variety and base-change action of A\_m}
\GSfield{Kind; source}{construction; \GSpage{P000989}}
\GSfield{The notes}{A chain \(V_1\to\cdots\to V_m\) is called a representation of \(A_m\), and the product of the general linear groups acts. Its last arrow is labelled \(A_m\) although there are only \(m-1\) arrows.}
\GSfield{Statement}{Fix finite-dimensional spaces \(V_i\) over an algebraically closed field \(k\), with \(d_i=\dim V_i\). For the equioriented quiver \(1\to2\to\cdots\to m\), put \(E(\mathbf d)=\prod_{i=1}^{m-1}\operatorname{Hom}(V_i,V_{i+1})\). The group \(G(\mathbf d)=\prod_iGL(V_i)\) acts by \((g\cdot A)_i=g_{i+1}A_i g_i^{-1}\). Its orbits are exactly the isomorphism classes of representations with this fixed dimension vector. The last arrow is \(A_{m-1}:V_{m-1}\to V_m\).}
\GSfield{Fate}{This is the setup of \citet[\S1.1, p.~403]{abeasisdelfrakraft1981geometry}, already in print at the time of this explicitly attributed exposition. The orientation is part of the problem, not a cosmetic choice: the particular composition-rank coordinates below use the linear orientation. The earlier finite-quiver discussion 1972-277.8 and GP-1972-301 is related background, rather than an additional unresolved problem here.}
\GSfield{Present status}{As of November 2026, this is an established result with the hypotheses specified above; the notebook occurrence is expository.}
\GSfield{Assessment}{literal statement: correct with missing hypotheses; publication match: exact; verification: read the result; development: annotated; confidence: high}
\end{GSitem}
\begin{GSitem}{1981-535.2}{E}{Finite orbit classification by interval multiplicities}
\GSfield{Kind; source}{theorem; \GSpage{P000989-P000990}}
\GSfield{The notes}{There are finitely many orbits, an explicit description is possible, and a second description uses multiplicities of occurrences.}
\GSfield{Statement}{Let \(I[p,q]\), \(1\le p\le q\le m\), have one-dimensional spaces at vertices \(p,\ldots,q\), zero spaces elsewhere, and identity maps between adjacent nonzero spaces. Every equioriented \(A_m\)-representation is uniquely a direct sum \(\bigoplus_{p\le q} I[p,q]^{\oplus m_{pq}}\), with \(m_{pq}\in\mathbb Z_{\ge0}\). For fixed \(\mathbf d\), the equations \(d_i=\sum_{p\le i\le q}m_{pq}\) permit finitely many such arrays, hence finitely many orbits. It is the fixed-dimension variety, not the collection over all possible dimensions, that has finitely many orbits.}
\GSfield{Reconstruction}{The interval modules and dimension equations unpack the notebook’s second classification without changing its fixed-dimension context.}
\GSfield{Fate}{The interval/multiplicity description is explicit in \citet[\S2.3, p.~404]{abeasisdelfrakraft1981geometry}. It is the type-A case of the pre-existing finite-type quiver classification. This item specifies what the notebook's abbreviated `occurrences' must mean; it does not attribute a new classification theorem to the reporting speaker.}
\GSfield{Present status}{As of November 2026, this is an established result with the hypotheses specified above; the notebook occurrence is expository.}
\GSfield{Assessment}{literal statement: correct with missing hypotheses; publication match: exact; verification: read the result; development: annotated; confidence: high}
\end{GSitem}
\begin{GSitem}{1981-535.3}{E}{Composition ranks, interval inversion and degeneration order}
\GSfield{Kind; source}{theorem; \GSpage{P000990}}
\GSfield{The notes}{The notes prescribe ranks \(r_{ij}\) of all possible compositions and say that a system of inequalities must hold.}
\GSfield{Statement}{Set \(r_{ii}=d_i\) and \(r_{ij}=\operatorname{rank}(A_{j-1}\cdots A_i)\) for \(i<j\), with ranks beyond the end vertices zero. Then \(r_{ij}=\sum_{p\le i\le j\le q}m_{pq}\), and \[m_{ij}=r_{ij}-r_{i-1,j}-r_{i,j+1}+r_{i-1,j+1}.\] Thus the ranks determine the orbit. Rank arrays of representations are precisely the arrays obtained from nonnegative integral \(m_{ij}\) by the first formula, subject to the prescribed diagonal. For two representations of the same dimension vector, \(\mathcal O_B\subseteq\overline{\mathcal O_A}\) if and only if \(r_{ij}(B)\le r_{ij}(A)\) for every \(i\le j\).}
\GSfield{Reconstruction}{The explicit finite-difference formula and the closure-order converse are read from the identified paper; the notebook only records the ranks and unspecified inequalities. They are labelled here as completion from that source.}
\GSfield{Fate}{\citet[\S2.1, Proposition 2.2 and \S2.3, p.~404]{abeasisdelfrakraft1981geometry} supplies the exact invariant and degeneration statements. Upper semicontinuity proves necessity of the rank inequalities; the converse is the nontrivial orbit-closure result, not an automatic consequence of writing down minors. The finite-difference inversion is also verified in the accompanying computation.}
\GSfield{Present status}{As of November 2026, this is an established result with the hypotheses specified above; the notebook occurrence is expository.}
\GSfield{Assessment}{literal statement: correct with missing hypotheses; publication match: exact; verification: read the result; development: annotated; confidence: high}
\end{GSitem}
\begin{GSitem}{1981-535.4}{E}{Normality, Cohen–Macaulayness and rational singularities}
\GSfield{Kind; source}{theorem; \GSpage{P000989}}
\GSfield{The notes}{The orbit closures are called good varieties: normal, Cohen--Macaulay, and so on.}
\GSfield{Statement}{Over an algebraically closed field of characteristic zero, every orbit closure in the representation space of the equioriented \(A_m\)-quiver is normal and Cohen--Macaulay, and has rational singularities. Here rational singularities means that for a resolution \(\pi:Y\to X\), \(\pi_*\mathcal O_Y=\mathcal O_X\) and \(R^i\pi_*\mathcal O_Y=0\) for \(i>0\). The claim concerns the reduced orbit closure, not an arbitrary scheme defined by equations with the same points.}
\GSfield{Reconstruction}{The field and characteristic are made explicit; rational singularities completes the openly abbreviated list from the paper being reported.}
\GSfield{Fate}{This is \citet[Theorem 0.2, p.~401; characteristic-zero convention in \S1.1]{abeasisdelfrakraft1981geometry}. The published theorem was visually checked in the author's full-text copy. `Rational singularities' is the identified paper's precise addition to the notebook's `etc.', rather than an independently recorded notebook claim. Later scheme-theoretic refinements should be cited separately.}
\GSfield{Present status}{As of November 2026, this is an established result with the hypotheses specified above; the notebook occurrence is expository.}
\GSfield{Assessment}{literal statement: correct with missing hypotheses; publication match: exact; verification: read the result; development: annotated; confidence: high}
\end{GSitem}
\begin{GSitem}{1981-535.5}{A}{The Schubert analogy and the Zelevinsky map}
\GSfield{Kind; source}{claim; \GSpage{P000990}}
\GSfield{The notes}{The quiver orbits are analogous to Schubert cells, with corresponding polynomials.}
\GSfield{Statement}{The exact geometric completion is a locally closed embedding of the equioriented quiver representation space into an appropriate partial flag variety, obtained from a block unitriangular matrix whose blocks are consecutive compositions. Under this map, each orbit closure identifies with a suitable open part of a Schubert variety, with an explicitly matched rank/double-coset index. Consequently local intersection-cohomology polynomials are type-A Kazhdan--Lusztig polynomials after the indexing and dimension shifts are fixed. An orbit itself need not be an affine cell: for one map \(\mathbb C\to\mathbb C\), the rank-one orbit is \(\mathbb C^*\).}
\GSfield{Reconstruction}{The map, flag variety and intersection-cohomology interpretation are the identified later realization of an analogy; they are not literally stated in the 1981 shorthand.}
\GSfield{Fate}{\citet[Theorem 1 and Corollary 1]{zelevinsky1985graded}, received 4 April 1983 and published in 1985, makes this analogy precise. \citet[\S1.3 and \S2.2]{lakshmibaimagyar1997degeneracy}, circulated in September 1997, strengthens the correspondence to determinantal schemes and explains what the earlier set-theoretic proof leaves to establish. Thus the later exact construction is credited to the speaker, while scheme refinements are credited separately.}
\GSfield{Present status}{As of November 2026, the Schubert-patch description and its scheme-theoretic strengthening are established. The elementary counterexample rules out reading “analogous” as a literal assertion that every quiver orbit is an affine Schubert cell.}
\GSfield{Assessment}{literal statement: correct with missing hypotheses; publication match: partial; verification: read the result; development: annotated; confidence: high}
\end{GSitem}
\begin{GSitem}{1981-535.6}{B}{Quiver intersection cohomology and p-adic representation multiplicities}
\GSfield{Kind; source}{claim; \GSpage{P000990}}
\GSfield{The notes}{The polynomials attached to these varieties are said to describe representations of p-adic groups, in analogy with classical-group representation theory and Schubert geometry. The notes name neither the group nor the multiplicity formula.}
\GSfield{Statement}{The securely identifiable nearby problem is Zelevinsky's 1981 conjecture for smooth finite-length complex representations of \(GL_n(F)\), with \(F\) a nonarchimedean local field. On one fixed cuspidal line, multisegments label standard induced modules \(\pi_a\), irreducibles \(\langle b\rangle\), and orbits \(\mathcal O_a,\mathcal O_b\) in a graded nilpotent/type-A quiver variety with the same dimension vector. With the original submodule convention, the conjecture gives \([\pi_a:\langle b\rangle]=\sum_j\dim\mathcal H^j(IC(\overline{\mathcal O_b}))_{x_a}\), with a consistent shift normalization and \(x_a\in\mathcal O_a\). The associated graded dimensions form a local intersection-cohomology polynomial, evaluated at one for the multiplicity. This is a statement about stalks and Jordan--H\"older factors, not ordinary global cohomology or a classification of all p-adic groups.}
\GSfield{Reconstruction}{The original and modern standard-module conventions are kept separate. The group GL(n), cuspidal-line reduction and stalk formula come from the identified 1981 article, not from the notebook line itself.}
\GSfield{Fate}{\citet[Introduction and \S1.9]{zelevinsky1981padic}, received 27 November 1980 and published April--June 1981, states the conjecture before this seminar; its original Russian full text was read. For the modern standard-quotient normalization, \citet[Theorem 3.9, p.~46, version 2]{balodis2025padic} proves the corresponding \(GL_n\) formula over a p-adic field, with explicit comparison of the geometric and representation-theoretic objects. This preprint first appeared 10 October 2025 and was revised 28 February 2026; it does not claim priority over all prior proofs. Its \S1.4 credits \citet[Theorem 5.4(b)]{solleveld2021graded} and the standard-module comparison \citet[Lemma 6.2]{solleveld2023submodules}. The two Solleveld preprints first circulated in June 2021 and September 2023; the relevant revisions include May 2022 and May 2024. Both current versions were checked, including January 2025 and April 2026 revisions respectively. The notebook's unspecific remark is therefore a partial match to this precise solved problem; the 1985 Schubert map alone does not prove the representation-theoretic equality.}
\GSfield{Present status}{As of November 2026, the identified GL(n) multiplicity statement is proved, with direct current theorem-level sources. No claim is made that all p-adic Kazhdan--Lusztig variants are settled, or that the first proof appeared in 2025.}
\GSfield{Assessment}{literal statement: correct with missing hypotheses; publication match: partial; verification: read the result; development: annotated; confidence: medium; problem \hyperref[prob:GP-1981-04]{GP-1981-04}}
\end{GSitem}
\begin{GSitem}{1981-535.7}{E}{A single arrow and determinantal varieties}
\GSfield{Kind; source}{construction; \GSpage{P000990}}
\GSfield{The notes}{For \(V_1\xrightarrow{A_1}V_2\), determinantal varieties arise.}
\GSfield{Statement}{For \(\dim V_1=p\), \(\dim V_2=q\), the \(GL(V_1)\times GL(V_2)\)-orbits in \(\operatorname{Hom}(V_1,V_2)\) are the rank strata \(\mathcal O_r\), \(0\le r\le\min(p,q)\). Their closures are \(D_r=\{A:\operatorname{rank}A\le r\}\), cut out by the \((r+1)\)-minors, of dimension \(r(p+q-r)\). Thus the general composition-rank description reduces to the ordinary rank invariant. The open rank-r orbit is smooth, but its lower-rank boundary can be singular.}
\GSfield{Fate}{This is the single-arrow special case in \citet[Introduction and \S2]{abeasisdelfrakraft1981geometry}. The dimension follows by choosing the image \(r\)-plane and a surjection onto it, giving \(r(q-r)+pr\). The geometric multiplicity calculations on determinantal varieties also occur in \citet[\S3]{zelevinsky1981padic}; that relation is more precise than identifying ordinary rank strata with affine cells.}
\GSfield{Present status}{As of November 2026, this is an established result with the hypotheses specified above; the notebook occurrence is expository.}
\GSfield{Assessment}{literal statement: correct; publication match: exact; verification: read the result; development: annotated; confidence: high}
\end{GSitem}

\GStalk{536}{1981-10-19}{A. A. Kirillov}{P000989-P000993}{Rank-two Serre presentations, free Lie dimensions and interlacing arrays}
Kirillov asks about the two-generator Lie algebra with two Serre relations, gives the multigraded Witt formula before relations, identifies a Catalan specialization, and passes to the enveloping algebra and a Hilbert-series logarithm. A bare lexicographic word list does not specify a reduction algorithm. Gelfand then asks for the number of integer interlacing arrays with four fixed top entries. The latter counts a basis of one irreducible representation, not distinct representations of that highest weight. The record ends on P000991; the next meeting starts on P000992.

\begin{GSitem}{1981-536.1}{E}{The two-generator Serre algebra and its dimension table}
\GSfield{Kind; source}{question; \GSpage{P000990}}
\GSfield{The notes}{For \(a,b\in\mathbb N\), ask about the Lie algebra generated by \(x,y\), with \((\operatorname{ad}x)^{a+1}y=(\operatorname{ad}y)^{b+1}x=0\). The table has dimensions \(3,4,6,\infty\), but question marks for \((a,b)=(3,2),(2,3)\).}
\GSfield{Statement}{Over characteristic zero and for positive integers \(a,b\), the presented algebra is the positive nilpotent subalgebra of the symmetrizable rank-two Kac--Moody algebra with Cartan matrix \(A=\bigl(\begin{smallmatrix}2&-a\\-b&2\end{smallmatrix}\bigr)\). Its total dimension is \(3\) for \(ab=1\), \(4\) for \(ab=2\), \(6\) for \(ab=3\), and infinite for \(ab\ge4\). The finite cases are \(A_2,B_2,G_2\), with those numbers of positive roots; \(ab=4\) is affine, and \(ab>4\) is indefinite. In particular both question marks in the table equal infinity.}
\GSfield{Reconstruction}{The positive-integer convention and identification with the positive Serre subalgebra are explicit; the dimensions concern that subalgebra, not the full Kac--Moody algebra.}
\GSfield{Fate}{The classification and Serre presentation were available before 1981; compare \citet[Chapter II, \S\S4--7]{kac1968graded} and the contemporaneous summary \citet[\S\S1--3]{macdonald1981affine}. The rank-two matrix is symmetrizable by \(D=\operatorname{diag}(b,a)\), and its symmetrized determinant has sign \(4-ab\), directly checking the boundary. The question marks are unresolved entries in these notes, not evidence of an open rank-two classification problem.}
\GSfield{Present status}{As of November 2026, this is an established result with the hypotheses specified above; the notebook occurrence is expository.}
\GSfield{Assessment}{literal statement: correct with missing hypotheses; publication match: exact; verification: read the result; development: annotated; confidence: high}
\end{GSitem}
\begin{GSitem}{1981-536.2}{E}{The bigraded Witt formula for a free Lie algebra}
\GSfield{Kind; source}{theorem; \GSpage{P000990}}
\GSfield{The notes}{Without the two relations, the number of commutators with \(k\) occurrences of \(x\) and \(\ell\) of \(y\) is given by a M\"obius sum.}
\GSfield{Statement}{For the free Lie algebra \(L(x,y)\) over a field of characteristic zero, and nonnegative integers \(k,\ell\) with \(k+\ell>0\), \[\dim L_{k,\ell}=\frac1{k+\ell}\sum_{d\mid\gcd(k,\ell)}\mu(d)\binom{(k+\ell)/d}{k/d}.\] The expression counts independent Lie elements, not all formal bracketings. It gives \(\dim L_{1,0}=\dim L_{0,1}=1\) and zero in pure-generator degrees above one. It does not compute the quotient dimensions after imposing Serre relations without further subtraction or reduction.}
\GSfield{Fate}{This is the classical multigraded Witt formula; \citet[proof of Theorem 2 and the free-Lie specialization, pp.~225--226]{bermanmoody1979multiplicities} derives the relevant logarithmic/M\"obius inversion before this seminar. In two letters it also counts primitive necklaces: M\"obius inversion removes proper repetitions, and each primitive necklace has \(k+\ell\) distinct cyclic rotations. The included finite enumeration checks that independent route against the formula.}
\GSfield{Present status}{As of November 2026, this is an established result with the hypotheses specified above; the notebook occurrence is expository.}
\GSfield{Assessment}{literal statement: correct; publication match: exact; verification: read the result; development: annotated; confidence: high}
\end{GSitem}
\begin{GSitem}{1981-536.3}{E}{The Catalan specialization and the missing product length}
\GSfield{Kind; source}{theorem; \GSpage{P000990}}
\GSfield{The notes}{Set \(k=n,\ell=n+1\); the notes identify the answer with the nth Catalan number and with parenthesizations of a nonassociative product, but do not write its number of factors.}
\GSfield{Statement}{For \(n\ge0\), coprimality gives \[\dim L_{n,n+1}=\frac1{2n+1}\binom{2n+1}{n}=\frac1{n+1}\binom{2n}{n}=C_n.\] The same integer counts full binary parenthesizations of a fixed word of length \(n+1\), rooted plane full binary trees with \(n\) internal vertices, and Dyck words with \(n\) pairs. This equality of counts by itself does not supply a basis of the Serre quotient or a canonical bijection from bracketings to Lie elements.}
\GSfield{Reconstruction}{The product length and the warning that a counting identity is not a quotient-basis theorem are supplied.}
\GSfield{Fate}{The assertion is an immediate specialization of the preceding pre-existing Witt formula \citep[pp.~225--226]{bermanmoody1979multiplicities}. The factorial simplification gives one check; the independently implemented Catalan recurrence \(C_0=1\), \(C_n=\sum_{j=0}^{n-1}C_jC_{n-1-j}\) gives another. The exact number \(n+1\) of factors is an editorial completion of the ellipsis.}
\GSfield{Present status}{As of November 2026, this is an established result with the hypotheses specified above; the notebook occurrence is expository.}
\GSfield{Assessment}{literal statement: correct with missing hypotheses; publication match: exact; verification: read the result; development: annotated; confidence: high}
\end{GSitem}
\begin{GSitem}{1981-536.4}{E}{The enveloping algebra and alternating associative Serre relations}
\GSfield{Kind; source}{construction; \GSpage{P000990}}
\GSfield{The notes}{The next question concerns the enveloping algebra. The two iterated brackets are expanded as alternating binomial sums of associative words.}
\GSfield{Statement}{If \(\mathfrak l=L(x,y)/(r_x,r_y)\) with \(r_x=(\operatorname{ad}x)^{a+1}y\), \(r_y=(\operatorname{ad}y)^{b+1}x\), then \[U(\mathfrak l)=\mathbb C\langle x,y\rangle\big/\left(\sum_{j=0}^{a+1}(-1)^j\binom{a+1}{j}x^{a+1-j}yx^j,\ \sum_{j=0}^{b+1}(-1)^j\binom{b+1}{j}y^{b+1-j}xy^j\right).\] These are two-sided ideal generators. The PBW theorem identifies the associated graded algebra for the PBW filtration with \(S(\mathfrak l)\); for the original positive root grading their Hilbert series also agree. It does not say that arbitrary unreduced words in \(x,y\) are a basis.}
\GSfield{Fate}{This is the universal property of enveloping algebras applied to a Lie presentation, together with PBW; the same Serre/PBW setting is reviewed in \citet[\S2]{macdonald1981affine}. The complete signs and upper limits here expand the notebook's ellipses rather than emend its two visible signs. The binomial identity follows because left and right multiplication by \(x\) commute, so \(\operatorname{ad}x=L_x-R_x\).}
\GSfield{Present status}{As of November 2026, this is an established result with the hypotheses specified above; the notebook occurrence is expository.}
\GSfield{Assessment}{literal statement: correct with missing hypotheses; publication match: exact; verification: read the result; development: annotated; confidence: high}
\end{GSitem}
\begin{GSitem}{1981-536.5}{E}{The symmetric-algebra Euler transform and divisor sums}
\GSfield{Kind; source}{theorem; \GSpage{P000990-P000991}}
\GSfield{The notes}{For a graded \(V\), the source first gives \(p(t)=\sum\dim V_k t^k\), then \(\log P(t)=\sum C_k t^k/k\), \(p(t)=\sum c_k t^k/k\), and \(C_k=\sum_{d\mid k}c_d\). The denominator in the latter expression for \(p\) is clearly written in the scan.}
\GSfield{Statement}{Assume \(V=\bigoplus_{k\ge1}V_k\) and every \(v_k=\dim V_k\) is finite. Then \[P(t)=\operatorname{Hilb}S(V)=\prod_{k\ge1}(1-t^k)^{-v_k},\qquad [t^n]\log P(t)=\frac1n\sum_{d\mid n}d v_d.\] Thus with \(c_k=kv_k\), the two notebook formulas for \(p(t)\) agree and \(C_n=\sum_{d\mid n}c_d\). Inversion is \(v_n=n^{-1}\sum_{d\mid n}\mu(d)C_{n/d}\). All identities are formal power-series identities; a nonzero degree-zero piece would invalidate finite coefficients of \(P(t)\).}
\GSfield{Reconstruction}{Positive grading, finite-dimensional pieces and the rescaled coefficient convention make the compressed formulas consistent.}
\GSfield{Fate}{This elementary Euler-transform computation is the commutative/PBW mechanism behind the logarithm and M\"obius inversion in \citet[proof of Theorem 2]{bermanmoody1979multiplicities}. In particular the source's \(c_k/k\) should \emph{not} be corrected to \(c_k\): its coefficient names have been rescaled. The notebook itself already allows a formal series after initially saying `Hilbert polynomial'.}
\GSfield{Present status}{As of November 2026, this is an established result with the hypotheses specified above; the notebook occurrence is expository.}
\GSfield{Assessment}{literal statement: correct with missing hypotheses; publication match: exact; verification: read the result; development: annotated; confidence: high}
\end{GSitem}
\begin{GSitem}{1981-536.6}{F}{The uncompleted lexicographic word method}
\GSfield{Kind; source}{method; \GSpage{P000991}}
\GSfield{The notes}{Words \(x^{k_1}y^{\ell_1}\cdots x^{k_s}y^{\ell_s}\), their two total degrees, and lexicographic order \(x<y\) are written, without a stated selection, bracketing or reduction rule.}
\GSfield{Statement}{No unique algorithm is specified. For a free Lie algebra, a possible completion would select Lyndon words and apply their standard recursive bracketing; for the Serre quotient one would additionally need reductions by the ideal, with a proof of spanning and independence. A total order on all associative words alone provides neither step, and the notebook does not state a termination or confluence claim.}
\GSfield{Reconstruction}{Lyndon bracketing is an explicitly labelled possible completion, not an attribution to Kirillov.}
\GSfield{Fate}{The preceding free-Lie dimension and enveloping-algebra discussion makes a Hall/Lyndon-type idea plausible, but neither the exact word selection nor a quotient normal-form theorem can be identified from the surviving three lines. Berman--Moody's multiplicity formula is only thematic context and does not identify the omitted algorithm. No modern basis construction is reported as a solution of an unspecified method.}
\GSfield{Present status}{As of November 2026, the precise proposed algorithm remains unidentified from the notes; no well-posed unresolved basis conjecture is extracted.}
\GSfield{Assessment}{literal statement: not determinable; publication match: none found; verification: unresolved lead; development: annotated; confidence: low}
\end{GSitem}
\begin{GSitem}{1981-536.7}{E}{Counting interlacing tables with four fixed top entries}
\GSfield{Kind; source}{problem; \GSpage{P000991}; I. M. Gelfand}
\GSfield{The notes}{Gelfand fixes \(m_1\le m_2\le m_3\le m_4\), asks how many integer triangular tables have each entry between its two upper neighbours including endpoints, and parenthetically calls this the number of representations of \(SL(4)\) with a given highest weight.}
\GSfield{Statement}{For fixed integers \(m_1\le\cdots\le m_4\), let \(a_{4,i}=m_i\) and impose \(a_{r+1,i}\le a_{r,i}\le a_{r+1,i+1}\) for \(1\le i\le r<4\). The number is \[N(m)=\prod_{1\le i<j\le4}\frac{m_j-m_i+j-i}{j-i}.\] Equivalently, \[\begin{aligned}12N(m)={}&(m_2-m_1+1)(m_3-m_1+2)(m_4-m_1+3)\\ &{}\cdot(m_3-m_2+1)(m_4-m_2+2)(m_4-m_3+1).\end{aligned}\] It is the dimension of the irreducible \(GL_4\)-module of highest weight \((m_4,m_3,m_2,m_1)\), or its irreducible restriction to \(SL_4\), because the patterns index a basis. There is one isomorphism class of irreducible representation for a fixed dominant integral \(SL_4\) highest weight; that number is not generally \(N(m)\).}
\GSfield{Reconstruction}{Integrality of every entry, reversal to the usual decreasing highest-weight convention, and the distinction between a basis count and a count of representations are explicit.}
\GSfield{Fate}{\citet[\S3, formula (3), reprint p.~655]{gelfandtsetlin1950unimodular} gives precisely the arbitrary-rank integer interlacing basis; its original article was submitted 16 February 1950 and published that year. The product is the finite type-A Weyl dimension specialization \citep[VII.8]{serre1987complex}. The notebook itself points to the Gelfand--Minlos--Shapiro appendix, but this identification is verified against the earlier primary paper. Direct recursive enumeration and the six-factor product agree for all top rows in the bounded computation supplied.}
\GSfield{Present status}{As of November 2026, this is a classical solved counting problem. The literal parenthetical identification in the notebook is false, although the intended dimension interpretation is exact.}
\GSfield{Assessment}{literal statement: false as written; publication match: exact; verification: read the result; development: annotated; confidence: high; problem \hyperref[prob:GP-1981-03]{GP-1981-03}}
\end{GSitem}

\GStalk{537}{1981-11-16}{A. S. Schwarz}{P000991-P000995}{Supergravity as the geometry of a real supersurface}
Schwarz treats real \((4|4)\)-dimensional supersurfaces in \(\mathbb C^{4|2}\), a holomorphic volume-preserving pseudogroup, the odd complex tangent distribution and its Levi form. The action is written as a delta-function surface integral; a preliminary exponent on P000993 is explicitly replaced by the Jacobian-corrected expression on P000994. An SCR structure retains a determinant normalization absent from ordinary CR data. The October 1981 announcement is identified bibliographically, and the detailed formulas are verified in \citet[\S3, pp.~44--48]{schwarz1982supergravity}, received 15 February 1982. The broader field-space-democracy work is not confused with this exact Levi determinant formula. The record ends before the 23 November heading on P000995.

\begin{GSitem}{1981-537.1}{E}{Superspace as a functor of Grassmann algebras}
\GSfield{Kind; source}{definition; \GSpage{P000992}}
\GSfield{The notes}{A real superspace is a row of even \(x^i\) and odd \(\theta^\alpha\) elements of a Grassmann algebra. The speaker insists on varying the algebra functorially instead of fixing it.}
\GSfield{Statement}{For a finite-dimensional real Grassmann algebra \(\Lambda=\Lambda_{\bar0}\oplus\Lambda_{\bar1}\), affine superspace has \(\mathbb R^{p|q}(\Lambda)=\Lambda_{\bar0}^{p}\times\Lambda_{\bar1}^{q}\). An even unital algebra homomorphism acts componentwise, giving a covariant functor of test algebras. Smooth supermaps are the natural maps induced by smooth functions in the even body variables and finite exterior polynomials in odd coordinates, with the appropriate smoothness structure on the functor. An arbitrary set-valued functor is not automatically a smooth supermanifold. Complexification introduces conjugation as additional real-structure data.}
\GSfield{Fate}{The local superfunction and integration framework predates the meeting: \citet[\S2, Definitions 1--4; \S3, Theorem 1]{berezinleites1975supermanifolds}, received 24 March 1975, was read in the original. Schwarz's exact functorial presentation is explained in \citet[\S1]{schwarz1982supergravity}. These are equivalent ways to organize the local model, but the 1975 sheaf definition is not falsely quoted as the same wording as the later functor exposition.}
\GSfield{Present status}{As of November 2026, this is an established result with the hypotheses specified above; the notebook occurrence is expository.}
\GSfield{Assessment}{literal statement: correct with missing hypotheses; publication match: exact; verification: read the result; development: annotated; confidence: high}
\end{GSitem}
\begin{GSitem}{1981-537.2}{E}{A real (4|4) supersurface and its odd complex tangent space}
\GSfield{Kind; source}{construction; \GSpage{P000992-P000993}}
\GSfield{The notes}{Four real even equations cut out \(\Omega\subset\mathbb C^{4,2}\), of real dimension \((4,4)\). Solving \(f^k_{x^a}dx^a+f^k_{\theta^\alpha}d\theta^\alpha=0\) gives a complex tangent space of dimension \((0,2)\).}
\GSfield{Statement}{Let \(z=(x^1,\ldots,x^4,\theta^1,\theta^2)\) be complex supercoordinates. Four real even functions \(f^k(z,\bar z)\) with independent real even body differentials define a real supersubmanifold of dimension \((4|4)\). On a generic chart assume additionally that the even \(4\times4\) matrix \(f_x=(\partial f^k/\partial x^a)\) has invertible body. Then \(T^{1,0}\Omega=\ker\partial f\) is locally free of rank \((0|2)\), with frame \[R_\alpha=\partial_{\theta^\alpha}-\sum_a(f_x^{-1}f_\theta)^a{}_{\alpha}\,\partial_{x^a}.\] Real regularity alone does not imply this complex genericity condition.}
\GSfield{Reconstruction}{The real regularity and separate invertibility of the even complex Jacobian are supplied. The explicit frame is obtained by the elimination written in the notes.}
\GSfield{Fate}{The regular-level-set mechanism is already \citet[\S4, Theorem 2]{berezinleites1975supermanifolds}; the generic supersurface and its complex tangent plane are treated explicitly in \citet[\S3, equations (1)--(3), pp.~44--46]{schwarz1982supergravity}. The notebook's linear elimination is valid on precisely this invertible chart. The dimension \((0|2)\) describes the tangent distribution, not the rank or dimension of the four-component Levi form.}
\GSfield{Present status}{As of November 2026, this is an established result with the hypotheses specified above; the notebook occurrence is expository.}
\GSfield{Assessment}{literal statement: correct with missing hypotheses; publication match: exact; verification: read the result; development: annotated; confidence: high}
\end{GSitem}
\begin{GSitem}{1981-537.3}{E}{Holomorphic transformations preserving the Berezin volume}
\GSfield{Kind; source}{definition; \GSpage{P000992-P000993}}
\GSfield{The notes}{Holomorphic coordinate transformations form \(L\); the notes then restrict to those satisfying \(\operatorname{Ber}(\partial\widetilde z/\partial z)=1\), and later again call the symmetry group \(L\).}
\GSfield{Statement}{Use \(L_{\mathrm{vol}}\) for the local holomorphic invertible coordinate changes preserving the chosen nonvanishing holomorphic Berezin volume. In a volume coordinate chart they satisfy \(\operatorname{Ber}J=1\). For an even block supermatrix \(J=\bigl(\begin{smallmatrix}A&B\\C&D\end{smallmatrix}\bigr)\), with \(D\) invertible, \(\operatorname{Ber}J=\det(A-BD^{-1}C)/\det D\). Its multiplicativity gives a consistent transformation law for Berezin integration. All local holomorphic transformations form a larger pseudogroup; a fixed-domain global group requires additional domain conditions.}
\GSfield{Reconstruction}{The volume-preserving subgroup notation separates the two successive meanings of L without changing the intended volume symmetry.}
\GSfield{Fate}{The Berezinian density rule is explicit in \citet[\S3, equation (3) and Theorem 1]{berezinleites1975supermanifolds}. The specific supergravity pseudogroup is used in \citet[\S3]{schwarz1982supergravity}. The passage from all holomorphic transformations to the volume-preserving subgroup is present in the source; the later symbol reuse is the ambiguity, not the determinant condition itself.}
\GSfield{Present status}{As of November 2026, this is an established result with the hypotheses specified above; the notebook occurrence is expository.}
\GSfield{Assessment}{literal statement: correct with missing hypotheses; publication match: exact; verification: read the result; development: annotated; confidence: high}
\end{GSitem}
\begin{GSitem}{1981-537.4}{E}{The graph superfield and field-space democracy}
\GSfield{Kind; source}{construction; \GSpage{P000993-P000994}}
\GSfield{The notes}{The supersurface is put in one-to-one correspondence with a field \(H^m(\xi,\theta,\bar\theta):\mathbb R^{4,4}\to\mathbb R^{4,0}\). The closing remark says that coordinates and field values mix under the ambient transformations.}
\GSfield{Statement}{Write complex even coordinates locally as \(x^m=\xi^m+i\eta^m\). If projection of \(\Omega\) to \((\xi,\theta,\bar\theta)\) is locally invertible, it has a unique graph \(\eta^m=H^m(\xi,\theta,\bar\theta)\). This is a local equivalence in a chosen projection chart; it is not a global statement for every embedded supersurface. A holomorphic transformation of \((x,\theta)\) can mix \(\xi\) and \(H\), so the field transforms by re-expressing the transformed surface as a graph wherever that projection remains regular.}
\GSfield{Reconstruction}{The projection condition, real/imaginary decomposition and local scope are supplied; the invariance concerns the surface description.}
\GSfield{Fate}{The geometric reformulation is part of Schwarz's earlier field-space-democracy program, bibliographically identified as \citet{schwarz1980equal}. The graph construction and its supergravity meaning are read directly in \citet[\S3]{schwarz1982supergravity}; the 1980 article itself was not read. This is a classical local implicit-function argument within the stated supergeometric model, not a claim of a global trivialization of all possible field configurations.}
\GSfield{Present status}{As of November 2026, this is an established result with the hypotheses specified above; the notebook occurrence is expository.}
\GSfield{Assessment}{literal statement: correct with missing hypotheses; publication match: exact; verification: read the result; development: annotated; confidence: high}
\end{GSitem}
\begin{GSitem}{1981-537.5}{A}{The vector-valued Levi form and its four-by-four determinant}
\GSfield{Kind; source}{construction; \GSpage{P000993-P000994}}
\GSfield{The notes}{Restrict \(\partial_A\partial_{\bar B}f^k\,dz^A d\bar z^B\) to the complex tangent plane. Its components are \(\Gamma^k_{\alpha\dot\beta}\); the index \(k\) has four values and the paired spinor index has \(2\times2=4\).}
\GSfield{Statement}{On the generic chart of 1981-537.2, evaluate the mixed Hessian of each \(f^k\) on \(R_\alpha,\overline{R_\beta}\), with a fixed convention for odd derivatives and any conventional overall factor of \(i\). This is the normal-valued Levi tensor. In the chosen normal and odd tangent frames it is a \(4\times4\) matrix \(\Gamma=(\Gamma^k_{\alpha\dot\beta})\), after ordering the four pairs \((\alpha,\dot\beta)\). Its determinant is a relative invariant of the frames, not a scalar independent of them. The supergravity density discussed next is defined on the locus where the body of this determinant is nonzero; fractional powers need a consistent real/complex branch.}
\GSfield{Reconstruction}{Invertibility, the paired-index ordering and branch conventions are supplied. Changing the conventional Levi factor changes only an overall normalization of the action.}
\GSfield{Fate}{\citet[\S3, equations (3)--(8), pp.~45--47]{schwarz1982supergravity} defines the Levi form, tracks changes of defining equations and tangent frame, and constructs the compensating invariant density. The October 1981 announcement \citep{schwarz1981complex} is an earlier identified publication in this line, but only its metadata was checked. Classification A refers to the detailed post-seminar treatment, not a claim that the Levi form itself originated in 1982.}
\GSfield{Present status}{As of November 2026, the construction and covariance law are explicit in the identified primary paper. The notebook does not determine a numerical choice of overall Levi normalization.}
\GSfield{Assessment}{literal statement: correct with missing hypotheses; publication match: exact; verification: read the result; development: annotated; confidence: high}
\end{GSitem}
\begin{GSitem}{1981-537.6}{A}{The invariant action density and the corrected exponents}
\GSfield{Kind; source}{theorem; \GSpage{P000993-P000994}}
\GSfield{The notes}{P000993 first writes \(A=|\det\Gamma|^{-1/2}\). P000994 then says \ru{и тогда надо писать} and gives \(A=(\det f_z)^{4/3}|\det\Gamma|^{-1/3}\). Both exponents and the added Jacobian are visible in the scans.}
\GSfield{Statement}{For a regular generic supersurface on the nondegenerate Levi locus, the invariant density in the natural odd-coordinate frame is, up to a constant and orientation conventions, \[A=|\det f_x|^{4/3}|\det\Gamma|^{-1/3},\qquad S(\Omega)=C\int A\,\prod_{k=1}^4\delta(f^k)\,Dz\,D\bar z.\] Here \(x^a\) denotes the four complex even coordinates, not all six supercoordinates; \(D\) is Berezin integration notation. Under replacement of equations by \(f'=\eta f\), the two determinants each acquire \(\det\eta\) on the surface, while the four-delta product acquires \(|\det\eta|^{-1}\); the exponents cancel. In a general tangent frame the first factor is Schwarz's compensator \(|W|^{4/3}\). Thus the action depends on the embedded surface and preserved ambient volume, and is invariant under \(L_{\mathrm{vol}}\).}
\GSfield{Reconstruction}{The first determinant is explicitly the even complex Jacobian; absolute values or an equivalent oriented-branch convention are supplied. The change of formula inside the notebook is retained as evidence of an internal correction.}
\GSfield{Fate}{This is \citet[\S3, equations (7)--(8), pp.~46--47]{schwarz1982supergravity}, received 15 February 1982; the author gives the defining-equation and tangent-frame covariance checks. The action had been announced in the October 1981 \emph{Yadernaya Fizika} paper \citep{schwarz1981complex}. The initial minus-one-half formula is not used as a second valid expression: under \(f\mapsto cf\) it has the wrong scaling in the displayed natural frame. The published formula also supplies absolute values/real-density conventions missing on the first determinant in the notebook.}
\GSfield{Present status}{As of November 2026, the corrected first/second-derivative action density is established in the cited primary source. The preliminary minus-one-half density is false under the explicitly chosen frames and equation-change covariance.}
\GSfield{Assessment}{literal statement: false as written; publication match: exact; verification: read the result; development: annotated; confidence: high}
\end{GSitem}
\begin{GSitem}{1981-537.7}{A}{Uniqueness within the stated first/second-derivative action class}
\GSfield{Kind; source}{claim; \GSpage{P000993}}
\GSfield{The notes}{The action is said to be the unique functional of the displayed kind, up to a multiplicative constant, invariant under \(L\). The following line about variation is unfinished and gives no field equation.}
\GSfield{Statement}{The precise identified claim restricts to local actions of generic nondegenerate \((4|4)\)-supersurfaces whose density depends on their defining functions through first and second derivatives, is independent of the choice of defining equations, and is invariant under holomorphic transformations preserving the ambient Berezin volume. Within that class the source asserts uniqueness up to a constant. At the determinant-monomial level, writing \(A=|W|^p|\Gamma|^q\), equation covariance and tangent-frame covariance impose \(p+q=1\) and \(p+4q=0\), hence \(p=4/3,q=-1/3\). This elementary exponent calculation does not, by itself, classify all admissible local functions of the jets.}
\GSfield{Reconstruction}{The derivative-order and geometric domain restrictions come from the identified paper. The two-exponent check is editorial; the stronger classification argument has not been independently read.}
\GSfield{Fate}{\citet[\S3, p.~48, following equation (9)]{schwarz1982supergravity} explicitly states the restricted uniqueness result, and at p.~58 refers its proof to \citet{gaydukromanovschwarz1981democracy}. The latter's bibliographic record and abstract were verified, but its proof was inaccessible in this search; therefore the full uniqueness proof is recorded as indirect verification. The 1982 paper's symmetry and determinant calculations were read. No uniqueness theorem for arbitrary derivative orders, extra matter fields or other supergravity models is inferred, and no Euler--Lagrange system is reconstructed from the unfinished sentence.}
\GSfield{Present status}{As of November 2026, the exact published uniqueness claim is identified with its stated action class; its original full proof remains an access limitation of this audit, not an open mathematical problem.}
\GSfield{Assessment}{literal statement: correct with missing hypotheses; publication match: exact; verification: indirect; development: annotated; confidence: medium}
\end{GSitem}
\begin{GSitem}{1981-537.8}{A}{SCR structure as CR data with ambient-volume normalization}
\GSfield{Kind; source}{construction; \GSpage{P000994}}
\GSfield{The notes}{A complex G-structure induces a CR structure on the surface. Here the ambient group is written \(SL(4,2\mid\mathbb C)\), and the induced SCR structure differs from CR by a determinant condition.}
\GSfield{Statement}{On a complex \((4|2)\)-supermanifold with a chosen holomorphic Berezin volume, the compatible ambient frames form an \(SL(4|2;\mathbb C)\)-reduction, where `SL' means Berezinian one. A generic real \((4|4)\)-supersurface inherits its odd complex tangent CR distribution together with the additional frame normalization determined by that volume. This is Schwarz's SCR structure. Forgetting the normalization yields the CR data; retaining it makes the determinant-compensated action intrinsic to the induced structure. The local automorphisms preserve both the holomorphic structure and the chosen volume.}
\GSfield{Reconstruction}{The modern vertical-bar notation and the distinction between ordinary determinant and Berezinian are supplied; the role of a fixed ambient volume is made explicit.}
\GSfield{Fate}{\citet[\S\S4--6, especially the SCR frame construction in \S6]{schwarz1982supergravity} gives the detailed G-structure formulation soon after the seminar. \citet[\S\S1--3]{konechnyschwarz1997geometry}, first circulated in September 1996 and published in 1997, subsequently uses these SCR superspaces for \(N=1\) super Yang--Mills theory in curved superspace. That is a concrete later application, not a purported first proof of the 1981 construction or a solution of an unstated conjecture.}
\GSfield{Present status}{As of November 2026, the SCR/G-structure construction and the cited later gauge-theory application are identified. The brief source does not assert a classification of all abstract SCR structures or a global embedding theorem.}
\GSfield{Assessment}{literal statement: correct with missing hypotheses; publication match: exact; verification: read the result; development: annotated; confidence: high}
\end{GSitem}

\GStalk{538}{1981-11-23}{A. V. Zelevinsky}{P000993-P000997}{Verma-module multiplicities and Kazhdan--Lusztig polynomials}
The actual section begins on P000995 and continues onto P000996, before Beilinson speaks. It recalls highest-weight modules and finite-length Verma modules, reduces composition multiplicities to a regular integral Weyl orbit, and reports the proved Kazhdan--Lusztig multiplicity conjecture. A subsequent coefficient interpretation anticipates the separate Jantzen-filtration discussion. The catalog page span includes earlier Schwarz material; it is not a mathematical boundary.

\begin{GSitem}{1981-538.1}{E}{Triangular decomposition and simple-root generators}
\GSfield{Kind; source}{definition; \GSpage{P000995}}
\GSfield{The notes}{The section begins with \(\mathfrak g\) complex semisimple, its Cartan subalgebra, \(\dim\mathfrak h=n\), generators \(e_i,h_i,f_i\), and schematic brackets identifying \(\mathfrak n_+\) and a Borel subalgebra.}
\GSfield{Statement}{Fix a Cartan subalgebra \(\mathfrak h\), positive roots, and simple roots \(\alpha_1,\ldots,\alpha_n\). Then \(\mathfrak g=\mathfrak n_-\oplus\mathfrak h\oplus\mathfrak n_+\), \(\mathfrak b=\mathfrak h\oplus\mathfrak n_+\), \(h_i=\alpha_i^\vee\), and \([e_i,f_j]=\delta_{ij}h_i\). The \(e_i\) generate \(\mathfrak n_+\) as a Lie algebra; their linear span need not equal it.}
\GSfield{Reconstruction}{The root system, positive system, and distinction between generating and spanning are supplied to make the schematic drawing precise.}
\GSfield{Fate}{This is the classical setup recalled in \citet[\S1]{brylinski1981differential}; the 1981 article reports a lecture delivered in August--September 1980.}
\GSfield{Present status}{As of 11 November 2026: standard structure theory, not a new seminar claim.}
\GSfield{Assessment}{literal statement: correct with missing hypotheses; publication match: exact; verification: read the result; development: annotated; confidence: high}
\end{GSitem}
\begin{GSitem}{1981-538.2}{E}{Highest vectors and highest weights}
\GSfield{Kind; source}{definition; \GSpage{P000995}}
\GSfield{The notes}{A highest vector is first characterized by \(e_i v=0\); the notes observe that \(h_jv\) has the same property. Later they impose \(hv=\chi(h)v\).}
\GSfield{Statement}{A highest-weight vector of weight \(\chi\in\mathfrak h^*\) is a nonzero vector killed by \(\mathfrak n_+\) and simultaneously satisfying \(hv=\chi(h)v\) for every \(h\in\mathfrak h\). The common kernel of the \(e_i\) is \(\mathfrak h\)-stable because \([e_i,h_j]\) is a scalar multiple of \(e_i\). In any nonzero finite-dimensional \(\mathfrak g\)-module this kernel contains a common \(\mathfrak h\)-eigenvector.}
\GSfield{Reconstruction}{Being killed by the raising operators alone does not select an eigenvector in an arbitrary module. Finite dimension, or the weight-module hypothesis, is needed for the stated extraction of a highest weight.}
\GSfield{Fate}{The standard highest-weight construction is recalled in \citet[\S1]{brylinski1981differential}. The kernel calculation here explains the notes\textquotesingle\ omitted transition.}
\GSfield{Present status}{As of 11 November 2026: correct after distinguishing a singular vector from a chosen highest-weight eigenvector.}
\GSfield{Assessment}{literal statement: correct with missing hypotheses; publication match: exact; verification: read the result; development: annotated; confidence: high}
\end{GSitem}
\begin{GSitem}{1981-538.3}{E}{Verma module and its unique simple quotient}
\GSfield{Kind; source}{construction; \GSpage{P000995}}
\GSfield{The notes}{The notes introduce a universal representation \(M(\chi)\) with highest vector of weight \(\chi\), and its unique irreducible quotient \(L(\chi)\).}
\GSfield{Statement}{Let \(\mathbb C_\chi\) be the one-dimensional \(\mathfrak b\)-module on which \(\mathfrak n_+\) acts by zero and \(h\) by \(\chi(h)\). Then \(M(\chi)=U(\mathfrak g)\otimes_{U(\mathfrak b)}\mathbb C_\chi\). It is generated by \(v_\chi\), and every module generated by a highest vector of this weight is a quotient. Its sum of proper submodules is proper, since none meets the one-dimensional top weight space, and hence it has a unique simple quotient \(L(\chi)\).}
\GSfield{Reconstruction}{The induction formula and the explanation of uniqueness are explicit reconstructions of the compressed universal-property statement.}
\GSfield{Fate}{Already part of highest-weight theory before the talk; see \citet[\S1]{brylinski1981differential} and the category construction in \citet[\S1]{bernsteingelfandgelfand1976category}.}
\GSfield{Present status}{As of 11 November 2026: established for every complex weight; finite-dimensionality of the simple quotient requires dominant integrality.}
\GSfield{Assessment}{literal statement: correct with missing hypotheses; publication match: exact; verification: read the result; development: annotated; confidence: high}
\end{GSitem}
\begin{GSitem}{1981-538.4}{E}{Finite length versus arbitrary subquotients}
\GSfield{Kind; source}{claim; \GSpage{P000995}}
\GSfield{The notes}{The notes say that \(M(\chi)\) has a finite composition series and that every subquotient has the form \(L(\psi)\).}
\GSfield{Statement}{For finite-dimensional complex semisimple \(\mathfrak g\), every Verma module has finite Jordan--Hölder length, and every simple composition factor is a highest-weight module \(L(\psi)\). An arbitrary subquotient need not be simple, and therefore need not be isomorphic to any \(L(\psi)\).}
\GSfield{Reconstruction}{The necessary emendation is ``simple subquotient\textquotesingle\ or ``composition factor\textquotesingle. It is a correction to the mathematical wording, not an alleged mistranscription.}
\GSfield{Fate}{Finite length is recalled explicitly in \citet[Introduction]{gabberjoseph1981towards}. The article appeared in 1981, and the result was standard before the seminar.}
\GSfield{Present status}{As of 11 November 2026: the finite-length claim is established; the unrestricted subquotient clause is false.}
\GSfield{Assessment}{literal statement: false as written; publication match: exact; verification: read the result; development: annotated; confidence: high}
\end{GSitem}
\begin{GSitem}{1981-538.5}{E}{Composition multiplicities in a Verma module}
\GSfield{Kind; source}{problem; \GSpage{P000995}}
\GSfield{The notes}{The first question asks for the multiplicity of \(M(\psi)\) in a composition series of \(M(\chi)\), but the next displayed definition uses \(L(\psi)\). Both the Russian and the scan contain this change of letter.}
\GSfield{Statement}{Determine \(m(\chi,\psi)=[M(\chi):L(\psi)]\), the Jordan--Hölder multiplicity of the simple module \(L(\psi)\) in \(M(\chi)\). A reducible Verma module cannot itself be a composition factor.}
\GSfield{Reconstruction}{The immediately following formula fixes the intended question unambiguously: the earlier \(M(\psi)\) must be emended to \(L(\psi)\).}
\GSfield{Fate}{This is a recurrence of GP-1979-41. The multiplicities in a regular integral block were determined by the Kazhdan--Lusztig theorem, whose proof was announced in 1980 and published in 1981; \citet[Introduction and \S1]{brylinski1981differential} records the proof context.}
\GSfield{Present status}{As of 11 November 2026: settled in the finite-dimensional semisimple setting; the source slip creates no separate problem.}
\GSfield{Assessment}{literal statement: false as written; publication match: exact; verification: read the result; development: annotated; confidence: high; problem \hyperref[prob:GP-1979-41]{GP-1979-41}}
\end{GSitem}
\begin{GSitem}{1981-538.6}{E}{Central-character linkage and the rho shift}
\GSfield{Kind; source}{theorem; \GSpage{P000995-P000996}}
\GSfield{The notes}{Nonzero multiplicity is said to require that \(\chi+\rho\) and \(\psi+\rho\) lie in one Weyl orbit; \(\rho(h_i)=1\). The subsequent definition of dominant integral weight changes the letter \(\chi\) to \(\lambda\).}
\GSfield{Statement}{Write \(w\cdot\chi=w(\chi+\rho)-\rho\), where \(\rho\) is half the sum of positive roots. The center of \(U(\mathfrak g)\) acts on \(M(\chi)\) by its Harish-Chandra character. Thus \([M(\chi):L(\psi)]\ne0\) implies \(\psi\in W\cdot\chi\). Dominant integral means \(\chi(h_i)\in\mathbb Z_{\ge0}\) for all simple coroots. Central-character equality alone is not sufficient for occurrence.}
\GSfield{Reconstruction}{The conclusion is explicitly a necessary condition. Renaming the variable in the integrality clause repairs notation without changing content.}
\GSfield{Fate}{The central-character parametrization and shifted Weyl action are stated in \citet[\S1, pp.~45--47]{brylinski1981differential}.}
\GSfield{Present status}{As of 11 November 2026: established; weights outside the regular integral case need their own block conventions.}
\GSfield{Assessment}{literal statement: correct with missing hypotheses; publication match: exact; verification: read the result; development: annotated; confidence: high}
\end{GSitem}
\begin{GSitem}{1981-538.7}{E}{Independence of the dominant integral base point}
\GSfield{Kind; source}{theorem; \GSpage{P000996}}
\GSfield{The notes}{For \(\chi=w_1(\chi_0+\rho)-\rho\) and \(\psi=w_2(\chi_0+\rho)-\rho\), the multiplicity is asserted to depend only on \(w_1,w_2\), when \(\chi_0\) is dominant integral.}
\GSfield{Statement}{For \(\lambda,\lambda^{\prime}\) dominant integral and \(x,y\in W\), one has \([M(x\cdot\lambda):L(y\cdot\lambda)]=[M(x\cdot\lambda^{\prime}):L(y\cdot\lambda^{\prime})]\). Here \(\lambda+\rho\) and \(\lambda^{\prime}+\rho\) are strictly dominant, so both infinitesimal characters are regular. Translation equivalences identify these regular integral blocks.}
\GSfield{Reconstruction}{The unshifted dominant weight may be on a chamber wall: regularity concerns \(\lambda+\rho\). This prevents an erroneous extra requirement that \(\lambda\) itself be strictly dominant.}
\GSfield{Fate}{This is the translation-principle setting treated in \citet[pp.~42--77]{jantzen1979moduln}; the table of contents and 1979 publication were checked, while the equality also follows from the explicit multiplicity theorem in \citet[\S3]{lusztig1983singularities}.}
\GSfield{Present status}{As of 11 November 2026: established in regular integral blocks; singular block analogues require parabolic modifications.}
\GSfield{Assessment}{literal statement: correct with missing hypotheses; publication match: exact; verification: read the result; development: annotated; confidence: high}
\end{GSitem}
\begin{GSitem}{1981-538.8}{E}{Kazhdan--Lusztig values give Verma multiplicities}
\GSfield{Kind; source}{theorem; \GSpage{P000996}}
\GSfield{The notes}{The notebook writes \(m(w_1,w_2)=P_{w_1,w_2}(1)\), notes nonnegative polynomial coefficients, and attributes the proof to Beilinson--Bernstein and Brylinski--Kashiwara.}
\GSfield{Statement}{For a finite Weyl group \(W\), ordinary Kazhdan--Lusztig polynomials normalized by \(P_{x,x}=1\), and a dominant integral \(\lambda\), \([M(x\cdot\lambda):L(y\cdot\lambda)]=P_{x,y}(1)\). The right side is zero unless \(x\le y\) in Bruhat order. This dominant-weight convention matters: formulas written with an antidominant base point or for characters involve changed indices and sometimes the longest element.}
\GSfield{Reconstruction}{The labels are fixed explicitly; one must not import a longest-element reversal from a differently normalized character formula.}
\GSfield{Fate}{The notes correctly report a proved theorem. \citet{beilinsonbernstein1981localisation} appeared in 1981 and \citet{brylinskikashiwara1981holonomic} in October 1981. Brylinski\textquotesingle s account derives from the August--September 1980 Toruń meeting. The precise character and multiplicity conventions are compared in \citet[\S3]{lusztig1983singularities}.}
\GSfield{Present status}{As of 11 November 2026: established well before 23 November 1981; this occurrence is expository rather than a new conjecture.}
\GSfield{Assessment}{literal statement: correct with missing hypotheses; publication match: exact; verification: read the result; development: annotated; confidence: high; problem \hyperref[prob:GP-1979-41]{GP-1979-41}}
\end{GSitem}
\begin{GSitem}{1981-538.9}{B}{What the coefficients remember beyond total multiplicity}
\GSfield{Kind; source}{question; \GSpage{P000996}}
\GSfield{The notes}{After \(P(q)=a_0+a_1q+\cdots\), the notes seek an interpretation of the individual \(a_i\). This question precedes the Beilinson heading.}
\GSfield{Statement}{The coefficient question asks for a natural filtration or geometric grading whose graded multiplicities recover \(P_{x,y}(q)\), rather than only its value at \(q=1\). In the representation-theoretic answer, Jantzen layers occur with a length-dependent shift and twice the polynomial degree; see 1981-540.6. Intersection cohomology gives a separate grading; see 1981-540.7.}
\GSfield{Reconstruction}{This is a question about a refinement of the already proved multiplicity theorem, not a second statement of the same conjecture. The notebook\textquotesingle s subsequent formula does not correctly normalize that refinement.}
\GSfield{Fate}{The filtration refinement was established by \citet{beilinsonbernstein1993jantzen}. The 1981 intermediate state is documented in \citet[Introduction]{gabberjoseph1981towards}; no verified prepublication date for the full 1993 theorem was found.}
\GSfield{Present status}{As of 11 November 2026: settled for ordinary Verma modules with the usual dominant deformation; it does not assert arbitrary-direction invariance.}
\GSfield{Assessment}{literal statement: correct with missing hypotheses; publication match: exact; verification: read the result; development: annotated; confidence: high; problem \hyperref[prob:GP-1979-41]{GP-1979-41}}
\end{GSitem}

\GStalk{539}{1981-11-23}{A. A. Beilinson}{P000995-P000997}{Category O and restricted contragredient duality}
The named intervention occupies P000996, between the coefficient question and S. Gelfand. It introduces category \(\mathcal O\), restricted duality, and the dual Verma module as preparation for the Jantzen filtration. The notebook confuses the sign of the locally nilpotent nilradical and does not spell out the anti-involution required for a duality preserving highest-weight category \(\mathcal O\).

\begin{GSitem}{1981-539.1}{E}{Category O and the sign of the nilradical}
\GSfield{Kind; source}{definition; \GSpage{P000996}}
\GSfield{The notes}{Beilinson defines a category of finitely generated modules with finite-dimensional Cartan behavior and local nilpotence of a nilradical. A marginal correction says \ru{Здесь только \(\mathfrak n_+\), а не \(\mathfrak n_-\), но это безразлично}.}
\GSfield{Statement}{For the chosen Borel, \(\mathcal O\) consists of finitely generated \(U(\mathfrak g)\)-modules that are \(\mathfrak h\)-semisimple and locally \(U(\mathfrak n_+)\)-finite. Their weight spaces are finite dimensional; positive root vectors act locally nilpotently. Reversing the Borel gives an equivalent category, but does not make local nilpotence of \(\mathfrak n_+\) and \(\mathfrak n_-\) equivalent on the same module.}
\GSfield{Reconstruction}{The marginal phrase is defensible only as a choice-of-Borel convention. On a generic Verma module, the lowering generators are not locally nilpotent.}
\GSfield{Fate}{The original definition and basic properties occur in \citet[\S1]{bernsteingelfandgelfand1976category}, received 22 December 1975 and published in 1976. Its full Russian text was inspected.}
\GSfield{Present status}{As of 11 November 2026: a standard abelian highest-weight category, with the chosen positive nilradical part of its definition.}
\GSfield{Assessment}{literal statement: correct with missing hypotheses; publication match: exact; verification: read the result; development: annotated; confidence: high}
\end{GSitem}
\begin{GSitem}{1981-539.2}{E}{Restricted contragredient duality}
\GSfield{Kind; source}{construction; \GSpage{P000996}}
\GSfield{The notes}{The notes discuss the dual of a Verma module, warn that the full dual can be too large, and select functions transforming finite-dimensionally under \(\mathfrak h\).}
\GSfield{Statement}{If \(M=\bigoplus_\mu M_\mu\in\mathcal O\), set \(dM=\bigoplus_\mu M_\mu^*\), not the product \(\prod_\mu M_\mu^*\). With the Chevalley anti-involution \(\tau(e_i)=f_i\), \(\tau(f_i)=e_i\), \(\tau(h)=h\), define \((x\phi)(m)=\phi(\tau(x)m)\). Then \(d\) is an exact contravariant involution of \(\mathcal O\), and \(dM(\chi)\) is the costandard module. An ordinary untwisted dual has different highest-weight conventions.}
\GSfield{Reconstruction}{The restriction to the direct sum of weight-space duals and the anti-involution are essential additions. The notebook does not define the action on its functional subspace.}
\GSfield{Fate}{This is the duality of \citet[\S1]{bernsteingelfandgelfand1976category}; it is precisely what permits the contravariant-form construction in the next intervention.}
\GSfield{Present status}{As of 11 November 2026: established; the full algebraic dual is generally outside category O.}
\GSfield{Assessment}{literal statement: correct with missing hypotheses; publication match: exact; verification: read the result; development: annotated; confidence: high}
\end{GSitem}

\GStalk{540}{1981-11-23}{S. I. Gelfand}{P000995-P000998}{Contravariant forms, Jantzen layers and deformation directions}
The named section starts at the foot of P000996 and ends above the new Zelevinsky heading on P000997. It describes the contravariant form, identifies Verma modules through PBW vectors, deforms the weight by \(t\rho\), and introduces the order-of-vanishing filtration. It asks whether changing \(\rho\) preserves layer multiplicities and links coefficients to Schubert geometry. Moving lifts, normalization of polynomial degrees, and the distinction between ordinary and intersection cohomology require explicit repair.

\begin{GSitem}{1981-540.1}{E}{The normalized contravariant Shapovalov form}
\GSfield{Kind; source}{construction; \GSpage{P000996-P000997}}
\GSfield{The notes}{A form is defined on monomials \(f_{i_1}\cdots f_{i_k}v\), with zero for different “unordered sets” of indices; the next page calls it \ru{инвар. отн. алг. Ли}.}
\GSfield{Statement}{The normalized form on \(M(\chi)\) satisfies \(S_\chi(v_\chi,v_\chi)=1\) and \(S_\chi(xm,n)=S_\chi(m,\tau(x)n)\), where \(\tau\) is the Chevalley anti-involution. Thus \[S_\chi(uv_\chi,wv_\chi)=\chi(\operatorname{HC}(\tau(u)w)).\] Here the Harish-Chandra projection is followed by evaluation on \(U(\mathfrak h)\). Distinct weight spaces are orthogonal. Since \(\tau\) reverses products, \(\tau(f_{i_1}\cdots f_{i_k})=e_{i_k}\cdots e_{i_1}\).}
\GSfield{Reconstruction}{“Invariant” must mean contravariant, not \(S(xm,n)+S(m,xn)=0\). The orthogonality test is equality of total root weight, hence of multiplicities of simple generators, not equality of sets with repetitions discarded.}
\GSfield{Fate}{The named form predates the talk: \citet{shapovalov1972bilinear}, received 15 March 1972. Its deformation use is explicit in \citet[\S1]{gabberjoseph1981towards}. The Shapovalov publisher record was checked; the Gabber--Joseph construction was read.}
\GSfield{Present status}{As of 11 November 2026: standard; the notebook omits the anti-involution and uses imprecise orthogonality and invariance terminology.}
\GSfield{Assessment}{literal statement: false as written; publication match: exact; verification: read the result; development: annotated; confidence: high}
\end{GSitem}
\begin{GSitem}{1981-540.2}{E}{PBW trivialization of the family of Verma modules}
\GSfield{Kind; source}{construction; \GSpage{P000997}}
\GSfield{The notes}{The notebook identifies \(f_{i_1}\cdots f_{i_l}v_\chi\) with the same monomial applied to \(v_{\chi^{\prime}}\).}
\GSfield{Statement}{PBW gives a vector-space isomorphism \(U(\mathfrak n_-)\simeq M(\chi)\), \(u\mapsto uv_\chi\), for every \(\chi\). Consequently all Verma modules can be put on one underlying vector space, while the action of \(\mathfrak h\oplus\mathfrak n_+\), and the contravariant form, vary with the parameter. This is not an isomorphism of \(\mathfrak g\)-modules when highest weights differ.}
\GSfield{Reconstruction}{The simple-root monomials span but are not a PBW basis without an ordering and composite-root vectors. PBW proves the proposed correspondence is well defined despite that presentation.}
\GSfield{Fate}{The deformation framework uses precisely this universal induced-module construction; see \citet[\S1]{gabberjoseph1981towards}.}
\GSfield{Present status}{As of 11 November 2026: valid as a vector-space trivialization.}
\GSfield{Assessment}{literal statement: correct with missing hypotheses; publication match: exact; verification: read the result; development: annotated; confidence: high}
\end{GSitem}
\begin{GSitem}{1981-540.3}{E}{Generic nondegeneracy along the rho deformation}
\GSfield{Kind; source}{claim; \GSpage{P000997}}
\GSfield{The notes}{For \(M(\chi+t\rho)\), the notes assert that the forms \((\ ,\ )_t\) are nondegenerate for small \(t\ne0\).}
\GSfield{Statement}{For a finite-dimensional semisimple \(\mathfrak g\), reducibility of \(M(\chi+t\rho)\) can occur only when \(\langle\chi+\rho+t\rho,\alpha^\vee\rangle\in\mathbb Z_{>0}\) for a positive root \(\alpha\). Since \(\langle\rho,\alpha^\vee\rangle>0\), the exceptional values form a locally discrete subset of \(\mathbb C\). Thus a sufficiently small punctured disk about zero contains none. Equivalently, the form is nondegenerate over \(\mathbb C(t)\).}
\GSfield{Reconstruction}{The formulation supplies the parameter field and the finite-root-system hypothesis. For a general direction \(\gamma\), it must at least be transverse to any reducibility hyperplane containing \(\chi\).}
\GSfield{Fate}{This is the determinant/deformation method underlying \citet{shapovalov1972bilinear} and \citet[\S1]{gabberjoseph1981towards}; it is not the later theorem about filtration multiplicities.}
\GSfield{Present status}{As of 11 November 2026: correct in the finite-dimensional setting with the stated regular deformation.}
\GSfield{Assessment}{literal statement: correct with missing hypotheses; publication match: exact; verification: read the result; development: annotated; confidence: high}
\end{GSitem}
\begin{GSitem}{1981-540.4}{E}{The Jantzen filtration requires parameter-dependent lifts}
\GSfield{Kind; source}{definition; \GSpage{P000997}}
\GSfield{The notes}{The literal definition is \(M_i=\{x\in M_\chi:(x,y)_t=O(t^i)\text{ for all }y\in M(\chi)\}\), with vectors regarded as fixed under the PBW identification. The next sentence says these are \(\mathfrak g\)-submodules.}
\GSfield{Statement}{Put \(A=\mathbb C[t]_{(t)}\) and let \(M_A\) be the deformed Verma module of highest weight \(\chi+t\gamma\). Define \(M_A^{(i)}=\{m\in M_A:S_A(m,M_A)\subset t^iA\}\) and \(J^iM=(M_A^{(i)}+tM_A)/tM_A\). Equivalently, a vector of \(M\) belongs to \(J^i\) if it admits a lift depending on \(t\) whose pairing is divisible by \(t^i\). Contravariance makes these \(\mathfrak g\)-submodules; \(J^1\) is the radical of the specialized form.}
\GSfield{Reconstruction}{Moving lifts are indispensable. A constant lift need not realize the maximal order of vanishing, so the notebook\textquotesingle s displayed set is not a complete definition of the Jantzen filtration. The verification subsection gives a two-dimensional bilinear-form example where a moving lift has order three and the constant lift only order one.}
\GSfield{Fate}{This is Jantzen\textquotesingle s pre-1981 construction, used in \citet[\S1]{gabberjoseph1981towards}; see also \citet[pp.~112--181]{jantzen1979moduln} for its original book placement.}
\GSfield{Present status}{As of 11 November 2026: established with the lattice-and-specialization definition; a fixed-vector reading can give the wrong filtration.}
\GSfield{Assessment}{literal statement: false as written; publication match: exact; verification: read the result; development: annotated; confidence: high}
\end{GSitem}
\begin{GSitem}{1981-540.5}{B}{Does changing the deformation direction preserve layer multiplicities?}
\GSfield{Kind; source}{question; \GSpage{P000997}}
\GSfield{The notes}{The exact question is \ru{Что если сменить \(\rho\)? Видимо, кратн. не будут меняться.} No cone or genericity condition on the replacement is written.}
\GSfield{Statement}{For a fixed Verma module, compare \([J^i_\gamma M/J^{i+1}_\gamma M:L(\psi)]\) as the regular direction \(\gamma\) varies. The dominant-direction theorem does not extend to every regular direction. In type \(A_3\), take highest weight \(-\alpha_1-\alpha_3\) and examine the weight lower by \(2\alpha_1+3\alpha_2+2\alpha_3\). Its dimension is 13. The layer dimensions for \(\gamma=\rho\) are \((7,3,2,1)\); for a regular \(\gamma\) satisfying \(\gamma(h_1+2h_2+h_3)=0\), they are \((7,2,4,0)\). For example, Dynkin coordinates \((1,-2,3)\) satisfy that relation and are nonzero on every positive coroot.}
\GSfield{Reconstruction}{The unrestricted question is reconstructed literally, while the positive answer under a dominant-cone restriction is separated from it. The explicit counterexample is from the later literature, not an unrecorded qualification attributed to S. Gelfand. The included exact PBW calculation reproduces the layer dimensions by two independent methods; see the verification subsection and its code/output.}
\GSfield{Fate}{\citet[\S\S1.4--1.6 and 8.4]{williamson2016local} proves dominant local Hodge theory and gives this counterexample. The preprint first circulated on 8 October 2014; the journal volume is dated 2016, with online publication on 7 July 2017. \citet{beilinsonbernstein1993jantzen} had proved the dominant-direction Jantzen conjectures.}
\GSfield{Present status}{As of 11 November 2026: the unrestricted expectation is false; the dominant-direction version is established. This is a settled problem, not a C/D opportunity.}
\GSfield{Assessment}{literal statement: false as written; publication match: exact; verification: read the result; development: annotated; confidence: high; problem \hyperref[prob:GP-1981-10]{GP-1981-10}}
\end{GSitem}
\begin{GSitem}{1981-540.6}{B}{Correct grading of Kazhdan--Lusztig coefficients in Jantzen layers}
\GSfield{Kind; source}{claim; \GSpage{P000997}}
\GSfield{The notes}{The notes identify \(P_{w_1,w_2}(q)\) with \(\sum_iq^i[L_\psi:M_i/M_{i+1}]\). The scan confirms the reversed order of the multiplicity brackets.}
\GSfield{Statement}{With \(\lambda\) dominant integral, \(x,y\in W\), the usual Jantzen filtration from a dominant regular deformation satisfies \[\sum_{j\ge0}[J^jM(x\cdot\lambda)/J^{j+1}M(x\cdot\lambda):L(y\cdot\lambda)]v^j=v^{\ell(y)-\ell(x)}P_{x,y}(v^{-2}).\] Thus a coefficient of \(q^i\) belongs to layer \(\ell(y)-\ell(x)-2i\). For \(\mathfrak{sl}_2\), \(P_{e,s}=1\), whereas the lower simple constituent lies in layer 1, already disproving the literal unshifted formula.}
\GSfield{Reconstruction}{Two separate repairs are required: put the module before the simple object in the multiplicity bracket, and include the length shift and doubled polynomial degree. The rank-one diagnostic follows immediately from the single singular vector.}
\GSfield{Fate}{The graded theorem is part of \citet{beilinsonbernstein1993jantzen}; \citet[\S1.4]{williamson2016local} states the equality with the grading filtration under dominant deformation. In November 1981 this refinement remained beyond the already proved ungraded multiplicity theorem.}
\GSfield{Present status}{As of 11 November 2026: established with these conventions. Summing layers gives 1981-538.8, but does not justify the notebook	extquotesingle s degree assignment.}
\GSfield{Assessment}{literal statement: false as written; publication match: exact; verification: read the result; development: annotated; confidence: high; problem \hyperref[prob:GP-1979-41]{GP-1979-41}}
\end{GSitem}
\begin{GSitem}{1981-540.7}{E}{Local intersection cohomology of Schubert closures}
\GSfield{Kind; source}{claim; \GSpage{P000997}}
\GSfield{The notes}{The coefficients are described as \(a_i=\dim H^i\) of one Schubert cell near another, with \ru{локальн. когомол. норм. сечения}; a later line writes ``cell contained in other cell\textquotesingle.}
\GSfield{Statement}{Let \(X_y=\overline{ByB/B}\subset G/B\), \(x\le y\), and \(p_x\in BxB/B\). Normalize the intersection complex \(IC^0(X_y)\) so that it restricts to the constant sheaf in degree zero on the smooth locus. Then \(P_{x,y}(q)=\sum_{i\ge0}\dim\mathcal H^{2i}(IC^0(X_y))_{p_x}q^i\); odd stalk groups vanish. A transverse normal slice gives the same local singularity data, with the appropriate intersection-complex normalization.}
\GSfield{Reconstruction}{Distinct Schubert cells are disjoint. The intended containment is a cell inside another cell\textquotesingle s closure, and the groups are intersection cohomology stalks, not arbitrary ordinary local cohomology. The degree is \(2i\), not \(i\).}
\GSfield{Fate}{This interpretation was already in \citet{kazhdanlusztig1980schubert}; the statement and normalization were checked in \citet[Theorem~4.1]{lusztig1983singularities}. It belongs with the earlier occurrence 1979-484.4.}
\GSfield{Present status}{As of 11 November 2026: established geometric interpretation; the missing word ``intersection\textquotesingle\ is mathematically decisive.}
\GSfield{Assessment}{literal statement: false as written; publication match: exact; verification: read the result; development: annotated; confidence: high; problem \hyperref[prob:GP-1979-41]{GP-1979-41}}
\end{GSitem}

\GStalk{541}{1981-11-23}{A. V. Zelevinsky}{P000996-P000999}{Finite-dimensional weight multiplicities and affine Weyl polynomials}
The section begins near the foot of P000997. Gelfand inserts the type-A alternant character formula and a Cauchy-determinant exercise. Zelevinsky then identifies weight multiplicities with values of affine Kazhdan--Lusztig polynomials, followed by an incompletely specified extension to all Kac--Moody algebras. The latter is kept apart from the finite-dimensional Verma-module multiplicity theorem.

\begin{GSitem}{1981-541.1}{E}{Weights and multiplicities in a finite-dimensional simple module}
\GSfield{Kind; source}{definition; \GSpage{P000997-P000998}}
\GSfield{The notes}{The source writes \(L(\chi)=\bigoplus_\lambda L(\chi)_\lambda\) but labels the summand as the weight space of weight \(\chi\).}
\GSfield{Statement}{For dominant integral highest weight \(\chi\), the finite-dimensional module \(L(\chi)\) decomposes into \(L(\chi)_\lambda=\{v:hv=\lambda(h)v\ \forall h\in\mathfrak h\}\). Its weight multiplicity is \(m(\chi,\lambda)=\dim L(\chi)_\lambda\). The highest weight labels the representation, whereas \(\lambda\) labels a summand.}
\GSfield{Reconstruction}{The annotation must say weight \(\lambda\), as the subscript and the formula on the next page independently show.}
\GSfield{Fate}{This is the finite-dimensional setup used in \citet[\S6]{lusztig1983singularities}, not the Verma composition multiplicity of record 538.}
\GSfield{Present status}{As of 11 November 2026: standard, with one faithful variable slip in the manuscript.}
\GSfield{Assessment}{literal statement: false as written; publication match: exact; verification: read the result; development: annotated; confidence: high}
\end{GSitem}
\begin{GSitem}{1981-541.2}{E}{Type-A alternants and the missing highest-weight shift}
\GSfield{Kind; source}{theorem; \GSpage{P000997-P000998}; I. M. Gelfand}
\GSfield{The notes}{Gelfand gives the ratio of determinants with numerator exponents \(a_i\), denominator exponents \(0,1,\ldots,n-1\), then identifies \(\chi=(a_1,\ldots,a_n)\).}
\GSfield{Statement}{For \(GL_n\), dominant integral \(\lambda_1\ge\cdots\ge\lambda_n\) gives \[\operatorname{ch}L(\lambda)=\frac{\det(x_j^{\lambda_i+n-i})_{i,j=1}^n}{\det(x_j^{n-i})_{i,j=1}^n}.\] Reversing rows in numerator and denominator yields the ascending-denominator convention of the manuscript, with \(a_i=\lambda_{n+1-i}+i-1\). The determinant quotient is a symmetric Laurent polynomial; setting \(\lambda=0\) gives 1.}
\GSfield{Reconstruction}{The \(\rho\)-shift cannot be omitted. If the \(a_i\) literally equal the unshifted coordinates of the highest weight, even the trivial representation gives a zero numerator when \(n>1\).}
\GSfield{Fate}{This is the classical alternant form of the Weyl character formula; its general form is recalled in \citet[\S6, (6.3)]{lusztig1983singularities}, and the determinant/Schur-function identification in \citet[\S2.1]{krattenthaler1999determinant}.}
\GSfield{Present status}{As of 11 November 2026: established; the displayed ratio requires shifted exponents and a stated row-order convention.}
\GSfield{Assessment}{literal statement: false as written; publication match: exact; verification: read the result; development: annotated; confidence: high}
\end{GSitem}
\begin{GSitem}{1981-541.3}{E}{Cauchy determinant exercise}
\GSfield{Kind; source}{problem; \GSpage{P000998}; I. M. Gelfand}
\GSfield{The notes}{Gelfand asks first-year students to calculate \(\det(1/(\lambda_i-\mu_j))\), under the sole restriction \(\lambda_i\ne\mu_j\).}
\GSfield{Statement}{For all cross-distinct parameters, \[\det\left(\frac1{\lambda_i-\mu_j}\right)_{i,j=1}^n=\frac{\prod_{i<j}(\lambda_j-\lambda_i)(\mu_i-\mu_j)}{\prod_{i,j}(\lambda_i-\mu_j)}.\] Coincidences within the \(\lambda\)-list or within the \(\mu\)-list are permitted and make both sides zero. Multiplying by the denominator gives an alternating polynomial in each list; divisibility by the two Vandermonde factors and comparison of one coefficient determine the formula.}
\GSfield{Reconstruction}{The solution and its sign are supplied as an annotation to an explicitly elementary exercise; no research problem is manufactured from it.}
\GSfield{Fate}{This classical double alternant appears in \citet[\S2.1, (2.7)]{krattenthaler1999determinant}, with plus-denominator parameters. Substituting \(Y_j=-\mu_j\) gives exactly the source convention.}
\GSfield{Present status}{As of 11 November 2026: long settled; the seminar wording is a pedagogical problem.}
\GSfield{Assessment}{literal statement: correct; publication match: exact; verification: read the result; development: annotated; confidence: high}
\end{GSitem}
\begin{GSitem}{1981-541.4}{B}{Finite weight multiplicities from affine Kazhdan--Lusztig polynomials}
\GSfield{Kind; source}{theorem; \GSpage{P000998}}
\GSfield{The notes}{Zelevinsky writes \(m(\lambda,\chi)=P_{n_\lambda,n_\chi}(1)\) and labels the indices as affine Weyl group elements, without defining them. This swaps the argument order used immediately above.}
\GSfield{Statement}{Let \(\mathfrak g\) be complex simple, \(P\) its weight lattice, and \(\lambda,\mu\) dominant integral with \(\mu\le\lambda\). Use the extended affine Weyl group \(W\ltimes P\) in Lusztig\textquotesingle s convention, and let \(n_\eta\) be the longest element of \(Wt_\eta W\). Then \(\dim L(\lambda)_\mu=P_{n_\mu,n_\lambda}(1)\). For an arbitrary integral weight, first replace it by its dominant Weyl-orbit representative. The root datum and translation lattice matter; ``affine Weyl group\textquotesingle\ alone is insufficient for arbitrary \(P\).}
\GSfield{Reconstruction}{The definition of \(n_\eta\), dominance conditions, and consistent highest-weight/weight argument order come from the identified publication, not from the notebook. Ordinary finite-Weyl KL polynomials from record 538 are not sufficient here.}
\GSfield{Fate}{An exact match is \citet[\S2 and Theorem~6.1]{lusztig1983singularities}. The volume records a Luminy conference of 6--10 July 1981, but that fact alone does not verify when this particular theorem first circulated. The November notebook demonstrates prepublication knowledge; the exact printed result is dated 1983. Class B records the later publication by another mathematician, not priority for the seminar.}
\GSfield{Present status}{As of 11 November 2026: established with the specified extended affine indices; it is a different multiplicity problem from the Verma theorem.}
\GSfield{Assessment}{literal statement: correct with missing hypotheses; publication match: exact; verification: read the result; development: annotated; confidence: high}
\end{GSitem}
\begin{GSitem}{1981-541.5}{F}{An unspecified extension to all Kac--Moody algebras}
\GSfield{Kind; source}{conjecture; \GSpage{P000998}}
\GSfield{The notes}{The next line says \ru{То же \(\forall\) алгебр Кац--Муди}; the entire subsequent expression \(\operatorname{ch}L_\lambda(K_\mu)=P_{\lambda,\mu}(q)\) is crossed out.}
\GSfield{Statement}{The surviving text supports only a proposed Kac--Moody generalization. It does not specify a symmetrizable root datum, the relevant category, the role of \(q\), the elements \(n_\lambda\), or whether it concerns weight multiplicities or highest-weight composition multiplicities. There is therefore no single well-defined proposition to test.}
\GSfield{Reconstruction}{It would be unjustified to replace the erased expression by a modern Kac--Moody character theorem. In infinite type there need not be a longest element of a finite-Weyl double coset of the kind used in 1981-541.4.}
\GSfield{Fate}{\citet[\S5 and \S6]{lusztig1983singularities} distinguishes several different extensions involving affine polynomials. The search found no exact identification of this erased formula. General Kac--Moody character results provide context only.}
\GSfield{Present status}{As of 11 November 2026: not identifiable at the level required for a problem-index entry; no conclusion of openness or nondevelopment is warranted.}
\GSfield{Assessment}{literal statement: not determinable; publication match: thematic only; verification: unresolved lead; development: annotated; confidence: low}
\end{GSitem}

\GStalk{542}{1981-11-23}{A. G. Elashvili}{P000997-P001000}{Primitive enveloping algebras and Frobenius semidirect products}
The actual intervention occupies the lower part of P000998. It asks Jacobson\textquotesingle s primitivity question, excludes nonzero semisimple and nilpotent examples, gives affine and binary-cubic semidirect products, and asserts a restricted classification with a further construction method. The exact 1982 announcement identifies the older Ooms criterion and supplies the missing reductive hypotheses; its full classification is broader than the notebook\textquotesingle s two examples.

\begin{GSitem}{1981-542.1}{E}{Primitivity of the universal enveloping algebra}
\GSfield{Kind; source}{definition; \GSpage{P000998}}
\GSfield{The notes}{For a linear algebraic group \(G\subset GL_n(\mathbb C)\), the notes define \(U(L(G))\) to be primitive when zero is a primitive ideal, equivalently the kernel of a faithful irreducible representation. ``Irreducible\textquotesingle\ is inserted above the line.}
\GSfield{Statement}{An associative unital algebra \(A\) is left primitive if it has a faithful simple left module; equivalently, \(0=\operatorname{Ann}_A(M)\) for a simple module \(M\). Applied to \(A=U(\mathfrak g)\), this requires faithfulness of the enveloping-algebra action, a stronger condition than faithfulness merely on the Lie subspace \(\mathfrak g\). The representation need not be finite dimensional.}
\GSfield{Reconstruction}{The insertion in the manuscript is mathematically essential. Every algebra has a faithful regular module, which would make the definition vacuous without simplicity.}
\GSfield{Fate}{This is the definition used in \citet[\S0]{elashvili1982frobenius}. That announcement was received on 30 November 1981, one week after the seminar, and published in 1982.}
\GSfield{Present status}{As of 11 November 2026: a standard ring-theoretic definition.}
\GSfield{Assessment}{literal statement: correct; publication match: exact; verification: read the result; development: annotated; confidence: high}
\end{GSitem}
\begin{GSitem}{1981-542.2}{E}{Jacobson's primitive-envelope question for algebraic Lie algebras}
\GSfield{Kind; source}{question; \GSpage{P000998}}
\GSfield{The notes}{The notes ask which \(G\) and \(L(G)\) have primitive enveloping algebra, attribute the question to Jacobson, and add ``last year\textquotesingle\ below his name.}
\GSfield{Statement}{For a finite-dimensional algebraic Lie algebra \(\mathfrak g=\operatorname{Lie}(G)\) over an algebraically closed field of characteristic zero, the following are equivalent: \(U(\mathfrak g)\) is primitive; the coadjoint action has a Zariski-open orbit; and some \(f\in\mathfrak g^*\) makes \(B_f(x,y)=f([x,y])\) nondegenerate. A Lie algebra satisfying the last condition is called Frobenius. One may replace \(G\) by its identity component without changing the criterion.}
\GSfield{Reconstruction}{The algebraicity hypothesis is retained: an open-orbit criterion for this problem must not silently be asserted for arbitrary nonalgebraic Lie algebras. The marginal date is not reliable evidence that Jacobson first posed the question in 1980.}
\GSfield{Fate}{\citet[\S0, theorem attributed to Ooms]{elashvili1982frobenius} explicitly dates the answer to 1976 and cites Ooms\textquotesingle s supplements in Proc. Amer. Math. Soc. 58, 67--72. The 1976 full text was not accessible in this run, so the theorem was verified through Elashvili\textquotesingle s exact statement. The question already occurs as 1964-37.1 in the supplied index.}
\GSfield{Present status}{As of 11 November 2026: settled before this seminar for algebraic Lie algebras; the bibliographic attribution is stronger than a topical match.}
\GSfield{Assessment}{literal statement: correct with missing hypotheses; publication match: exact; verification: read the result; development: annotated; confidence: high; problem \hyperref[prob:GP-1981-12]{GP-1981-12}}
\end{GSitem}
\begin{GSitem}{1981-542.3}{E}{Why nonzero semisimple and nilpotent Lie algebras are excluded}
\GSfield{Kind; source}{claim; \GSpage{P000998}}
\GSfield{The notes}{The source says semisimple algebras, and an inserted parenthetical ``nilpotent ones\textquotesingle, are unsuitable.}
\GSfield{Statement}{A nonzero complex semisimple Lie algebra is not Frobenius: every coadjoint stabilizer has dimension at least its positive rank. A nonzero nilpotent Lie algebra has nonzero center, which lies in the radical of every \(B_f\). Hence neither class has primitive enveloping algebra under the algebraic criterion. The zero Lie algebra is an exception: \(U(0)=\mathbb C\) is primitive.}
\GSfield{Reconstruction}{The qualification ``nonzero\textquotesingle\ repairs the trivial edge case. The two exclusions have different mechanisms; nilpotence by itself is not being confused with associative nilpotence of \(U(\mathfrak g)\).}
\GSfield{Fate}{These deductions are immediate applications of the Ooms criterion printed in \citet[\S0]{elashvili1982frobenius}; no new classification is asserted.}
\GSfield{Present status}{As of 11 November 2026: established with the zero-algebra exception.}
\GSfield{Assessment}{literal statement: correct with missing hypotheses; publication match: exact; verification: read the result; development: annotated; confidence: high}
\end{GSitem}
\begin{GSitem}{1981-542.4}{E}{The affine Lie algebra is Frobenius}
\GSfield{Kind; source}{construction; \GSpage{P000998}}
\GSfield{The notes}{The first example is the Lie algebra of the affine group, represented by matrices whose last column is zero.}
\GSfield{Statement}{For \(n\ge1\), \(\mathfrak{aff}_n=\mathfrak{gl}_n\ltimes\mathbb C^n\), with the standard action and abelian translation ideal, is Frobenius and therefore has primitive enveloping algebra. The usual column-vector convention realizes it by matrices \(\left(\begin{smallmatrix}A&b\\0&0\end{smallmatrix}\right)\). The notebook\textquotesingle s last-column-zero matrices give an isomorphic Lie algebra via minus transpose, so they are not a mathematical error.}
\GSfield{Reconstruction}{The row/column convention is stated explicitly. For the Frobenius assertion, the open orbit of \(GL_n\) on nonzero covectors has stabilizer algebra isomorphic to \(\mathfrak{aff}_{n-1}\); iteration reduces to the zero algebra.}
\GSfield{Fate}{This is the first reduced family in \citet[\S2, table]{elashvili1982frobenius}; the stabilizer criterion is its \S1 lemma. The two-dimensional base case was already noted in 1964-37.2.}
\GSfield{Present status}{As of 11 November 2026: established for all positive n; the source	extquotesingle s matrix convention is retained as valid.}
\GSfield{Assessment}{literal statement: correct with missing hypotheses; publication match: exact; verification: read the result; development: annotated; confidence: high}
\end{GSitem}
\begin{GSitem}{1981-542.5}{E}{The binary-cubic semidirect product}
\GSfield{Kind; source}{construction; \GSpage{P000998}}
\GSfield{The notes}{The second example is \(\mathfrak{gl}_2\ltimes V_4\), with abelian \(V_4\), called the unique four-dimensional irreducible representation.}
\GSfield{Statement}{Take \(V_4=\operatorname{Sym}^3\mathbb C^2\), with the natural \(\mathfrak{gl}_2\)-action. Then \(\mathfrak{gl}_2\ltimes V_4\) is Frobenius. A generic binary cubic has three distinct projective roots; its stabilizer in \(GL_2\) is finite, so the infinitesimal orbit has dimension four and is open in \(V_4\) (or its dual). The abelian semidirect-product criterion then applies with zero generic stabilizer.}
\GSfield{Reconstruction}{The irreducible four-dimensional \(\mathfrak{sl}_2\)-module is unique, but extensions to \(\mathfrak{gl}_2\) depend on the central character. A zero central action leaves a central ideal and cannot be Frobenius. Nonzero scalar central actions yield isomorphic Lie semidirect products after rescaling the central generator; the natural cubic action is one choice.}
\GSfield{Fate}{This is the second reduced family in \citet[\S2, table]{elashvili1982frobenius}. The publication provides an exact example match, while the notebook\textquotesingle s uniqueness language needs the stated qualification.}
\GSfield{Present status}{As of 11 November 2026: established with nontrivial action of the scalar center.}
\GSfield{Assessment}{literal statement: false as written; publication match: exact; verification: read the result; development: annotated; confidence: high}
\end{GSitem}
\begin{GSitem}{1981-542.6}{A}{The asserted two-example classification and its missing reductive hypotheses}
\GSfield{Kind; source}{theorem; \GSpage{P000998}}
\GSfield{The notes}{The notebook claims there are no other examples for \(L=R\ltimes V\), with \(R\) simple, \(V\) abelian and irreducible. The word preceding ``simple\textquotesingle\ is crossed out.}
\GSfield{Statement}{The literal hypothesis cannot include either displayed \(\mathfrak{gl}_n\) example, since \(\mathfrak{gl}_n\) is not simple. The identified paper classifies reduced algebraic Frobenius triples \((R,V,\varphi)\) with \(R\) reductive and with either irreducible \(V\), or abelian \(V\) and simple derived algebra \([R,R]\). Its table has ten families. Under the simultaneous narrow requirements of irreducible \(V\) and simple derived algebra, the first two reduced families are the affine standard module and the binary cubic; the \(n=1\) affine edge case has zero derived algebra.}
\GSfield{Reconstruction}{This is a comparison with the exact later announcement, not a claim that the erased word can be read as ``reductive\textquotesingle. Reduction and the distinction between \(R\) and \([R,R]\) are substantive missing hypotheses.}
\GSfield{Fate}{\citet[\S0 and \S2]{elashvili1982frobenius}, received 30 November 1981 and printed in 1982, is a close chronological and mathematical match. Its broader ten-family theorem prevents the seminar\textquotesingle s compressed ``no others\textquotesingle\ from being quoted without qualification.}
\GSfield{Present status}{As of 11 November 2026: the relevant reduced classification is published; the literal simple-R formulation is false/incompatible with its examples.}
\GSfield{Assessment}{literal statement: false as written; publication match: partial; verification: read the result; development: annotated; confidence: high}
\end{GSitem}
\begin{GSitem}{1981-542.7}{A}{Castling as a method for generating further Frobenius examples}
\GSfield{Kind; source}{method; \GSpage{P000998}}
\GSfield{The notes}{The intervention ends: there is a method for constructing examples beyond these. No method is named or written out.}
\GSfield{Statement}{A plausible precise identification is the castling operation in the ensuing publication. If \(\dim V=m\) and \(0<n<m\), compare \((R\oplus\mathfrak{sl}_n,V\otimes\mathbb C^n,\varphi\otimes\mathrm{std})\) and \((R\oplus\mathfrak{sl}_{m-n},V^*\otimes\mathbb C^{m-n},\varphi^*\otimes\mathrm{std})\). The associated abelian semidirect products are simultaneously Frobenius. A reduced triple minimizes the module dimension in its castling class.}
\GSfield{Reconstruction}{The notebook supplies only the existence of a method. Identification with castling is supported by the speaker\textquotesingle s paper, but remains a reconstruction; the displayed transformation is not transcribed from the seminar.}
\GSfield{Fate}{\citet[\S2, lemma and definition of \ru{приведённая} triples]{elashvili1982frobenius} gives exactly this operation and uses it for the classification. The paper explicitly traces it to Sato--Kimura transformations.}
\GSfield{Present status}{As of 11 November 2026: a published construction; the connection to the unnamed seminar remark is probable, not demonstrable from that remark alone.}
\GSfield{Assessment}{literal statement: correct with missing hypotheses; publication match: partial; verification: read the result; development: annotated; confidence: medium}
\end{GSitem}
\begin{GSitem}{1981-542.8}{F}{An unidentified gamma function associated with the cubic example}
\GSfield{Kind; source}{remark; \GSpage{P000998}}
\GSfield{The notes}{Next to \(\mathfrak{gl}_2\ltimes V_4\), a marginal note reads that another interesting \(\Gamma\)-function arises here.}
\GSfield{Statement}{No function, integral, functional equation, normalization, or parameter is recorded. Possible connections to relative invariants and local zeta integrals cannot determine a unique gamma function from this line.}
\GSfield{Reconstruction}{The symbol is visibly \(\Gamma\), but neither the manuscript nor the short Frobenius announcement supplies enough information to reconstruct the object.}
\GSfield{Fate}{The full \citet{elashvili1982frobenius} was inspected; it classifies Frobenius algebras and mentions Sato--Kimura, but contains no matching gamma-function formula. Searches combining Elashvili, Frobenius, binary cubics and gamma functions produced no exact identification.}
\GSfield{Present status}{As of 11 November 2026: unidentified, not an established open problem.}
\GSfield{Assessment}{literal statement: not determinable; publication match: none found; verification: unresolved lead; development: annotated; confidence: low}
\end{GSitem}

\GStalk{543}{1981-11-30}{I. M. Gelfand}{P000998-P001000}{Nonstandard analysis, relaxation oscillations and surreal sign sequences}
The section begins on P000999 and continues to the Krichever heading on P001000. It reports Cartier\textquotesingle s November 1981 Bourbaki exposition of Reeb\textquotesingle s methods, gives the van der Pol singular limit and its Rayleigh form, and sketches Cartier\textquotesingle s random-element quotient using a zero-one finitely additive measure. It then introduces surreal numbers through arrow sequences. The ultrapower and surreal field are distinct constructions; the open-interval probability codomain and finite-arrow halving rule require mathematical emendation.

\begin{GSitem}{1981-543.1}{E}{The singular limit of the van der Pol equation}
\GSfield{Kind; source}{problem; \GSpage{P000999}}
\GSfield{The notes}{Gelfand writes \(\varepsilon\ddot x+(x^2-1)\dot x+x=0\) and asks what happens when \(\varepsilon\to0\). The heading names Reeb and, underneath, Cartier\textquotesingle s November 1981 Bourbaki survey.}
\GSfield{Statement}{Take \(\varepsilon>0\) and let it decrease to zero. The second-order equation loses one derivative when \(\varepsilon=0\), so naive substitution cannot accommodate two arbitrary initial data or the rapid transitions. A precise version asks for the limiting attracting periodic orbit and the matching of its slow segments across fast jumps; 1981-543.3 states this version. A claim of uniform convergence for all initial data on all times is not supported.}
\GSfield{Reconstruction}{The notes ask a broad singular-limit question. The periodic-orbit formulation is selected because the sketches show relaxation jumps and because it is the actual question treated in the identified source.}
\GSfield{Fate}{\citet[\S1 and \S5, nos.~21--23]{cartier1982perturbations} is an exact match, delivered in November 1981 and printed in 1982. Cartier reports methods of Reeb\textquotesingle s school; Reeb was not the speaker of Bourbaki exposé 580.}
\GSfield{Present status}{As of 11 November 2026: the classical relaxation limit in this formulation is established; the seminar reports a contemporary exposition of existing and recent methods.}
\GSfield{Assessment}{literal statement: correct with missing hypotheses; publication match: exact; verification: read the result; development: annotated; confidence: high; problem \hyperref[prob:GP-1981-13]{GP-1981-13}}
\end{GSitem}
\begin{GSitem}{1981-543.2}{E}{Integrating the reduced first-order equation}
\GSfield{Kind; source}{construction; \GSpage{P000999}}
\GSfield{The notes}{At \(\varepsilon=0\), the notes write \(dx/dt=-x/(x^2-1)\), separate variables, and obtain \(t+C=\log|x|-x^2/2\).}
\GSfield{Statement}{On an interval where \(x\ne0,\pm1\), the reduced equation is \(\dot x=x/(1-x^2)\), hence \(d(\log|x|-x^2/2)/dt=1\). The equilibrium \(x\equiv0\) must be restored after division by \(x\). A finite derivative cannot satisfy the reduced equation at \(x=\pm1\); those turning points require a fast transition in the full equation.}
\GSfield{Reconstruction}{The division conditions and omitted equilibrium are supplied. The implicit formula parametrizes individual slow branches, not a complete periodic solution.}
\GSfield{Fate}{The same reduction is given in \citet[\S5, no.~21, equation~(17)]{cartier1982perturbations}. The derivative check is elementary and confirms the signs in both transcriptions.}
\GSfield{Present status}{As of 11 November 2026: correct with the excluded points and equilibrium explicitly recorded.}
\GSfield{Assessment}{literal statement: correct with missing hypotheses; publication match: exact; verification: read the result; development: annotated; confidence: high}
\end{GSitem}
\begin{GSitem}{1981-543.3}{E}{Relaxation jumps and the limiting periodic cycle}
\GSfield{Kind; source}{claim; \GSpage{P000999}}
\GSfield{The notes}{Two sketches have nearly vertical segments, accompanied by the statement that limiting solutions have such jumps.}
\GSfield{Statement}{Write \(F(x)=x^3/3-x\) and \(y=\varepsilon\dot x+F(x)\), giving \(\varepsilon\dot x=y-F(x)\), \(\dot y=-x\). As \(\varepsilon\downarrow0\), the attracting periodic orbit approaches the slow arcs \(y=F(x)\), \(x\in[1,2]\) and \(x\in[-2,-1]\), connected by fast segments at \(y=-2/3\) and \(y=2/3\). In \(x(t)\), the jumps are \(1\to-2\) and \(-1\to2\). Convergence of the phase-plane cycle is understood geometrically; time-parametrized convergence is local away from jumps after fixing phase.}
\GSfield{Reconstruction}{The sketch does not label the endpoints or specify the topology of convergence. These details are reconstructed from the polynomial slow manifold and the matched-cycle analysis, not from the drawing alone.}
\GSfield{Fate}{\citet[\S5, nos.~22--23 and Figure~5]{cartier1982perturbations} gives this Liénard-plane construction and the attracting limit cycle. The November 1981 exposition precedes the Moscow report; the 1982 print date does not make it a new Moscow conjecture.}
\GSfield{Present status}{As of 11 November 2026: established relaxation asymptotics; uniform convergence of continuous trajectories to a discontinuous graph across a jump is not the asserted theorem.}
\GSfield{Assessment}{literal statement: correct with missing hypotheses; publication match: exact; verification: read the result; development: annotated; confidence: high; problem \hyperref[prob:GP-1981-13]{GP-1981-13}}
\end{GSitem}
\begin{GSitem}{1981-543.4}{E}{Rayleigh and van der Pol equations are linked by differentiation}
\GSfield{Kind; source}{construction; \GSpage{P000999}}
\GSfield{The notes}{The second displayed equation is \(\varepsilon\ddot u+(\dot u^3/3-\dot u)+u=0\), called Rayleigh\textquotesingle s equation, and is said to reduce to van der Pol.}
\GSfield{Statement}{If \(u\) solves Rayleigh\textquotesingle s equation and \(x=\dot u\), differentiation gives \(\varepsilon\ddot x+(x^2-1)\dot x+x=0\). Conversely, from a van der Pol solution define \(u=-\varepsilon\dot x-x^3/3+x\). Then the van der Pol equation gives \(\dot u=x\), and substitution verifies Rayleigh\textquotesingle s equation. For smooth solutions with \(\varepsilon\ne0\), these formulas give the precise correspondence without an arbitrary integration constant.}
\GSfield{Reconstruction}{The substitution and converse are supplied to make ``reduces\textquotesingle\ exact. Cartier subsequently changes a sign in his Liénard coordinate convention; that does not change this displayed differentiation identity.}
\GSfield{Fate}{The two equations and the derivative relation are written in \citet[\S5, no.~21]{cartier1982perturbations}.}
\GSfield{Present status}{As of 11 November 2026: an established elementary equivalence, with the variable transformation explicit.}
\GSfield{Assessment}{literal statement: correct with missing hypotheses; publication match: exact; verification: read the result; development: annotated; confidence: high}
\end{GSitem}
\begin{GSitem}{1981-543.5}{E}{Using an infinitesimal parameter in singular perturbation theory}
\GSfield{Kind; source}{method; \GSpage{P000999}}
\GSfield{The notes}{The notebook emphasizes that equation (1) is studied by nonstandard analysis.}
\GSfield{Statement}{In a suitable nonstandard enlargement of the real field and its function spaces, choose positive infinitesimal \(\varepsilon\), analyze the internal differential equation, and take standard parts on limited slow portions. Different rescalings resolve fast portions. Transfer and the precise internal/external distinction are required; the mere existence of a field containing infinitesimals does not prove a singular-limit theorem.}
\GSfield{Reconstruction}{This is a method, not a claim that setting an ordinary parameter equal to zero is legitimate. The later surreal-number sketch is a different construction and does not itself provide the enlargement or transfer principle.}
\GSfield{Fate}{\citet[\S1, no.~2; \S3; \S5]{cartier1982perturbations} explains the infinitesimal method and its application to relaxation oscillations. The report is specifically about Reeb\textquotesingle s school, as presented by Cartier in November 1981.}
\GSfield{Present status}{As of 11 November 2026: an established framework for these asymptotics; its foundations must not be conflated with arbitrary ordered-field extensions.}
\GSfield{Assessment}{literal statement: correct with missing hypotheses; publication match: exact; verification: read the result; development: annotated; confidence: high}
\end{GSitem}
\begin{GSitem}{1981-543.6}{E}{The zero-one probability space and ultrapower quotient}
\GSfield{Kind; source}{construction; \GSpage{P000999}}
\GSfield{The notes}{The scan and both transcriptions give \(p:2^\Omega\to(0,1)\), finite additivity, and equivalence of \(x^{\prime},x^{\prime\prime}:\Omega\to X\) when \(P(x^{\prime}=x^{\prime\prime})=1\).}
\GSfield{Statement}{The intended construction uses a normalized finitely additive measure \(p:2^\Omega\to\{0,1\}\), with \(p(\Omega)=1\). Then \(\mathcal U=\{A:p(A)=1\}\) is an ultrafilter, and \(X^\Omega/\mathcal U\) identifies maps equal on a \(\mathcal U\)-large set. For \(X=\mathbb R\), operations descend pointwise. Taking \(\Omega=\mathbb N\) and a nonprincipal ultrafilter produces a nontrivial ultrapower in which \([1/n]\) is positive infinitesimal.}
\GSfield{Reconstruction}{The codomain is inferred from Cartier\textquotesingle s exact construction, not visually read as braces: the notebook truly has round parentheses. Finite additivity forces \(p(\varnothing)=0\), and the equivalence relation requires the value 1, so the open interval is impossible. Replacing it merely by \([0,1]\) would allow an ordinary probability algebra whose quotient need not be a field; the zero-one condition is decisive.}
\GSfield{Fate}{\citet[\S2, no.~3, assumptions (a)--(c), and no.~4]{cartier1982perturbations} explicitly requires values zero or one and quotienting random elements by almost-sure equality. This exact source resolves the manuscript error.}
\GSfield{Present status}{As of 11 November 2026: the ultrapower construction is established after this emendation; normalization and a nonprincipal choice are additionally needed for the intended infinitesimals.}
\GSfield{Assessment}{literal statement: false as written; publication match: exact; verification: read the result; development: annotated; confidence: high}
\end{GSitem}
\begin{GSitem}{1981-543.7}{E}{Surreal numbers represented by transfinite sign sequences}
\GSfield{Kind; source}{construction; \GSpage{P001000}}
\GSfield{The notes}{Under the heading \ru{Числа Крускала}, the notes call the construction a simple field extending \(\mathbb R\), identify the empty word with 0, up with 1, down with \(-1\), and give \(n\) ups, up-down \(=1/2\), and up-down-up \(=3/4\).}
\GSfield{Statement}{These are surreal sign expansions: a surreal number is a function from an ordinal to \(\{+,-\}\), ordered lexicographically with \(-<\text{undefined}<+\) at the first differing position. The empty sequence is zero. Constant finite sequences give integers, and finite mixed sequences give dyadic rationals. With Conway\textquotesingle s recursively defined operations, the full class \(\mathbf{No}\) is a real-closed ordered field in the proper-class sense, containing copies of the reals and ordinals.}
\GSfield{Reconstruction}{The identification with surreal numbers is strong; ``Kruskal\textquotesingle\ records the seminar\textquotesingle s attribution, not the standard origin of the field. The proper-class qualification and arithmetic are absent from the notebook. This is not the same object as the preceding set-sized ultrapower.}
\GSfield{Fate}{The field was published in \citet{conway1976numbers}, before this talk; \citet[Chs.~2--4]{gonshor1986surreal} systematically develops sign expansions. The readable modern verification is \citet[\S2.1--2.2]{mantovamatusinski2017survey}. The earlier occurrence 1978-446.22 concerns a different use of surreal/omnific arithmetic.}
\GSfield{Present status}{As of 11 November 2026: identified as classical surreal-number theory; the Kruskal label is retained as historical evidence.}
\GSfield{Assessment}{literal statement: correct with missing hypotheses; publication match: exact; verification: read the result; development: annotated; confidence: high}
\end{GSitem}
\begin{GSitem}{1981-543.8}{E}{The finite-arrow halving rule needs correction}
\GSfield{Kind; source}{claim; \GSpage{P001000}}
\GSfield{The notes}{The explanatory parenthesis says that the step is halved when the direction changes.}
\GSfield{Statement}{For a finite sign sequence beginning with \(k\ge1\) identical signs of value \(s\in\{1,-1\}\), its value is \(sk\) if it stops there. If the first change occurs after that initial run, each subsequent sign contributes successively \(\pm2^{-1},\pm2^{-2},\ldots\), irrespective of whether it changes direction again. Thus \(+-=1/2\), \(+-+=3/4\), and \(+--=1/4\). Halving only on each direction change would instead give the incorrect value zero for \(+--\).}
\GSfield{Reconstruction}{The supplied three examples do not distinguish the two rules. A fourth example does, so the correction concerns a faithful mathematical shorthand rather than a transcription error.}
\GSfield{Fate}{The sign-expansion construction is treated in \citet[Chs.~2--4]{gonshor1986surreal} and recalled in \citet[\S2.1]{mantovamatusinski2017survey}. The corrected finite rule follows from the simplest-number-between-neighbors recursion.}
\GSfield{Present status}{As of 11 November 2026: the literal operational rule is false; the corrected rule recovers all finite surreal sign expansions.}
\GSfield{Assessment}{literal statement: false as written; publication match: exact; verification: read the result; development: annotated; confidence: high}
\end{GSitem}
\begin{GSitem}{1981-543.9}{E}{Infinite arrow blocks and the distinction between omega and infinity}
\GSfield{Kind; source}{claim; \GSpage{P001000}}
\GSfield{The notes}{An infinitely repeated up arrow is called \(\infty\); appending a down arrow gives \(\infty-1\). An up block followed by a down block is called \(\tfrac12\infty\), their ratio is 2, a double hat gives \(\infty^2\), and one up followed by infinitely many downs gives \(\varepsilon\).}
\GSfield{Statement}{Interpreting each single infinite block as length \(\omega\), the identities are \(+^\omega=\omega\), \(+^\omega-=\omega-1\), \(+^\omega-^\omega=\omega/2\), and \(+-^\omega=\omega^{-1}\). In particular \(\omega/(\omega/2)=2\). The double-hat identity becomes correct if it means a constant positive sequence of length \(\omega^2\). Here \(\omega\) is a positive infinite surreal, not the extended-real endpoint \(+\infty\).}
\GSfield{Reconstruction}{The ordinal lengths and arithmetic convention are essential reconstructions. A vague ``infinite repetition of an infinite repetition\textquotesingle\ is not a definition. The nearby ``no predecessor\textquotesingle\ can refer to ordinal \(\omega\), but every surreal \(x\) has \(x-1\) and there are no immediate neighbors in an ordered field.}
\GSfield{Fate}{These are consequences of the classical surreal order and arithmetic described in \citet{conway1976numbers} and \citet[Chs.~4--5]{gonshor1986surreal}; \citet[\S2]{mantovamatusinski2017survey} supplies the modern sign and ordinal conventions. The exact infinite-block formulas were checked by the canonical-cut interpretation, rather than treated as ordinary real limits.}
\GSfield{Present status}{As of 11 November 2026: correct after assigning ordinal lengths and replacing the overloaded infinity notation.}
\GSfield{Assessment}{literal statement: correct with missing hypotheses; publication match: exact; verification: read the result; development: annotated; confidence: high}
\end{GSitem}

\GStalk{544}{1981-11-30}{I. M. Krichever}{P000999-P001000}{An elliptic matrix-function fragment}
Only the bottom of P001000 is headed Krichever. It names joint work with Brazovskii and Dzyaloshinskii, writes an unspecified ascending union, and displays a sigma-product matrix function on an elliptic curve with a divisor-sum condition. Krichever\textquotesingle s 1983 survey, \S1.3, identifies a closely matching algebra of matrix Baker--Akhiezer functions; the filtration remains a reconstruction. The 1982 joint Peierls paper confirms the collaboration, but supplies only thematic context for the surviving fragment.

\begin{GSitem}{1981-544.1}{A}{An ascending union plausibly identified with a filtered elliptic matrix algebra}
\GSfield{Kind; source}{construction; \GSpage{P001000}}
\GSfield{The notes}{After naming joint work with Brozovskii, Dzyaloshinskii and Krichever, the source writes only \(G=\bigcup_nG_n\), \(G_n\subset G_{n+1}\). No group law, algebra structure, or definition of \(G_n\) is given.}
\GSfield{Statement}{A close later construction is the algebra of matrix functions on a punctured elliptic curve whose entries near the puncture have prescribed exponential factors \(e^{(x_i-x_j)/z}\) times Laurent series with finitely many negative powers. Bounding the pole order after a local diagonal exponential gauge gives finite-dimensional pieces \(G_n\), with \(G_nG_m\subseteq G_{n+m}\). This supplies an algebra filtration, not a claim that each \(G_n\) is a group.}
\GSfield{Reconstruction}{The identification is a reconstruction supported by the immediately following sigma formula and by Krichever\textquotesingle s later notation \(G(\Gamma,z)\). The handwritten letter \(G\) alone does not authorize calling the union a group. The collaborator is almost certainly S. A. Brazovskii; the notebook\textquotesingle s first vowel is faithfully ``o\textquotesingle.}
\GSfield{Fate}{\citet[\S1.3, (1.91)--(1.94)]{krichever1983elliptic} defines the matching matrix-function framework; the original Russian survey is dated 1983 and the inspected English translation 1985. The joint \citet{brazovskiidzyaloshinskiikrichever1982peierls}, submitted 9 February 1982, confirms the collaboration but is only a thematic match to this isolated fragment.}
\GSfield{Present status}{As of 11 November 2026: partially identified with an established filtered-algebra construction; no complete seminar theorem can be recovered.}
\GSfield{Assessment}{literal statement: not determinable; publication match: partial; verification: read the result; development: annotated; confidence: medium}
\end{GSitem}
\begin{GSitem}{1981-544.2}{A}{Sigma-product entries and cancellation of elliptic multipliers}
\GSfield{Kind; source}{construction; \GSpage{P001000}}
\GSfield{The notes}{The source gives matrix entries \(u^{ij}(z)=\prod_{k=1}^n\sigma(z-a_k^{ij})\,\sigma(z)^{-n}\exp(\zeta(z)(x_i-x_j))\), with \(\sum_ka_k^{ij}=x_i-x_j\), with \(\Gamma\) labelled as a torus.}
\GSfield{Statement}{Fix a lattice \(\Lambda\subset\mathbb C\), its Weierstrass sigma and zeta functions, and chosen complex lifts satisfying the stated sum exactly. For a period \(2\omega\), the sigma quotient acquires multiplier \(e^{-2\eta\sum_ka_k^{ij}}\), whereas the exponential acquires \(e^{2\eta(x_i-x_j)}\); they cancel. Thus the displayed entry is single valued on \((\mathbb C/\Lambda)\setminus\{0\}\). Near zero it is \(e^{(x_i-x_j)/z}\) times a meromorphic germ of pole order at most \(n\). If \(x_i\ne x_j\), the puncture is generally essential, not a meromorphic pole.}
\GSfield{Reconstruction}{The lattice, choice of lifts, and local nature of the exponential gauge are supplied. The gauge-transformed entry need not itself be a global elliptic function. Scalar prefactors may be added in a general factorization, but are not written in the notebook.}
\GSfield{Fate}{The analytic framework and sigma construction in \citet[\S1.3, (1.91)--(1.94)]{krichever1983elliptic} give a partial publication match. The period-cancellation calculation verifies exactly what the surviving formula supports; no unrecorded Lax equation, Peierls Hamiltonian, or claimed completeness theorem is assigned to it.}
\GSfield{Present status}{As of 11 November 2026: the displayed quasi-periodicity cancellation is correct under the stated choices; the intended global theorem remains unspecified.}
\GSfield{Assessment}{literal statement: correct with missing hypotheses; publication match: partial; verification: read the result; development: annotated; confidence: medium}
\end{GSitem}

\GStalk{545}{}{}{P001001-P001002}{Year 1982}
The catalog's undated year heading marks the transition to calendar year 1982 within the academic year 1981--82. It supplies no independent mathematical assertion. Its catalog page range overlaps Müller's dated lecture and must not create a duplicate talk or duplicate occurrence.


\GStalk{546}{1982-02-08}{Werner Müller}{P001001-P001005}{Lattice gravity and curvature of piecewise flat spaces}
Müller reports joint work with Jeff Cheeger and Robert Schrader. A formal path-integral motivation leads to Regge's hinge curvature, nondegenerate smooth triangulations, convergence of curvature measures, higher Lipschitz--Killing curvatures, finite-dimensional metric spaces of simplices, and an integral-geometric proof idea. Gelfand asks which curvature actions have lower-order field equations. The source explicitly separates curvature convergence from the still unclear convergence of the equations. The catalog range extends through P001007; the next dated heading actually starts at P001006. Several faithful mathematical errors require emendation, especially the heat sign, the Euler-form trace, and the edge-length Jacobian.

\begin{GSitem}{1981-546.1}{E}{Schrödinger evolution and the formal Feynman kernel}
\GSfield{Kind; source}{construction; \GSpage{P001001}}
\GSfield{The notes}{For \(H(p,q)=p^2/(2m)+V(q)\), the page sends \(q\) to multiplication and \(p\) to \(+i\hbar\partial_x\), writes evolution \(e^{-it\widehat H/\hbar}\), and calls its position kernel a transition probability. Its formal path expression has exponent \(-iS\), with \(S=\int(\dot x^2-V)\,dt\).}
\GSfield{Statement}{For real bounded \(V\), one elementary rigorous setting is the self-adjoint operator \(\widehat H=-(\hbar^2/2m)\partial_x^2+V\) on \(H^2(\mathbb R)\), whose unitary evolution solves \(i\hbar\partial_t f=\widehat Hf\). The kernel is an amplitude (generally a distribution), not a probability density. The conventional action is \(\int(m\dot x^2/2-V)\,dt\), with formal phase \(e^{iS/\hbar}\). A product of Lebesgue measures over all times is not a defined countably additive path-space measure.}
\GSfield{Reconstruction}{The self-adjoint domain and bounded-potential hypothesis are a sufficient modern reconstruction, not the most general theorem. Changing the momentum sign alone leaves this quadratic Hamiltonian unchanged; the displayed phase and omitted action constants require separate correction.}
\GSfield{Fate}{This is the pre-existing path-integral motivation, not the new result of the talk. The book cited in the notebook is Konopleva--Popov; a separately verified historical reference is \citet[book and errata information]{feynmanhibbs1965path}. No claim of a rigorous real-time measure is attributed to that citation.}
\GSfield{Present status}{As of November 2026, unitary Schrödinger evolution is standard. Equations (2), (4), and (5) are not mutually normalized as written; the formal path integral must remain explicitly formal.}
\GSfield{Assessment}{literal statement: false as written; publication match: thematic only; verification: publisher metadata checked; development: annotated; confidence: high}
\end{GSitem}
\begin{GSitem}{1981-546.2}{E}{The heat-semigroup interpretation of the path integral}
\GSfield{Kind; source}{theorem; \GSpage{P001002}}
\GSfield{The notes}{The notes credit Kac with giving meaning to the integral after \(i\mapsto1\). They write \(\hbar\partial_t f=+\widehat Hf\), followed by \(e^{-t\widehat H/\hbar}\), and put \(x(\tau)\), rather than the potential, inside the time integral in the exponential.}
\GSfield{Statement}{For \(m,\hbar>0\) and real bounded continuous \(V\), Brownian motion with generator \((\hbar/2m)\partial_x^2\) gives \((e^{-t\widehat H/\hbar}f)(x)=\mathbb E_x[e^{-\hbar^{-1}\int_0^t V(B_s)ds}f(B_t)]\). This solves \(\hbar\partial_t f=-\widehat Hf\). A fixed-endpoint kernel uses Brownian-bridge expectation multiplied by the free heat kernel; unconditioned Wiener measure is not by itself an endpoint kernel.}
\GSfield{Reconstruction}{The displayed bounded continuous hypothesis avoids the stronger integrability assumptions needed for singular potentials. The missing minus sign, missing \(V\), and endpoint normalization are emendations of a faithfully copied formula.}
\GSfield{Fate}{Kac's 1949 article is the historical source identified in the notebook; its original DOI was attempted but full access failed in this run, so that entry is marked unverified rather than presented as a theorem read in the original. The formula above is checked directly against the heat generator and is a standard reconstruction, not a new claim by Müller.}
\GSfield{Present status}{As of November 2026, the Feynman--Kac formula is established. The real-time-to-imaginary-time slogan does not by itself justify analytic continuation or a gravitational measure.}
\GSfield{Assessment}{literal statement: false as written; publication match: exact; verification: indirect; development: annotated; confidence: high}
\end{GSitem}
\begin{GSitem}{1981-546.3}{E}{The Einstein--Hilbert action has second-order field equations}
\GSfield{Kind; source}{theorem; \GSpage{P001002}}
\GSfield{The notes}{The scalar-curvature action contains second derivatives of the metric, so fourth-order equations seem possible; a local noncovariant divergence term is said to remove them.}
\GSfield{Statement}{On a closed smooth manifold, the first variation of \(\int_M\operatorname{Scal}_g\,d\operatorname{vol}_g\) with respect to an inverse-metric variation has interior coefficient \(\operatorname{Ric}_g-\tfrac12\operatorname{Scal}_g g\). The remaining terms are divergences. The resulting vacuum equation is second order in \(g\); adding a volume term adds a multiple of \(g\). On a manifold with boundary one must specify boundary conditions or the appropriate boundary action.}
\GSfield{Reconstruction}{Closedness or boundary conditions are not stated. The local first-derivative expression need not be globally covariant, whereas its Euler--Lagrange tensor is.}
\GSfield{Fate}{This is the classical Einstein--Hilbert variational identity. It is also the first nonconstant member of the natural-tensor family classified by Lovelock; see the precise modern formulation in \citet[Theorem 2.6]{navarronavarro2011lovelock}.}
\GSfield{Present status}{As of November 2026, this order reduction is established. The manuscript's historical aside about what Einstein initially believed is not independently verified and is not used as mathematical evidence.}
\GSfield{Assessment}{literal statement: correct with missing hypotheses; publication match: partial; verification: read the result; development: annotated; confidence: high}
\end{GSitem}
\begin{GSitem}{1981-546.4}{E}{Which curvature actions have lower-order Euler--Lagrange equations?}
\GSfield{Kind; source}{question; \GSpage{P001002}; I. M. Gelfand}
\GSfield{The notes}{Gelfand asks \ru{Найти все комб. кривизн, где получ. ур-е меньшего порядка}; the pencilled answer is Einstein and Pontryagin classes.}
\GSfield{Statement}{A precise established version fixes dimension and metric signature and asks for all natural symmetric divergence-free tensors depending on the metric through its two-jet. Their real vector space is spanned by the Lovelock tensors \(L_0,\ldots,L_{\lfloor(n-1)/2\rfloor}\). The associated Lovelock scalar densities yield second-order metric equations. In their top dimension Euler and Pontryagin characteristic numbers have identically zero interior variation; they are null Lagrangians, not additional nonzero field equations.}
\GSfield{Reconstruction}{The notebook neither fixes dimension nor defines lower order, polynomiality, covariance, or equivalence modulo divergences. The natural-tensor statement is a carefully bounded reconstruction. It is not a classification of every conceivable curvature-dependent action. In dimensions at least five, the nontrivial Gauss--Bonnet Lovelock tensor already shows why an unqualified Einstein-only nontrivial answer is incomplete.}
\GSfield{Fate}{The relevant classification predates the seminar: Lovelock's original paper appeared in 1971. Its exact tensor theorem and a proof were read in \citet[Theorem 2.6]{navarronavarro2011lovelock}, first circulated in May 2010 and published in 2011. The 1971 publisher page was not fully accessible and is separately marked unverified.}
\GSfield{Present status}{As of November 2026, the reconstructed natural second-order tensor problem is solved. The unrestricted wording of the notebook cannot be assigned that theorem without the stated restrictions, and its proposed answer needs amendment.}
\GSfield{Assessment}{literal statement: false as written; publication match: partial; verification: read the result; development: annotated; confidence: high; problem \hyperref[prob:GP-1981-20]{GP-1981-20}}
\end{GSitem}
\begin{GSitem}{1981-546.5}{F}{A formal gravitational integral and finite triangulation measures}
\GSfield{Kind; source}{method; \GSpage{P001002,P001004-P001005}}
\GSfield{The notes}{The notes integrate \(e^{-S[g]-\kappa\operatorname{Vol}(g)}\) over metrics modulo diffeomorphisms and formally write \(\prod_x(\det g(x))^{5/2}\prod_{\mu\leq\nu}dg_{\mu\nu}(x)\). They later propose finite-dimensional measures on compatible simplex metrics.}
\GSfield{Statement}{For a fixed finite triangulation \(P\), admissible edge lengths form an open subset of a finite-dimensional real vector space. A specified smooth Riemannian metric on that subset defines an ordinary volume measure. A continuum integral on \(\operatorname{Met}(M)/\operatorname{Diff}(M)\) additionally requires a measure, gauge treatment, integrability, and a limiting prescription, none of which follows from the notation \(\prod_xdg(x)\). For covariant positive-definite \(n\times n\) matrices the congruence-invariant density is \((\det g)^{-(n+1)/2}\prod_{i\leq j}dg_{ij}\), so the four-dimensional exponent is negative.}
\GSfield{Reconstruction}{The finite-dimensional construction is recoverable; a continuum existence theorem is not stated. The sign diagnosis assumes the covariant matrix entries actually written, not a silent change to contravariant variables.}
\GSfield{Fate}{The finite Regge program is identifiable in the 1981 symposium announcement published as \citet{cheegermullerschrader1982lattice}. Neither that bibliographic match nor the curvature theorem of \citet{cheegermullerschrader1984curvature} certifies the formal continuum measure in (9)--(10).}
\GSfield{Present status}{As of November 2026, no exact continuum-measure theorem is identifiable from this fragment. This F classification concerns the unquantified formal integral, not the well-defined simplex measures indexed separately.}
\GSfield{Assessment}{literal statement: not determinable; publication match: thematic only; verification: indirect; development: annotated; confidence: medium}
\end{GSitem}
\begin{GSitem}{1981-546.6}{E}{Piecewise flat metrics and Regge hinge curvature}
\GSfield{Kind; source}{definition; \GSpage{P001003}}
\GSfield{The notes}{A metric is encoded by positive edge lengths \(\ell:P^{(1)}\to\mathbb R_+\). At a codimension-two simplex the angle defect is \(2\pi\) minus the sum of incident dihedral angles, and \(R(P)\) is the defect times hinge volume summed over hinges.}
\GSfield{Statement}{For a finite triangulated closed \(n\)-manifold with compatible nondegenerate Euclidean simplex metrics, define \(\delta_h=2\pi-\sum_{\sigma\supset h}\theta_{h,\sigma}\) and \(\mathcal R(P)=\sum_h\delta_h\operatorname{vol}_{n-2}(h)\). Positive edge lengths must satisfy the simplex Gram positivity conditions, not merely positivity. With the convention \(\operatorname{Scal}=2K\) on surfaces, the scalar-curvature integral corresponding to this hinge sum has factor \(2\mathcal R(P)\); other normalizations absorb this factor.}
\GSfield{Reconstruction}{Closedness excludes boundary hinges, whose defect convention differs. The notes call \(\mathcal R\) scalar curvature without fixing the factor of two; this is made explicit rather than silently identifying inequivalent conventions.}
\GSfield{Fate}{Regge's original paper appeared in 1961, long before the talk; the publisher record \citet{regge1961relativity} verifies its title, pages and chronology. The hinge construction is the starting point of \citet[Introduction]{cheegermullerschrader1984curvature}.}
\GSfield{Present status}{As of November 2026, the construction is established. It defines a curvature measure concentrated on hinges, not a smooth pointwise curvature tensor.}
\GSfield{Assessment}{literal statement: correct with missing hypotheses; publication match: exact; verification: read the result; development: annotated; confidence: high}
\end{GSitem}
\begin{GSitem}{1981-546.7}{A}{Convergence of Regge curvature under nondegenerate refinement}
\GSfield{Kind; source}{theorem; \GSpage{P001003}}
\GSfield{The notes}{After \(\operatorname{mesh}(P_j)\to0\), the notes announce convergence of total Regge curvature to the smooth integral, then add that angles must stay bounded away from zero and a smooth triangulation is used.}
\GSfield{Statement}{Let \((M,g)\) be a closed smooth Riemannian manifold and let secant piecewise flat approximations arise from admissible smooth triangulations with mesh tending to zero and a uniform positive fatness bound. The corresponding Lipschitz--Killing curvature measures converge weakly to the smooth ones; in particular their total scalar-curvature terms converge with consistent normalization. Uniform fatness controls every simplex face and excludes slivers; mesh convergence alone is insufficient.}
\GSfield{Reconstruction}{The detailed approximation and fatness hypotheses are supplied from the published theorem, not all from the shorthand angle condition. Normalizations must agree with the preceding item. No convergence of derivatives of the action is inferred.}
\GSfield{Fate}{The seminar is a documented 1982 circulation point, after the 1981 Heisenberg Symposium announcement (published 1982). The full result was read in \citet[Theorem 5.1, pp. 427--428]{cheegermullerschrader1984curvature}, published in 1984.}
\GSfield{Present status}{As of November 2026, this is an established curvature-measure convergence theorem. It does not establish convergence of critical points of the discrete actions.}
\GSfield{Assessment}{literal statement: correct with missing hypotheses; publication match: exact; verification: read the result; development: annotated; confidence: high}
\end{GSitem}
\begin{GSitem}{1981-546.8}{C}{Does curvature convergence imply convergence of the Regge equations?}
\GSfield{Kind; source}{question; \GSpage{P001003}}
\GSfield{The notes}{Immediately after the curvature theorem the page says \ru{Пока неясно, сходятся ли ур-я} (it is still unclear whether the equations converge).}
\GSfield{Statement}{For a specified family of fat triangulations approximating a smooth metric, one must distinguish convergence of the Regge action, convergence of its first variation or Einstein tensor in a distributional norm, and convergence of solutions of the resulting discrete boundary-value problem. The strongest literal reading asks about the equations; it supplies no gauge, boundary conditions, topology of solution convergence, or uniqueness class. A precise bounded subproblem is convergence of the distributional Einstein tensor of Regge interpolants of a fixed smooth metric.}
\GSfield{Reconstruction}{The three-way distinction and the interpolant subproblem are editorial reconstructions. The question is not promoted to an unqualified conjecture that arbitrary stationary Regge metrics converge to a prescribed Einstein metric.}
\GSfield{Fate}{For the subproblem, \citet[Theorem 4.1 and Corollary 4.3, pp. 16--17 of the preprint]{gawlikneunteufel2023einstein} prove convergence in an \(H^{-2}\) norm of the distributional densitized Einstein tensor for suitable positive-degree Regge finite-element interpolants on shape-regular meshes, in dimension at least three. The preprint first appeared in October 2023 and the journal version online in March 2025 (2026 volume). This is a partial match, not a theorem about convergence of all discrete solutions; degree-zero interpolation is not covered by the same positive conclusion.}
\GSfield{Present status}{As of November 2026, distributional tensor approximation has rigorous positive results under explicit mesh and degree hypotheses. The notes' unspecified convergence-of-equations question, and especially convergence of actual stationary metrics under a definite gauge and boundary problem, is not resolved by the sources verified here.}
\GSfield{Assessment}{literal statement: not determinable; publication match: partial; verification: read the result; development: annotated; confidence: medium; problem \hyperref[prob:GP-1981-21]{GP-1981-21}}
\GSfield{Worth developing}{Worth a further narrowly specified investigation: an explicit stable boundary-value problem could separate consistency of the discrete curvature from convergence of its solutions.}
\end{GSitem}
\begin{GSitem}{1981-546.9}{E}{Fat smooth triangulations and the danger of repeated barycentric subdivision}
\GSfield{Kind; source}{method; \GSpage{P001003}}
\GSfield{The notes}{The notes discourage barycentric subdivision because angles become small, illustrate a better planar subdivision, and refer to Müller's Ray--Singer torsion work for smooth triangulations.}
\GSfield{Statement}{A useful approximation family must combine mesh tending to zero with scale-invariant lower bounds on simplex shape. One standard fatness condition bounds \(\operatorname{vol}_j(\tau)/\operatorname{diam}(\tau)^j\) away from zero for every \(j\)-face \(\tau\). Repeated refinement must be chosen to preserve such bounds. The planar midpoint refinement illustrates shape preservation for a triangle; it is not by itself a construction in all dimensions.}
\GSfield{Reconstruction}{The explicit all-face fatness definition makes the notebook's angle requirement precise. The torsion paper is a bibliographic pointer, not an assertion that every smooth triangulation or every barycentric subdivision sequence works.}
\GSfield{Fate}{The geometric condition is defined and used in \citet[Section 1]{cheegermullerschrader1984curvature}. Müller's pre-talk torsion article was published in 1978; its title and date were checked in his academy bibliography, while no exact lemma in it is claimed here.}
\GSfield{Present status}{As of November 2026, nondegeneracy remains an essential hypothesis. The translation unnecessary weakens the Russian practical warning not to use such a refinement when it destroys the angle bound.}
\GSfield{Assessment}{literal statement: correct with missing hypotheses; publication match: partial; verification: read the result; development: annotated; confidence: high}
\end{GSitem}
\begin{GSitem}{1981-546.10}{A}{Higher piecewise flat curvature from flags and solid angles}
\GSfield{Kind; source}{construction; \GSpage{P001003-P001004}}
\GSfield{The notes}{Formula (18) sums hinge volumes in codimension \(2k\), multiplied by alternating products of normalized angles along nested even-codimension face flags. The margin says norm.; the next page normalizes full solid angle to one.}
\GSfield{Statement}{A piecewise flat manifold carries intrinsic Lipschitz--Killing curvature measures in even curvature degrees. Their local weights are universal alternating angle expressions attached to face flags, multiplied by the corresponding face volumes. Once the empty flag, angular normalization, and universal dimensional constants are fixed, these give the discrete counterparts of the smooth Lipschitz--Killing densities.}
\GSfield{Reconstruction}{The manuscript does not state all flag ranges or normalization constants. Its displayed sum is retained as evidence, while the exact published definition is the reference for a usable formula; arbitrary constants are not reconstructed from the margin.}
\GSfield{Fate}{The matching detailed definition is in \citet[Sections 2--3]{cheegermullerschrader1984curvature}, following the jointly announced lattice-gravity program of 1981--82. This is the curvature family actually treated there, not an arbitrary collection of Pontryagin forms.}
\GSfield{Present status}{As of November 2026, the construction is established. The notebook's abbreviated flag sum is a partial transcription of the definition, insufficient for numerical use without the missing conventions.}
\GSfield{Assessment}{literal statement: correct with missing hypotheses; publication match: partial; verification: read the result; development: annotated; confidence: high}
\end{GSitem}
\begin{GSitem}{1981-546.11}{A}{Higher curvature convergence and the Euler-form correction}
\GSfield{Kind; source}{theorem; \GSpage{P001004}}
\GSfield{The notes}{The notes write \(R_k(P_j)\to\int_M\operatorname{tr}(\Omega^k)\) and, for \(k=n/2\), identify that integral with \(\chi(M)\). They later mention multiplication by a standard form.}
\GSfield{Statement}{The convergent quantities are consistently normalized Lipschitz--Killing curvature measures, with smooth densities formed by the appropriate alternating curvature contractions and coframes. In even dimension \(n=2m\), the top density is the Euler form \(\operatorname{Pf}(\Omega/2\pi)\), with the standard orientation convention, whose integral is \(\chi(M)\). It is not generally \(\operatorname{tr}(\Omega^m)\). A trace of a matrix-valued curvature power has degree \(2m\) but represents a different invariant polynomial; in dimension two the orthogonal curvature matrix has trace zero.}
\GSfield{Reconstruction}{The replacement of the trace by the Lipschitz--Killing density is substantive mathematical emendation, not a corrected reading of the handwriting. The dimension-two countercheck is our elementary verification of the mismatch.}
\GSfield{Fate}{\citet[Introduction and Theorem 5.1]{cheegermullerschrader1984curvature} explicitly identify the top Lipschitz--Killing measure with Euler characteristic and establish curvature-measure convergence. The paper does not support the trace notation literally.}
\GSfield{Present status}{As of November 2026, the corrected theorem is established; the literal Euler trace identity is false, even for a round two-sphere.}
\GSfield{Assessment}{literal statement: false as written; publication match: partial; verification: read the result; development: annotated; confidence: high}
\end{GSitem}
\begin{GSitem}{1981-546.12}{E}{An angle-flag Gauss--Bonnet formula for Euler characteristic}
\GSfield{Kind; source}{theorem; \GSpage{P001004}}
\GSfield{The notes}{Formula (20) expresses \(\chi(M)\) as an alternating sum of products of normalized solid angles along flags starting at vertices and passing through even-dimensional faces.}
\GSfield{Statement}{For a closed even-dimensional piecewise flat manifold, the total normalized top Lipschitz--Killing curvature equals its Euler characteristic. A flag-angle expression computes the vertex weights. The empty-chain contribution and the convention for normalized full solid angle are necessary for the formula to hold, including in dimension zero.}
\GSfield{Reconstruction}{The topological conclusion is identifiable; the exact unprinted summation conventions are not. This item separates Gauss--Bonnet from the erroneous matrix-trace expression preceding it.}
\GSfield{Fate}{Combinatorial and polyhedral Gauss--Bonnet formulas predate the talk. \citet[Section 3]{cheegermullerschrader1984curvature} supply the particular intrinsic angle framework used in this announcement.}
\GSfield{Present status}{As of November 2026, the normalized identity is established. The notebook's bare summation should not be used without the empty-flag convention.}
\GSfield{Assessment}{literal statement: correct with missing hypotheses; publication match: partial; verification: read the result; development: annotated; confidence: high}
\end{GSitem}
\begin{GSitem}{1981-546.13}{F}{The unidentified earlier Brin formula}
\GSfield{Kind; source}{remark; \GSpage{P001004}}
\GSfield{The notes}{The note mentions a similar work of Brin presented in the seminar around a date transcribed as 1946, adding that its limiting meaning is unknown.}
\GSfield{Statement}{No exact author identity, publication, formula, or limit problem can be fixed from this sentence. The date and surname must remain diplomatic evidence pending identification; they do not define a separate mathematical conjecture.}
\GSfield{Reconstruction}{The handwriting was inspected. It does not justify silently changing the decade or assigning the remark to a particular Brin solely from a name resemblance.}
\GSfield{Fate}{Searches combining Brin, polyhedral curvature, solid angles, Euler characteristic, and the stated seminar chronology did not yield an exact publication match. The nearby Cheeger--Müller--Schrader formula is thematic context only.}
\GSfield{Present status}{As of November 2026, this provenance lead remains unresolved. No problem-index entry is created from it.}
\GSfield{Assessment}{literal statement: not determinable; publication match: none found; verification: unresolved lead; development: annotated; confidence: low}
\end{GSitem}
\begin{GSitem}{1981-546.14}{E}{Simplex metrics as positive-definite Gram matrices}
\GSfield{Kind; source}{construction; \GSpage{P001004}}
\GSfield{The notes}{For an \(n\)-simplex based at vertex zero, \(g_{ij}=\tfrac12(\ell_{0i}^2+\ell_{0j}^2-\ell_{ij}^2)\), and the metric space is identified with \(GL(n,\mathbb R)/O(n)\).}
\GSfield{Statement}{The admissible squared edge lengths of a labelled nondegenerate Euclidean \(n\)-simplex are in bijection with \(\operatorname{Sym}_n^+\) by the Gram formula. Positivity is equivalent to existence of linearly independent edge vectors with those lengths. The congruence action of \(GL(n,\mathbb R)\) is transitive with orthogonal stabilizer, giving \(\operatorname{Sym}_n^+\cong GL(n,\mathbb R)/O(n)\).}
\GSfield{Reconstruction}{Labels and a distinguished base vertex remove combinatorial quotient ambiguities. The reconstruction is elementary linear algebra, not a new theorem of the lecture.}
\GSfield{Fate}{This standard parametrization is used in the lattice program identified by \citet{cheegermullerschrader1982lattice}. Its formula is checked directly: the polarization identity recovers all pairwise distances from the Gram matrix.}
\GSfield{Present status}{As of November 2026, the identification is correct for nondegenerate labelled simplices. Arbitrary positive edge assignments need not lie in its domain.}
\GSfield{Assessment}{literal statement: correct with missing hypotheses; publication match: exact; verification: indirect; development: annotated; confidence: high}
\end{GSitem}
\begin{GSitem}{1981-546.15}{E}{The invariant metric measure and its missing edge-length Jacobian}
\GSfield{Kind; source}{construction; \GSpage{P001004}}
\GSfield{The notes}{The notes give the following two expressions and call them the same measure: \[\begin{aligned}\text{(22)}\quad &(\det g)^{-(n+1)/2}\prod_{i\leq j}dg_{ij},\\ \text{(23)}\quad &\operatorname{vol}(\sigma,\ell)^{-(n+1)}\prod d\ell_{ij}.\end{aligned}\]}
\GSfield{Statement}{Squared edge lengths \(s_{ij}=\ell_{ij}^2\) are linear coordinates on symmetric matrices, and \(\det g=(n!\operatorname{vol}\sigma)^2\). The congruence-invariant density on \(\operatorname{Sym}_n^+\) therefore has the coordinate expressions \[\begin{aligned}d\mu&=C(\det g)^{-(n+1)/2}\prod_{i\leq j}dg_{ij}\\ &=C'\operatorname{vol}(\sigma)^{-(n+1)}\prod_{i<j}\ell_{ij}\,d\ell_{ij}\\ &=C''\operatorname{vol}(\sigma)^{-(n+1)}\prod_{i<j}ds_{ij}.\end{aligned}\] The edge-coordinate products extend over all simplex edges, including those from vertex zero.}
\GSfield{Reconstruction}{The factor \(\prod\ell_{ij}\) is a required Jacobian if the notation means lengths, as explicitly defined on P001003. The independent finite-dimensional derivation and exact small-dimensional check are supplied in the code and output; no new invariant-measure claim is attributed to the notes.}
\GSfield{Fate}{The invariant density follows from the determinant of congruence on symmetric matrices, \(|\det A|^{n+1}\), and the elementary coordinate change. The lattice-gravity reference provides context; the correction itself is proved and checked here.}
\GSfield{Present status}{As of November 2026, (22) is correct up to normalization. Equation (23) is false for length coordinates; already for one edge the invariant measure is \(d\ell/\ell\), not \(d\ell/\ell^2\).}
\GSfield{Assessment}{literal statement: false as written; publication match: partial; verification: indirect; development: annotated; confidence: high}
\end{GSitem}
\begin{GSitem}{1981-546.16}{A}{Compatible simplex metrics and the induced finite-dimensional volume}
\GSfield{Kind; source}{construction; \GSpage{P001005}}
\GSfield{The notes}{The space of metrics on a triangulation embeds into the product of its simplex metric spaces. The notes impose gluing conditions and induce a metric, hence a volume measure, from invariant symmetric-space metrics.}
\GSfield{Statement}{For a fixed finite simplicial complex of pure dimension \(n\), take one \(\operatorname{Sym}_n^+\) for every top simplex and impose agreement of all shared edge lengths. The admissible compatible metrics form a smooth open cone in the independent squared-edge coordinates. Pulling back the product affine-invariant Riemannian metric yields a positive Riemannian metric and its volume density on this cone. The induced density depends on this choice and is not automatically the product of single-simplex densities with duplicated edge variables deleted.}
\GSfield{Reconstruction}{The distinction between restricting a Riemannian metric and naively restricting a product measure removes an ambiguity in the scribbled line. No continuum or probabilistic normalization is supplied by this finite-dimensional construction.}
\GSfield{Fate}{This is part of the jointly announced lattice-gravity program, identified bibliographically in \citet{cheegermullerschrader1982lattice}. The exact finite-dimensional construction is mathematically valid as stated here; the accessible metadata alone does not verify a later continuum theorem.}
\GSfield{Present status}{As of November 2026, the fixed-triangulation construction is sound. Convergence of these measures as the mesh tends to zero is an additional problem not made precise in the source.}
\GSfield{Assessment}{literal statement: correct with missing hypotheses; publication match: partial; verification: publisher metadata checked; development: annotated; confidence: high}
\end{GSitem}
\begin{GSitem}{1981-546.17}{E}{Chern's kinematic formula as the curvature-convergence mechanism}
\GSfield{Kind; source}{method; \GSpage{P001005}}
\GSfield{The notes}{For an isometric embedding \(M^n\subset\mathbb R^N\), the notes integrate \(\chi(M\cap E^k)\) over affine \(k\)-planes, with \(k=N-n+2\ell\), to recover the corresponding curvature integral. They propose the analogous polyhedral identity and eventual almost-everywhere agreement of the intersection Euler characteristics.}
\GSfield{Statement}{With invariant measures normalized consistently, affine Crofton formulas express the total \(2\ell\)-curvature term as a dimensional constant times the average Euler characteristic of generic \((N-n+2\ell)\)-plane sections, for smooth compact embedded manifolds and their suitable polyhedral counterparts. Recovering convergence from pointwise stabilization additionally requires integrable uniform control; almost-everywhere agreement alone is not a dominated-convergence proof.}
\GSfield{Reconstruction}{The curvature density is the Lipschitz--Killing density, not a raw trace of a curvature power. Transversality, the generic-section convention, normalization, and domination are the necessary reconstructed hypotheses.}
\GSfield{Fate}{The integral-geometric ingredients predate the talk. Their use in the piecewise-flat convergence argument is verified in \citet[Sections 3--5]{cheegermullerschrader1984curvature}; the notebook supplies a proof idea, not a complete proof of the limiting step.}
\GSfield{Present status}{As of November 2026, the corrected integral-geometric method is established. The source's higher-degree trace notation and bare almost-everywhere argument must not be copied as a complete theorem.}
\GSfield{Assessment}{literal statement: correct with missing hypotheses; publication match: partial; verification: read the result; development: annotated; confidence: high}
\end{GSitem}

\GStalk{547}{1982-02-15}{V. L. Ginzburg [catalog attribution]}{P001006}{Hamiltonian actions and the moment map}
A short announcement records the definition of a Hamiltonian Lie-group action and its moment map, including infinitesimal equivariance and Poisson compatibility. The manuscript heading says only Ginzburg; the initials are those of the supplied catalog and are not independently established by the page. No additional theorem or unpublished program is recoverable. This content ends before the dated Arnold heading on P001007.

\begin{GSitem}{1981-547.1}{E}{Hamiltonian Lie-group actions and the global exactness condition}
\GSfield{Kind; source}{definition; \GSpage{P001006}}
\GSfield{The notes}{The group \(A\) preserves the symplectic form, and contraction with an infinitesimal action field is written as globally (?) exact. Hamiltonians \(a\mapsto H_a\) are required to form a Lie-algebra homomorphism.}
\GSfield{Statement}{For a smooth action of a Lie group \(A\) on \((M,\omega)\), each fundamental field \(a_M\) is symplectic exactly when \(d(\iota_{a_M}\omega)=0\). A Hamiltonian lift requires a linear assignment \(H:\mathfrak a\to C^\infty(M)\) with \(dH_a=\iota_{a_M}\omega\), after fixing conventions, and the compatible Poisson identity \(\{H_a,H_b\}=H_{[a,b]}\). Closedness gives local primitives; global exactness and compatibility of the additive constants are extra requirements.}
\GSfield{Reconstruction}{The ambiguous word was inspected in the scan. A Hamiltonian action needs global primitives; reading the word locally would give only an infinitesimally symplectic action. Signs in the fundamental-field and Poisson conventions must be coordinated.}
\GSfield{Fate}{This is established pre-1982 Hamiltonian geometry, not a new result announced here. The distinction and obstruction framework were read in \citet[Sections 4.6--4.9]{cannas2007introduction}; the cited text records a 2002 course revised in 2007, used as a modern verification source, not as the origin.}
\GSfield{Present status}{As of November 2026, these definitions are standard. No theorem specific to this Ginzburg announcement can be recovered beyond them.}
\GSfield{Assessment}{literal statement: correct with missing hypotheses; publication match: exact; verification: read the result; development: annotated; confidence: high}
\end{GSitem}
\begin{GSitem}{1981-547.2}{E}{The equivariant Poisson moment map}
\GSfield{Kind; source}{construction; \GSpage{P001006}}
\GSfield{The notes}{The moment map is \(\mu(m)(a)=H_a(m)\); it is said to commute with the \(A\)-action and be compatible with the Poisson bracket.}
\GSfield{Statement}{Let \(\mu:M\to\mathfrak a^*\) be defined by a Hamiltonian Lie-algebra lift. With the compatible Lie--Poisson convention it is a Poisson map. If \(A\) is connected, infinitesimal equivariance integrates to \(A\)-equivariance. For disconnected \(A\), equivariance under the remaining components must be separately assumed or verified; a Lie-algebra condition alone does not control them.}
\GSfield{Reconstruction}{Connectedness is supplied to justify the full group statement. No reduction theorem, convexity assertion, or existence theorem for moment maps is present on this page.}
\GSfield{Fate}{\citet[Sections 4.7 and 4.9]{cannas2007introduction} verify the moment-map formulation and the connected-group qualification. The earlier occurrence 1980-524.2 concerns orbit symplectic forms and reduction; it is a thematic link, not recurrence of a definite question.}
\GSfield{Present status}{As of November 2026, the stated construction and Poisson property are established under the indicated conventions and equivariance hypothesis.}
\GSfield{Assessment}{literal statement: correct with missing hypotheses; publication match: exact; verification: read the result; development: annotated; confidence: high}
\end{GSitem}

\GStalk{549}{1982-02-22}{V. I. Arnold}{P001007}{Integral geometry and lattice analogues}
Arnold compares the parallel-body polynomial, curvature integrals, projection averages, and Euler-characteristic Crofton formulas, then proposes their lattice analogues and refers to McMullen's Vancouver congress lecture. The lattice discussion states normalized-volume and quotient-lattice ideas, not a fully quantified discrete replacement for every continuous formula. This heading precedes Gelfand's Macdonald-book intervention on the physical page, although the catalog numbers put record 548 first.

\begin{GSitem}{1981-549.1}{E}{The parallel-body polynomial and the curvature-index shift}
\GSfield{Kind; source}{theorem; \GSpage{P001007}}
\GSfield{The notes}{Arnold writes the volume of the \(\varepsilon\)-neighborhood of a convex body as a degree-\(n\) polynomial and expresses its coefficients as boundary integrals of elementary symmetric functions of the principal curvatures. The polynomial uses inconsistent indices, and the coefficient and curvature-polynomial indices are not shifted.}
\GSfield{Statement}{For a convex body \(D\subset\mathbb R^n\), \(\operatorname{vol}_n(D+\varepsilon B^n)\) is a polynomial in \(\varepsilon\geq0\). If \(\partial D\) is \(C^2\) and principal curvatures are taken nonnegative on a convex sphere, its constant coefficient is \(\operatorname{vol}_nD\), while the coefficient of \(\varepsilon^j\), \(1\leq j\leq n\), is \(j^{-1}\int_{\partial D}\sigma_{j-1}(\kappa_1,\ldots,\kappa_{n-1})\,dS\). Indeed, the outer normal Jacobian is \(\prod_i(1+t\kappa_i)\), whose integral in \(t\) gives this expression.}
\GSfield{Reconstruction}{The explicit coefficient formula is our normalization of the notes' unspecified constants. It requires \(n-1\) principal curvatures and the index shift. General convex bodies follow through the usual continuous extension of intrinsic volumes, not through an assumed classical second fundamental form.}
\GSfield{Fate}{This is the classical Steiner--Minkowski framework, already standard in the pre-talk integral geometry referenced by the notebook. \citet[pp. 491--492]{mcmullen1975metrical} discusses mixed volumes, intrinsic volumes, and their angle formulas in the 1974 congress lecture, published in 1975. Its full five-page text was read. The displayed smooth Jacobian derivation is the present reconstruction of the notebook's coefficient normalization.}
\GSfield{Present status}{As of November 2026, the corrected polynomial and smooth boundary formula are established. The formula with \(\sigma_j\) for the coefficient of \(\varepsilon^j\) is false as written.}
\GSfield{Assessment}{literal statement: false as written; publication match: partial; verification: read the result; development: annotated; confidence: high}
\end{GSitem}
\begin{GSitem}{1981-549.2}{E}{Projection averages and Crofton averages of Euler characteristic}
\GSfield{Kind; source}{theorem; \GSpage{P001007}}
\GSfield{The notes}{The same coefficients are expressed through average projection volumes and through \(\int\chi(h\cap D)\,dh\) over affine planes. Buffon's needle and Santaló are mentioned, and constants are to be calculated by balls.}
\GSfield{Statement}{Write \(\operatorname{vol}_n(D+\varepsilon B^n)=\sum_{j=0}^n\kappa_{n-j}V_j(D)\varepsilon^{n-j}\), where \(\kappa_r\) is the unit \(r\)-ball volume. Kubota's formula gives \(V_j(D)\) as a positive dimensional constant times the mean \(j\)-volume of orthogonal projections onto \(j\)-planes. Crofton expresses the same \(V_j\) as a dimensional constant times the invariant affine \((n-j)\)-plane integral of \(\chi(D\cap E)\). For convex compact sections the Euler characteristic is one if nonempty and zero otherwise; normalization constants depend on the chosen Haar measure.}
\GSfield{Reconstruction}{The notes mix the degree in \(\varepsilon\), the intrinsic-volume index, and the dimension of the averaging planes. The reconstruction fixes them explicitly. The Buffon reference is an analogy for averaging intersections, not a distinct formula for arbitrary nonconvex sets.}
\GSfield{Fate}{These are pre-talk integral-geometric representations. McMullen's congress survey supplies the convex-valuation context; the smooth and polyhedral Euler-section mechanism is also documented in \citet[Sections 3--5]{cheegermullerschrader1984curvature}. The latter is verification of the related method, not the first publication of Crofton or Kubota.}
\GSfield{Present status}{As of November 2026, the convex-body statements are established. Extensions to nonconvex sets require a specified valuation category and integrability assumptions.}
\GSfield{Assessment}{literal statement: correct with missing hypotheses; publication match: partial; verification: read the result; development: annotated; confidence: high}
\end{GSitem}
\begin{GSitem}{1981-549.3}{E}{Normalized lattice volume and height in a quotient lattice}
\GSfield{Kind; source}{construction; \GSpage{P001007}}
\GSfield{The notes}{In the lattice version, volume is related to counting integer points; an elementary parallelepiped is drawn, volume counts such cells, and quotienting by a direction is proposed as a lattice analogue of height.}
\GSfield{Statement}{Let \(\Lambda\subset V\) be a full lattice and normalize Lebesgue measure so that a fundamental parallelepiped has volume one. For independent \(v_1,\ldots,v_n\in\Lambda\), the half-open parallelepiped contains \([\Lambda:\sum_i\mathbb Zv_i]\) lattice points and has that normalized volume. If \(W\) is a rational subspace, the saturated sublattice \(\Lambda\cap W\) gives the quotient lattice \(\Lambda/(\Lambda\cap W)\) in \(V/W\). The primitive quotient step supplies the correct lattice height; an arbitrary unsaturated subgroup may introduce torsion instead.}
\GSfield{Reconstruction}{Half-open cells prevent boundary overcounting. The saturated quotient-lattice condition and normalization are supplied to make the diagram a precise statement; neither a Euclidean height nor a lattice-point count of an arbitrary closed parallelepiped is automatically this index.}
\GSfield{Fate}{The construction is elementary lattice geometry and belongs to the framework surveyed by \citet{mcmullen1975metrical}. The cited Vancouver lecture is a real 1974 contribution published in the 1975 proceedings; the notebook does not identify an additional theorem beyond this analogy.}
\GSfield{Present status}{As of November 2026, the index-volume identity is established. It should not be replaced by an equality between volume and the number of all boundary-inclusive lattice points.}
\GSfield{Assessment}{literal statement: correct with missing hypotheses; publication match: partial; verification: read the result; development: annotated; confidence: high}
\end{GSitem}
\begin{GSitem}{1981-549.4}{E}{The proposed replacement of volume by lattice-point count}
\GSfield{Kind; source}{remark; \GSpage{P001007}}
\GSfield{The notes}{Arnold says to transfer the formulas to the lattice case and replace volume by the number of integer points, referring to McMullen's Vancouver lecture for literature.}
\GSfield{Statement}{One precise nearby established theorem is Ehrhart's: for a lattice polytope \(P\) in its affine lattice, \(L_P(t)=\#(tP\cap\Lambda)\), for nonnegative integer \(t\), is a polynomial of degree \(\dim P\), with leading coefficient its normalized Euclidean volume. For rational polytopes the corresponding function is generally a quasipolynomial. This provides a discrete polynomial theory; it does not assert that every continuous curvature or Crofton formula remains unchanged after replacing volumes by counts.}
\GSfield{Reconstruction}{Ehrhart's theorem is a modern precise interpretation of the programmatic sentence, not a formula actually stated in the notes. The match to the program is thematic rather than an exact recorded claim.}
\GSfield{Fate}{\citet[pp. 493--494]{mcmullen1975metrical} proves polynomiality for lattice valuations and explicitly applies it to lattice-point counts and Ehrhart reciprocity. This full-text match confirms the notebook's bibliographic pointer. No new lattice analogue with exact quantifiers is recoverable from this line, so no new open problem is manufactured.}
\GSfield{Present status}{As of November 2026, lattice-point polynomiality is established under the stated arithmetic hypotheses. The scope of Arnold's intended further analogies remains unspecified.}
\GSfield{Assessment}{literal statement: not determinable; publication match: thematic only; verification: read the result; development: annotated; confidence: medium}
\end{GSitem}

\GStalk{548}{1982-02-22}{I. M. Gelfand}{P001007-P001008}{A partition complement identity and Guillemin's symplectic cones}
Gelfand gives a combinatorial exercise from Macdonald's book and its generating-function corollary, then introduces symplectic cones from work of Guillemin on Weyl asymptotics. The partition identity is identifiable exactly; its written corollary omits a conjugation and an inverse argument. The cone discussion introduces the subject resumed under Arnold's heading in record 550. The surviving notes state no final new spectral asymptotic theorem.

\begin{GSitem}{1981-548.1}{E}{The complementary beta-numbers of a partition and its conjugate}
\GSfield{Kind; source}{theorem; \GSpage{P001007-P001008}}
\GSfield{The notes}{For a partition inside an \(m\)-by-\(n\) rectangle, the notes ask to show that the numbers \(\lambda_i+n-i\), \(1\leq i\leq n\), and \(n-1+j-\lambda'_j\), \(1\leq j\leq m\), permute \(0,\ldots,m+n-1\). The example is \(\lambda=(5,4,4,1)\), \(\lambda'=(4,3,3,3,1)\).}
\GSfield{Statement}{Let \(m,n\geq0\), let \(m\geq\lambda_1\geq\cdots\geq\lambda_n\geq0\) be integers padded by zeros, and set \(\lambda'_j=\#\{i:\lambda_i\geq j\}\). Then the two displayed sets are disjoint and their union is exactly \(\{0,1,\ldots,m+n-1\}\). The vertical and horizontal steps of the rectangle-boundary path provide the two sets of labels.}
\GSfield{Reconstruction}{Nonnegativity and zero-padding are made explicit. This is an exercise with a known proof, not an unresolved research question. The empty rectangle cases are interpreted with empty sums and sets.}
\GSfield{Fate}{The exact match is \citet[Chapter I, Section 1, formula (1.7), pp. 3--4]{macdonald1995symmetric}. The opened second edition identifies the first edition as 1979; the notebook expressly attributes the exercise to that book. Earlier occurrence 1976-415.1 supplies the Young-diagram background but is not the same question.}
\GSfield{Present status}{As of November 2026, the identity is elementary and established. The supplied code checks all partitions in a bounded collection of rectangles by two independent enumeration methods; that finite check supplements the general boundary-path proof.}
\GSfield{Assessment}{literal statement: correct with missing hypotheses; publication match: exact; verification: read the result; development: annotated; confidence: high}
\end{GSitem}
\begin{GSitem}{1981-548.2}{E}{The corrected generating-function identity for complementary partitions}
\GSfield{Kind; source}{claim; \GSpage{P001008}}
\GSfield{The notes}{For \(f_{\lambda,n}(t)=\sum_{i=1}^n t^{\lambda_i+n-i}\), the notes write \(f_{\lambda,n}(t)+t^{m+n-1}f_{\lambda,m}(t)=(1-t^{m+n})/(1-t)\), with no conjugate partition and no inverse argument in the second term.}
\GSfield{Statement}{Under the rectangle hypotheses of the preceding item, the polynomial identity is \[f_{\lambda,n}(t)+t^{m+n-1}f_{\lambda',m}(t^{-1})=\sum_{r=0}^{m+n-1}t^r.\] The left side is a polynomial after the indicated exponent shift. At \(t=1\), the right side is \(m+n\); at \(t=0\), use the polynomial form rather than evaluating an isolated Laurent expression.}
\GSfield{Reconstruction}{Both the prime on \(\lambda'\) and the inverse in \(t^{-1}\) are mathematical emendations: the inspected scan omits them. The result follows by summing the complementary exponents, not from an independent reciprocity conjecture.}
\GSfield{Fate}{\citet[Chapter I, Section 1, formula (1.7'), p. 4]{macdonald1995symmetric} gives the exact corrected formula. The source's preceding set identity provides an independent local check of the emendation.}
\GSfield{Present status}{As of November 2026, the corrected identity is established. The displayed notebook corollary is false, for example for a one-row rectangle containing its full partition.}
\GSfield{Assessment}{literal statement: false as written; publication match: exact; verification: read the result; development: annotated; confidence: high}
\end{GSitem}
\begin{GSitem}{1981-548.3}{E}{Symplectic cones and their two scaling examples}
\GSfield{Kind; source}{definition; \GSpage{P001008}}
\GSfield{The notes}{Gelfand introduces \(\rho_a^*\omega=a\omega\), \(a>0\), from work of Guillemin on new conditions for Weyl asymptotics. The examples are simultaneous half-degree scaling of \((x,\xi)\in\mathbb R^{2n}\) and fiber scaling of \(T^*X\).}
\GSfield{Statement}{A symplectic cone for the subsequent quotient construction is a symplectic manifold with a free proper \(\mathbb R_{>0}\)-action satisfying \(\rho_a^*\omega=a\omega\). The model examples are \(\mathbb R^{2n}\setminus\{0\}\), with \(\rho_a(x,\xi)=(a^{1/2}x,a^{1/2}\xi)\), and \(T^*X\setminus0_X\), with \(\rho_a(x,\xi)=(x,a\xi)\). The removed loci are fixed by the actions on the full spaces.}
\GSfield{Reconstruction}{Homogeneity alone is the literal definition. Freeness, properness, and removal of zero are added only when making the contact-quotient claim of record 550 precise. No new Weyl remainder estimate is stated on this page.}
\GSfield{Fate}{The examples and symplectization geometry are standard pre-talk material. The later operator framework is identified indirectly by \citet[Introduction and Section 1]{ammannlauternistorvasy2004complex}, which explicitly traces its Weyl algebras to \citet{guillemin1985weyl}.}
\GSfield{Present status}{As of November 2026, the cone examples are valid with these quotient hypotheses. The surviving phrase new conditions for Weyl asymptotics is a topic label, not enough to reconstruct a particular improved asymptotic theorem.}
\GSfield{Assessment}{literal statement: correct with missing hypotheses; publication match: partial; verification: indirect; development: annotated; confidence: high}
\end{GSitem}

\GStalk{550}{1982-02-22}{V. I. Arnold}{P001008-P001009}{Contact cones, Weyl quantization and axiomatic operator algebras}
Arnold's intervention identifies the contact quotient of a symplectic cone. The ensuing exposition gives a Weyl kernel and axioms for complex-order operator classes: symbolic composition, asymptotic completeness, boundedness, compactness, trace-class thresholds, detection of principal symbols, splitting, and holomorphic order families. The connection with Guillemin's later axiomatic Weyl-algebra paper is identifiable, but the original 1985 article was not available in full during this audit. Local uniqueness, global quantization, and spectral conclusions must be kept separate.

\begin{GSitem}{1981-550.1}{E}{The contact quotient and symplectization of a free symplectic cone}
\GSfield{Kind; source}{theorem; \GSpage{P001008}}
\GSfield{The notes}{Arnold says that cones are locally the same and that their quotient is a contact manifold, writing symplectic cones equivalent to contact manifolds. He stresses the role of an action one-form.}
\GSfield{Statement}{For a free proper symplectic cone let \(E\) generate logarithmic dilation. Then \(\alpha=\iota_E\omega\) satisfies \(d\alpha=\omega\) and \(\rho_a^*\alpha=a\alpha\). Its kernel descends to a cooriented contact hyperplane distribution on the quotient. Conversely a cooriented contact manifold \((Y,\ker\beta)\) has symplectization \((Y\times\mathbb R_{>0},d(r\beta))\). Local equivalence is Darboux equivalence of contact structures, not a global classification of cones.}
\GSfield{Reconstruction}{The proof follows from Cartan's formula and nondegeneracy; the hypotheses are supplied from the preceding definition. A contact form is a choice of scale in the coorientation line, not globally determined by the hyperplane distribution alone.}
\GSfield{Fate}{This is classical symplectic/contact geometry, not a later solution of a seminar problem. The standard differential-form identities can be checked directly; the role of the cone in operator quantization is treated in the later framework discussed by \citet[Introduction]{ammannlauternistorvasy2004complex}.}
\GSfield{Present status}{As of November 2026, the correspondence is established under the free proper and coorientation assumptions. The zero loci in the preceding examples cannot be included in a smooth contact quotient.}
\GSfield{Assessment}{literal statement: correct with missing hypotheses; publication match: partial; verification: indirect; development: annotated; confidence: high}
\end{GSitem}
\begin{GSitem}{1981-550.2}{F}{The unqualified existence and uniqueness assertion for quantization}
\GSfield{Kind; source}{claim; \GSpage{P001008}}
\GSfield{The notes}{The notes claim that quantization is always possible and the Weyl algebra is uniquely constructed, then qualify this as microlocal uniqueness and a local classical Weyl algebra.}
\GSfield{Statement}{A Darboux contact chart admits the usual local microlocal operator model. Any uniqueness assertion must specify the category of algebras or algebroids, permitted equivalences, order filtration, symbol identification, and whether it is local or global. The quoted line does not supply those data, so it does not determine a global existence-and-uniqueness theorem.}
\GSfield{Reconstruction}{The local model is identifiable; a global quantization functor or a globally unique Hilbert-space representation is not reconstructed. A theorem about formal deformation quantization would address a different category unless that identification were separately justified.}
\GSfield{Fate}{Guillemin's 1985 paper is an identifiable later lead, cited explicitly by \citet[Introduction]{ammannlauternistorvasy2004complex}. The original article's full text could not be checked, and the latter source does not establish the notebook's unqualified global wording.}
\GSfield{Present status}{As of November 2026, the exact intended uniqueness theorem is undetermined from the notes. This item is F rather than an invented open quantization problem.}
\GSfield{Assessment}{literal statement: not determinable; publication match: thematic only; verification: unresolved lead; development: annotated; confidence: medium}
\end{GSitem}
\begin{GSitem}{1981-550.3}{E}{Weyl quantization and the momentum-sign inconsistency}
\GSfield{Kind; source}{construction; \GSpage{P001008}}
\GSfield{The notes}{The page sends \(\xi_k\) to \(+i\partial_{x_k}\) but uses the kernel \[(Pf)(x)=(2\pi)^{-n}\int p\!\left(\frac{x+y}{2},\xi\right)e^{i\xi\cdot(x-y)}u(y)\,dy\,d\xi.\] The source labels the left side with \(f\) and the input on the right with \(u\).}
\GSfield{Statement}{For a suitable symbol \(p\), the displayed oscillatory integral defines Weyl quantization \(\operatorname{Op}^w(p)\) on Schwartz functions. With its positive exponential phase, \(\operatorname{Op}^w(x_k)\) is multiplication and \(\operatorname{Op}^w(\xi_k)=-i\partial_{x_k}\). Choosing \(+i\partial\) instead requires the opposite Fourier phase. The same input function must occur on both sides.}
\GSfield{Reconstruction}{Oscillatory interpretation and a symbol class are required for unbounded symbols. The sign follows by differentiating the exponential and integrating distributionally, so it is independent of any bibliographic identification.}
\GSfield{Fate}{This is the standard pre-talk Weyl prescription, represented by the model operator calculi in \citet[Section 1, examples]{ammannlauternistorvasy2004complex}. The literal mismatch is preserved in the diplomatic source and corrected here.}
\GSfield{Present status}{As of November 2026, the quantization formula is established with a consistent Fourier convention. The scan confirms the two inconsistent signs and the input-function slip.}
\GSfield{Assessment}{literal statement: false as written; publication match: exact; verification: read the result; development: annotated; confidence: high}
\end{GSitem}
\begin{GSitem}{1981-550.4}{E}{Classical symbols on the isotropic cone and the principal-order quotient}
\GSfield{Kind; source}{definition; \GSpage{P001008}}
\GSfield{The notes}{Symbols have an expansion \(p=\sum_{j\geq0}p_j\), with \(p_j\) homogeneous of ordinary degree \(2\mu-2j\). The notes call this order \(\mu\), define \(W^\mu\), and take the symbol in \(W^\mu/W^{\mu-1}\).}
\GSfield{Statement}{In the global isotropic classical calculus on \(\mathbb R^{2n}\), use an asymptotic, not necessarily convergent, expansion outside a compact set, with terms homogeneous of degree \(2\mu-2j\) and derivative remainder estimates of the corresponding degrees. Under \((x,\xi)\mapsto a^{1/2}(x,\xi)\), the leading term has weight \(\mu\). The principal symbol is the class modulo order \(\mu-1\). For nonintegral complex order the full indexed family is not one totally ordered filtration; integer downward shifts are the stated inclusions.}
\GSfield{Reconstruction}{Global derivative estimates and the asymptotic meaning of the series are supplied. The notation for finite sums of different complex orders in the source is not used to infer extra nesting for arbitrary unrelated complex numbers.}
\GSfield{Fate}{The standard symbolic construction is encoded in the later extended-Weyl examples of \citet[Section 1]{ammannlauternistorvasy2004complex}. It is background to the framework attributed to Guillemin, rather than a new symbolic calculus created at this lecture.}
\GSfield{Present status}{As of November 2026, the classical symbol interpretation is established. Homogeneity without estimates is insufficient to obtain the operator properties listed on the next page.}
\GSfield{Assessment}{literal statement: correct with missing hypotheses; publication match: partial; verification: read the result; development: annotated; confidence: high}
\end{GSitem}
\begin{GSitem}{1981-550.5}{B}{An axiomatic complex-order Weyl algebra}
\GSfield{Kind; source}{construction; \GSpage{P001008-P001009}}
\GSfield{The notes}{A common dense domain in a Hilbert space carries operator classes \(\underline W^\mu\). Axioms A1--A5 require integer-order inclusions, additive order under products, an order drop for commutators, conjugation of order under adjoints, and an order-zero identity.}
\GSfield{Statement}{One seeks a filtered algebra of operators on a common invariant dense domain, with \(W^{\mu-1}\subset W^\mu\), \(W^\mu W^\nu\subset W^{\mu+\nu}\), \([W^\mu,W^\nu]\subset W^{\mu+\nu-1}\), \((W^\mu)^*\subset W^{\bar\mu}\), and \(I\in W^0\). Appropriate symbol spaces, topologies, and a smoothing ideal are part of a complete operator calculus; these inclusions alone do not prove a spectral law.}
\GSfield{Reconstruction}{This is an axiom system, not a theorem claiming that arbitrary operators satisfy it. The domain and topology qualifications prevent formal manipulations with unbounded operators from being read as unrestricted Hilbert-space equalities.}
\GSfield{Fate}{The later named framework is Guillemin's Weyl algebra, published in 1985. \citet[Introduction and Section 1]{ammannlauternistorvasy2004complex}, first circulated in November 2002 and published in 2004, explicitly identify and extend it. Because Guillemin is neither this record's catalogued speaker nor Gelfand, B is used for the later published framework; the original theorem is an indirect match.}
\GSfield{Present status}{As of November 2026, axiomatic Weyl calculi and their extensions are established. The source gives no complete list of spectral hypotheses sufficient to identify an exact new Weyl asymptotic result.}
\GSfield{Assessment}{literal statement: correct with missing hypotheses; publication match: partial; verification: indirect; development: annotated; confidence: high}
\end{GSitem}
\begin{GSitem}{1981-550.6}{E}{Asymptotic completeness of the operator filtration}
\GSfield{Kind; source}{construction; \GSpage{P001009}}
\GSfield{The notes}{Axiom A6 asks that for \(A_j\in W^{\mu-j}\) one can find \(A\in W^\mu\) with \(A-\sum_{j<N}A_j\in W^{\mu-N}\) for every \(N\).}
\GSfield{Statement}{In a classical pseudodifferential calculus with its usual symbol estimates, every prescribed sequence of successively lower-order symbol terms admits an asymptotic sum, unique only modulo smoothing operators after all symbol terms are fixed. Quantization transfers this Borel-type summation to the operator filtration. An ordinary norm-convergent series of the \(A_j\) is not asserted.}
\GSfield{Reconstruction}{The topology and smoothing ambiguity are supplied. A6 itself is a formal requirement on an abstract algebra; its verification belongs to each concrete calculus.}
\GSfield{Fate}{Asymptotic summation is pre-talk pseudodifferential background. It is explicitly part of the construction revisited in \citet[Section 1 and Section 2]{ammannlauternistorvasy2004complex}.}
\GSfield{Present status}{As of November 2026, the standard classical calculus satisfies this property. The notebook's summation should consistently be read asymptotically.}
\GSfield{Assessment}{literal statement: correct with missing hypotheses; publication match: exact; verification: read the result; development: annotated; confidence: high}
\end{GSitem}
\begin{GSitem}{1981-550.7}{E}{Boundedness, compactness and trace-class thresholds}
\GSfield{Kind; source}{theorem; \GSpage{P001009}}
\GSfield{The notes}{Axioms B require boundedness at nonpositive real order, compactness at negative real order, and trace class below some threshold \(-s_0\).}
\GSfield{Statement}{For scalar classical pseudodifferential operators on a closed \(d\)-manifold, order at most zero gives \(L^2\)-boundedness, negative order gives compactness, and real order below \(-d\) gives trace class. For the global isotropic \(\mathbb R^{2n}\) calculus normalized by weight \(\mu\), decay in both position and momentum gives trace class when \(\operatorname{Re}\mu<-n\). These conclusions require the complete symbol estimates; local negative order on a noncompact base alone does not imply compactness.}
\GSfield{Reconstruction}{The notes state only the existence of \(s_0\). The two concrete thresholds are separate models explaining that the order normalization and phase-space dimension matter; they are not claimed to have been stated in the lecture.}
\GSfield{Fate}{The standard mapping properties predate the talk. \citet[Introduction and Section 1]{ammannlauternistorvasy2004complex} specifically explain why their noncompact extension drops Guillemin's negative-order compactness axiom. That distinction verifies the necessary global qualification.}
\GSfield{Present status}{As of November 2026, these results hold in the indicated calculi. A translation-invariant negative-order operator on \(\mathbb R^d\) supplies the familiar obstruction to interpreting compactness as a purely local assertion.}
\GSfield{Assessment}{literal statement: correct with missing hypotheses; publication match: partial; verification: read the result; development: annotated; confidence: high}
\end{GSitem}
\begin{GSitem}{1981-550.8}{E}{Compactness detects a vanishing classical principal symbol}
\GSfield{Kind; source}{theorem; \GSpage{P001009}}
\GSfield{The notes}{Axiom C says that a compact member of \(W^0\) belongs to \(W^{-1}\).}
\GSfield{Statement}{For scalar classical order-zero pseudodifferential operators on a closed manifold, compactness forces the principal symbol to vanish, so the operator has order at most \(-1\). Conversely negative-order operators are compact. The word classical and the unit-spaced asymptotic expansion are essential: in a broader order-zero class, decay of a symbol such as \(1/\log(2+|\xi|)\) can yield compactness without order \(-1\).}
\GSfield{Reconstruction}{The counterexample is interpreted as a Fourier multiplier on a torus, with a harmless smooth modification near the origin; its eigenvalues tend to zero but do not satisfy an \(O(|\xi|^{-1})\) bound. This elementary check shows why the hidden symbol-class restriction matters.}
\GSfield{Fate}{The axiom belongs to the classical Weyl setting traced by \citet[Introduction]{ammannlauternistorvasy2004complex} to Guillemin. The standard compact-principal-symbol criterion predates that abstract organization.}
\GSfield{Present status}{As of November 2026, the classical assertion is established. The literal axiom is appropriate only after the class \(W^0\) has been fixed as in the preceding items.}
\GSfield{Assessment}{literal statement: correct with missing hypotheses; publication match: partial; verification: read the result; development: annotated; confidence: high}
\end{GSitem}
\begin{GSitem}{1981-550.9}{E}{Splitting the principal-symbol sequence}
\GSfield{Kind; source}{construction; \GSpage{P001009}}
\GSfield{The notes}{Axiom D displays \(0\to W^{\mu-1}\to W^\mu\to W^\mu/W^{\mu-1}\to0\) and asks that it split nicely.}
\GSfield{Statement}{The displayed sequence is algebraically exact by definition of the quotient. The meaningful analytic assertion is existence of a continuous linear quantization section from the designated principal-symbol space to \(W^\mu\), for specified locally convex topologies, giving a topological direct-sum decomposition as vector spaces. Such a section is generally noncanonical and need not be multiplicative.}
\GSfield{Reconstruction}{The word nicely is not a complete topology or norm estimate. The reconstruction describes the standard requirement rather than attributing an unspecified canonical algebra splitting to the notebook.}
\GSfield{Fate}{Principal-symbol exactness and quantization sections are standard. \citet[Section 1, condition (sigma)]{ammannlauternistorvasy2004complex} make explicit the kernel condition needed in an abstract extension.}
\GSfield{Present status}{As of November 2026, appropriate continuous sections exist in the usual classical models. An arbitrary quotient of topological vector spaces need not have one.}
\GSfield{Assessment}{literal statement: correct with missing hypotheses; publication match: partial; verification: read the result; development: annotated; confidence: high}
\end{GSitem}
\begin{GSitem}{1981-550.10}{E}{Holomorphic order families and complex powers}
\GSfield{Kind; source}{construction; \GSpage{P001009}}
\GSfield{The notes}{Axiom E requires \(A(\mu)\in W^\mu\) analytic in \(\mu\), with \(A(0)-I\in W^{-1}\) and multiplicativity modulo \(W^{\mu+\nu-1}\).}
\GSfield{Statement}{A standard model uses a strictly positive invertible elliptic operator \(P\) of normalized order one. Its complex powers satisfy \(P^\mu\in W^\mu\), \(P^0=I\), and \(P^\mu P^\nu=P^{\mu+\nu}\), in the appropriate holomorphic family calculus. If a kernel is present, it must be removed or a finite-rank correction specified. The notebook's axiom is weaker, requiring these identities only at principal order.}
\GSfield{Reconstruction}{Strict positivity, invertibility, the normalized order, and the topology of holomorphic families are supplied to give a rigorous concrete model. No claim that axiom E alone constructs complex powers of every elliptic element is made.}
\GSfield{Fate}{Seeley's complex-power construction predates the talk. The source's Guillemin-style alternative and its noncompact extension are read in \citet[Introduction, Sections 7--8]{ammannlauternistorvasy2004complex}; the paper explicitly distinguishes those analytic hypotheses.}
\GSfield{Present status}{As of November 2026, the concrete elliptic complex-power result and suitable abstract extensions are established. The source lists the axiom but does not state an ensuing spectral theorem.}
\GSfield{Assessment}{literal statement: correct with missing hypotheses; publication match: exact; verification: read the result; development: annotated; confidence: high}
\end{GSitem}

\GStalk{552}{1982-03-01}{Robert MacPherson}{P001010}{Proposed topics: perverse sheaves and bivariant theories}
The handwritten English explicitly introduces titles of a paper he would speak about. One diagram joins perverse sheaves, holonomic differential modules, and characteristic varieties; another proposes a theory attached to maps with homology and cohomology as its endpoints. These are substantive identifiable announcements, but the page does not document a completed lecture or proofs. The second topic matches the already published Fulton--MacPherson bivariant theory.

\begin{GSitem}{1981-552.1}{E}{The regular-holonomic Riemann--Hilbert correspondence in MacPherson's proposed title}
\GSfield{Kind; source}{claim; \GSpage{P001010}}
\GSfield{The notes}{A proposed talk-title diagram writes perverse sheaves approximately equivalent to holonomic \(\mathcal D\)-modules. No regularity hypothesis or functor is written.}
\GSfield{Statement}{On a complex manifold \(X\) of complex dimension \(d\), the analytic Riemann--Hilbert correspondence identifies regular holonomic left \(\mathcal D_X\)-modules with complex-constructible perverse sheaves through the de Rham functor, with the de Rham complex placed so that \(\operatorname{DR}(\mathcal O_X)=\mathbb C_X[d]\). The solution functor is the corresponding contravariant formulation. Arbitrary irregular holonomic modules are not recovered by ordinary constructible perverse sheaves.}
\GSfield{Reconstruction}{Regularity, analytic constructibility, and the degree convention are required reconstructions. The page records an announced subject, not a proof or a priority assertion by MacPherson.}
\GSfield{Fate}{\citet[main theorem, XIX.1]{kashiwara1980constructible}, a seminar lecture dated 22 April 1980, already announces the derived regular-holonomic correspondence. The detailed paper appeared in 1984. Thus the 1982 topic is pre-existing even though later expositions became definitive.}
\GSfield{Present status}{As of November 2026, the regular statement is established. The source's omission of regular is mathematically consequential; modern irregular correspondences require additional sheaf-theoretic structure and do not validate the literal unrestricted diagram.}
\GSfield{Assessment}{literal statement: correct with missing hypotheses; publication match: exact; verification: read the result; development: annotated; confidence: high}
\end{GSitem}
\begin{GSitem}{1981-552.2}{B}{What corresponds to a characteristic variety on the sheaf side?}
\GSfield{Kind; source}{question; \GSpage{P001010}}
\GSfield{The notes}{The diagram places characteristic variety below holonomic modules and a dashed question-mark arrow from perverse sheaves to it.}
\GSfield{Statement}{A precise realization assigns the microsupport \(\operatorname{SS}(F)\subset T^*X\) to a constructible complex \(F\). For a coherent differential module \(\mathcal M\), the microsupport of its derived solution complex equals \(\operatorname{Char}(\mathcal M)\), under the standard identification of complex cotangent covectors with real covectors. In the holonomic constructible case this is a conic complex Lagrangian set. A characteristic cycle adds multiplicities and should not be conflated with its support.}
\GSfield{Reconstruction}{The notebook asks only for the missing arrow. Microsupport is a precise later implementation, but the page does not determine whether MacPherson intended the set, the cycle, or a specific index formula. No further equality of cycle conventions is asserted.}
\GSfield{Fate}{\citet[Corollary 9.3.3, p. 166; Theorem 10.1.1, p. 174]{kashiwaraschapira1985microlocal} proves the equality, first in the regular Riemann--Hilbert setting and then for a bounded coherent differential-module complex. Both statements and their proof contexts were read in the full text. This B classification identifies a definitive published realization, not a claim that the geometric idea had no pre-1982 predecessors; Remark 9.3.4 itself points to earlier vanishing-cycle formulations.}
\GSfield{Present status}{As of November 2026, microsupport provides the required characteristic-set construction. The intended multiplicity-level statement in this brief proposed title remains unspecified.}
\GSfield{Assessment}{literal statement: not determinable; publication match: partial; verification: read the result; development: annotated; confidence: medium}
\end{GSitem}
\begin{GSitem}{1981-552.3}{E}{Bivariant homology attached to maps and its composition product}
\GSfield{Kind; source}{construction; \GSpage{P001010}}
\GSfield{The notes}{Another proposed title assigns \(H(X\to Y)\) to a map, with \(H(X\to\mathrm{pt})=H_*(X)\), \(H(X\xrightarrow{\mathrm{id}}X)=H^*(X)\), and a product covering composition of maps.}
\GSfield{Statement}{A graded bivariant theory associates \(\mathbb B^*(X\xrightarrow fY)\) to suitable morphisms, together with a product \(\mathbb B^i(X\to Y)\otimes\mathbb B^j(Y\to Z)\to\mathbb B^{i+j}(X\to Z)\), proper pushforwards, and pullbacks in designated Cartesian squares, satisfying compatibility axioms. Its covariant theory is \(\mathbb B_*(X)=\mathbb B^{-*}(X\to\mathrm{pt})\), and its contravariant theory is \(\mathbb B^*(X)=\mathbb B^*(X\xrightarrow{\mathrm{id}}X)\). Topological examples use the corresponding homology theory, often Borel--Moore homology on noncompact spaces.}
\GSfield{Reconstruction}{The grading reversal, properness, pullback squares, and choice of homology are omitted in the source and supplied from the formalism. Ordinary singular homology and Borel--Moore homology must not be interchanged without a compactness assumption.}
\GSfield{Fate}{This exactly identifies the framework of \citet[Sections 1--2]{fultonmacpherson1981categorical}, published in 1981 before this announced title. The publisher description and table of contents were verified; no unpublished 1982 theorem is inferred from the announcement.}
\GSfield{Present status}{As of November 2026, the bivariant framework and its homology/cohomology specializations are established. This source is evidence of dissemination, not of its first publication.}
\GSfield{Assessment}{literal statement: correct with missing hypotheses; publication match: exact; verification: publisher metadata checked; development: annotated; confidence: high}
\end{GSitem}

\GStalk{553}{1982-03-01}{Dusa McDuff}{P001010-P001011}{Canonical classes of volume-preserving diffeomorphism groups}
McDuff constructs the homotopy fiber of the discrete-to-topological comparison, the universal transverse volume foliation, and its compactly supported cohomology classes. She records a torus test, the parity of the top class for spheres, and the opposite parity of a secondary class on Euclidean space. These claims match her 1983 paper, whose author bibliography reproduces its abstract; the full proofs were not obtained. Fuchs's interjection occurs inside this discussion and has its own catalog record.

\begin{GSitem}{1981-553.1}{E}{The discrete-to-topological homotopy fiber for volume-preserving diffeomorphisms}
\GSfield{Kind; source}{definition; \GSpage{P001010}}
\GSfield{The notes}{The group \(\mathcal G=\operatorname{Diff}^{c}_{\omega,0}(M)\) has its smooth topology, while \(G\) is the same group made discrete. The source defines \(\bar B\mathcal G\) as the homotopy fiber of \(BG\to B\mathcal G\), explicitly instructing one first to replace the map by a fibration.}
\GSfield{Statement}{Let \((M^n,\omega)\) be a smooth oriented manifold without boundary with a nowhere-vanishing volume form. Let \(\mathcal G\) be the identity component of its compactly supported volume-preserving smooth diffeomorphisms, with the standard support-controlled smooth direct-limit topology when \(M\) is noncompact. Put \(\bar B\mathcal G=\operatorname{hofib}(B\mathcal G^\delta\to B\mathcal G)\). The source \(B\mathcal G^\delta\) is a \(K(\mathcal G^\delta,1)\); the fiber need not be one.}
\GSfield{Reconstruction}{The topology at infinity is made explicit so that compact parameter families have controlled supports. The notation is disentangled from the ordinary space \(B\mathcal G\) and from a set-theoretic quotient.}
\GSfield{Fate}{The construction is standard pre-talk classifying-space machinery. \citet[Section 1, p. 610]{mcduff1982local} gives exactly this notation and the direct-limit compact-support convention (p. 612); its full text was read.}
\GSfield{Present status}{As of November 2026, the construction is established. Its comparison with foliation classifying spaces is extra structure, not part of the definition.}
\GSfield{Assessment}{literal statement: correct with missing hypotheses; publication match: exact; verification: read the result; development: annotated; confidence: high}
\end{GSitem}
\begin{GSitem}{1981-553.2}{E}{Local homology depends on the group germ at the identity}
\GSfield{Kind; source}{claim; \GSpage{P001010}}
\GSfield{The notes}{The notes say that \(\bar B\mathcal G\) depends only on a neighborhood of the identity and call it a topological version of a Lie algebra in some sense.}
\GSfield{Statement}{The relevant locality is the germ of the topological group, including multiplication and inversion wherever defined near the identity. The local-homology construction is invariant under isomorphisms of such local-group germs. A homeomorphism between underlying neighborhoods without compatibility with the local group law is insufficient. The analogy with a Lie algebra is heuristic, not an equivalence of categories or an identification with Chevalley--Eilenberg cohomology.}
\GSfield{Reconstruction}{The group-law qualification supplies the precise meaning of neighborhood. No new assertion about integration of arbitrary infinite-dimensional Lie algebras is reconstructed.}
\GSfield{Fate}{\citet[Section 1, p. 610, the parenthetical definition of local homology]{mcduff1982local} explicitly states this locality and attributes it to van Est. The result is background to McDuff's new classes, not a separate new theorem of this talk.}
\GSfield{Present status}{As of November 2026, the local-group interpretation is established in the indicated classifying-space framework. The Lie-algebra comparison remains only the source's motivational analogy.}
\GSfield{Assessment}{literal statement: correct with missing hypotheses; publication match: exact; verification: read the result; development: annotated; confidence: high}
\end{GSitem}
\begin{GSitem}{1981-553.3}{A}{The universal flat foliation and its transverse volume form}
\GSfield{Kind; source}{construction; \GSpage{P001010-P001011}}
\GSfield{The notes}{The universal \(M\)-bundle becomes flat after pullback to the discrete classifying space and becomes topologically trivial after further pullback to the homotopy fiber. The resulting product carries a transverse foliation and a closed \(n\)-form \(\Omega\) restricting to \(\omega\) on every \(M\)-fiber; \(\Omega-\pi^*\omega\) vanishes near infinity.}
\GSfield{Statement}{On finite smooth parameter models for \(\bar B\mathcal G\), the pulled-back flat bundle admits the specified trivialization and a horizontal foliation transverse to its \(M\)-fibers. Since the holonomy preserves \(\omega\), the vertical volume forms extend to a closed transverse form \(\Omega\), annihilating the foliation directions. In the compact-support model the foliation is the product near infinity, so the difference \(\Omega-\pi_M^*\omega\) is supported in a compact part of \(M\) over each compact parameter family. Universal forms can be expressed in a compatible simplicial de Rham model.}
\GSfield{Reconstruction}{Ordinary differential forms on an unspecified infinite-dimensional classifying space are not assumed. The finite-parameter or simplicial formulation and support control make the source's construction rigorous.}
\GSfield{Fate}{\citet[Section 1, pp. 610--612]{mcduff1982local} gives the same pullback, trivialization, foliation, and support construction. The canonical-class paper followed in January 1983. The March 1982 lecture documents circulation before those publications.}
\GSfield{Present status}{As of November 2026, this construction is established. Its classifying map and its induced homology equivalence are stronger statements than the notebook records, so those additional theorems are not inserted as new occurrences.}
\GSfield{Assessment}{literal statement: correct with missing hypotheses; publication match: exact; verification: read the result; development: annotated; confidence: high}
\end{GSitem}
\begin{GSitem}{1981-553.4}{A}{The canonical cohomology classes obtained by fiber integration}
\GSfield{Kind; source}{construction; \GSpage{P001011}}
\GSfield{The notes}{The page defines \(c_k(M)\in H^k(\bar B\mathcal G;H_c^{n-k}(M;\mathbb R))\), by associating to a \(k\)-cycle \(K\) the closed \((n-k)\)-form written as \(x\mapsto\int_{K\times x}\Omega\).}
\GSfield{Statement}{For \(1\leq k\leq n\), the compact-in-\(M\) class represented by \(\Omega-\pi_M^*\omega\) has a Künneth or slant component defining \(c_k\) with the stated coefficients. Evaluation on a parameter cycle is integration in its \(k\) directions and leaves an \((n-k)\)-form on \(M\). The identity component acts trivially on \(H_c^*(M;\mathbb R)\), explaining the constant coefficient group in this setting. For \(k>0\), the pullback \(\pi_M^*\omega\) has zero parameter-degree component.}
\GSfield{Reconstruction}{The integral notation is interpreted as fiber integration, not as integration of an \(n\)-form over a lower-dimensional cycle to obtain a number. The isolated superscript \(k\) on \(\operatorname{Diff}\) in the cycle line is inconsistent with the compact-support superscript \(c\) in the surrounding definitions.}
\GSfield{Fate}{The exact family is identified in \citet[abstract, reproduced in the author's complete bibliography]{mcduff1983canonical}, published January 1983, pp. 345--356. The defining geometric mechanism is independently read in \citet[Section 1]{mcduff1982local}. The 1983 proof text was not obtained.}
\GSfield{Present status}{As of November 2026, these are established canonical classes. The geometric construction verifies the intended notation, but this audit does not claim to have read every independence-of-choices proof in the 1983 article.}
\GSfield{Assessment}{literal statement: correct with missing hypotheses; publication match: exact; verification: indirect; development: annotated; confidence: high}
\end{GSitem}
\begin{GSitem}{1981-553.5}{A}{The torus addition cycle detects the top class}
\GSfield{Kind; source}{construction; \GSpage{P001011}}
\GSfield{The notes}{For \(M=T^n\), the notes use addition \(a:T^n\times T^n\to T^n\), \((x,y)\mapsto x+y\), to produce a foliation and evaluate \(c_n(T^n)\) on the parameter torus.}
\GSfield{Statement}{Take the standard translation-invariant volume form \(\omega\) on \(T^n\). The fibers of addition define a foliation transverse to the second-factor torus fibers, with transverse form \(\Omega=a^*\omega\). For fixed \(y\), addition restricts to the translation \(x\mapsto x+y\), hence \(\int_{T^n\times\{y\}}\Omega=\int_{T^n}\omega\ne0\). This evaluates the represented cycle nontrivially and detects \(c_n(T^n)\).}
\GSfield{Reconstruction}{Translation invariance is a necessary hypothesis for the explicit formula as written; it is available for the standard torus model. The integral is independent of \(y\), as stated in the notes. The direct calculation is checked here and is not presented as new research.}
\GSfield{Fate}{The example belongs to the family announced in \citet{mcduff1983canonical}. Its publication match is partial because the exact example's location in the inaccessible full paper was not verified; the calculation itself is complete from the source data.}
\GSfield{Present status}{As of November 2026, the torus test is valid and supplies a concrete nonzero pairing. The test does not imply the sphere parity theorem.}
\GSfield{Assessment}{literal statement: correct with missing hypotheses; publication match: partial; verification: indirect; development: annotated; confidence: high}
\end{GSitem}
\begin{GSitem}{1981-553.6}{A}{Parity of the top canonical class on spheres}
\GSfield{Kind; source}{theorem; \GSpage{P001011}}
\GSfield{The notes}{Proposition 1 says \(c_n(S^n)=0\) for even \(n\) and \(c_n(S^n)\ne0\) for odd \(n\). A short explanation mentions a transverse foliation of \(K^n\times S^n\) and integration over a horizontal \ru{слой}.}
\GSfield{Statement}{For the canonical class on the homotopy fiber of the identity component of the volume-preserving diffeomorphism group of \(S^n\), with \(n\geq1\), the stated even/odd vanishing dichotomy holds. The parameter-cycle pairing is integration over a horizontal slice or section \(K\times\{x\}\), whereas a transverse volume form already annihilates tangent directions to an actual foliation leaf. Those two meanings of \ru{слой} cannot be interchanged in the argument.}
\GSfield{Reconstruction}{The precise class is the preceding \(c_n\). The recorded explanation is too short to supply the parity proof, and the ambiguous horizontal word is retained as an interpretive issue instead of being turned into a proof valid in all dimensions.}
\GSfield{Fate}{The exact parity assertion is explicitly stated in the abstract of \citet{mcduff1983canonical}, read in the author's bibliography; the paper was published in January 1983. Full proof access was not obtained, so the verification field deliberately records an indirect check.}
\GSfield{Present status}{As of November 2026, the parity theorem is established. The notebook statement is an exact match; its sketch is not a complete independent proof.}
\GSfield{Assessment}{literal statement: correct; publication match: exact; verification: indirect; development: annotated; confidence: high}
\end{GSitem}
\begin{GSitem}{1981-553.7}{A}{The secondary class on Euclidean space and its opposite parity}
\GSfield{Kind; source}{theorem; \GSpage{P001011}}
\GSfield{The notes}{The notes say that another class \(s\) is constructed and give Proposition 2: \(s(\mathbb R^n)\ne0\) for even \(n\), and zero for odd \(n\). They do not define \(s\) or its degree.}
\GSfield{Statement}{When \(H_c^i(M;\mathbb R)=0\) for \(0\leq i<n\), the primary classes vanish and McDuff's construction gives a secondary class \(s(M)\in H^{n-1}(\bar B\operatorname{Diff}^{c}_{\omega,0}(M);\mathbb R)\). For \(M=\mathbb R^n\), \(n\geq1\), it vanishes for odd \(n\) and is nonzero for even \(n\). In dimension two its normalization is twice the Calabi invariant in the cited construction.}
\GSfield{Reconstruction}{The cohomological hypotheses, degree, and normalization are recovered from the publication abstract, not from the shorthand source. No explicit secondary cochain formula is invented.}
\GSfield{Fate}{\citet[abstract]{mcduff1983canonical} matches the exact parity and supplies the missing hypotheses and degree; it attributes the nonvanishing argument to a relationship with Hurder's classes. That abstract was read, but the full secondary-class proof was not.}
\GSfield{Present status}{As of November 2026, the stated result is established with McDuff's normalization. There is no unresolved existence problem in the notebook's omission of the definition.}
\GSfield{Assessment}{literal statement: correct with missing hypotheses; publication match: exact; verification: indirect; development: annotated; confidence: high}
\end{GSitem}

\GStalk{551}{1982-03-01}{D. B. Fuchs}{P001010}{A subgroup-fibration analogy during McDuff's talk}
Fuchs contributes the familiar subgroup fibration and a heuristic about a very dense subgroup having an infinitesimal fundamental domain. The homotopy fiber is the mathematically reliable object here; the naive quotient for the identity map from a group with discrete topology cannot serve as its replacement. This is an intervention, not a separate developed lecture. Its source position follows the beginning of McDuff's record 553.

\begin{GSitem}{1981-551.1}{E}{The subgroup fibration and the homotopy-quotient qualification}
\GSfield{Kind; source}{theorem; \GSpage{P001010}}
\GSfield{The notes}{Fuchs writes \(G/H\to BH\to BG\), marked usually, in response to the homotopy-fiber construction.}
\GSfield{Statement}{If \(H\) is a closed subgroup of a Lie group \(G\), the map \(EG/H\to EG/G\) is a fiber bundle with fiber \(G/H\); since \(EG/H\) models \(BH\), this realizes the stated homotopy-fiber sequence. More generally suitable local-section hypotheses are needed. For the map \(G^\delta\to G\), the underlying image equals all of \(G\), but the source topology differs, so the naive quotient is a point and cannot replace \(\operatorname{hofib}(BG^\delta\to BG)\).}
\GSfield{Reconstruction}{The familiar closed-subgroup case is the rigorous content; it is not applied literally to a topologically different copy of the same group. The relevant replacement in the ongoing discussion is the homotopy fiber already defined by McDuff.}
\GSfield{Fate}{This is classical classifying-space theory. The exact comparison-fiber usage can be checked in \citet[Section 1, p. 610]{mcduff1982local} and \citet[Section 1]{milnor1983discrete}; neither source identifies the fiber with the naive quotient for the discrete identity map.}
\GSfield{Present status}{As of November 2026, the subgroup fibration is established under its local triviality hypotheses. The intervention is explanatory, not a new separate fibration theorem.}
\GSfield{Assessment}{literal statement: correct with missing hypotheses; publication match: partial; verification: read the result; development: annotated; confidence: high}
\end{GSitem}
\begin{GSitem}{1981-551.2}{F}{The infinitesimal fundamental-domain analogy for a dense subgroup}
\GSfield{Kind; source}{remark; \GSpage{P001010}}
\GSfield{The notes}{Fuchs says that if \(H\) is very dense in \(G\), a fundamental domain is an infinitesimal neighborhood of the identity.}
\GSfield{Statement}{Density alone neither produces a measurable fundamental domain nor equips a quotient with a manifold neighborhood. The phrase can motivate the local-group homotopy fiber, but it specifies no topology, measure, or limiting process from which a literal infinitesimal-domain statement follows.}
\GSfield{Reconstruction}{The word very is an insertion visible in the scan. It is not replaced by a technical density condition that the speaker did not state.}
\GSfield{Fate}{No exact independent publication or theorem was identifiable from this metaphor. The neighboring local-homology construction is identifiable, but that does not turn the metaphor into a theorem.}
\GSfield{Present status}{As of November 2026, this remains a heuristic interpretation of the discussion, with no definite new problem to index.}
\GSfield{Assessment}{literal statement: not determinable; publication match: none found; verification: unresolved lead; development: annotated; confidence: high}
\end{GSitem}

\GStalk{554}{1982-03-01}{John Milnor}{P001011-P001012}{Homology of Lie groups made discrete}
Milnor asks how the homology of a Lie group changes after forgetting its topology. The mod-prime isomorphism conjecture is accompanied by rational vanishing for compact groups, a large integral homology obstruction, solvable cases, finite-coefficient split surjectivity, integral cohomological injectivity, and low-dimensional results. The 1983 article was read and provides exact theorem-level identifications. Its proved partial statements do not settle the general Friedlander--Milnor conjecture.

\begin{GSitem}{1981-554.1}{E}{The comparison map after forgetting the topology of a Lie group}
\GSfield{Kind; source}{construction; \GSpage{P001011-P001012}}
\GSfield{The notes}{Milnor begins with a Lie group \(G\) and its underlying discrete group, writes \(BG^\delta\to BG\), and asks what happens to homology. He lists algebraic \(K\)-theory, characteristic classes of foliations, and cutting polyhedra as applications.}
\GSfield{Statement}{For a finite-dimensional real Lie group \(G\), the continuous identity homomorphism \(G^\delta\to G\) induces \(i:BG^\delta\to BG\). Here \(BG^\delta\) is a \(K(G^\delta,1)\), whereas \(BG\) is the classifying space for the original topology. The comparison uses singular homology with constant coefficients. Integral, rational, and finite coefficients yield substantially different questions; the coefficient group must be fixed before stating an isomorphism or vanishing assertion.}
\GSfield{Reconstruction}{The untwisted-coefficient convention is recovered from Milnor's publication. The named applications are motivations, not separately stated equivalences or theorems about scissors congruence.}
\GSfield{Fate}{\citet[Introduction and Section 1, pp. 72--73]{milnor1983discrete} gives this construction and the same applications. The machinery predates the talk. The 1 March 1982 notes verify circulation of this particular account before its 1983 publication; no earlier dated version was verified.}
\GSfield{Present status}{As of November 2026, this comparison is a standard construction. It should be read alongside McDuff's homotopy-fiber construction in 1981-553.1, without identifying a finite-dimensional Lie group with a diffeomorphism group.}
\GSfield{Assessment}{literal statement: correct; publication match: exact; verification: read the result; development: annotated; confidence: high}
\end{GSitem}
\begin{GSitem}{1981-554.2}{C}{The Friedlander--Milnor finite-coefficient isomorphism conjecture}
\GSfield{Kind; source}{conjecture; \GSpage{P001012}}
\GSfield{The notes}{The notes label as a conjecture that \(H_*(BG^\delta;\mathbb Z/p)\to H_*(BG;\mathbb Z/p)\) is an isomorphism, then list proved special cases.}
\GSfield{Statement}{For every finite-dimensional real Lie group \(G\), prime \(p\), and integer \(q\geq0\), does the identity comparison induce an isomorphism \[H_q(BG^\delta;\mathbb F_p)\longrightarrow H_q(BG;\mathbb F_p)?\] Equivalently, one asks for all constant finite abelian coefficient groups, by coefficient exact sequences, or for the corresponding mod-prime cohomology comparison. The conjecture is not an integral or rational isomorphism assertion. Solvable identity component is a proved case; finite-component split surjectivity leaves a possible kernel.}
\GSfield{Reconstruction}{The all-degrees and all-primes quantifiers, constant coefficients, and finite-dimensional Lie setting are explicit here. Milnor's original wording permits arbitrary component groups; the finite-component form in Friedlander's later survey is not substituted for the broader source statement.}
\GSfield{Fate}{The exact conjecture is in \citet[Introduction, p. 72]{milnor1983discrete}. \citet[Conjecture 32.4 and p. 92]{friedlander2008conjecture} records the finite-component formulation and stable progress, which does not establish the fixed-group assertion in all degrees. A renewed proof claim occurs in \citet[abstract and version history]{amrani2015friedlander}: first posting 2015, withdrawn version in 2016, new version in July 2021. No independently verified proof or journal acceptance for that claim was obtained. A 2025 publisher extract, \citet[Conjecture 1.1]{dranishnikov2025generalized}, still states the Lie-group assertion as a conjecture; its full article was not accessible and is marked unverified.}
\GSfield{Present status}{As of November 2026, the verified literature supports substantial partial results, not a certified general resolution. The index therefore retains C. Recent abelian results for more general topological groups and counterexamples outside the Lie category do not answer this exact question. Search failure is not presented as proof of current impossibility or of the invalidity of a particular claimed proof.}
\GSfield{Assessment}{literal statement: correct; publication match: exact; verification: read the result; development: annotated; confidence: medium; problem \hyperref[prob:GP-1981-22]{GP-1981-22}}
\GSfield{Worth developing}{Worth developing only through a sharply delimited group and degree or an independently audited proof claim: the precise comparison is central, while the general claim exceeds a bounded companion project.}
\end{GSitem}
\begin{GSitem}{1981-554.3}{E}{Positive-degree rational homology vanishes under the comparison for compact groups}
\GSfield{Kind; source}{theorem; \GSpage{P001012}}
\GSfield{The notes}{For compact \(G\), the notes say that the induced map in reduced rational homology is zero and label the assertion Chern--Weyl. The reduced-homology tildes already exclude degree zero.}
\GSfield{Statement}{If \(G\) is a compact Lie group and \(q>0\), then \(i_*:H_q(BG^\delta;\mathbb Q)\to H_q(BG;\mathbb Q)\) is zero; the same holds with real coefficients. Indeed, Chern--Weil classes generate real cohomology of \(BG\), and all positive-degree curvature polynomials vanish on a flat bundle. Pairing with homology proves the statement. The map in degree zero is the identity, not zero.}
\GSfield{Reconstruction}{The positive-degree formulation unfolds the source's reduced-homology notation. The curvature-polynomial argument is supplied from the publication, not written out on this page. This does not assert that rational homology of the discrete group itself vanishes, nor that every integral comparison is zero without further information on torsion.}
\GSfield{Fate}{\citet[Lemma 8, p. 78, and its proof, p. 80]{milnor1983discrete} supplies the exact statement and Chern--Weil argument. The underlying flat-bundle vanishing mechanism is pre-talk theory; the paper presents it as one of the known tools of the subject.}
\GSfield{Present status}{As of November 2026, the statement is established. It sharply distinguishes the rational comparison from the finite-coefficient conjecture and does not contradict integral cohomological injectivity.}
\GSfield{Assessment}{literal statement: correct; publication match: exact; verification: read the result; development: annotated; confidence: high}
\end{GSitem}
\begin{GSitem}{1981-554.4}{E}{The large integral second homology of a Lie group made discrete}
\GSfield{Kind; source}{theorem; \GSpage{P001012}}
\GSfield{The notes}{The page explicitly assumes \(G\ne1\) connected, contrasts finite generation of \(H_2(BG;\mathbb Z)\) with the discrete group, and credits Sah--Wagoner for an epimorphism onto an uncountable sum of copies of \(\mathbb Q\). The sum symbol is written as a large disjoint-union sign.}
\GSfield{Statement}{If \(G\) is a nontrivial connected finite-dimensional Lie group, \(H_2(BG^\delta;\mathbb Z)\) has a quotient which is an uncountable rational vector space. In contrast \(H_2(BG;\mathbb Z)\cong\pi_1(G)\) is finitely generated. This explains why a finite-coefficient comparison cannot be extrapolated to integral homology. Divisible rational summands disappear after tensoring with \(\mathbb F_p\), though one must also account for the adjacent torsion term in the coefficient sequence.}
\GSfield{Reconstruction}{Connectedness and nontriviality are actually explicit in the scan and are retained. Interpreting the handwritten union symbol as a vector-space direct sum is a mathematical emendation. The statement concerns a quotient, not a complete calculation of \(H_2\) or a canonical splitting.}
\GSfield{Fate}{\citet[Introduction, p. 72, and the Sah--Wagoner reference]{milnor1983discrete} explicitly attributes this pre-talk result to Sah and Wagoner. Milnor's statement was read; the original 1977 article was not obtained, so the historical attribution is checked indirectly rather than its original proof being claimed as read.}
\GSfield{Present status}{As of November 2026, the quoted obstruction is established in the indicated connected nontrivial setting. The notebook's hypotheses should be retained rather than extending this to arbitrary disconnected groups.}
\GSfield{Assessment}{literal statement: correct with missing hypotheses; publication match: exact; verification: indirect; development: annotated; confidence: high}
\end{GSitem}
\begin{GSitem}{1981-554.5}{A}{The isomorphism conjecture for groups with solvable identity component}
\GSfield{Kind; source}{theorem; \GSpage{P001012}}
\GSfield{The notes}{The first proved statement says that the conjecture holds when the Lie algebra \(\mathfrak g\) is solvable and calls the proof elementary.}
\GSfield{Statement}{If the identity component \(G_0\) of a finite-dimensional Lie group \(G\) is solvable, then \(BG^\delta\to BG\) induces isomorphisms in all homology and cohomology degrees with every constant finite abelian coefficient group. No finiteness assumption on the component group is needed. Locally isomorphic groups have the same relevant local homology; a simply connected solvable group is handled by induction using a quotient onto \(\mathbb R\) and a lower-dimensional simply connected solvable kernel.}
\GSfield{Reconstruction}{The coefficient convention and lack of a finite-component restriction distinguish this lemma from the next theorem. The induction uses finite-coefficient acyclicity of uniquely divisible discrete abelian groups, not integral acyclicity of \(\mathbb R^\delta\).}
\GSfield{Fate}{\citet[Lemmas 1--3, pp. 73--74]{milnor1983discrete} gives this exact argument. The 1982 talk precedes the 1983 publication; Milnor expressly describes the methods as largely known, so A denotes the near-contemporary publication of the reported result and does not assert a wholly new proof method.}
\GSfield{Present status}{As of November 2026, this case is established. It includes abelian Lie groups and disconnected groups with solvable identity component, but not a general semisimple identity component.}
\GSfield{Assessment}{literal statement: correct; publication match: exact; verification: read the result; development: annotated; confidence: high}
\end{GSitem}
\begin{GSitem}{1981-554.6}{A}{Finite-coefficient split surjectivity for finitely many components}
\GSfield{Kind; source}{theorem; \GSpage{P001012}}
\GSfield{The notes}{The second proved statement assumes finitely many components and asserts that the mod-prime homology comparison is a split epimorphism.}
\GSfield{Statement}{For a Lie group \(G\) with finitely many components, finite abelian group \(A\), and \(q\geq0\), the map \(H_q(BG^\delta;A)\to H_q(BG;A)\) is a split surjection. The associated cohomology map is a split injection. A maximal compact subgroup \(K\), the normalizer \(N\) of a maximal torus of \(K\), and \(\chi(K/N)=1\) give a transfer splitting; the solvable-identity-component case applies to \(N\).}
\GSfield{Reconstruction}{The extension from the written mod-prime coefficients to every finite abelian coefficient group is supplied from the published theorem. Splitting is as abelian groups, and no canonical or multiplicative splitting is inferred. The possible complementary summand in the discrete homology is exactly what a general isomorphism theorem would have to eliminate.}
\GSfield{Fate}{\citet[Theorem 1, p. 75]{milnor1983discrete} is an exact full-text match. The paper credits the Becker--Gottlieb transfer and describes its statement as generalizing a Bott--Heitsch result. The source verifies circulation by March 1982; publication followed in 1983.}
\GSfield{Present status}{As of November 2026, this split-surjectivity theorem is established. It is a proper partial result toward 1981-554.2, not a proof of that conjecture.}
\GSfield{Assessment}{literal statement: correct; publication match: exact; verification: read the result; development: annotated; confidence: high}
\end{GSitem}
\begin{GSitem}{1981-554.7}{A}{Integral cohomology injects after discretization}
\GSfield{Kind; source}{theorem; \GSpage{P001012}}
\GSfield{The notes}{The notes append an arrow between the integral cohomology groups; the handwritten arrow shape does not safely distinguish an isomorphism from an injection.}
\GSfield{Statement}{For a Lie group \(G\) with finitely many components, \(i^*:H^q(BG;\mathbb Z)\hookrightarrow H^q(BG^\delta;\mathbb Z)\) is injective for every \(q\). If a class maps to zero, finite-coefficient injectivity forces it into \(mH^q(BG;\mathbb Z)\) for every positive integer \(m\); finite generation of \(H^q(BG;\mathbb Z)\) makes their intersection zero.}
\GSfield{Reconstruction}{The direction is contravariant, and the conclusion is injectivity, not an integral cohomology isomorphism. Rational vanishing is compatible with this statement because the discrete classifying space need not be of finite type: rationalizing integral cohomology is not automatically the same as taking rational cohomology there.}
\GSfield{Fate}{\citet[Corollary 1 and its proof, p. 76]{milnor1983discrete} settles the arrow's mathematical meaning exactly. This is a substantive clarification of an ambiguous arrow, not a claim that the scan contains the modern hook symbol.}
\GSfield{Present status}{As of November 2026, integral injectivity holds under the finite-component hypothesis. An analogous integral homology surjectivity statement is explicitly false in the examples discussed by Milnor.}
\GSfield{Assessment}{literal statement: not determinable; publication match: exact; verification: read the result; development: annotated; confidence: high}
\end{GSitem}
\begin{GSitem}{1981-554.8}{A}{The finite-coefficient comparison in degree one}
\GSfield{Kind; source}{theorem; \GSpage{P001012}}
\GSfield{The notes}{The notes say that the conjecture always holds in dimension one.}
\GSfield{Statement}{For every Lie group \(G\) and constant finite abelian group \(A\), the map \(H_1(BG^\delta;A)\to H_1(BG;A)\) is an isomorphism. For connected \(G\), generation by one-parameter subgroups makes its abstract abelianization divisible, so its tensor product with a finite group vanishes; \(H_1(BG;A)=0\) as well. For arbitrary components, the component group gives the same degree-one contribution on both sides.}
\GSfield{Reconstruction}{Dimension means homological degree, not the dimension of the Lie group. The finite-coefficient qualification is essential: a connected abelian Lie group need not have zero integral \(H_1(BG^\delta;\mathbb Z)\).}
\GSfield{Fate}{\citet[Introduction, p. 72; Section 1 and Section 3]{milnor1983discrete} explicitly records the all-groups degree-one result. The elementary argument above spells out the coefficient reason and is not attributed as a new seminar discovery.}
\GSfield{Present status}{As of November 2026, the degree-one assertion is established. It supplies no general higher-degree isomorphism.}
\GSfield{Assessment}{literal statement: correct with missing hypotheses; publication match: exact; verification: read the result; development: annotated; confidence: high}
\end{GSitem}
\begin{GSitem}{1981-554.9}{E}{The degree-two comparison for real and complex special linear groups}
\GSfield{Kind; source}{theorem; \GSpage{P001012}}
\GSfield{The notes}{For degree two the notes list the interesting examples \(SL(n,\mathbb R)\) and \(SL(n,\mathbb C)\), without an argument.}
\GSfield{Statement}{For \(G=SL(n,\mathbb R)\) or \(SL(n,\mathbb C)\), \(n\geq2\), the comparison \(H_2(BG^\delta;\mathbb F_p)\to H_2(BG;\mathbb F_p)\) is an isomorphism for every prime \(p\), and likewise with constant finite abelian coefficients. Milnor's proof framework separates the finitely generated \(\pi_1G\) contribution from the second integral homology of the comparison homotopy fiber; the latter is uniquely divisible for these Chevalley examples.}
\GSfield{Reconstruction}{The homological degree and coefficient qualifications are restored. A bar on the homotopy-fiber notation in the published argument is essential: it is not the assertion that all of \(H_2(BG;\mathbb Z)\) or \(H_2(BG^\delta;\mathbb Z)\) is uniquely divisible.}
\GSfield{Fate}{\citet[Lemmas 6--7 and examples, pp. 77--78]{milnor1983discrete} identifies these exact cases and refers to the earlier Sah--Wagoner proof based on Steinberg, Moore, and Matsumoto. The exact statement in Milnor was read; no new independent verification of those deeper original proofs is claimed.}
\GSfield{Present status}{As of November 2026, these low-degree cases are established. A fixed-rank assertion in all homological degrees is substantially stronger and is not inferred.}
\GSfield{Assessment}{literal statement: correct with missing hypotheses; publication match: exact; verification: read the result; development: annotated; confidence: high}
\end{GSitem}

\GStalk{555}{1982-03-01}{André Haefliger}{P001012-P001013}{Smooth coinvariants for the modular action on the circle}
Haefliger asks whether every smooth function on the projective circle is a sum of the two generator coboundaries for the modular group. Exact smooth solvability is called unknown; smooth density is asserted separately. The source links the question to foliations with minimal leaves but gives no theorem making the connection precise. The companion proves exact continuous solvability, independently checked, while leaving smooth regularity unresolved. Kirillov's next heading is on P001014, outside this record.

\begin{GSitem}{1981-555.1}{C}{Exact smooth coinvariants for the modular action on the projective circle}
\GSfield{Kind; source}{question; \GSpage{P001012-P001013}}
\GSfield{The notes}{For \(\Gamma=PSL(2,\mathbb Z)\) acting on the upper half-plane and its projective-circle boundary by linear-fractional transformations, Haefliger asks whether every smooth \(f\) can be written \(g_1-g_1(x+1)+g_2-g_2(-1/x)\). The notes say that existence of such smooth \(g_1,g_2\) is unknown.}
\GSfield{Statement}{Let \(X=\mathbb {RP}^1\) with its standard smooth structure, \(T(x)=x+1\), \(S(x)=-1/x\), and \(D(u,v)=u-u\circ T+v-v\circ S\). Is the continuous linear map \(D:C^\infty(X;\mathbb R)^2\to C^\infty(X;\mathbb R)\) surjective? The transformations are smooth on the projective circle, and the central element of \(SL(2,\mathbb Z)\) acts trivially. All smooth right sides, including constants, are asked for; no zero-average condition is imposed.}
\GSfield{Reconstruction}{The real-valued convention is harmless because real and imaginary parts separate. The compact smooth structure at infinity is necessary: smoothness on the affine line alone is a weaker requirement. The formula uses two fixed generator coboundaries, not an unrestricted infinite sum.}
\GSfield{Fate}{The exact modular-action equation was not found in the checked primary literature. Haefliger's 1980 minimal-leaf paper is an unverified related lead; \citet[abstract]{abouqatebelkacimi2000invariant} treats invariant functionals and basic currents, but its compact-group theorem does not answer this noncompact discrete action. The development below proves, with a complete independent audit, that \(D:C(X)^2\to C(X)\) is surjective and has a uniform existence bound. It supplies no derivative estimate for its continuous solutions.}
\GSfield{Present status}{As of November 2026, the strongest result verified in this companion is exact continuous solvability. Exact smooth solvability remains unproved here and no exact published resolution was verified. This is a bounded partial result, not an assertion of historical novelty or a definitive world-literature claim of openness.}
\GSfield{Assessment}{literal statement: correct; publication match: none found; verification: unresolved lead; development: developed in this chapter; confidence: medium; problem \hyperref[prob:GP-1981-23]{GP-1981-23}}
\GSfield{Worth developing}{Worth a focused further investigation because the continuous theorem removes measure obstructions, leaving a concrete smooth closed-range or derivative-estimate question for a single explicit operator.}
\end{GSitem}
\begin{GSitem}{1981-555.2}{C}{The asserted smooth density is distinct from exact solvability}
\GSfield{Kind; source}{claim; \GSpage{P001013}}
\GSfield{The notes}{After calling the exact decomposition unknown, the last page says that the displayed sums approximate every smooth function in the \(C^\infty\) topology and that this is easy.}
\GSfield{Statement}{The strongest literal interpretation is that \(\operatorname{ran}D\) is dense in the usual Fr\'echet topology of \(C^\infty(\mathbb {RP}^1)\). By continuous-dual separation, this is equivalent to absence of a nonzero distribution invariant under both \(S\) and \(T\). Even if this assertion holds, a dense range in a Fr\'echet space need not be the whole space. Absence of invariant finite measures only addresses order-zero distributions and does not prove the stated smooth density.}
\GSfield{Reconstruction}{The topology is read literally from the source's \(C^\infty\), not silently weakened to the uniform norm. The continuous theorem proved in this companion implies uniform density of smooth coboundary sums by smooth approximation of continuous solutions, but it does not establish density in every derivative seminorm.}
\GSfield{Fate}{No exact independently verified reference or proof for the notebook's smooth-density assertion was obtained. The invariant-functional terminology in \citet[abstract]{abouqatebelkacimi2000invariant} is thematic only. The development identifies the missing distinction and reduces exact smooth solvability to surjectivity of the averaged difference operator \(\tfrac12(I+C_S)(I-C_T)\) onto the \(S\)-invariant smooth functions.}
\GSfield{Present status}{As of November 2026, smooth density is retained as an assertion in the notes, not promoted to a proved result of this companion. Exact \(C^0\) solvability and the elementary averaging reduction are proved; the higher-order distribution and regularity questions remain outside those proofs.}
\GSfield{Assessment}{literal statement: not determinable; publication match: none found; verification: unresolved lead; development: developed in this chapter; confidence: medium; problem \hyperref[prob:GP-1981-23]{GP-1981-23}}
\GSfield{Worth developing}{Worth checking before attempting a smooth inverse: an explicit invariant-distribution classification would test the notebook's density claim and identify any remaining obstruction accurately.}
\end{GSitem}
\begin{GSitem}{1981-555.3}{F}{The connection with foliations having minimal leaves}
\GSfield{Kind; source}{remark; \GSpage{P001013}}
\GSfield{The notes}{The page connects the problem with foliations on a Riemannian manifold whose leaves are minimal submanifolds, without giving a construction or a theorem relating the two.}
\GSfield{Statement}{A suspension of a group action is one standard route from action coinvariants to transverse foliation data. The notebook supplies neither a suspension manifold, a metric-existence criterion, nor an equivalence between its equation and a stated minimal-leaf theorem. No such exact bridge can therefore be recovered from this sentence.}
\GSfield{Reconstruction}{The suspension interpretation is a motivation only. The Russian explicitly mentions a Riemannian manifold and minimal surfaces, so minimality is geometric here, not the dynamical assertion that every leaf is dense. That distinction is determined by the source, rather than left as an English-word ambiguity.}
\GSfield{Fate}{The title of \citet{haefliger1980minimal} is a natural pre-talk lead, and \citet{abouqatebelkacimi2000invariant} provides later related invariant-functional terminology. The first original was inaccessible and the second has different hypotheses. Neither was verified as an exact publication of this modular equation.}
\GSfield{Present status}{As of November 2026, the source's subject link is recorded, while its precise mathematical bridge remains unidentified. No additional problem-index entry is created from it.}
\GSfield{Assessment}{literal statement: not determinable; publication match: thematic only; verification: unresolved lead; development: annotated; confidence: low}
\end{GSitem}

\GStalk{556}{1982-03-15}{A. A. Kirillov}{P001013-P001018}{Orbit geometry, multiplicities, and student research topics}
The actual heading and mathematical discussion occupy P001014--P001016; the wider catalogue range also includes Haefliger and Givental. Kirillov proposes topics on multiplicities, illustrates the geometry of an \(SU(3)\) orbit and its diagonal projection, and relates a recent Guillemin--Sternberg preprint to an index on a reduced space. Gelfand adds questions about duality, difference calculus and hypergeometric functions. The final remarks connect torus convexity with Newton polyhedra. The missing quotient in the multiplicity formula is mathematically consequential.

\begin{GSitem}{1981-556.1}{B}{Geometric computation of compact-group weight multiplicities}
\GSfield{Kind; source}{problem; \GSpage{P001014}}
\GSfield{The notes}{Weight multiplicities and restriction to subgroups are proposed as student topics; Guillemin--Sternberg are said to give an answer in principle.}
\GSfield{Statement}{For a compact connected Lie group \(G\), an integral dominant weight \(\lambda\), and a maximal torus \(T\), determine \(m_\lambda(\nu)=\dim\operatorname{Hom}_T(\mathbb C_\nu,V_\lambda)\), and interpret it as a quantization of the symplectic reduction of the corresponding prequantized orbit.}
\GSfield{Reconstruction}{The choice of a highest weight, prequantum bundle and reduction is our precise version of a topic, not a claim that the notebook states the general singular theorem.}
\GSfield{Fate}{The notebook supplies circulation evidence for a preprint by March 1982. \citet{guilleminsternberg1982multiplicities} was published in October 1982; the Kähler result and the later general symplectic/singular extensions must be distinguished. \citet[Introduction and \S2.4]{meinrenkensjamaar1999singular} give the latter, with a July 1997 preprint preceding the 1999 publication.}
\GSfield{Present status}{As of 11 November 2026: The compact prequantizable Hamiltonian case has an index-theoretic solution, including singular reductions in the specified Riemann--Roch sense. No new open problem is inferred from this student topic.}
\GSfield{Assessment}{literal statement: correct with missing hypotheses; publication match: partial; verification: read the result; development: annotated; confidence: high; problem \hyperref[prob:GP-1981-41]{GP-1981-41}}
\end{GSitem}
\begin{GSitem}{1981-556.2}{F}{The unspecified hyperbolic multiplicity programme}
\GSfield{Kind; source}{remark; \GSpage{P001014}}
\GSfield{The notes}{Hyperbolic algebras are described by a Cartan matrix having a negative eigenvalue; multiplicities and Macdonald are mentioned.}
\GSfield{Statement}{A symmetrizable generalized Cartan matrix is hyperbolic only if it is indecomposable of indefinite type and every proper connected principal submatrix is of finite or affine type. A negative eigenvalue alone does not impose these conditions. No matrix, module, weight, or proposed multiplicity formula is specified.}
\GSfield{Fate}{This is not enough to identify a single problem or match a later theorem. The general root-multiplicity discussion earlier this year and the compact-group problem are distinct. Search by the phrase and the cited surname yielded no uniquely recoverable target.}
\GSfield{Present status}{As of 11 November 2026: The programme is identifiable by subject only; an exact fate cannot be assigned. No claim that hyperbolic multiplicities are generally uncomputed is made.}
\GSfield{Assessment}{literal statement: not determinable; publication match: none found; verification: unresolved lead; development: annotated; confidence: low}
\end{GSitem}
\begin{GSitem}{1981-556.3}{E}{The regular SU(3) orbit and its hexagonal diagonal image}
\GSfield{Kind; source}{construction; \GSpage{P001014-P001015}}
\GSfield{The notes}{The skew-Hermitian trace-zero matrix \(X\) has fixed \(\operatorname{tr}X^2=c_2\), \(\operatorname{tr}X^3=c_3\). Its orbit is six-dimensional and its diagonal projection is drawn as a hexagon.}
\GSfield{Statement}{For three distinct real eigenvalues \(\lambda_1>\lambda_2>\lambda_3\), with sum zero, the orbit of \(i\operatorname{diag}(\lambda_1,\lambda_2,\lambda_3)\) is \(SU(3)/T^2\), of real dimension six. Its diagonal image is the convex hull of the six permutations of \(\lambda\). For repeated eigenvalues the orbit dimension and polygon degenerate.}
\GSfield{Reconstruction}{Regularity and the ordered real eigenvalues make the sketch precise.}
\GSfield{Fate}{\citet[Introduction and Theorem 1]{atiyah1982convexity} explicitly relate the general torus theorem to the older Schur--Horn and coadjoint-orbit cases. The notebook attributes the orbit picture to \citet{berezingelfand1956spherical}; that attribution is retained as a partial historical match, not evidence that every later quantization assertion is in the 1956 paper.}
\GSfield{Present status}{As of 11 November 2026: the stated qualified result is established; the literal notebook shorthand should be read with the hypotheses above.}
\GSfield{Assessment}{literal statement: correct with missing hypotheses; publication match: partial; verification: read the result; development: annotated; confidence: high}
\end{GSitem}
\begin{GSitem}{1981-556.4}{E}{Dimensions and singular values of the torus moment map}
\GSfield{Kind; source}{claim; \GSpage{P001015}}
\GSfield{The notes}{A regular fibre of the diagonal projection \(p\) has dimension four; the diagram marks internal walls where the fibre becomes singular.}
\GSfield{Statement}{For the Hamiltonian action of the effective two-torus on a regular six-dimensional orbit, a regular level \(\mu^{-1}(\nu)\) has dimension four and its quotient by a locally free torus action has dimension two. Values with positive-dimensional isotropy are critical, including the boundary of the moment polygon; an interior-wall drawing is not a complete classification by itself.}
\GSfield{Reconstruction}{The quotient and the qualification about the diagram are editorial additions.}
\GSfield{Fate}{This is the regular-value theorem together with Hamiltonian reduction. \citet[\S2.1]{meinrenkensjamaar1999singular} distinguish the level set from its quotient and explain the stratification when regularity fails.}
\GSfield{Present status}{As of 11 November 2026: the stated qualified result is established; the literal notebook shorthand should be read with the hypotheses above.}
\GSfield{Assessment}{literal statement: correct with missing hypotheses; publication match: partial; verification: read the result; development: annotated; confidence: high}
\end{GSitem}
\begin{GSitem}{1981-556.5}{E}{SU(3) weight hexagons and multiplicity layers}
\GSfield{Kind; source}{claim; \GSpage{P001015}}
\GSfield{The notes}{Irreducible representations are indexed by two integers; a hexagonal weight diagram has increasing multiplicities towards a central triangle.}
\GSfield{Statement}{For \(SU(3)\), dominant integral highest weights are \(a\omega_1+b\omega_2\), \(a,b\in\mathbb Z_{\ge0}\). Weight multiplicities are Weyl-invariant, have the familiar successive-layer pattern, and are bounded by \(\min(a,b)+1\). Only weights in the appropriate root-lattice coset occur; a drawn polygon does not make every lattice point a weight.}
\GSfield{Fate}{The compact semisimple character formula in \citet{kirillov1968characters} provides the established representation-theoretic setting. The diagram is expository and does not constitute a new multiplicity theorem.}
\GSfield{Present status}{As of 11 November 2026: The highest-weight classification and the type A2 multiplicity pattern are classical. The sketch is not a specification of a new formula.}
\GSfield{Assessment}{literal statement: correct with missing hypotheses; publication match: partial; verification: publisher metadata checked; development: annotated; confidence: medium}
\end{GSitem}
\begin{GSitem}{1981-556.6}{B}{Riemann–Roch multiplicity requires a reduced space}
\GSfield{Kind; source}{claim; \GSpage{P001015}}
\GSfield{The notes}{The notes write \(X=p^{-1}(y_0)\) and \(\langle e^B\operatorname{Td},X\rangle=\text{multiplicity}\), for large parameters.}
\GSfield{Statement}{At a regular integral weight \(\nu\), with a free effective torus action and the descended prequantum line bundle \(L_\nu\), the relevant index is \(\int_{M_\nu}\operatorname{ch}(L_\nu)\operatorname{Td}(TM_\nu)\), where \(M_\nu=\mu^{-1}(\nu)/T\). Appropriate positivity identifies this index with the relevant holomorphic multiplicity. Orbifold stabilizers require an orbifold index, and singular levels require a separately defined singular Riemann--Roch number.}
\GSfield{Reconstruction}{We supply \(M_\nu\), the descended line bundle, and the index interpretation. The notebook does not justify using an ordinary Todd integral on a singular quotient.}
\GSfield{Fate}{The 1982 Guillemin--Sternberg result and \citet[\S2.4]{meinrenkensjamaar1999singular} support the reduced-space interpretation. The scanned formula genuinely omits the quotient, so this is a faithful mathematical error rather than a translation slip.}
\GSfield{Present status}{As of 11 November 2026: the stated qualified result is established; the literal notebook shorthand should be read with the hypotheses above.}
\GSfield{Assessment}{literal statement: false as written; publication match: partial; verification: read the result; development: annotated; confidence: high; problem \hyperref[prob:GP-1981-41]{GP-1981-41}}
\end{GSitem}
\begin{GSitem}{1981-556.7}{F}{An unspecified duality between groups and representations}
\GSfield{Kind; source}{question; \GSpage{P001015}; I. M. Gelfand}
\GSfield{The notes}{Gelfand calls the duality between groups and their representations the main question.}
\GSfield{Statement}{No category of groups, dual object, topology, or universal property is specified. Pontryagin, Tannakian and orbit-method dualities are different possible interpretations.}
\GSfield{Fate}{No exact theorem can be selected from this sentence without importing a programme into the record.}
\GSfield{Present status}{As of 11 November 2026: The question remains a historical programme statement, not an indexed mathematical conjecture.}
\GSfield{Assessment}{literal statement: not determinable; publication match: none found; verification: unresolved lead; development: annotated; confidence: low}
\end{GSitem}
\begin{GSitem}{1981-556.8}{E}{Falling factorials as a basis for finite differences}
\GSfield{Kind; source}{method; \GSpage{P001015}; I. M. Gelfand}
\GSfield{The notes}{Gelfand contrasts \(d(x^n)/dx=nx^{n-1}\) with \(\Delta x^{\underline n}=n x^{\underline{n-1}}\).}
\GSfield{Statement}{For \(\Delta f(x)=f(x+1)-f(x)\) and \(x^{\underline n}=x(x-1)\cdots(x-n+1)\), \(\Delta x^{\underline n}=n x^{\underline{n-1}}\) for \(n\ge1\); equivalently \(\Delta\binom{x}{n}=\binom{x}{n-1}\).}
\GSfield{Fate}{This follows immediately by factoring the two falling factorials. It is the classical Newton-series basis underlying \citet{mahler1958interpolation}, not a theorem originating in this talk.}
\GSfield{Present status}{As of 11 November 2026: the stated qualified result is established; the literal notebook shorthand should be read with the hypotheses above.}
\GSfield{Assessment}{literal statement: correct; publication match: exact; verification: publisher metadata checked; development: annotated; confidence: high}
\end{GSitem}
\begin{GSitem}{1981-556.9}{E}{Orbit volume and representation dimension need a normalization and a shift}
\GSfield{Kind; source}{claim; \GSpage{P001016}}
\GSfield{The notes}{The invariant orbit measure is unique up to scale and can be normalized so that its volume is the representation dimension.}
\GSfield{Statement}{On a compact homogeneous orbit the invariant measure is unique up to a positive scalar. Rescaling it to any prescribed positive total volume is tautologically possible. With the standard Kirillov--Kostant Liouville normalization, the dimension formula instead uses the orbit through \(\lambda+\rho\), not generally the unshifted orbit through \(\lambda\).}
\GSfield{Fate}{The character formula of \citet{kirillov1968characters} is the relevant established source. The note is correct as a freedom to choose scale, but cannot by itself identify that scale with the canonical symplectic normalization.}
\GSfield{Present status}{As of 11 November 2026: the stated qualified result is established; the literal notebook shorthand should be read with the hypotheses above.}
\GSfield{Assessment}{literal statement: correct with missing hypotheses; publication match: partial; verification: publisher metadata checked; development: annotated; confidence: medium}
\end{GSitem}
\begin{GSitem}{1981-556.10}{E}{Fractional differentiation and the Euler hypergeometric integral}
\GSfield{Kind; source}{method; \GSpage{P001016}; I. M. Gelfand}
\GSfield{The notes}{The displayed integral for the \(p\)-th derivative of \(x^k(1-x)^\ell\) has kernel \((x-t)^{-p-1}\) and no gamma prefactor.}
\GSfield{Statement}{Initially for \(0<x<1\), \(\operatorname{Re}k>-1\) and \(\operatorname{Re}p<0\), define \[D^p f(x)=\frac1{\Gamma(-p)}\int_0^x(x-t)^{-p-1}f(t)\,dt.\] Then \[D^p[x^k(1-x)^\ell]=\frac{\Gamma(k+1)}{\Gamma(k-p+1)}x^{k-p}{}_2F_1(-\ell,k+1;k-p+1;x),\] with consistent branches and meromorphic continuation in the parameters. In particular \(D^p x^k=\Gamma(k+1)x^{k-p}/\Gamma(k-p+1)\).}
\GSfield{Reconstruction}{The branch, convergence domain and normalized fractional operator are supplied here.}
\GSfield{Fate}{The substitution \(t=xu\) reduces this directly to the classical Euler integral; compare \citet[Eq.~15.6.1]{nist2026hypergeometric}. The missing \(1/\Gamma(-p)\) is visible in the scan. Ordinary derivatives at nonnegative integral \(p\) require continuation or the corresponding fractional-calculus definition, not a divergent literal integral.}
\GSfield{Present status}{As of 11 November 2026: the stated qualified result is established; the literal notebook shorthand should be read with the hypotheses above.}
\GSfield{Assessment}{literal statement: false as written; publication match: exact; verification: read the result; development: annotated; confidence: high}
\end{GSitem}
\begin{GSitem}{1981-556.11}{F}{The programme of multivariable hypergeometric functions}
\GSfield{Kind; source}{remark; \GSpage{P001016}; I. M. Gelfand}
\GSfield{The notes}{There is said to be no satisfactory theory of hypergeometric functions of several variables.}
\GSfield{Statement}{This is an assessment of the state of a subject, without a specified class of differential systems or a criterion of completeness. A precise later framework is the \(A\)-hypergeometric system attached to an integer matrix and a complex parameter.}
\GSfield{Fate}{\citet{gelfandzelevinskykapranov1989hypergeometric} developed a major framework, preceded by the 1987--1988 announcements listed in that paper. This is a thematic successor, not proof that every intended multivariable question was solved.}
\GSfield{Present status}{As of 11 November 2026: Several precise multivariable theories are established; the unrestricted evaluative sentence has no single true-or-false resolution.}
\GSfield{Assessment}{literal statement: not determinable; publication match: thematic only; verification: publisher metadata checked; development: annotated; confidence: medium}
\end{GSitem}
\begin{GSitem}{1981-556.12}{E}{Atiyah convexity and the torus approach to representations}
\GSfield{Kind; source}{theorem; \GSpage{P001016}}
\GSfield{The notes}{Atiyah is cited for the convexity of the momentum image of a torus action; Khovanskii and Kazarnovskii are mentioned in relation to Newton polyhedra.}
\GSfield{Statement}{For a Hamiltonian torus acting on a compact connected symplectic manifold, the moment-map image is a convex polytope, the convex hull of the images of the fixed-point components, and its fibres are connected. Compactness, connectedness and a genuine Hamiltonian moment map are essential.}
\GSfield{Reconstruction}{The precise theorem is separated from the unformulated application programme.}
\GSfield{Fate}{\citet[Theorems 1--2]{atiyah1982convexity}, received 27 March 1981 and published in 1982, gives the result; the previous-year item 1980-508.15 already records its circulation. The names of Khovanskii and Kazarnovskii identify an adjacent Newton-polyhedron programme, but no specific enumerative formula is stated here.}
\GSfield{Present status}{As of 11 November 2026: the stated qualified result is established; the literal notebook shorthand should be read with the hypotheses above.}
\GSfield{Assessment}{literal statement: correct with missing hypotheses; publication match: exact; verification: read the result; development: annotated; confidence: high}
\end{GSitem}

\GStalk{557}{1982-03-29}{A. B. Givental}{P001016-P001020}{Relative symplectic normal forms and Lagrangian planes}
The manuscript heading is on P001017 and the talk ends on P001018. A relative Darboux theorem leads to local symplectic extension, the symplectic geometry of rays and obstacle problems, and the invariant alternating form on an odd symmetric power of the standard \(SL_2\)-module. An intervention attributed to Goldsheid asks for uniform approximation of a Lagrangian plane by a product of lines in fixed coordinate symplectic planes. That last question has much stronger quantifiers than the coordinatewise observation following it.

\begin{GSitem}{1981-557.1}{A}{Relative Darboux theorem for a smooth submanifold}
\GSfield{Kind; source}{theorem; \GSpage{P001017}}
\GSfield{The notes}{If \(\omega_0|_{TN}=\omega_1|_{TN}\), then near each point of \(N\) a diffeomorphism fixing \(N\) pulls \(\omega_1\) back to \(\omega_0\).}
\GSfield{Statement}{Let \(N\) be a smooth embedded submanifold of a smooth manifold and let \(\omega_0,\omega_1\) be symplectic forms near \(x\in N\), with equal pullbacks to \(N\). There are neighbourhood germs and a diffeomorphism germ \(g\), identity on \(N\), such that \(g^*\omega_1=\omega_0\). No constant-rank assumption on the restricted form is needed for this local statement.}
\GSfield{Reconstruction}{We interpret the displayed U-to-U map as an equivalence of neighbourhood germs; one fixed invariant open set is unnecessary.}
\GSfield{Fate}{\citet[Chapter 2, \S1.3, Theorem I and proof, pp.~24--25]{arnoldgivental1990symplectic} state and prove this result. The Russian chapter appeared in 1985 and the English version in 1990. The March 1982 notebook is earlier circulation evidence, not proof that the 1985 chapter was the first announcement.}
\GSfield{Present status}{As of 11 November 2026: Established as the relative Darboux–Givental theorem; the cited proof covers precisely equality of pullbacks.}
\GSfield{Assessment}{literal statement: correct with missing hypotheses; publication match: exact; verification: read the result; development: annotated; confidence: high}
\end{GSitem}
\begin{GSitem}{1981-557.2}{E}{Moser’s deformation argument relative to a submanifold}
\GSfield{Kind; source}{method; \GSpage{P001017}}
\GSfield{The notes}{Choose \(g_t^*\omega_t=\omega_0\), differentiate, and solve an ordinary differential equation; a point gives Darboux.}
\GSfield{Statement}{First align the two bilinear forms at the base point by a linear map fixing \(T_xN\). On a sufficiently small neighbourhood, the interpolating forms are nondegenerate. Choose a primitive of \(\omega_1-\omega_0\) that vanishes as an ambient covector along \(N\); solve \(\iota_{V_t}\omega_t=-\alpha\), so \(V_t|_N=0\), and integrate. Equality of pullbacks alone does not say the unadjusted straight-line path is nondegenerate.}
\GSfield{Reconstruction}{The linear adjustment and vanishing of the primitive as an ambient covector are supplied proof details.}
\GSfield{Fate}{The full relative argument is given in \citet[Chapter 2, \S1.3]{arnoldgivental1990symplectic}. The homotopy method is older; the cited source explains the preliminary linear adjustment omitted in the notebook.}
\GSfield{Present status}{As of 11 November 2026: the stated qualified result is established; the literal notebook shorthand should be read with the hypotheses above.}
\GSfield{Assessment}{literal statement: correct with missing hypotheses; publication match: partial; verification: read the result; development: annotated; confidence: high}
\end{GSitem}
\begin{GSitem}{1981-557.3}{A}{When does a two-form extend to an ambient symplectic form?}
\GSfield{Kind; source}{question; \GSpage{P001017}}
\GSfield{The notes}{The question asks when a two-form is the restriction of a symplectic form; closedness is named. Gelfand asks first for the vector-space structure.}
\GSfield{Statement}{Let \(N\subset M\) be a submanifold germ with \(\dim M\) even and let \(\beta\) be a smooth two-form on \(N\). It extends locally near \(x\) to a symplectic form on \(M\) if and only if \(d\beta=0\) near \(x\) and \(\dim\ker\beta_x\le\operatorname{codim}_M N\).}
\GSfield{Reconstruction}{The ambient dimension and corank inequality are absent from the question and supplied from the later theorem.}
\GSfield{Fate}{\citet[Chapter 2, \S1.3, Extension Theorem, p.~25]{arnoldgivental1990symplectic} gives the exact criterion. The pointwise algebraic question omits closedness legitimately, whereas a smooth symplectic extension cannot.}
\GSfield{Present status}{As of 11 November 2026: The local smooth extension problem is solved with this criterion; it does not assert a global embedding theorem.}
\GSfield{Assessment}{literal statement: correct with missing hypotheses; publication match: exact; verification: read the result; development: annotated; confidence: high; problem \hyperref[prob:GP-1981-42]{GP-1981-42}}
\end{GSitem}
\begin{GSitem}{1981-557.4}{E}{The symplectic space of oriented rays}
\GSfield{Kind; source}{construction; \GSpage{P001017}}
\GSfield{The notes}{The space of straight lines in \(\mathbb R^3\) is symplectic and is identified in Arnold’s explanation with \(T^*\mathbb R^2\).}
\GSfield{Statement}{The space of oriented affine lines in \(\mathbb R^3\) is globally \(T^*S^2\): direction \(u\in S^2\) and closest point \(v\perp u\). Reduction of the unit-speed geodesic flow supplies its canonical symplectic form. A direction chart gives \(T^*\mathbb R^2\) only locally. Reversing the line orientation is anti-symplectic, so the unoriented quotient is not endowed with this same symplectic form.}
\GSfield{Reconstruction}{Orientation and the local meaning of the affine cotangent chart are necessary qualifications.}
\GSfield{Fate}{\citet[Example 1]{givental1983lagrangian} explicitly identifies the reduced space as \(T^*S^{n-1}\). The result is standard before the talk; the later paper verifies the exact geometric context.}
\GSfield{Present status}{As of 11 November 2026: the stated qualified result is established; the literal notebook shorthand should be read with the hypotheses above.}
\GSfield{Assessment}{literal statement: correct with missing hypotheses; publication match: exact; verification: read the result; development: annotated; confidence: high}
\end{GSitem}
\begin{GSitem}{1981-557.5}{A}{Extremals and the obstacle example}
\GSfield{Kind; source}{claim; \GSpage{P001017}}
\GSfield{The notes}{The space of extremals of any variational problem is said to be symplectic; a point and a sphere produce a two-dimensional family.}
\GSfield{Statement}{For a regular finite-dimensional first-order variational problem, the Legendre transform identifies local Cauchy data with a symplectic phase space. On a regular fixed-energy level one may reduce by the nonvanishing Hamiltonian flow. Families generated by a fixed initial point can then give Lagrangian subvarieties; obstacle tangencies can create singularities. Degenerate Lagrangians or singular quotients require a separate treatment.}
\GSfield{Reconstruction}{Regularity, reduction, and genericity are explicit modern hypotheses not supplied by the notebook.}
\GSfield{Fate}{\citet[Examples 1--2, Theorem and Corollary]{givental1983lagrangian}, received 19 July 1982 and published in 1983, connects generic obstacle singularities to open swallowtails. This supplies a concrete development of the example, not a proof of the unrestricted word “any”.}
\GSfield{Present status}{As of 11 November 2026: the stated qualified result is established; the literal notebook shorthand should be read with the hypotheses above.}
\GSfield{Assessment}{literal statement: false as written; publication match: partial; verification: read the result; development: annotated; confidence: high}
\end{GSitem}
\begin{GSitem}{1981-557.6}{E}{The symplectic odd binary-form representation}
\GSfield{Kind; source}{construction; \GSpage{P001017-P001018}}
\GSfield{The notes}{The space of binary forms of degree \(2n-1\) has Darboux coefficients \(q_j,p_j\) and an \(SL_2(\mathbb R)\)-invariant \(\omega=\sum dp_j\wedge dq_j\).}
\GSfield{Statement}{Write the dehomogenized polynomial as \[F(w)=\sum_{j=1}^n q_j\frac{w^{2n-j}}{(2n-j)!}+\sum_{j=1}^n(-1)^{n-j}p_j\frac{w^{j-1}}{(j-1)!}.\] The irreducible module \(\operatorname{Sym}^{2n-1}\mathbb R^2\) has an invariant nondegenerate alternating form, unique up to scalar; the displayed coefficient normalization makes it \(\sum dp_j\wedge dq_j\).}
\GSfield{Reconstruction}{The alternating signs in the second sum fix the overwritten sign on P001017. Different Hamiltonian/sign conventions in the publication are equivalent after a coordinate sign change.}
\GSfield{Fate}{\citet[Example 2]{givental1983lagrangian} publishes this representation and its use in singularity models. The representation-theoretic alternating form itself is classical; the application is developed in the subsequent paper.}
\GSfield{Present status}{As of 11 November 2026: the stated qualified result is established; the literal notebook shorthand should be read with the hypotheses above.}
\GSfield{Assessment}{literal statement: correct with missing hypotheses; publication match: exact; verification: read the result; development: annotated; confidence: high}
\end{GSitem}
\begin{GSitem}{1981-557.7}{A}{Translation Hamiltonian and reduction by q1}
\GSfield{Kind; source}{construction; \GSpage{P001018}}
\GSfield{The notes}{\(w\mapsto w+t\) is generated by \(h=q_n^2/2-q_{n-1}p_n-\cdots-q_1p_2\); \(q_1\) is a first integral and \(q_1=1,p_1=1\) are imposed.}
\GSfield{Statement}{With \(\omega=\sum dp_j\wedge dq_j\) and \(\iota_{X_h}\omega=dh\), the translation flow has \(\dot q_1=0\), \(\dot q_j=q_{j-1}\) for \(j>1\), \(\dot p_n=q_n\), and \(\dot p_j=-p_{j+1}\) for \(j<n\). The hypersurface \(q_1=1\) has characteristic direction \(\partial_{p_1}\); its quotient is symplectic of dimension \(2n-2\). Fixing \(p_1=1\) chooses a slice, rather than a second conserved quantity.}
\GSfield{Reconstruction}{The Hamiltonian sign convention and the quotient meaning of the slice are supplied.}
\GSfield{Fate}{\citet[Example 2]{givental1983lagrangian} gives the corresponding Hamiltonian representation and reduction. Direct coefficient differentiation of \(F(w+t)\) checks the notebook signs under the stated convention.}
\GSfield{Present status}{As of 11 November 2026: The construction is established; no claim that p1 is preserved by the translation Hamiltonian is justified.}
\GSfield{Assessment}{literal statement: correct with missing hypotheses; publication match: exact; verification: read the result; development: annotated; confidence: high}
\end{GSitem}
\begin{GSitem}{1981-557.8}{C}{Goldsheid’s uniform approximation by coordinate-product Lagrangians}
\GSfield{Kind; source}{question; \GSpage{P001018}; I. Ya. Goldsheid}
\GSfield{The notes}{Goldsheid asks whether every Lagrangian \(L\subset\mathbb R^{2n}\) admits \(L_1=\operatorname{span}(\alpha_je_j+\beta_jf_j)\) with \(\|P_L-P_{L_1}\|\le1/\sqrt2\); even any constant less than one independent of \(n\) would suffice.}
\GSfield{Statement}{Use the standard Euclidean and symplectic structures, with \(Je_j=f_j\). Must every Lagrangian \(n\)-plane admit one nonzero vector in each prescribed coordinate plane \(\operatorname{span}(e_j,f_j)\), whose direct sum \(K\) satisfies \(\|P_L-P_K\|_{2\to2}\le2^{-1/2}\)? The weaker question asks for a universal \(c<1\) for all \(n\).}
\GSfield{Fate}{The scan unambiguously gives an operator norm and a dimension-independent requirement. Searches under Goldsheid/Goldsheyd/\ru{Гольдшейд}, Lagrangian projections, symmetric-unitary diagonal approximation, Math-Net author records, arXiv and MathOverflow located no exact resolution. \citet[Theorem 3.4]{mehrmannpoloni2012doubling} bounds individual entries of a permuted Lagrangian graph basis; that does not supply the requested uniform operator norm. The development below proves a Fourier-matrix counterexample to the sharp constant at \(n=6\), a universal lower obstruction \(\cos(\pi/8)\), and a dimension-dependent upper bound.}
\GSfield{Present status}{As of 11 November 2026: We disprove the proposed \(1/\sqrt2\) constant in dimension \(n=6\). More generally, any universal constant must be at least \(\cos(\pi/8)\). We establish \(\cos(\pi/(2n))<1\) for each fixed \(n\), but do not resolve the existence of a universal \(c\in[\cos(\pi/8),1)\). No exact later resolution of that weaker question was located in this bounded search.}
\GSfield{Assessment}{literal statement: correct; publication match: partial; verification: read the result; development: developed in this chapter; confidence: medium; problem \hyperref[prob:GP-1981-40]{GP-1981-40}}
\GSfield{Worth developing}{Worth developing because the exact finite-dimensional optimization is explicit and admits certified partial bounds while preserving the unresolved uniform quantifier.}
\end{GSitem}
\begin{GSitem}{1981-557.9}{E}{Coordinatewise closeness does not control the gap of the spans}
\GSfield{Kind; source}{claim; \GSpage{P001018}}
\GSfield{The notes}{For every \(j\), one can choose \(\varphi_j=\alpha_je_j+\beta_jf_j\) with \(\sin\angle(\varphi_j,L)\le1/\sqrt2\).}
\GSfield{Statement}{The compression of \(P_L\) to each coordinate \(J\)-plane has trace one, so its largest eigenvalue is at least \(1/2\), proving the coordinatewise claim. This supplies no bound for all linear combinations of the chosen vectors. Indeed for \(n\ge2\), \(v=(1,\ldots,1)/\sqrt n\) and \(L=\mathbb Rv\oplus i v^\perp\) give \(\operatorname{dist}(f_j,L)=1/\sqrt n\), yet \(\|P_L-P_{i\mathbb R^n}\|=1\).}
\GSfield{Reconstruction}{The trace proof and the example separating the quantifiers are ours.}
\GSfield{Fate}{The observation in the notebook is correct. The explicit counterexample to the inference about the span is proved and independently checked in our development; it is not a counterexample to Goldsheid’s existential question.}
\GSfield{Present status}{As of 11 November 2026: The elementary coordinatewise statement is verified; its use as a proof of uniform subspace approximation would be invalid.}
\GSfield{Assessment}{literal statement: correct; publication match: none found; verification: read the result; development: annotated; confidence: high; problem \hyperref[prob:GP-1981-40]{GP-1981-40}}
\end{GSitem}

\GStalk{558}{1982-04-05}{Mariusz Wodzicki}{P001018-P001023}{Local index densities and spectral boundary conditions}
The actual notes run from the heading on P001019 through P001021. Wodzicki reviews spectral zeta functions, local heat coefficients and the signature theorem, then explains the nonlocal boundary correction and spectral boundary conditions. Gelfand asks about additive invariants. The notebook mixes a closed-manifold formula with a boundary hypothesis and sketches a normal, rather than necessarily self-adjoint, tangential operator; these scopes are kept separate.

\begin{GSitem}{1981-558.1}{E}{The index as an integral of a local density}
\GSfield{Kind; source}{theorem; \GSpage{P001019}}
\GSfield{The notes}{For an elliptic differential system \(D\), possibly on a manifold with boundary, \(\operatorname{ind}D=\int_M\alpha_0\) is attributed to Atiyah--Singer.}
\GSfield{Statement}{For an elliptic differential operator \(D:C^\infty(E)\to C^\infty(F)\) on a closed compact manifold, the Fredholm index equals the integral of the constant coefficient in the local heat supertrace. On a manifold with boundary one must first specify an elliptic boundary realization; its formula generally includes boundary terms.}
\GSfield{Reconstruction}{We allow different source and target bundles and separate the closed case from boundary realizations.}
\GSfield{Fate}{This is an established index/heat-kernel theorem; \citet[Chapter 1, \S1.7]{gilkey1984invariance} gives the analytic formulation, while \citet[Theorem (3.10)]{atiyahpatodisinger1975spectral1} exhibit the boundary correction. A bare differential expression on a manifold with boundary does not have the asserted Fredholm index.}
\GSfield{Present status}{As of 11 November 2026: the stated qualified result is established; the literal notebook shorthand should be read with the hypotheses above.}
\GSfield{Assessment}{literal statement: false as written; publication match: exact; verification: read the result; development: annotated; confidence: high}
\end{GSitem}
\begin{GSitem}{1981-558.2}{E}{Meromorphic continuation of a positive elliptic zeta function}
\GSfield{Kind; source}{theorem; \GSpage{P001019}}
\GSfield{The notes}{For \(A=D^*D\) or \(DD^*\), poles are listed at \((n-j)/m\), with none at \(0,-1,-2,\ldots\).}
\GSfield{Statement}{Let \(A\) be a positive self-adjoint elliptic differential operator of order \(m>0\) on a closed \(n\)-manifold, with zero eigenvalues omitted. Then \(\zeta_A(s)=\sum_{\lambda>0}\lambda^{-s}\) converges for \(\operatorname{Re}s>n/m\), extends meromorphically with at most simple poles at \((n-j)/m\), and is regular at nonpositive integers. Some candidate poles vanish. General classical pseudodifferential operators need a more qualified statement.}
\GSfield{Reconstruction}{Closedness, differential order, positivity, omission of the kernel and ``possible poles are essential additions.}
\GSfield{Fate}{This is classical complex-power/heat asymptotics: \citet{seeley1968powers} is the 1968 Russian translation of earlier work, and \citet[\S1.10]{gilkey1984invariance} supplies the heat formulation. It continues the expository item 1972-296.1.}
\GSfield{Present status}{As of 11 November 2026: the stated qualified result is established; the literal notebook shorthand should be read with the hypotheses above.}
\GSfield{Assessment}{literal statement: correct with missing hypotheses; publication match: exact; verification: read the result; development: annotated; confidence: high}
\end{GSitem}
\begin{GSitem}{1981-558.3}{E}{The shifted-zeta formula for the index}
\GSfield{Kind; source}{claim; \GSpage{P001019}}
\GSfield{The notes}{\(\operatorname{ind}D=\zeta_{1+D^*D}(0)-\zeta_{1+DD^*}(0)\).}
\GSfield{Statement}{For a closed elliptic Fredholm realization, the positive eigenvalues of \(D^*D\) and \(DD^*\) agree with multiplicity. Therefore for every \(\varepsilon>0\), \[\zeta_{\varepsilon+D^*D}(s)-\zeta_{\varepsilon+DD^*}(s)=\varepsilon^{-s}\operatorname{ind}D\] first in a convergent half-plane and then by continuation; evaluation at zero gives the notebook formula.}
\GSfield{Reconstruction}{The displayed all-s identity is our explanatory derivation, not a second formula claimed to be in the notes.}
\GSfield{Fate}{The equality is an elementary spectral-pairing consequence of the heat-index method, as in \citet[\S1.7]{gilkey1984invariance}. It avoids the missing-kernel problem of unshifted zeta functions.}
\GSfield{Present status}{As of 11 November 2026: the stated qualified result is established; the literal notebook shorthand should be read with the hypotheses above.}
\GSfield{Assessment}{literal statement: correct with missing hypotheses; publication match: exact; verification: read the result; development: annotated; confidence: high}
\end{GSitem}
\begin{GSitem}{1981-558.4}{E}{Local zeta values from a parameter-dependent parametrix}
\GSfield{Kind; source}{method; \GSpage{P001020}}
\GSfield{The notes}{Equation (5) integrates \(b^\circ_{-m-n}(x,\xi,\lambda)\) over \(|\xi|=1\) and a ray, with a minus sign and factor \(1/m\); it is a density.}
\GSfield{Statement}{For an invertible positive elliptic differential \(A\), if \(\operatorname{tr}K_A(t;x,x)\sim\sum_j a_j(x)t^{(j-n)/m}\), then \(\zeta_A(0;x)=a_n(x)\). If zero modes are omitted, subtract the kernel-projection density. The coefficient is computed from the homogeneous degree \(-m-n\) resolvent symbol, using an admissible ray and a fixed convention for \((A-\lambda)^{-1}\), Fourier factors and contour orientation.}
\GSfield{Reconstruction}{We supply a convention-independent heat statement and preserve the resolvent formula as a separate, convention-sensitive notebook assertion.}
\GSfield{Fate}{The heat version is verified in \citet[\S1.10]{gilkey1984invariance}. The parametrix expression belongs to the Seeley/Wodzicki calculus; the notebook alone does not fix enough conventions to certify a sign error in Seeley. The stronger assertion of cut-independent pointwise density is not valid for arbitrary pseudodifferential operators; compare \citet[Remark 4.4]{ponge2006spectral}.}
\GSfield{Present status}{As of 11 November 2026: the stated qualified result is established; the literal notebook shorthand should be read with the hypotheses above.}
\GSfield{Assessment}{literal statement: correct with missing hypotheses; publication match: partial; verification: read the result; development: annotated; confidence: medium}
\end{GSitem}
\begin{GSitem}{1981-558.5}{E}{The signature operator and its chirality splitting}
\GSfield{Kind; source}{construction; \GSpage{P001020}}
\GSfield{The notes}{The notes introduce \(D=(d+d^*)|_{\Omega_+}\), with a phase times the Hodge star, and identify its index with the signature.}
\GSfield{Statement}{On a closed oriented Riemannian manifold of dimension \(4k\), the Hodge-star phase defines the standard chirality involution on complex differential forms; \(d+d^*\) interchanges its eigensubbundles. The index of its positive-to-negative part, with the compatible chirality convention, is the signature of the middle-dimensional intersection pairing.}
\GSfield{Reconstruction}{The notebook phase uses an insufficiently specified n; we do not silently replace it by a different exponent. The modern statement specifies the standard chirality by its operator properties.}
\GSfield{Fate}{This is the classical signature operator. \citet[Chapter 3, \S3.1]{gilkey1984invariance} fixes the chirality convention; \citet[\S4]{atiyahpatodisinger1975spectral1} explain the boundary modification. The preceding scanned hypothesis \(\partial M\ne\varnothing\) cannot accompany the uncorrected integral equality.}
\GSfield{Present status}{As of 11 November 2026: the stated qualified result is established; the literal notebook shorthand should be read with the hypotheses above.}
\GSfield{Assessment}{literal statement: correct with missing hypotheses; publication match: exact; verification: read the result; development: annotated; confidence: high}
\end{GSitem}
\begin{GSitem}{1981-558.6}{E}{The div–grad–curl example and its missing sign}
\GSfield{Kind; source}{claim; \GSpage{P001020}; I. M. Gelfand}
\GSfield{The notes}{Gelfand writes \(\bigl(\begin{smallmatrix}0&\operatorname{div}\\\operatorname{grad}&\operatorname{rot}\end{smallmatrix}\bigr)^2=\Delta I\), with a parenthetical \(\mathbb R^4\).}
\GSfield{Statement}{Using the usual Euclidean operators on \(\mathbb R^3\), the literal square is \(\operatorname{diag}(\Delta,2\operatorname{grad}\operatorname{div}-\Delta I_3)\). A self-adjoint Dirac-type version is \[B=\begin{pmatrix}0&-\operatorname{div}\\\operatorname{grad}&\operatorname{curl}\end{pmatrix},\qquad B^2=-\Delta I_4.\] The four refers to a scalar plus three vector components, not a four-dimensional base.}
\GSfield{Reconstruction}{We explicitly fix the conventional Euclidean signs instead of treating the shorthand as a literal self-adjoint matrix identity.}
\GSfield{Fate}{\citet[Introduction, p.~45]{atiyahpatodisinger1975spectral1} discuss the same operator only up to sign conventions. Our vector-calculus calculation and the repeated \(\mathbb R^3\) wording on P001024 settle the notebook ambiguity; the accompanying symbolic check is included.}
\GSfield{Present status}{As of 11 November 2026: the stated qualified result is established; the literal notebook shorthand should be read with the hypotheses above.}
\GSfield{Assessment}{literal statement: false as written; publication match: partial; verification: read the result; development: annotated; confidence: high}
\end{GSitem}
\begin{GSitem}{1981-558.7}{E}{Local signature and Todd densities}
\GSfield{Kind; source}{theorem; \GSpage{P001020}}
\GSfield{The notes}{Atiyah--Bott--Patodi identify \(\alpha_0\) with \(L_k(p)\), in dimension four \(p_1/3\); for \(\bar\partial\), a Kähler assumption is stated for the Todd density.}
\GSfield{Statement}{For the geometric signature operator the constant heat supertrace density is the top-degree Hirzebruch \(L\)-form. For the Dolbeault complex with the usual Kähler metric and compatible Hermitian connection, the local index density is the appropriate Todd form (and the Chern character for coefficients). Equality of integrated characteristic numbers does not make arbitrary pointwise representatives identical.}
\GSfield{Fate}{\citet{atiyahbottpatodi1973heat} is the 1973 publication named by the talk. \citet[Chapter 3, \S\S3.1 and 3.6--3.7]{gilkey1984invariance} provides the precise geometric versions. The notebook correctly warns against dropping the Kähler qualification in its stated pointwise formula.}
\GSfield{Present status}{As of 11 November 2026: the stated qualified result is established; the literal notebook shorthand should be read with the hypotheses above.}
\GSfield{Assessment}{literal statement: correct with missing hypotheses; publication match: exact; verification: read the result; development: annotated; confidence: high}
\end{GSitem}
\begin{GSitem}{1981-558.8}{E}{The signature theorem with its eta boundary correction}
\GSfield{Kind; source}{theorem; \GSpage{P001020}}
\GSfield{The notes}{\(\operatorname{sign}M=\int_M L_k(p)-\eta(\partial M)\), for a product metric near the boundary.}
\GSfield{Statement}{For compact oriented Riemannian \(M^{4k}\), product near \(Y=\partial M\), define the signature on the image of \(H^{2k}(M,Y;\mathbb R)\to H^{2k}(M;\mathbb R)\). With the APS odd-signature operator on even forms and its boundary orientation, \[\operatorname{sign}M=\int_M L-\eta(B_Y;0).\] This normalization is different from writing \((\eta+h)/2\) for a general Dirac boundary problem.}
\GSfield{Fate}{\citet[Theorem (4.14)]{atiyahpatodisinger1975spectral1} is the exact result. Its announcement appeared in 1973, the full paper was received 25 March 1974, and publication followed in 1975; the notebook’s ``1974 should not replace those distinct dates.}
\GSfield{Present status}{As of 11 November 2026: the stated qualified result is established; the literal notebook shorthand should be read with the hypotheses above.}
\GSfield{Assessment}{literal statement: correct with missing hypotheses; publication match: exact; verification: read the result; development: annotated; confidence: high}
\end{GSitem}
\begin{GSitem}{1981-558.9}{F}{Is signature the only additive invariant?}
\GSfield{Kind; source}{question; \GSpage{P001020}; I. M. Gelfand}
\GSfield{The notes}{Gelfand asks whether signature is the unique additive invariant of manifolds with boundary.}
\GSfield{Statement}{“Additive” has no specified gluing category or normalization here. Under gluing compact even-dimensional manifolds along complete closed odd-dimensional boundary components, Euler characteristic is also additive, since \(\chi(M\cup_YN)=\chi(M)+\chi(N)-\chi(Y)\) and \(\chi(Y)=0\). It is independent of signature: \(\chi(D^{4k})=1\), \(\sigma(D^{4k})=0\). If orientation-oddness or cobordism invariance is imposed, that is a different question.}
\GSfield{Reconstruction}{The Euler-characteristic example refutes only the plainly stated broad interpretation; no extra axiom is attributed to Gelfand.}
\GSfield{Fate}{\citet[Chapter 1, Corollary (1.4)]{karraskreckneumannossa1973cutting} already classify the corresponding real-valued cut-and-paste invariants of closed oriented manifolds by Euler characteristic and signature. This is useful context, not an exact match to an unstated boundary gluing law.}
\GSfield{Present status}{As of 11 November 2026: The broad uniqueness interpretation is false, while the intended restricted axioms cannot be recovered from this line.}
\GSfield{Assessment}{literal statement: not determinable; publication match: thematic only; verification: read the result; development: annotated; confidence: medium}
\end{GSitem}
\begin{GSitem}{1981-558.10}{E}{Filling independence and nonlocality of the signature defect}
\GSfield{Kind; source}{claim; \GSpage{P001020}}
\GSfield{The notes}{Doubling shows that \(\eta(\partial M)\) is independent of the filling; lens spaces show it is nonlocal because it fails multiplicativity under coverings.}
\GSfield{Statement}{For two product-metric fillings of the same oriented metric boundary, glue one to the orientation reverse of the other and apply the closed signature theorem and signature additivity. Their defects \(\int L-\sigma\) agree. A universal natural local boundary-density integral would multiply under finite local-isometry coverings; the signature eta invariant need not do so.}
\GSfield{Reconstruction}{Two possibly different fillings, common boundary metric and orientation reversal make the gluing argument precise.}
\GSfield{Fate}{\citet[Introduction, pp.~43--45]{atiyahpatodisinger1975spectral1} give exactly this motivation and the sphere/lens-space obstruction. Filling independence does not say eta is independent of the boundary metric.}
\GSfield{Present status}{As of 11 November 2026: the stated qualified result is established; the literal notebook shorthand should be read with the hypotheses above.}
\GSfield{Assessment}{literal statement: correct with missing hypotheses; publication match: exact; verification: read the result; development: annotated; confidence: high}
\end{GSitem}
\begin{GSitem}{1981-558.11}{E}{Spectral boundary projections and the normal-operator qualification}
\GSfield{Kind; source}{construction; \GSpage{P001021}}
\GSfield{The notes}{Near the boundary \(D=\partial_u+A\), with \(A\) normal and spectrum outside a cone near \(i\mathbb R\); \(P_A\) projects to \(\operatorname{Re}\lambda\ge0\). Local elliptic boundary conditions need not exist.}
\GSfield{Statement}{The classical APS construction assumes \(D=\sigma(\partial_u+A)\) in a product collar, where \(\sigma:E|_Y\to F|_Y\) is an isomorphism and \(A\) is self-adjoint elliptic on \(E|_Y\). Requiring \(P_{\ge0}(A)(f|_Y)=0\) gives an elliptic Fredholm realization. More general sectorial projections require symbol cuts and resolvent estimates; normality plus a bare spectral drawing is not a complete analytic hypothesis list.}
\GSfield{Reconstruction}{The modern statement deliberately separates the proved classical boundary theorem from the notebook’s larger normal-operator claim.}
\GSfield{Fate}{\citet[\S2 and Theorem (3.10)]{atiyahpatodisinger1975spectral1} verify the self-adjoint case. \citet[\S3 and Appendix A]{ponge2006spectral} provide the modern spectral-cut framework, not by themselves a proof of this entire normal-boundary index assertion.}
\GSfield{Present status}{As of 11 November 2026: the stated qualified result is established; the literal notebook shorthand should be read with the hypotheses above.}
\GSfield{Assessment}{literal statement: not determinable; publication match: partial; verification: read the result; development: annotated; confidence: medium}
\end{GSitem}
\begin{GSitem}{1981-558.12}{E}{APS index formula and the role of the zero modes}
\GSfield{Kind; source}{theorem; \GSpage{P001021}}
\GSfield{The notes}{\(\operatorname{ind}D=\int\alpha_0-(h_0+\eta(A))/2\); \(h_0\) counts purely imaginary eigenvalues, while the notes explicitly say APS used self-adjoint \(A\).}
\GSfield{Statement}{Under the classical self-adjoint product-collar assumptions, \[\begin{aligned}\operatorname{ind}D_{\mathrm{APS}}&=\int_M\alpha_0-\frac{\eta_A(0)+\dim\ker A}{2},\\ \eta_A(s)&=\sum_{\lambda\ne0}\operatorname{sign}(\lambda)|\lambda|^{-s}.\end{aligned}\] The choice to exclude nonnegative boundary eigenmodes determines the kernel term; changing the zero-mode boundary subspace changes the index.}
\GSfield{Reconstruction}{The precise theorem assumes self-adjointness. The manuscript’s distinction is retained.}
\GSfield{Fate}{\citet[Theorem (3.10)]{atiyahpatodisinger1975spectral1} proves this formula. The normal-operator expression using \(\operatorname{Re}\lambda\) is a proposed extension with only a partial publication match here. It is not silently substituted into the original APS theorem.}
\GSfield{Present status}{As of 11 November 2026: the stated qualified result is established; the literal notebook shorthand should be read with the hypotheses above.}
\GSfield{Assessment}{literal statement: not determinable; publication match: partial; verification: read the result; development: annotated; confidence: medium}
\end{GSitem}
\begin{GSitem}{1981-558.13}{E}{Variation of the signature eta invariant on a cylinder}
\GSfield{Kind; source}{claim; \GSpage{P001021}}
\GSfield{The notes}{\(\int_{N\times I}L_k(p)=\eta_{p_1}-\eta_{p_2}\) is written below a cylinder diagram.}
\GSfield{Statement}{For an oriented cylinder of dimension \(4k\), with product metrics at both ends, its signature is zero. Applying the signature boundary theorem yields that its \(L\)-integral equals the signed difference of the two boundary eta invariants, with signs determined by outward boundary orientation. For a general family of self-adjoint elliptic operators, an analogous formula requires the spectral-flow/kernel terms.}
\GSfield{Reconstruction}{The theorem is restricted to the geometric signature situation actually illustrated; the indices p or rho in the manuscript are not interpreted as additional data.}
\GSfield{Fate}{The signature case follows directly from \citet[Theorem (4.14)]{atiyahpatodisinger1975spectral1}. The source says only “elliptic operator” at this point; a general family formula cannot be inferred by omitting integer jumps.}
\GSfield{Present status}{As of 11 November 2026: the stated qualified result is established; the literal notebook shorthand should be read with the hypotheses above.}
\GSfield{Assessment}{literal statement: correct with missing hypotheses; publication match: partial; verification: read the result; development: annotated; confidence: high}
\end{GSitem}

\GStalk{559}{1982-04-12}{Mariusz Wodzicki}{P001021-P001025}{Eta regularity, spectral cuts, and local zeta invariants}
The actual heading is on P001022. The first part repeats the spectral boundary problem and discusses global versus local residues. Gelfand asks why self-adjointness and a filling should be needed. Wodzicki then compares zeta functions defined by two spectral cuts, relates their integer values to local residues, and writes a shifted-zeta formula with a missing power of the shift. The final p-adic analogy is not sufficiently specified to identify an automorphic theorem.

\begin{GSitem}{1981-559.1}{E}{The collar operator and boundary bundle identification}
\GSfield{Kind; source}{construction; \GSpage{P001022}}
\GSfield{The notes}{\(D:E\to F\) is written as \(\partial_u+A\), while \(A:E|_Y\to F|_Y\) is then assigned eigenspaces.}
\GSfield{Statement}{Factor the invertible normal symbol: \(D=\sigma(\partial_u+A)\). Then \(A\) is an operator on the single boundary bundle \(E|_Y\), after identifying the target by \(\sigma^{-1}\). Its nonnegative spectral projection can act on \(f|_Y\). Without this identification, an operator between distinct bundles has no intrinsic spectrum.}
\GSfield{Reconstruction}{The missing normal-symbol isomorphism is supplied, not read from the scan.}
\GSfield{Fate}{This factorization is explicit in \citet[Theorem (3.10)]{atiyahpatodisinger1975spectral1}. The notebook suppresses \(\sigma\), a harmless shorthand only if the bundle identification is supplied.}
\GSfield{Present status}{As of 11 November 2026: the stated qualified result is established; the literal notebook shorthand should be read with the hypotheses above.}
\GSfield{Assessment}{literal statement: correct with missing hypotheses; publication match: exact; verification: read the result; development: annotated; confidence: high}
\end{GSitem}
\begin{GSitem}{1981-559.2}{E}{The repeated APS formula and the choice of projection}
\GSfield{Kind; source}{theorem; \GSpage{P001022}}
\GSfield{The notes}{\(P_+f|_Y=0\) and \(\operatorname{ind}D=\int\alpha_0-(h+\eta)/2\) are repeated for normal \(A\).}
\GSfield{Statement}{For self-adjoint \(A\), the displayed \(-h/2\) formula corresponds to \(P_{\ge0}\), not the strictly positive projection alone. For \(P_{>0}\), allowing the kernel changes the index by \(h\). The normal case requires an explicitly chosen imaginary-axis spectral subspace and additional analytic hypotheses.}
\GSfield{Reconstruction}{We distinguish positive from nonnegative and do not claim to have proved the full normal-operator extension.}
\GSfield{Fate}{\citet[Theorem (3.10)]{atiyahpatodisinger1975spectral1} verifies the nonnegative convention. This record is a genuine repeated occurrence of 1981-558.11--12; it should remain linked rather than counted as a second independent discovery.}
\GSfield{Present status}{As of 11 November 2026: the stated qualified result is established; the literal notebook shorthand should be read with the hypotheses above.}
\GSfield{Assessment}{literal statement: not determinable; publication match: partial; verification: read the result; development: annotated; confidence: medium}
\end{GSitem}
\begin{GSitem}{1981-559.3}{E}{Heat supertrace as the basis of the boundary index proof}
\GSfield{Kind; source}{method; \GSpage{P001023}}
\GSfield{The notes}{Equation (9) writes \(\operatorname{ind}D=\operatorname{Tr}e^{-tD^*D}-\operatorname{Tr}e^{-tDD^*}\).}
\GSfield{Statement}{For a closed densely defined Fredholm operator with compact-resolvent nonnegative realizations \(D^*D\) and \(DD^*\), positive eigenvalues pair under \(D\), so their heat-trace difference equals the index for every \(t>0\). For a boundary problem, \(D^*\) must be the adjoint of the chosen APS realization, with its corresponding complementary boundary condition.}
\GSfield{Reconstruction}{The operator domains are indispensable supplied hypotheses.}
\GSfield{Fate}{\citet[\S3]{atiyahpatodisinger1975spectral1} use this identity in the boundary calculation. The formula is older than the talk and is not valid for two independently chosen boundary domains.}
\GSfield{Present status}{As of 11 November 2026: the stated qualified result is established; the literal notebook shorthand should be read with the hypotheses above.}
\GSfield{Assessment}{literal statement: correct with missing hypotheses; publication match: exact; verification: read the result; development: annotated; confidence: high}
\end{GSitem}
\begin{GSitem}{1981-559.4}{E}{Positive spectral zeta: a zero global residue can have a nonzero local density}
\GSfield{Kind; source}{theorem; \GSpage{P001023}}
\GSfield{The notes}{\(\zeta_+(s)=\sum_{\lambda>0}\lambda^{-s}=\operatorname{Tr}(P_+A^{-s})\) has a finite value at zero, although its local density may have a residue whose integral is zero.}
\GSfield{Statement}{For a self-adjoint classical elliptic operator of positive order on a closed manifold, \(\zeta_+(s)=\operatorname{Tr}(P_+|A|^{-s})\) is regular at zero. The meromorphic diagonal trace density can have a nonzero pointwise residue; its integral vanishes. Thus regularity of a global trace does not establish pointwise holomorphy.}
\GSfield{Reconstruction}{We use absolute values on the spectral powers and specify closedness and classical order.}
\GSfield{Fate}{\citet[Theorem 0.1 and Introduction]{gilkey1981residue} establish global eta regularity in all dimensions and discuss nonzero local residues. Together with zeta regularity this gives the positive part. \citet[Theorem 4.3 and Remarks 4.4--4.5]{ponge2006spectral} explain the residue of projections. Gilkey’s 1981 publication precedes this talk.}
\GSfield{Present status}{As of 11 November 2026: the stated qualified result is established; the literal notebook shorthand should be read with the hypotheses above.}
\GSfield{Assessment}{literal statement: correct with missing hypotheses; publication match: exact; verification: read the result; development: annotated; confidence: high}
\end{GSitem}
\begin{GSitem}{1981-559.5}{E}{The sum and difference of the two spectral halves}
\GSfield{Kind; source}{definition; \GSpage{P001023}}
\GSfield{The notes}{\(\zeta(A)=\zeta_+(A)+\zeta_-(A)\), \(\eta(A)=\zeta_+(A)-\zeta_-(A)\).}
\GSfield{Statement}{For real nonzero spectrum define \(\zeta_\pm(s)=\sum_{\pm\lambda>0}|\lambda|^{-s}\). Then \(\zeta_{|A|}(s)=\zeta_+(s)+\zeta_-(s)\) and \(\eta_A(s)=\zeta_+(s)-\zeta_-(s)\). A spectral-cut zeta function of \(A\), rather than \(|A|\), has a phase on the negative part. At zero those phases can be discarded only after checking the regularity just discussed.}
\GSfield{Reconstruction}{We separate functions from their regularized values at zero and make the hidden absolute-value convention explicit.}
\GSfield{Fate}{The distinction is explicit in \citet[\S4, equations (4.1)--(4.2)]{ponge2006spectral}. The notebook writes values at zero; extending its first equality to all \(s\) without \(|A|\) would be false.}
\GSfield{Present status}{As of 11 November 2026: the stated qualified result is established; the literal notebook shorthand should be read with the hypotheses above.}
\GSfield{Assessment}{literal statement: correct with missing hypotheses; publication match: exact; verification: read the result; development: annotated; confidence: high}
\end{GSitem}
\begin{GSitem}{1981-559.6}{F}{Boundary values for unspecified general characteristic classes}
\GSfield{Kind; source}{question; \GSpage{P001023}; I. M. Gelfand}
\GSfield{The notes}{Gelfand cites Gabrielov, Gelfand and Fuks for algebraic reasons for nonlocality and says the analogous task for general classes has not been done.}
\GSfield{Statement}{The class of characteristic classes, geometric structures, equivalence relation on boundary data and desired invariant are unspecified. The sentence does not determine a unique generalization of the signature eta invariant.}
\GSfield{Fate}{The adjacent Pontryagin discussion suggests characteristic-class transgression, but is not enough to identify a theorem. The well-known Gabrielov--Gelfand--Losik papers must not be silently substituted for the name \ru{Фукс} actually visible in the scan; the attribution remains unresolved.}
\GSfield{Present status}{As of 11 November 2026: No precise new open problem or exact later solution can be assigned without additional source evidence.}
\GSfield{Assessment}{literal statement: not determinable; publication match: none found; verification: unresolved lead; development: annotated; confidence: low}
\end{GSitem}
\begin{GSitem}{1981-559.7}{E}{Stredder and Epstein: an identifiable bibliographic pointer}
\GSfield{Kind; source}{remark; \GSpage{P001023}}
\GSfield{The notes}{Two papers by Stredder and Epstein are listed as J. Differential Geometry 10 (1975), 645, about natural differential operators.}
\GSfield{Statement}{The identifiable works are Epstein, “Natural tensors on Riemannian manifolds”, pp.~631--645, and Stredder, “Natural differential operators on Riemannian manifolds and representations of the orthogonal and special orthogonal groups”, pp.~647--660. No particular classification theorem is stated in the seminar notes.}
\GSfield{Fate}{The two titles and page ranges are independently recorded in the reference list of \citet{fegansteer2005second}. The original DOI/Project Euclid routes did not open in this run, so the two original entries are marked unverified and the verification level remains indirect. Page 645 is the end of Epstein’s paper, not a common starting page.}
\GSfield{Present status}{As of 11 November 2026: The bibliographic pointer is identified; no full-text verification of either original classification is claimed.}
\GSfield{Assessment}{literal statement: correct with missing hypotheses; publication match: exact; verification: indirect; development: annotated; confidence: medium}
\end{GSitem}
\begin{GSitem}{1981-559.8}{E}{Is self-adjointness necessary for spectral asymmetry?}
\GSfield{Kind; source}{question; \GSpage{P001023-P001024}; I. M. Gelfand}
\GSfield{The notes}{Gelfand asks \ru{зачем самосопряженность A?}; the subsequent construction uses complex spectral cuts.}
\GSfield{Statement}{For a positive-order classical elliptic operator on a closed manifold, complex powers and a zeta function can be defined without self-adjointness when a ray avoids both the principal-symbol spectrum and the operator spectrum and satisfies the required parameter-dependent resolvent estimates. Replacing a real-sign eta sum requires specifying the cut or sector; it is not the same definition without choices.}
\GSfield{Reconstruction}{We answer the closed-operator functional-calculus question precisely. The broader normal boundary index assertion in 1981-559.2 is not proved by this answer.}
\GSfield{Fate}{Seeley’s complex-power theory predates the talk; \citet{seeley1968powers} is an early published translation. \citet[\S\S2--4]{ponge2006spectral} give the later precise asymmetry comparison. Wodzicki’s relevant papers appeared in February 1982 and February 1984; the closed-manifold construction is not an unresolved question of principle.}
\GSfield{Present status}{As of 11 November 2026: Self-adjointness is unnecessary for admissible-cut complex powers. The choice dependence and boundary-domain issues must be specified separately.}
\GSfield{Assessment}{literal statement: correct with missing hypotheses; publication match: partial; verification: read the result; development: annotated; confidence: high; problem \hyperref[prob:GP-1981-44]{GP-1981-44}}
\end{GSitem}
\begin{GSitem}{1981-559.9}{E}{Does eta regularity require a filling?}
\GSfield{Kind; source}{question; \GSpage{P001023}}
\GSfield{The notes}{An eta invariant is described through a filling and the notes say the case without a filling has not been proved.}
\GSfield{Statement}{For every self-adjoint classical elliptic operator of positive order on a closed manifold, in any dimension, the eta function has a meromorphic continuation regular at zero. Its spectral definition needs no bounding manifold. A geometric filling can give an index formula for its value but is not required for existence.}
\GSfield{Reconstruction}{The ambiguity between self-adjoint regularity and the normal extension is retained rather than resolved by attribution.}
\GSfield{Fate}{\citet[Theorem 0.1]{gilkey1981residue} proves the general regularity, following the earlier odd-dimensional APS result. The notebook sentence is false if read as a claim about all such self-adjoint operators; it may instead refer to the unspecified nonself-adjoint extension. Earlier occurrences are 1972-296.2--3 and 1972-296.5.}
\GSfield{Present status}{As of 11 November 2026: The self-adjoint no-filling problem was solved before this talk. The intended broader hypothesis set is not recoverable.}
\GSfield{Assessment}{literal statement: false as written; publication match: partial; verification: read the result; development: annotated; confidence: medium; problem \hyperref[prob:GP-1981-43]{GP-1981-43}}
\end{GSitem}
\begin{GSitem}{1981-559.10}{E}{A shifted local zeta representative of the index density}
\GSfield{Kind; source}{claim; \GSpage{P001023}}
\GSfield{The notes}{\(\alpha_0(x)=\zeta(D^*D+\varepsilon,x)-\zeta(DD^*+\varepsilon,x)\).}
\GSfield{Statement}{For a closed elliptic operator, the difference of shifted local zeta values at zero integrates to its index. It equals the constant term in \(e^{-\varepsilon t}\) times the local heat supertrace. If negative-power supertrace coefficients vanish pointwise, as in the classical geometric Dirac case, this is the usual local index density; for a general elliptic operator the chosen local representative can depend on the shift.}
\GSfield{Reconstruction}{We distinguish the shifted representative from a separately fixed canonical density and supply the suppressed evaluation at s=0.}
\GSfield{Fate}{The heat-to-zeta relation in \citet[\S\S1.7 and 1.10]{gilkey1984invariance} and the spectral pairing in 1981-558.3 give this exact interpretation. Equality of integrals should not be strengthened into equality of all local representatives.}
\GSfield{Present status}{As of 11 November 2026: the stated qualified result is established; the literal notebook shorthand should be read with the hypotheses above.}
\GSfield{Assessment}{literal statement: correct with missing hypotheses; publication match: partial; verification: read the result; development: annotated; confidence: high}
\end{GSitem}
\begin{GSitem}{1981-559.11}{E}{The corrected ray integral and the alleged Seeley sign error}
\GSfield{Kind; source}{method; \GSpage{P001024}}
\GSfield{The notes}{Equation (15) has \(+1/m\); equation (16), labelled correct, has \(-e^{i\theta}/m\) and \(b^\circ_{-m-n}(x,\xi,e^{i\theta}\lambda)\). The margin says Seeley has the wrong sign.}
\GSfield{Statement}{The substitution from a ray variable \(z=e^{i\theta}\lambda\) contributes the factor \(e^{i\theta}\). A meaningful comparison must also fix whether the resolvent is \((A-z)^{-1}\) or \((z-A)^{-1}\), the orientation, trace/Fourier normalization and any kernel term. The scan supports the corrected minus sign and a zero first argument of the local zeta value, but does not by itself establish a publication error.}
\GSfield{Reconstruction}{Only the ray substitution is convention-independent; the full formula is not silently normalized to agree with a modern source.}
\GSfield{Fate}{\citet[\S2]{ponge2006spectral} state explicit complex-power and residue conventions. The original Seeley formula was not verified at this level in this run; the attribution of an error is consequently an unresolved lead rather than an endorsed correction to his paper.}
\GSfield{Present status}{As of 11 November 2026: the stated qualified result is established; the literal notebook shorthand should be read with the hypotheses above.}
\GSfield{Assessment}{literal statement: not determinable; publication match: partial; verification: indirect; development: annotated; confidence: medium}
\end{GSitem}
\begin{GSitem}{1981-559.12}{E}{The repeated div–grad–curl square}
\GSfield{Kind; source}{claim; \GSpage{P001024}; I. M. Gelfand}
\GSfield{The notes}{The all-positive vector-calculus matrix is now explicitly on \(\mathbb R^3\), but is again said to square to \(\Delta I\).}
\GSfield{Statement}{With standard Euclidean conventions, its vector block squares to \(2\operatorname{grad}\operatorname{div}-\Delta I_3\), not \(\Delta I_3\). Replacing the upper-right entry by \(-\operatorname{div}\) gives the self-adjoint elliptic matrix whose square is \(-\Delta I_4\).}
\GSfield{Reconstruction}{The correction fixes a declared convention; the scan is transcribed faithfully.}
\GSfield{Fate}{This is the same occurrence as the example in 1981-558.6, repeated a week later. The algebraic calculation is included with its independent polynomial-symbol check; \citet[Introduction]{atiyahpatodisinger1975spectral1} is the historical geometric context.}
\GSfield{Present status}{As of 11 November 2026: the stated qualified result is established; the literal notebook shorthand should be read with the hypotheses above.}
\GSfield{Assessment}{literal statement: false as written; publication match: partial; verification: read the result; development: annotated; confidence: high}
\end{GSitem}
\begin{GSitem}{1981-559.13}{A}{Differences of spectral-cut zeta functions at integer points}
\GSfield{Kind; source}{theorem; \GSpage{P001024-P001025}}
\GSfield{The notes}{For two rays of minimal growth, \(\rho(s)=\zeta_\theta(A,s)-\zeta_{\theta^{\prime}}(A,s)\) is meromorphic; its residues vanish microlocally at integral candidate poles and its integral-point values are local.}
\GSfield{Statement}{Let \(P\) be a positive integer-order \(m\) classical elliptic operator on a closed manifold with two admissible spectral cuts, and let \(\Pi_{\theta,\theta^{\prime}}(P)\) be the associated sectorial projection. With arguments in \((\theta-2\pi,\theta)\), \[\rho(s)=(1-e^{-2\pi i s})\operatorname{TR}(\Pi_{\theta,\theta^{\prime}}P_\theta^{-s}).\] The factor cancels the possible simple pole at any integer \(k\), and \[\rho(k)=\frac{2\pi i}{m}\operatorname{Res}(\Pi_{\theta,\theta^{\prime}}P^{-k}).\] The residue is an integral of a local symbol density.}
\GSfield{Reconstruction}{The canonical-trace notation and orientation convention are modern; the notebook’s sign in its symbol-residue formula is not compared without matching conventions.}
\GSfield{Fate}{\citet[Theorem 4.2, equations (4.13)--(4.14)]{ponge2006spectral} prove this precise formulation and trace the method to Wodzicki. The 1982 paper was already published before the talks; the fuller 1984 treatment followed them. The primary Wodzicki publisher records were opened, but the exact original theorem text was unavailable, so priority between those two articles is not inferred from the notebook.}
\GSfield{Present status}{As of 11 November 2026: The integer-value residue formula is established in the stated classical calculus; the microscopic cancellation and global integral are distinct statements.}
\GSfield{Assessment}{literal statement: correct with missing hypotheses; publication match: partial; verification: read the result; development: annotated; confidence: medium}
\end{GSitem}
\begin{GSitem}{1981-559.14}{A}{Cut independence at zero and the nonpositive-integer extension}
\GSfield{Kind; source}{claim; \GSpage{P001025}}
\GSfield{The notes}{Part (d) says \(\rho(\ell)=0\) for \(\ell=0,-1,-2,\ldots\), while the following hint mentions integers \(>0\).}
\GSfield{Statement}{For admissible positive-order classical elliptic operators, \(\rho(0)=0\): the noncommutative residue of a pseudodifferential projection vanishes. At other integers, the exact criterion is \(\operatorname{Res}(\Pi P^{-k})=0\); it is not permissible to erase this condition for arbitrary classical pseudodifferential \(P\). The blanket nonpositive-integer assertion can fail even for differential operators: in odd dimensions, negative-integer poles of the absolute zeta or eta functions can give nonzero cut differences. Additional parity or residue-vanishing hypotheses are needed.}
\GSfield{Reconstruction}{The valid zero-value theorem is separated from the false general extension. Regularity of a cut zeta function at a negative integer does not imply equality of its values for different cuts.}
\GSfield{Fate}{\citet[Theorem 4.3 and Remark 4.5]{ponge2006spectral} verify zero-value cut independence. \citet[\S4]{loyamoroianu2012singularities} verify admissible negative poles for Dirac-type differential operators in odd dimensions; together with the cut-difference identity these rule out the unrestricted extension. An independent explicit check in \hyperref[dev:1981-negative-cut]{the spectral verification} constructs an invertible massive Dirac operator on the flat three-torus: its two cut zeta values at \(-1\) are \(\mp i\pi^2\mu^4\). The positive-integer hint conflicts with the displayed range and remains a source ambiguity.}
\GSfield{Present status}{As of 11 November 2026: Cut independence at zero is established. The stated extension to every nonpositive integer is false without further restrictions, including within differential-operator settings.}
\GSfield{Assessment}{literal statement: false as written; publication match: partial; verification: read the result; development: annotated; confidence: high}
\end{GSitem}
\begin{GSitem}{1981-559.15}{A}{The shifted-zeta polynomial and its missing powers}
\GSfield{Kind; source}{method; \GSpage{P001025}}
\GSfield{The notes}{For order one, equation (20) gives a binomial sum plus \[t^k\sum_{\mu=1}^n\frac{(-1)^\mu\delta_\mu}{\mu\binom{k+\mu}{k}},\] without \(t^\mu\) in its summands.}
\GSfield{Statement}{For an invertible elliptic differential operator \(A\) of order one on a closed \(n\)-manifold, with a fixed admissible cut and a sufficiently small scalar shift preserving invertibility and the cut, write \(\delta_\mu=\operatorname*{res}_{s=\mu}\zeta_A(s)\). Then for \(k\ge0\), \[\begin{split}\zeta_{A+t}(-k)={}&\sum_{\nu=0}^{k}\binom{k}{\nu}t^\nu\zeta_A(\nu-k)\\&+\sum_{\mu=1}^{n}\frac{(-1)^\mu\delta_\mu\,t^{k+\mu}}{\mu\binom{k+\mu}{k}}.\end{split}\] Kernel terms require separate bookkeeping when invertibility is dropped.}
\GSfield{Reconstruction}{The complete corrected coefficient derivation is included as an annotation in the development file, not claimed as a new theorem of spectral theory.}
\GSfield{Fate}{The correction follows by expanding \((A+t)^{-s}\) and taking the limit at \(s=-k\); the coefficient is independently checked symbolically in the included code. Wodzicki’s local-invariant programme is the historical context, but a direct published equation-by-equation match to (20) was not read. The literal omission is already contradicted by setting \(t=0\) when a residue is nonzero.}
\GSfield{Present status}{As of 11 November 2026: The corrected local shift identity is verified under these hypotheses; the notebook formula without the additional powers is false.}
\GSfield{Assessment}{literal statement: false as written; publication match: partial; verification: publisher metadata checked; development: annotated; confidence: medium}
\end{GSitem}
\begin{GSitem}{1981-559.16}{F}{The Mahler analogy and the unspecified p-adic continuation}
\GSfield{Kind; source}{remark; \GSpage{P001025}}
\GSfield{The notes}{The first finite sum is called a Mahler transform; an automorphic function for it is said to continue to a p-adic domain.}
\GSfield{Statement}{Mahler’s classical theorem expands a continuous \(f:\mathbb Z_p\to\mathbb Q_p\) in the binomial basis, with coefficients tending to zero. To turn the spectral shift identity into a p-adic interpolation statement one would need a specified p-adic coefficient field, sequence, continuity/congruence estimates and meaning of automorphy. None is supplied here.}
\GSfield{Fate}{\citet{mahler1958interpolation} is a genuine antecedent for the binomial basis; the further automorphic/p-adic sentence cannot be matched to a precise theorem from the notes. No general p-adic continuation of spectral zeta values is asserted.}
\GSfield{Present status}{As of 11 November 2026: The Mahler basis theorem is established; the proposed spectral/automorphic continuation is not identifiable from the recorded data.}
\GSfield{Assessment}{literal statement: not determinable; publication match: thematic only; verification: publisher metadata checked; development: annotated; confidence: low}
\end{GSitem}

\section{Developments}\label{dev:1981}
\subsection{Editorial verification: what the root-multiplicity question actually asks}

No item in records 530--537 is assigned class C or D. The following proofs and computations are annotations checking the reconstruction and the literal errors; they are not claimed as new research. In particular, an unspecified wish for an illuminating multiplicity formula is not turned into an open demand for a positive formula. The symmetrizable formula in \citet[Theorem 2]{bermanmoody1979multiplicities} was already available in 1979. The elementary argument below covers the finite computational interpretation of 1981-533.4 without importing that paper's symmetrizability hypothesis.

\begin{proposition}[Recursive lowering determines every specified positive degree]\ours
Let \(A\) be a finite generalized Cartan matrix over characteristic zero. Use a Cartan realization with independent simple roots, and define the contragredient algebra by factoring the Chevalley cross-relation algebra by its largest ideal disjoint from the Cartan subalgebra. If \(\alpha\in Q_+\) has height greater than one, then
\[
 \mathfrak g_\alpha\longrightarrow
 \bigoplus_{i:\,\alpha-\alpha_i\in Q_+\setminus\{0\}}
 \mathfrak g_{\alpha-\alpha_i},\qquad
 v\longmapsto([f_i,v])_i
\]
is injective. Consequently a prescribed root-space dimension and independent representing words are computable by finitely many exact linear operations, inductively on height.
\end{proposition}
\begin{proof}
The positive part of the algebra before the maximal quotient is the free Lie algebra on the \(e_i\); its component of degree \(\alpha\) is finite-dimensional. The triangular decomposition and maximal-ideal construction are those of \citet[Chapter II, \S1]{kac1968graded}. The simple positive components have dimension one.

Suppose \(v\in\mathfrak g_\alpha\) is killed by all \(f_i\). Let \(J\) be the span of all iterated brackets
\(
 [e_{i_1},[e_{i_2},\ldots,[e_{i_s},v]\ldots]]
\), including \(v\). It is stable under the \(e_i\) and the Cartan subalgebra. It is stable under each \(f_j\) as well: move \(\operatorname{ad}f_j\) through the iterated positive adjoints, using
\([f_j,e_i]=-\delta_{ij}h_i\), and then move each Cartan adjoint through the remaining positive adjoints. Every resulting term is a scalar multiple of a shorter positive iterated bracket on \(v\); the terminal \([f_j,v]\) vanishes. This induction also covers an arbitrary linear combination of such brackets. Hence \(J\) is an ideal of \(\mathfrak g\).

Every weight of \(J\) has degree \(\alpha+\beta\) for \(\beta\in Q_+\), so \(J\cap\mathfrak h=0\). The maximal-quotient definition prohibits any nonzero such ideal: its preimage would strictly enlarge the largest ideal disjoint from \(\mathfrak h\). Therefore \(v=0\), proving injectivity. Terms with a negative coefficient in \(\alpha-\alpha_i\) are automatically zero; height greater than one excludes a Cartan-valued target. Removing the redundant central directions of Rudakov's enlarged realization does not change these positive components.

For effectiveness, choose a finite basis or spanning list of free Lie words of degree \(\alpha\). By induction suppose quotient coordinate maps are known in every smaller positive degree. The brackets \([f_i,w]\) are calculated from the cross relations and expressed in those already computed quotient coordinates. Stack the resulting matrices. Its kernel is exactly the degree-\(\alpha\) part of the maximal ideal, by the injection just proved. Its rank is the desired dimension, and a set of pivot columns supplies independent word representatives. All coefficients are obtained from the integer entries of \(A\), so exact rational arithmetic suffices for a generalized Cartan matrix. The calculation at every fixed height terminates because all free Lie components in question are finite-dimensional. No complexity bound, canonical choice of pivots or cancellation-free formula is asserted.
\end{proof}

\begin{remark}[Independent check of the lowering argument]\oursComputed
The supplied \texttt{editorial\_checks.py} computes the same quotient dimensions by two different procedures: the lowering-map recursion above, and a degree-by-degree Lie ideal generated by the two explicit Serre relations. Both use an exact free-Lie basis, but the second procedure never uses a lowering operator. For the symmetrizable rank-two pairs \((a,b)=(1,1),(2,1),(3,1),(2,2),(3,2)\), they agree through height eight. The first three height profiles are respectively
\[
 (2,1,0,0,0,0,0,0),\quad
 (2,1,1,0,0,0,0,0),\quad
 (2,1,1,1,1,0,0,0).
\]
These independently recover total dimensions \(3,4,6\). The affine pair \((2,2)\) gives \((2,1,2,1,2,1,2,1)\); the indefinite pair \((3,2)\) gives \((2,1,2,2,3,3,6,7)\). Finite output does not prove infinite-dimensionality; that assertion uses the previously established Cartan classification. The second procedure is used only for symmetrizable rank-two matrices, where the Serre presentation is justified. It is not substituted for the maximal-ideal definition in the general statement above.
\end{remark}

\subsection{Editorial counterexamples for the transformation-group notes}

For 1981-534.1--5 write \(G_{(S)}\) for the pointwise stabilizer and \(G_{\{S\}}\) for the setwise stabilizer. Then the source's closure is
\(
 \operatorname{cl}(S)=\operatorname{Fix}(G_{(S)}).
\)
This distinction is also needed for the restriction homomorphisms in the notebook's fourth axiom. The exact match of those axioms to Seifert's 1984 theorem remains unverified; the following corrections concern elementary statements, independent of that publication-identification gap.

\begin{proposition}[Singleton closure need not mean equal stabilizers]\ours
For every group action,
\[
 \operatorname{cl}\{x\}=\{y:G_x\subseteq G_y\}.
\]
Even in a faithful transitive action the containment can be strict. If the transitive set is finite, it is necessarily equality.
\end{proposition}
\begin{proof}
The displayed equivalence is just the definition that every element fixing \(x\) also fixes \(y\). For strictness let \(G=\mathbb Z[1/2]\rtimes\langle t\rangle\), with \(tbt^{-1}=2b\) on the additive translation subgroup, let \(H=\mathbb Z\), and let \(G\) act on left cosets \(G/H\). At \(x=H\), the stabilizer is \(H\). At \(y=t^{-1}H\), it is \(t^{-1}Ht=\tfrac12\mathbb Z\), strictly larger. Therefore \(y\in\operatorname{cl}\{x\}\) while \(G_y\ne G_x\). The action is transitive by construction and faithful: its kernel lies in \(\bigcap_{n\ge0}2^n\mathbb Z=\{0\}\), hence is trivial. Thus faithfulness does not repair the notebook's statement.

In a finite transitive set, \(G_x\) and \(G_y\) have the same finite index in \(G\). If \(G_x\subseteq G_y\), the index multiplication formula implies \([G_y:G_x]=1\). This is a separate verification of the missing finite-action hypothesis, and does not require \(G\) itself to be finite.
\end{proof}

\begin{example}[The degree-two exception to the double-transitivity remark]\ours
Let \(C_2\) act regularly on its two elements. It is transitive on the two ordered pairs of distinct elements, hence doubly transitive. Every point stabilizer is trivial, so \(\operatorname{cl}\{x\}=X\), strictly larger than \(\{x\}\). Thus the notebook's proposed obstruction requires a restriction. For a doubly transitive set with at least three points, \(G_x\) acts transitively on \(X\setminus\{x\}\), a set of at least two elements, and no point there is fixed by all of \(G_x\). In that range \(\operatorname{cl}\{x\}=\{x\}\) does hold.
\end{example}

\subsection{Independent finite checks of combinatorial formulas}

\begin{proposition}[Interval inversion]\ours
For interval multiplicities \(m_{pq}\) of an equioriented type-A representation,
\[
 r_{ij}=\sum_{p\le i\le j\le q}m_{pq},\qquad
 m_{ij}=r_{ij}-r_{i-1,j}-r_{i,j+1}+r_{i-1,j+1},
\]
where ranks beyond the ends are zero.
\end{proposition}
\begin{proof}
An interval summand contributes one to the rank of the composition from vertex \(i\) to vertex \(j\) exactly when it contains both endpoints, proving the first formula. The difference \(r_{ij}-r_{i-1,j}\) retains exactly intervals starting at \(i\) and ending at or beyond \(j\). Subtracting the corresponding difference with endpoint \(j+1\) retains exactly the intervals ending at \(j\). This proves the second formula. It verifies the indexing in the source completion from \citet[\S2.3]{abeasisdelfrakraft1981geometry}, not the deeper converse degeneration theorem.
\end{proof}

\begin{remark}[Independent enumeration results]\oursComputed
The code constructs explicit direct-sum arrow matrices for all 120 choices of three distinct intervals on four vertices, with independently varied multiplicities one or two. Matrix products and exact ranks agree with the finite-difference reconstruction in every case. For the free Lie formula, direct enumeration of primitive binary necklaces agrees with the M\"obius formula in all 65 bidegrees of total degree one through ten. The independent binary-tree recurrence agrees with the Catalan specialization for \(n=0,\ldots,12\).

For the Gelfand--Tsetlin question, recursive enumeration of every intervening integer row agrees with the Weyl product for all 210 weakly increasing top rows drawn from \(\{-2,-1,0,1,2,3,4\}\). Examples give counts \(1\) for \((0,0,0,0)\), \(4\) for \((0,0,0,1)\), \(64\) for \((0,1,2,3)\), and \(630\) for the additional row \((-2,-1,1,4)\). The general theorem is supplied by the interlacing basis \citep[\S3]{gelfandtsetlin1950unimodular} and the Weyl dimension formula; bounded enumeration is an independent check on indexing, endpoints and normalization.

Finally, multiplying a symmetric-algebra Euler product and separately expanding its formal logarithm through degree fourteen confirms \(C_n=\sum_{d\mid n}d\dim V_d\). This checks that the notebook's rescaled coefficients \(c_k=k\dim V_k\), including the denominator \(k\) in \(p(t)\), are internally consistent. Complete parameter ranges and output are retained beside the code.
\end{remark}

\subsection{Two sign and scaling checks}

\begin{example}[A direct test of the BGG rho shift]\ours
With \([e,f]=h\), \([h,f]=-2f\), let \(M(0)\) be the \(\mathfrak{sl}_2\) Verma module generated by \(v_0\), with \(ev_0=hv_0=0\). The nonzero vector \(fv_0\) has weight \(-2\), and \(e(fv_0)=hv_0=0\). The universal property gives a nonzero map \(M(-2)\to M(0)\); on the polynomial PBW bases it is multiplication by \(f\), so it is injective. Since \(\rho(h)=1\), the weights satisfy \(-2+1=-(0+1)\). They do not satisfy \(-2-1=\pm(0-1)\). This verifies the plus-rho convention against both an explicit module calculation and the shifted definition in \citet[\S3]{bernsteingelfandgelfand1971structure}.
\end{example}

\begin{proposition}[Determinant weights of the supergravity density]\ours
In the determinant-monomial ansatz \(A=|W|^p|\Gamma|^q\), the defining-equation and odd-frame covariance laws used by \citet[\S3, equation (8)]{schwarz1982supergravity} force \(p=4/3\), \(q=-1/3\).
\end{proposition}
\begin{proof}
Replacing the four real defining equations by \(f'=\eta f\) multiplies each determinant by \(\det\eta\), while the product of delta functions has weight \(-1\). Thus \(p+q=1\). Changing the complex odd tangent frame by a \(2\times2\) matrix \(S\) gives \(|W|\) weight \(|\det S|\) and the paired-spinor determinant \(|\Gamma|\) weight \(|\det S|^4\); hence \(p+4q=0\). The two equations have the unique solution stated.

As a separate check, rescale all defining equations by a positive real constant \(c\). Both four-by-four determinants acquire \(c^4\), and the delta product acquires \(c^{-4}\). The final density has total exponent \(16/3-4/3-4=0\), while the preliminary expression \(|\det\Gamma|^{-1/2}\prod\delta(f^k)\) has exponent \(-2-4=-6\). This confirms the source's internal correction. The argument determines exponents within this ansatz; it is not a replacement for the stronger local uniqueness classification attributed to \citet{gaydukromanovschwarz1981democracy}, whose full proof was not accessible here.
\end{proof}

\subsection{Verification of the Jantzen emendations}

This subsection checks two annotations to record 540. It does not claim a new solution of a C/D problem: the unrestricted direction question has the published negative answer in \citet[\S8.4]{williamson2016local}. The checks below provide independently reproducible evidence that the hypotheses and grading in the companion have not merely been copied from an abstract.

\begin{example}[Why moving lifts matter]\ours
Let \(A=\mathbb Q[t]_{(t)}\), and give \(A^2\) the symmetric bilinear form with matrix
\[
 S(t)=\begin{pmatrix}1&t\\t&t^2+t^3\end{pmatrix}.
\]
Its elementary divisors have orders \(0,3\). The class of \(e_2\) in \(A^2/tA^2\) therefore belongs to the third Jantzen-type step. The constant lift \(e_2\) has pairing \((t,t^2+t^3)\), whose order is only one. The moving lift \(e_2-te_1\) instead has pairing \((0,t^3)\), as required. This supplies an explicit reason why the literal fixed-vector definition on P000997 is incomplete. Direct multiplication and the determinant \(\det S(t)=t^3\) independently check the calculation.
\end{example}

\begin{example}[Exact reproduction of the type-A3 direction counterexample]\oursComputed
The file \texttt{verify\_annotations.py} builds the contravariant form from the matrix-unit relations
\[
 [E_{ij},E_{kl}]=\delta_{jk}E_{il}-\delta_{li}E_{kj}
\]
on a PBW basis of the weight space of highest weight \(-\alpha_1-\alpha_3\), lowered by \(2\alpha_1+3\alpha_2+2\alpha_3\). There are thirteen basis monomials. All arithmetic is exact. Dynkin coordinates of the highest weight are \((-2,2,-2)\); the three deformation directions are displayed below.
\[
\begin{array}{c|c|c}
\gamma(h_1),\gamma(h_2),\gamma(h_3)&
(\dim \operatorname{gr}^0,\dim \operatorname{gr}^1,
 \dim \operatorname{gr}^2,\dim \operatorname{gr}^3)&
(N(1),N(2),N(3),N(4))\\\hline
(1,1,1)&(7,3,2,1)&(6,9,10,10)\\
(1,-2,3)&(7,2,4,0)&(6,10,10,10)\\
(3,-2,1)&(7,2,4,0)&(6,10,10,10)
\end{array}
\]
Here \(N(k)\) is the nullity, over \(\mathbb Q\), of multiplication by the form matrix modulo \(t^k\). Every positive-coroot value of each direction is nonzero. The two non-dominant directions satisfy \(\gamma(h_1+2h_2+h_3)=0\); the third row independently checks the reflected choice.

Two different exact procedures agree. First, elimination over rational power series modulo \(t^5\) gives elementary-divisor orders. Second, fraction-free integer row elimination gives the ranks of the block matrices modulo \(t,t^2,t^3,t^4\). If the elementary-divisor orders are \(a_1,\ldots,a_{13}\), then
\[
 N(k)=\sum_{r=1}^{13}\min(k,a_r),\qquad
 \dim\operatorname{gr}^j=\#\{r:a_r=j\}.
\]
Stabilization between \(k=3\) and \(k=4\) also rules out hidden higher orders. The saved \texttt{verification\_output.json} and \texttt{verification\_run.log} record the full results. They reproduce the published counterexample and do not constitute a priority claim.
\end{example}

The same executable independently checks the signs in the Cauchy determinant at sizes 1 through 5, including a repeated-row edge case; the differentiated Rayleigh equation and its inverse transformation; an explicit nondegenerate binary-cubic Kirillov pairing; and all 511 finite surreal sign sequences of length at most eight by both canonical dyadic cuts and the corrected halving formula. These finite checks support particular emendations. They are not presented as proofs of an unbounded numerical conjecture.

\subsection{Haefliger's modular-group coinvariant equation: a complete continuous result}
\label{dev:1981-haefliger}

The question in 1981-555.1 is exact smooth solvability on the projective circle
\(X=\mathbb {RP}^{1}=\mathbb R\cup\{\infty\}\): for every real
\(f\in C^\infty(X)\), do there exist real \(u,v\in C^\infty(X)\) such that
\[
 f=u-u\circ T+v-v\circ S,
 \qquad T(x)=x+1,\qquad S(x)=-1/x?
\]
Both maps extend smoothly across the apparent affine-coordinate singularities.
The notebook calls exact solvability unknown and asserts density of the image
in the \(C^\infty\) topology. Those are two separate assertions. In particular,
density in a Fr\'echet space does not imply surjectivity.

\paragraph{Literature and the scope of this development.}
The search documented in the research log located Haefliger's 1980 paper on
minimal leaves, and the later invariant-functional framework of
\citet[abstract]{abouqatebelkacimi2000invariant}. The latter's
compact-group result does not cover this noncompact discrete group action.
Neither lead supplied a verified theorem answering the displayed smooth
equation. The subject connection with minimal leaves is recorded in the notes;
no historical explanation of the question's subsequent neglect is inferred.
The following result is a bounded contribution proved here, with no claim of
historical priority. It solves the exact equation on continuous functions. It
does not prove a regularity theorem for its solutions.

\begin{lemma}[No invariant probability measure]\ours
There is no Borel probability measure on \(X\) invariant under both \(T\)
and \(S\).
\end{lemma}
\begin{proof}
If \(\nu\) is \(T\)-invariant, the pairwise disjoint Borel sets
\([k,k+1)\), \(k\in\mathbb Z\), have the same measure. Finiteness forces
each measure to be zero, and countable additivity then gives
\(\nu(\mathbb R)=0\). Thus \(\nu=\delta_\infty\). Since
\(S(\infty)=0\), this probability is not \(S\)-invariant.
\end{proof}

\begin{lemma}[A uniform defect of invariance]\ours
Let \(\mathcal M(X)\) be the real finite signed regular Borel measures, with
total-variation norm \(\|\cdot\|_{\rm TV}\), and let \(g_*\mu\) denote
pushforward. There is a number \(\epsilon>0\) such that every
\(\mu\in\mathcal M(X)\) satisfies
\[
 \|\mu-T_*\mu\|_{\rm TV}+\|\mu-S_*\mu\|_{\rm TV}
 \ \geq\ \epsilon\|\mu\|_{\rm TV}.
 \tag{H1}
\]
\end{lemma}
\begin{proof}
The probabilities on compact \(X\) form a weak-* compact subset of
\(C(X)^*\). Pushforward by a homeomorphism is weak-* continuous and the
total-variation norm is weak-* lower semicontinuous. Consequently
\[
 \epsilon:=\inf_{\nu\ \mathrm{probability}}
 \bigl(\|\nu-T_*\nu\|_{\rm TV}+\|\nu-S_*\nu\|_{\rm TV}\bigr)
\]
is attained. If it were zero, a minimizer would contradict the preceding
lemma. Thus \(\epsilon>0\).

For signed measures \(\alpha,\beta\),
\(\||\alpha|-|\beta|\|_{\rm TV}\leq\|\alpha-\beta\|_{\rm TV}\).
Indeed, taking Radon--Nikodym derivatives with respect to
\(|\alpha|+|\beta|\) reduces this assertion to
\(\bigl||a|-|b|\bigr|\leq|a-b|\) and integration. A homeomorphism also
satisfies \(|g_*\mu|=g_*|\mu|\). Apply the positive-measure estimate to
\(|\mu|/\|\mu\|_{\rm TV}\), when \(\mu\ne0\), and use these two facts
for \(g=T,S\). The zero measure causes no exception.
\end{proof}

\begin{theorem}[Exact solvability for continuous functions]\ours
For every real \(f\in C(X)\), there are real \(u,v\in C(X)\) with
\[
 f=u-u\circ T+v-v\circ S.
 \tag{H2}
\]
If \(0<\epsilon'<\epsilon\), the solutions can be chosen to satisfy
\[
 \max(\|u\|_\infty,\|v\|_\infty)
 \leq \frac{2}{\epsilon'}\|f\|_\infty.
 \tag{H3}
\]
Here \(\epsilon\) is the positive constant in (H1); the theorem does not
claim an explicit numerical value for it or a linear choice of solutions.
\end{theorem}
\begin{proof}
Equip \(E=C(X)\oplus C(X)\) with the maximum norm and define the bounded
operator \(D:E\to C(X)\) by the right side of (H2). Under the Riesz
identification \(C(X)^*=\mathcal M(X)\), its adjoint is
\[
 D^*\mu=(\mu-T_*\mu,\mu-S_*\mu).
\]
The dual norm on \(E^*\) is the sum of the two total-variation norms.
Thus (H1) states precisely that \(D^*\) is bounded below by \(\epsilon\).

Let \(B_E\) be the closed unit ball and let
\(K=\overline{D(B_E)}\), with norm closure. For every \(\mu\),
\[
 \sup_{h\in K}\langle\mu,h\rangle
 =\|D^*\mu\|\geq\epsilon\|\mu\|_{\rm TV}.
\]
The Hahn--Banach separation theorem for closed convex sets implies
\(\epsilon' B_{C(X)}\subset K\) for every \(0<\epsilon'<\epsilon\).
To see the strict-inequality point explicitly, a hypothetical separating
functional for an element of \(\epsilon' B_{C(X)}\setminus K\) would
have supremum on \(K\) at least \(\epsilon\|\mu\|\), while its value
on that element is at most \(\epsilon'\|\mu\|\), a contradiction.

Set \(r_0=f\). Given \(r_j\ne0\), the inclusion just proved permits
a pair \(z_j\in E\) with
\[
 \|z_j\|_E\leq\|r_j\|_\infty/\epsilon',\qquad
 \|r_j-Dz_j\|_\infty\leq\tfrac12\|r_j\|_\infty.
\]
Put \(r_{j+1}=r_j-Dz_j\); if a residual is zero, take all subsequent
pairs zero. Then \(\|r_j\|_\infty\leq2^{-j}\|f\|_\infty\) and
\(\sum_j\|z_j\|_E\leq2\|f\|_\infty/\epsilon'\). Completeness gives
\(z=\sum_j z_j=(u,v)\in E\), and continuity of \(D\) gives
\(Dz=f\). This also proves (H3).
\end{proof}

\begin{proposition}[An exact reduction of the smooth question]\ours
For \(C_gw=w\circ g\), put \(A=(I+C_S)/2\). A smooth function \(f\)
has a smooth decomposition (H2) if and only if
\[
 Af\in\operatorname{ran}\bigl(A(I-C_T):C^\infty(X)\to
 C^\infty(X)^S\bigr),
 \tag{H4}
\]
where \(C^\infty(X)^S\) denotes the \(S\)-invariant functions.
\end{proposition}
\begin{proof}
Since \(S^2=I\) on \(\mathbb {RP}^{1}\), \(A\) is a projection and
\(A(I-C_S)=0\). Applying it to (H2) proves necessity. Conversely, choose
\(u\) with \(A(I-C_T)u=Af\). The remainder
\(r=f-(I-C_T)u\) satisfies \(Ar=0\), hence \(C_Sr=-r\).
Putting \(v=r/2\) gives \((I-C_S)v=r\) and proves sufficiency.
\end{proof}

\paragraph{Second check and remaining obstruction.}
An independent audit re-derived the measure argument and checked all norms.
It also gave a separate functional-analytic check: (H1) makes
\(\operatorname{ran}D^*\) norm closed and \(\ker D^*=0\); the Banach
closed-range theorem makes \(\operatorname{ran}D\) closed, while the
annihilator identity makes it dense, so \(D\) is onto. This agrees with
the constructive Hahn--Banach argument above. The use of total variation
of a signed measure is essential: mere absence of invariant signed
measures would show density, not a quantitative lower bound or
surjectivity.

The proof controls only uniform norms. Uniformly convergent series of
continuous solutions need not have bounded derivatives, and smoothing
does not commute with the two nonlinear coordinate actions in a way that
this argument estimates. In the smooth setting the dual obstructions are
distributions, not only measures. The notebook's claim of smooth density,
even if independently established, would still leave a closed-range or
regularity estimate to prove. Thus the strongest verified result in this
development is exact \(C^0\) solvability; the strongest formulation of the
original problem remains exact \(C^\infty\) solvability, retained as an
unresolved question rather than promoted to a theorem.

One further bounded investigation can target (H4): establish a tame
estimate for the averaged parabolic difference operator near the fixed
point \(\infty\), with a finite and explicitly controlled derivative
loss, or construct an obstruction to such an estimate. That is a precise
analytic task. The present argument supplies neither outcome, and the
absence of a verified resolution in this search is not evidence that no
resolution exists.

\subsection{Independent finite checks of two source emendations}

These checks support annotations 1981-548.1--2 and 1981-546.15; they are not
additional open problems. The partition theorem itself is the known identity
in \citet[Chapter I, Section 1, (1.7)--(1.7')]{macdonald1995symmetric}.

\begin{example}[Exact partition enumeration]\oursComputed
The supplied Python program enumerates partitions in every rectangle
\(0\leq m,n\leq6\) in two ways: weakly decreasing row lengths, and
\(n\)-element subsets of \(\{0,\ldots,m+n-1\}\), inverted as beta numbers.
Both enumerations agree and contain \(\binom{m+n}{n}\) partitions for each
rectangle. For all 3431 partitions across these 49 rectangles, the two
complementary exponent sets are disjoint and fill that interval. A separate
run on \(0\leq m,n\leq4\), with 251 partitions, agrees. The negative control
\(m=n=1\), \(\lambda=(1)\), makes the written generating-function left side
\(t+t^2\), whereas the required right side is \(1+t\). The corrected inverse
argument and conjugation give \(1+t\). These finite checks detect the source
error; they are not the proof for arbitrary rectangles.
\end{example}

\begin{proposition}[The simplex edge-length Jacobian]\ours
Order the independent Gram entries by diagonals first and then off-diagonals.
Write \(s_{ij}=\ell_{ij}^2\) for all \(n(n+1)/2\) edges of a labelled
\(n\)-simplex, with vertex zero distinguished. The absolute Jacobians are
\[
 \left|\det\frac{\partial(g_{ij})_{1\leq i\leq j\leq n}}
                         {\partial(s_{ij})_{0\leq i<j\leq n}}\right|
 =2^{-n(n-1)/2},
 \qquad
 \left|\det\frac{\partial(g_{ij})}{\partial(\ell_{ij})}\right|
 =2^n\prod_{0\leq i<j\leq n}\ell_{ij}.
\]
Consequently the density in equation (23) needs the product of lengths if
its differentials are length differentials.
\end{proposition}
\begin{proof}
The diagonal formula is \(g_{ii}=s_{0i}\). Each off-diagonal coordinate is
\(g_{ij}=(s_{0i}+s_{0j}-s_{ij})/2\). In the stated ordering the Jacobian is
block triangular, with diagonal blocks the identity and
\(-\tfrac12 I_{n(n-1)/2}\). This proves the first identity. The change from
lengths to squared lengths has diagonal Jacobian entries \(2\ell_{ij}\).
Their product gives the second identity. Finally
\(\det g=(n!\operatorname{vol}\sigma)^2\), so multiplying the invariant
matrix density by this Jacobian gives, up to a fixed constant,
\(\operatorname{vol}(\sigma)^{-(n+1)}\prod\ell_{ij}\,d\ell_{ij}\).
\end{proof}

\begin{example}[Second check of the Jacobian]\oursComputed
Independent symbolic differentiation of the full coordinate map in dimensions
\(n=1,2,3,4,5\) gives signed squared-length Jacobians
\(1,-1/2,-1/8,1/64,1/1024\), respectively. Division of the absolute
length-coordinate determinant by the product of all edge lengths gives
\(2,4,8,16,32\). An independent positive-definite test uses based orthonormal
edge vectors, whose pairwise squared distances are one or two, and recovers
the identity Gram matrix exactly in every dimension. For a single edge,
\(g=\ell^2\) gives \(g^{-1}dg=2\,d\ell/\ell\), while the written formula
would give \(d\ell/\ell^2\). The complete code and its output are included.
\end{example}

\subsection{Goldsheid's coordinate-product approximation: a counterexample and quantitative bounds}
\label{dev:1981-goldsheid}

\paragraph{The exact question and its two quantifiers.}
The intervention on P001018 (1981-557.8; GP-1981-40) concerns the standard
Euclidean symplectic space \(\mathbb R^{2n}\), with an orthonormal symplectic
basis \(e_1,\ldots,e_n,f_1,\ldots,f_n\). For every Lagrangian \(n\)-plane
\(L\), can one choose a nonzero vector in each prescribed plane
\(E_j=\operatorname{span}(e_j,f_j)\), so that their span \(K\) satisfies
\[
 \|P_L-P_K\|_{2\to2}\le 1/\sqrt2? \tag{G1}
\]
The notebook separately says that any constant \(c<1\), independent of
\(n\), would suffice. All projections here are Euclidean orthogonal
projections. Changing to a different, unquantified notion of distance would
change the problem.

\paragraph{Literature and historical limits.}
The searches in the accompanying log found no exact answer under the
speaker's name or the equivalent matrix formulation below. This is a
bounded negative retrieval result, not a proof that the problem has been
absent from the literature. A precise related result is
\citet[Theorem 3.4]{mehrmannpoloni2012doubling}: after symplectic coordinate
permutations, a Lagrangian has a graph matrix whose diagonal entries have
absolute value at most one and whose off-diagonal entries have absolute
value at most \(\sqrt2\). The theorem controls entries, and hence gives
dimension-dependent norm estimates; it does not state (G1) or a uniform
operator-gap estimate. We found no source explaining why Goldsheid's
specific question was dropped or whether it was later reformulated.

The results below are proved in this task, without a claim of historical
priority. They disprove (G1), show that any universal constant must be at
least \(\cos(\pi/8)\), and give an upper bound for every fixed dimension.
The still plausible weaker formulation is the existence of some
\(c\in[\cos(\pi/8),1)\) working in every dimension; no particular value
in that interval is asserted to work.

\begin{lemma}[Symmetric-unitary formulation]\ours
Identify \(\mathbb R^{2n}=\mathbb C^n\), with \(f_j=i e_j\), and let
\(Cz=\overline z\). Every real Lagrangian \(L\) is \(U\mathbb R^n\)
for a unitary matrix \(U\). Set \(S=UU^T\). Then \(S\) is symmetric
unitary and, as real-linear operators,
\[
 P_L=\tfrac12(I+SC).
\]
For \(K=D_\theta\mathbb R^n\), where
\(D_\theta=\operatorname{diag}(e^{i\theta_j})\),
\[
 \|P_L-P_K\|=\tfrac12\|S-D_\theta^2\|. \tag{G2}
\]
Consequently (G1) is equivalent to choosing a diagonal unitary \(D\)
such that \(\operatorname{Re}(S^*D)\) is positive semidefinite, where
\(\operatorname{Re}B=(B+B^*)/2\).
\end{lemma}
\begin{proof}
Choose a real orthonormal basis of \(L\). Its Hermitian Gram matrix is
the identity: the real part is the Euclidean Gram matrix and the imaginary
part vanishes because \(L\) is isotropic. Thus its columns form a unitary
\(U\). Reflection across \(L\) is the real-linear involution
\(UCU^*=(UU^T)C\). Averaging it with the identity gives the projection.
The same formula for \(D_\theta\) gives (G2), because conjugation is an
isometry and the norm of a complex-linear matrix is the same whether
regarded as a real or complex operator. Every diagonal unitary has a
diagonal unitary square root, so no choices are lost. Finally
\[
 (S-D)^*(S-D)=2I-S^*D-D^*S
             =2I-2\operatorname{Re}(S^*D),
\]
which proves the last equivalence.
\end{proof}

\begin{theorem}[An explicit six-plane refutes the proposed constant]\ours
Let \(\omega=e^{2\pi i/3}\) and
\[
 V=\frac1{\sqrt3}
 \begin{pmatrix}1&1&1\\1&\omega&\omega^2\\1&\omega^2&\omega\end{pmatrix}.
\]
The real six-dimensional subspace
\[
 L=\{(V\overline z,z):z\in\mathbb C^3\}\subset\mathbb C^6
       =\mathbb R^{12}
\]
is Lagrangian, and every coordinate-product Lagrangian \(K\) satisfies
\(\|P_L-P_K\|>1/\sqrt2\).
\end{theorem}
\begin{proof}
The Fourier matrix \(V\) is unitary. For \(z,w\in\mathbb C^3\), the
Hermitian inner product of the corresponding vectors in \(L\) is
\(\overline{z^*w}+z^*w\), which is real. The parametrization is real
linear and injective, so \(L\) is isotropic of half the ambient real
dimension and hence Lagrangian. An explicit real orthonormal basis is
\[
 \frac{(Ve_j,e_j)}{\sqrt2},\qquad
 \frac{(-iVe_j,ie_j)}{\sqrt2},\qquad j=1,2,3.
\]
Its reflection matrix is
\[
 S=\begin{pmatrix}0&V\\V^T&0\end{pmatrix}.
\]
Indeed \(S\) is symmetric unitary, and \(SC\) fixes exactly the displayed
plane: \(V^T\overline V=I\) supplies the second block equation.

Suppose a diagonal unitary \(D=\operatorname{diag}(d_1,\ldots,d_6)\)
made \(\|S-D\|\le\sqrt2\). The preceding lemma says
\(H=(S^*D+D^*S)/2\ge0\). Since every diagonal entry of \(S\) is zero,
every diagonal entry of \(H\) is zero. A positive semidefinite Hermitian
matrix obeys \(|H_{ij}|^2\le H_{ii}H_{jj}\), so \(H=0\).
For a nonzero entry \(S_{ij}\), this implies
\[
 d_i d_j=-\frac{S_{ij}}{\overline{S_{ij}}}. \tag{G3}
\]
The first row and first column of \(V\) are positive real. Their
equations force \(d_1=d_2=d_3=a\), \(d_4=d_5=d_6=b\), and \(ab=-1\).
The entry \(V_{22}=\omega/\sqrt3\) instead forces
\(ab=-\omega^2\), a contradiction. Thus the inequality cannot hold.
Compactness of the diagonal-unitary torus also shows that the minimum
gap for this one plane is strictly larger than \(1/\sqrt2\), not merely
that the infimum fails to be attained.
\end{proof}

\begin{theorem}[A Fourier-family obstruction to universal constants]\ours
For every odd integer \(m\ge3\), put
\(V_{jk}=m^{-1/2}e^{2\pi i jk/m}\), \(0\le j,k<m\), and form the
Lagrangian \(L_m=\{(V\overline z,z):z\in\mathbb C^m\}\subset\mathbb C^{2m}\).
Every coordinate-product Lagrangian \(K\) satisfies
\[
 \|P_{L_m}-P_K\|\ \geq\
 \frac12\sqrt{\,2+\sqrt{2-2/\sqrt m}\,}. \tag{G4}
\]
In particular any dimension-independent constant in Goldsheid's weaker
question must satisfy \(c\ge\cos(\pi/8)\).
\end{theorem}
\begin{proof}
As above, \(S=\bigl(\begin{smallmatrix}0&V\\V^T&0\end{smallmatrix}\bigr)\).
Write a diagonal unitary as \(D=\operatorname{diag}(A,B)\), with diagonal
entries \(a_j,b_k\). Then
\[
 \operatorname{Re}(S^*D)=\frac12
 \begin{pmatrix}0&C_0\\C_0^*&0\end{pmatrix},\qquad
 C_0=\overline V B+A^*V.
\]
The eigenvalues of the displayed off-diagonal Hermitian matrix are the
positive and negative singular values of \(C_0\), divided by two.
It follows exactly, without an estimate at this step, that
\[
 \|S-D\|^2=2+\|C_0\|. \tag{G5}
\]
Let \(W=(V_{jk}^2)\), the entrywise square, not the matrix square.
Because \(m\) is odd, multiplication by two permutes the residue classes
modulo \(m\). Consequently \(W\) is \(m^{-1/2}\) times a unitary
Fourier matrix with permuted columns, so \(\|W\|=m^{-1/2}\).
Expanding the Frobenius norm gives
\[
 \|C_0\|_F^2
 =2m+2\operatorname{Re}\sum_{j,k}\overline{a_j}\overline{b_k}V_{jk}^2
 \ge2m-2\sqrt m,
\]
because \(\|a\|=\|b\|=\sqrt m\) and
\(|a^*W\overline b|\le\|a\|\|W\|\|b\|=\sqrt m\).
There are \(m\) singular values, so
\(\|C_0\|\ge\|C_0\|_F/\sqrt m\ge\sqrt{2-2/\sqrt m}\).
Equations (G2) and (G5) prove (G4). Its right side tends to
\(\tfrac12\sqrt{2+\sqrt2}=\cos(\pi/8)\) as odd \(m\to\infty\).
Every \(c\) strictly below that limit is therefore defeated by some
finite odd \(m\). The endpoint itself is not excluded by this estimate.
\end{proof}

\begin{proposition}[A fixed-dimension upper bound]\ours
Every Lagrangian \(L\subset\mathbb C^n\) admits a common-phase product
plane \(K=e^{i\theta}\mathbb R^n\) with
\[
 \|P_L-P_K\|\le\cos\!\left(\frac{\pi}{2n}\right). \tag{G6}
\]
In particular the original bound is true for \(n=1,2\).
\end{proposition}
\begin{proof}
The eigenvalues of the unitary \(S\) lie on the unit circle. Of the \(n\)
circular gaps, including the wraparound gap and allowing repetitions,
one has length at least \(2\pi/n\). All eigenvalues therefore lie in a
closed complementary arc of length at most \(2\pi-2\pi/n\).
Let \(z=e^{i\phi}\) be its midpoint. Their angular distances from \(z\)
are at most \(\pi-\pi/n\), so, by the spectral theorem,
\[
 \|S-zI\|=\max_{\lambda\in\operatorname{Spec}S}|\lambda-z|
 \le2\sin\!\left(\frac\pi2-\frac{\pi}{2n}\right).
\]
Choose \(\theta=\phi/2\) and use (G2). If \(n=1\), the complementary
arc is a point and the bound is zero, as required.
\end{proof}

\begin{example}[Sharpness in dimension two]\ours
For \(S=\bigl(\begin{smallmatrix}0&1\\1&0\end{smallmatrix}\bigr)\),
each column of \(S-D\) has norm \(\sqrt2\) for every diagonal unitary
\(D\), so every product plane has gap at least \(1/\sqrt2\).
Taking \(D=iI\) gives equality, since \(S\) is self-adjoint unitary.
The associated Lagrangian is spanned over \(\mathbb R\) by
\((e_1+e_2)/\sqrt2\) and \((f_1-f_2)/\sqrt2\).
Thus (G6) is sharp at \(n=2\), although (G1) fails in higher dimension.
\end{example}

\begin{proposition}[The notebook's coordinatewise observation]\ours
For every coordinate plane \(E_j\), there is a unit vector
\(\varphi_j\in E_j\) with \(\operatorname{dist}(\varphi_j,L)\le1/\sqrt2\).
This does not bound the gap of the span of the chosen vectors.
\end{proposition}
\begin{proof}
Since \(L\) is Lagrangian, \(JL=L^\perp\), hence
\(P_L+JP_LJ^{-1}=I\). Compressing to the \(J\)-invariant plane \(E_j\)
and taking traces shows that the compression of \(P_L\) has trace one.
Its largest eigenvalue is at least \(1/2\), which proves the first
assertion by choosing a unit eigenvector.

For the second assertion, take \(n\ge2\),
\(v=(1,\ldots,1)/\sqrt n\in\mathbb R^n\), and
\(L=\mathbb Rv\oplus i v^\perp\). Its Hermitian pairings have zero
imaginary part and its real dimension is \(n\), so it is Lagrangian.
Every \(f_j=ie_j\) is at distance \(1/\sqrt n\) from \(L\), but
\(iv\in i\mathbb R^n\cap L^\perp\). Thus the span
\(K=i\mathbb R^n\) has \(\|P_L-P_K\|=1\). The equality follows also
from the fact that differences of orthogonal projections have norm at
most one. This example refutes an inference from a particular selection
of vectors, whereas the Fourier example refutes the existential (G1).
\end{proof}

\begin{remark}[Independent checks and numerical scope]\oursComputed
The independent mathematical audit re-derived the antiunitary reflection,
the positive-semidefinite zero-diagonal argument, the Fourier cross-term
estimate, the arc bound including \(n=1\), and the coordinatewise example.
It confirmed every hypothesis and emphasized that the endpoint
\(\cos(\pi/8)\) is not ruled out. A second representation of each Fourier
plane by the explicit real orthonormal basis checked isotropy directly.

The included \texttt{verify\_root.py} uses two independent random seeds,
five dimensions and four samples per dimension to compare the real
projection matrices with (G2); the maximum discrepancy was below
\(7\times10^{-16}\). It also checks Fourier examples for
\(m=3,5,7,11\). The rigorous lower bound for \(m=3\) is approximately
\(0.854312835963\), already above \(1/\sqrt2\). Two exploratory
optimizations found feasible gaps near \(0.866026\); neither met its
convergence criterion, and neither is reported as a global optimum.
The proofs of (G4)--(G6) do not depend on floating-point evidence. The
earlier random exploration, whose sampled gaps were all below
\(1/\sqrt2\), is also included: it illustrates why that evidence could
not justify the proposed universal bound.
\end{remark}

\paragraph{What remains and a bounded next investigation.}
Define
\[
 c_*:=\sup_{n\ge1}\ \sup_{S=S^T,\ S^*S=I}
       \inf_{D\ \mathrm{diagonal\ unitary}}\frac12\|S-D\|.
\]
The proofs establish \(\cos(\pi/8)\le c_*\le1\), with a strictly
subunit bound in each fixed dimension. Goldsheid's weaker question is
exactly whether \(c_*<1\). The original stronger assertion
\(c_*\le1/\sqrt2\) is false. A further bounded research run could
determine the exact optimum of the explicit \(m=3\) Fourier example by
an algebraic or interval certificate, or improve the Fourier-family
lower bound using higher singular-value moments. Either target is
concrete. A numerical optimizer alone would not decide the supremum over
all dimensions, and no mechanism for a uniform upper bound has been
proved here.

\subsection{Verification annotations for the spectral formulas}
\label{dev:1981-spectral-checks}

This subsection verifies corrections to established spectral theory. It
does not promote the Wodzicki exposition to a newly open C/D problem.

\begin{proposition}[Correct coefficients in the shifted-zeta identity]\ours
Let \(A\) be an invertible first-order elliptic differential operator on a
closed \(n\)-manifold, with a fixed admissible spectral cut. For a scalar
\(t\) sufficiently small that the cut and invertibility persist, put
\(\delta_\mu=\operatorname*{res}_{s=\mu}\zeta_A(s)\). For \(k\ge0\),
\[
\begin{split}
 \zeta_{A+t}(-k)={}&\sum_{\nu=0}^{k}\binom{k}{\nu}t^\nu\zeta_A(\nu-k)\\
 &+\sum_{\mu=1}^{n}
 \frac{(-1)^\mu\delta_\mu}{\mu\binom{k+\mu}{k}}\,t^{k+\mu}.
\end{split}
\tag{Z1}
\]
The second sum in P001025 (20) omits \(t^\mu\).
\end{proposition}
\begin{proof}
For sufficiently small \(t\) the complex-power binomial expansion,
initially in a trace-class half-plane, is
\[
 \zeta_{A+t}(s)=\sum_{j=0}^\infty
       \frac{(-1)^j(s)_j}{j!}t^j\zeta_A(s+j),
 \qquad (s)_j=s(s+1)\cdots(s+j-1).
\]
One may alternatively take a finite expansion with a remainder of
arbitrarily negative operator order, which is holomorphic near the point
under consideration. For \(0\le j\le k\), the coefficient at
\(s=-k\) is \(\binom{k}{j}\). For \(j=k+\mu\), \(\mu\ge1\),
the Pochhammer product has a simple zero and
\[
 \left.\frac{d}{ds}\frac{(-1)^{k+\mu}(s)_{k+\mu}}{(k+\mu)!}
 \right|_{s=-k}
 =\frac{(-1)^\mu k!(\mu-1)!}{(k+\mu)!}
 =\frac{(-1)^\mu}{\mu\binom{k+\mu}{k}}.
\]
It contributes only when \(\zeta_A(s+j)\) has a pole at \(\mu\).
For a differential operator of order one the positive integer poles are
confined to \(1,\ldots,n\), and the nonpositive integer values are
regular. Multiplying the zero by the pole gives (Z1). The factor \(t^j\)
has remained \(t^{k+\mu}\) throughout. The holomorphic remainder makes
no additional contribution, since its binomial coefficient vanishes at
\(-k\) once the expansion is taken beyond \(k+n\).
\end{proof}

\begin{remark}[Second check and the limit of the inference]\oursComputed
Direct differentiation of the finite Pochhammer product agrees with the
displayed coefficient in all 42 exact rational checks
\(0\le k\le6\), \(1\le\mu\le6\). Setting \(t=0\) independently
detects the notebook's missing powers, whenever the residue combination
is nonzero. Kernel-crossing cases are excluded from the proposition;
adding \(h(A)\) at zero without tracking which zero modes move is not a
substitute for that hypothesis.

The check also prevents an invalid inference about spectral cuts.
Cut independence at \(s=0\) for every shifted operator does not imply
cut independence at every negative integer. In the difference of (Z1)
for two cuts, the case \(k=0\) contains no negative values from which to
start such an induction. The general valid identity is instead
\citet[Theorem 4.2]{ponge2006spectral},
\(\rho(k)=(2\pi i/m)\operatorname{Res}(\Pi P^{-k})\).
Moreover, negative poles of eta and absolute zeta functions occur for
Dirac-type differential operators in odd dimensions
\citep[\S4]{loyamoroianu2012singularities}. For self-adjoint \(P\), the
two cut functions have the form
\(\zeta_+(s)+e^{\pm i\pi s}\zeta_-(s)\); their difference at an
integer is the limit of \(2i\sin(\pi s)\zeta_-(s)\).
A nonzero residue of \(\zeta_-\) therefore gives a nonzero difference,
even though each cut function is regular there. This is why the blanket
claim in P001025(d) is recorded as false without further hypotheses.
\end{remark}

\begin{theorem}[An explicit first-order differential operator with cut-dependent negative value]\ours
\label{dev:1981-negative-cut}
Let \(\mathbb T^3=(\mathbb R/2\pi\mathbb Z)^3\) have its flat metric,
let \(\mu>0\), and let \(\sigma_1,\sigma_2,\sigma_3\) be the Pauli
matrices. On the trivial Hermitian bundle of rank four set
\[
 \alpha_j=\sigma_1\otimes\sigma_j\quad(1\le j\le3),
 \qquad \beta=\sigma_3\otimes I_2,
 \qquad
 P=\sum_{j=1}^3\alpha_j(-i\partial_{x_j})+\mu\beta.
\]
Then \(P\) is an invertible self-adjoint elliptic differential operator
of order one and of Dirac type. Define \(Z_+(s)\) using argument
\(+\pi\) for negative eigenvalues, and \(Z_-(s)\) using argument
\(-\pi\); both branches give argument zero to positive eigenvalues.
The respective branch ranges are \((-\pi/2,3\pi/2)\) and
\((-3\pi/2,\pi/2)\), so the cuts are the negative and positive
imaginary axes. Both cut zeta functions are regular at \(s=-1\), but
\[
 \boxed{Z_+(-1)=-i\pi^2\mu^4,\qquad
        Z_-(-1)=+i\pi^2\mu^4.} \tag{Z2}
\]
In particular, their difference there is \(-2i\pi^2\mu^4\ne0\).
\end{theorem}
\begin{proof}
The four matrices \(\alpha_1,\alpha_2,\alpha_3,\beta\) are
Hermitian, square to the identity, and anticommute in distinct pairs.
Therefore
\[
 P^2=(-\Delta+\mu^2)I_4.
\]
The principal symbol squares to \(|\xi|^2I_4\), proving ellipticity
and the Dirac-type assertion. On the Fourier mode indexed by
\(k\in\mathbb Z^3\), the Hermitian matrix
\(\sum_j\alpha_jk_j+\mu\beta\) has square
\((|k|^2+\mu^2)I_4\) and trace zero. Its eigenvalues are
\(\pm\sqrt{|k|^2+\mu^2}\), each with multiplicity two. Fourier
expansion supplies the self-adjoint realization with domain \(H^1\),
and \(\mu>0\) excludes a kernel. Both stated cuts are admissible for
the operator and its principal symbol.

For \(\operatorname{Re}s>3\), put
\[
 Z_{\mathrm{abs}}(s)=\operatorname{Tr}|P|^{-s}
 =4\sum_{k\in\mathbb Z^3}(|k|^2+\mu^2)^{-s/2}.
\]
The multiplicities and branch conventions give exactly
\[
 Z_\pm(s)=\frac{1+e^{\mp i\pi s}}2\,Z_{\mathrm{abs}}(s). \tag{Z3}
\]
These identities continue meromorphically. To determine the pole of
\(Z_{\mathrm{abs}}\), use its Mellin representation
\[
 Z_{\mathrm{abs}}(s)
 =\frac1{\Gamma(s/2)}\int_0^\infty
 t^{s/2-1}\operatorname{Tr}(e^{-tP^2})\,dt.
\]
Gaussian Poisson summation on this torus yields
\[
\begin{split}
 \operatorname{Tr}(e^{-tP^2})
 &=4e^{-\mu^2t}\sum_{k\in\mathbb Z^3}e^{-t|k|^2}\\
 &=4\pi^{3/2}t^{-3/2}
 \left(1-\mu^2t+\frac{\mu^4t^2}{2}+O(t^3)\right)
   +O\!\left(t^{-3/2}e^{-c/t}\right)
\end{split}
\]
for some \(c>0\) as \(t\downarrow0\). The coefficient of
\(t^{1/2}\) is \(a=2\pi^{3/2}\mu^4\). The corresponding Mellin
integral contributes \(2a/(s+1)\) before division by
\(\Gamma(s/2)\). Since \(\Gamma(-1/2)=-2\sqrt\pi\),
\[
 \operatorname*{res}_{s=-1}Z_{\mathrm{abs}}(s)
 =\frac{2a}{\Gamma(-1/2)}=-2\pi\mu^4. \tag{Z4}
\]
After the displayed heat terms are subtracted, the small-time integral
is holomorphic near \(-1\); the large-time integral is entire because
\(\mu>0\). Thus the pole computation is justified without a formal
exchange of divergent series.

Each prefactor in (Z3) has a simple zero at \(-1\), with
\[
 \left.\frac{d}{ds}\frac{1+e^{\mp i\pi s}}2\right|_{s=-1}
 =\pm\frac{i\pi}{2}.
\]
Multiplying this derivative by the residue in (Z4) proves regularity
and the two values (Z2). The holomorphic part of
\(Z_{\mathrm{abs}}\) contributes zero at that point.
\end{proof}

\begin{remark}[Independent algebraic and normalization checks]\oursComputed
The separate executable \texttt{verify\_negative\_cut.py} checks all
Clifford anticommutators, Hermitian symmetry, the square and the exact
characteristic polynomial, the Mellin residue, and both branch limits.
The heat normalization was independently checked as
\(4\operatorname{vol}(\mathbb T^3)/(4\pi)^{3/2}=4\pi^{3/2}\).
A second mathematical audit reproduced the residue and phase signs.
The symmetric spectrum also gives \(\eta(P;s)=0\) identically:
this example uses the nonzero residue of the absolute zeta function.
It is an explicit verification of the notebook correction, with no
historical priority claim.
\end{remark}

\begin{remark}[The vector-calculus sign check]\oursComputed
For conventional Euclidean operators,
\(\operatorname{div}\operatorname{curl}=0\),
\(\operatorname{curl}\operatorname{grad}=0\), and
\(\operatorname{curl}\operatorname{curl}
=\operatorname{grad}\operatorname{div}-\Delta\).
Direct multiplication gives the corrected square in 1981-558.6 and
1981-559.12. An independent symbolic polynomial calculation replaces
the derivatives by three commuting indeterminates and reproduces both
the erroneous matrix's actual square and the corrected square exactly.
The output is included with the coefficient checks above.
\end{remark}

\section{Problem index}\label{problems:1981}
Earlier-year occurrence identifiers are retained as text; current-year identifiers link to their entries.

\subsection{GP-1981-01: Root-space dimensions and independent Lie words}\label{prob:GP-1981-01}
For the generalized Cartan matrices and maximal contragredient quotient specified in records 532--533, compute a prescribed root-space dimension and independent representing words. No efficiency or canonical-basis requirement is present in the notebook.
\par\noindent\textit{Occurrences:} \hyperref[item:1981-533.4]{1981-533.4}.
\par\noindent\textit{Status (2026-11):} As of November 2026: computationally determined degree by degree. Berman--Moody's 1979 formula covers symmetrizable matrices; the larger literal class is handled by finite recursive lowering maps as explained in the annotation. This is not recorded as an open demand for a uniform positive combinatorial formula. Earlier root-system and denominator occurrences are thematic predecessors, not exact previous formulations of this problem.

\subsection{GP-1981-02: Three classical double-coset spaces}\label{prob:GP-1981-02}
Describe \(U(n)\backslash GL_n(\mathbb C)/U(n)\), \(B\backslash GL_n(\mathbb C)/B\), and \(U(n)\backslash GL_n(\mathbb C)/B\), with \(U(n)\) unitary and \(B\) the invertible upper triangular subgroup; the unitary interpretation is explicitly reconstructed.
\par\noindent\textit{Occurrences:} \hyperref[item:1981-534.9]{1981-534.9}, \hyperref[item:1981-534.10]{1981-534.10}, \hyperref[item:1981-534.11]{1981-534.11}.
\par\noindent\textit{Status (2026-11):} As of November 2026: classical solved exercises. The invariants are respectively ordered positive singular values, a permutation in \(S_n\), and no invariant (one double coset). The three distinct questions share this problem entry. The earlier Bruhat occurrence 1979-483.1 is a related theorem, not an identical three-part question.

\subsection{GP-1981-03: Integer Gelfand–Tsetlin arrays with four fixed top entries}\label{prob:GP-1981-03}
Count integer triangular interlacing arrays with fixed weakly increasing top row \(m_1,m_2,m_3,m_4\), including equality at every interlacing bound.
\par\noindent\textit{Occurrences:} \hyperref[item:1981-536.7]{1981-536.7}.
\par\noindent\textit{Status (2026-11):} As of November 2026: solved by the type-A Weyl dimension product displayed in 1981-536.7 and the 1950 Gelfand--Tsetlin basis. The earlier 1976-415.5 and GP-1976-24 concern standard tableaux/Young-graph paths, a related but different branching graph; no exact recurrence of this fixed-integer-row question is present in the supplied index.

\subsection{GP-1981-04: The type-A quiver formula for p-adic GL(n) multiplicities}\label{prob:GP-1981-04}
For standard induced representations of \(GL_n(F)\) on a fixed cuspidal line, express their irreducible Jordan--H\"older multiplicities through local intersection cohomology of the corresponding graded nilpotent/type-A quiver orbit closures. The notebook occurrence is an abbreviated allusion, identified through Zelevinsky's 1981 paper.
\par\noindent\textit{Occurrences:} \hyperref[item:1981-535.6]{1981-535.6}.
\par\noindent\textit{Status (2026-11):} As of November 2026: proved in the specified GL(n) setting. Current direct sources are Solleveld's Theorem 5.4(b), his standard-module comparison Lemma 6.2, and Balodis's version-2 Theorem 3.9. First circulation and revision dates are separated in the occurrence. GP-1979-41 concerns the earlier Verma-module/Schubert Kazhdan--Lusztig problem; it is a related precursor, not this p-adic multiplicity question.

\subsection{GP-1979-41: Kazhdan--Lusztig multiplicity and cohomological conjectures}\label{prob:GP-1979-41}
For ordinary Verma modules of a finite-dimensional complex semisimple Lie algebra, identify Jordan--Hölder multiplicities with Kazhdan--Lusztig values, and determine the coefficient refinements from local intersection cohomology and the Jantzen filtration with its correct grading.
\par\noindent\textit{Occurrences:} 1979-483.6, 1979-484.4, \hyperref[item:1981-538.5]{1981-538.5}, \hyperref[item:1981-538.8]{1981-538.8}, \hyperref[item:1981-538.9]{1981-538.9}, \hyperref[item:1981-540.6]{1981-540.6}, \hyperref[item:1981-540.7]{1981-540.7}.
\par\noindent\textit{Status (2026-11):} As of November 2026, the ungraded theorem and geometric interpretation were already established before the November 1981 talk; the dominant-direction filtration refinement was proved in the 1993 Beilinson--Bernstein publication. Arbitrary-direction invariance is separately indexed as GP-1981-10, because it is false and has different quantifiers.

\subsection{GP-1981-10: Dependence of Jantzen-layer multiplicities on deformation direction}\label{prob:GP-1981-10}
For a fixed regular-integral Verma module, are the multiplicities in each layer of the Jantzen filtration independent of the regular direction used to deform its highest weight?
\par\noindent\textit{Occurrences:} \hyperref[item:1981-540.5]{1981-540.5}.
\par\noindent\textit{Status (2026-11):} As of November 2026: false without a dominant-cone restriction. Williamson, \S8.4, gives an explicit type-A3 example with a 13-dimensional weight space whose layer dimensions change from \((7,3,2,1)\) to \((7,2,4,0)\). First arXiv circulation was 8 October 2014; the journal volume is dated 2016 and its online release 7 July 2017. Dominant deformations satisfy the established Jantzen theorem. No unresolved residue is asserted for the literal seminar question.

\subsection{GP-1981-12: Primitivity of enveloping algebras of algebraic Lie algebras}\label{prob:GP-1981-12}
Characterize finite-dimensional algebraic Lie algebras over an algebraically closed field of characteristic zero whose universal enveloping algebra has a faithful simple module.
\par\noindent\textit{Occurrences:} 1964-37.1, 1964-37.2, \hyperref[item:1981-542.2]{1981-542.2}.
\par\noindent\textit{Status (2026-11):} As of November 2026: already settled by the 1976 Ooms criterion, quoted exactly in Elashvili, \S0. Primitivity is equivalent to an open coadjoint orbit, equivalently a nondegenerate Kirillov form. Elashvili\textquotesingle s 1982 paper further classifies specified reduced semidirect-product families. The earlier supplied index gives no existing GP identifier for 1964-37.1, so a new problem ID is assigned while retaining its earlier occurrence.

\subsection{GP-1981-13: Van der Pol relaxation limit and matching of slow branches}\label{prob:GP-1981-13}
Describe the limit as \(\varepsilon\downarrow0\) of the attracting periodic orbit of \(\varepsilon\ddot x+(x^2-1)\dot x+x=0\), including the connection of its slow branches by rapid transitions.
\par\noindent\textit{Occurrences:} \hyperref[item:1981-543.1]{1981-543.1}, \hyperref[item:1981-543.3]{1981-543.3}.
\par\noindent\textit{Status (2026-11):} As of November 2026: established in the classical relaxation setting. Cartier\textquotesingle s November 1981 Bourbaki exposition, printed in 1982, gives a nonstandard-analysis treatment and the matched cycle. The notebook does not specify an additional unresolved initial-value, rate, or canard problem; those are not silently added.

\subsection{GP-1981-20: Curvature actions with lower-order field equations}\label{prob:GP-1981-20}
Gelfand asks for curvature combinations whose metric Euler--Lagrange equations have lower order. A bounded precise version classifies natural symmetric divergence-free metric two-tensors of differential order at most two in fixed dimension and signature, modulo the separate issue of null Lagrangians. The unrestricted source wording is broader and lacks a specified class of actions.
\par\noindent\textit{Occurrences:} \hyperref[item:1981-546.4]{1981-546.4}.
\par\noindent\textit{Status (2026-11):} As of November 2026, the reconstructed natural-tensor problem is solved by Lovelock's pre-1982 classification; \citet[Theorem 2.6]{navarronavarro2011lovelock} gives a verified exact formulation. The handwritten Einstein-and-Pontryagin answer is not a complete all-dimensional classification. Earlier Einstein-action discussions are thematic antecedents, not this definite classification question.

\subsection{GP-1981-21: Convergence of the equations of Regge gravity}\label{prob:GP-1981-21}
For nondegenerate piecewise flat approximations of a smooth metric, determine a precise topology and geometric hypotheses under which the first variations or discrete Einstein equations converge. This must be separated from convergence of the curvature action and from convergence of actual solutions with fixed gauge and boundary data; the source specifies neither of the latter.
\par\noindent\textit{Occurrences:} \hyperref[item:1981-546.8]{1981-546.8}.
\par\noindent\textit{Status (2026-11):} As of November 2026, \citet[Theorem 4.1 and Corollary 4.3 of the preprint]{gawlikneunteufel2023einstein} establish a rigorous positive-degree interpolant result for the distributional densitized Einstein tensor in \(H^{-2}\). This is a partial resolution of a specified consistency problem. No general convergence theorem for arbitrary stationary Regge metrics was verified. Curvature-measure convergence alone does not imply convergence of variational equations.

\subsection{GP-1981-22: The Friedlander--Milnor comparison with finite coefficients}\label{prob:GP-1981-22}
For every finite-dimensional Lie group \(G\), prime \(p\), and degree \(q\geq0\), is \(H_q(BG^\delta;\mathbb F_p)\to H_q(BG;\mathbb F_p)\) an isomorphism? All coefficients are constant, and the topology on \(G\) on the right is its Lie topology. The corresponding statement for every finite abelian coefficient group is equivalent.
\par\noindent\textit{Occurrences:} \hyperref[item:1981-554.2]{1981-554.2}.
\par\noindent\textit{Status (2026-11):} As of November 2026, retain C: solvable identity components, finite-component split surjectivity, and specified low-degree cases are verified in \citet{milnor1983discrete}; stable developments are recorded in \citet{friedlander2008conjecture}. The renewed 2021 preprint claim \citep{amrani2015friedlander} was located but not certified by an independent proof check. The 2025 publisher extract \citep{dranishnikov2025generalized} still labels the general Lie statement a conjecture. No accepted general resolution was verified in this audit.

\subsection{GP-1981-23: Smooth modular-group coinvariants on the projective circle}\label{prob:GP-1981-23}
For every \(f\in C^\infty(\mathbb {RP}^1;\mathbb R)\), do there exist smooth \(u,v\) such that \(f(x)=u(x)-u(x+1)+v(x)-v(-1/x)\), with the two fractional-linear maps interpreted smoothly at infinity? The notebook also asserts density of these sums in the smooth topology; this is a separate claim.
\par\noindent\textit{Occurrences:} \hyperref[item:1981-555.1]{1981-555.1}, \hyperref[item:1981-555.2]{1981-555.2}.
\par\noindent\textit{Status (2026-11):} As of November 2026, this companion proves exact solvability in \(C^0\), with a uniform norm bound and an independent proof audit, and gives an exact involution-averaging reduction of the smooth problem. Neither a derivative estimate nor a proof of the notebook's smooth-density assertion is supplied. No exact later resolution of the \(C^\infty\) equation was verified; the located minimal-leaf and invariant-functional sources remain related leads.

\subsection{GP-1981-40: Uniform coordinate-product approximation of a Lagrangian plane}\label{prob:GP-1981-40}
For every Lagrangian \(L\subset\mathbb R^{2n}\), choose one line in each fixed coordinate symplectic plane so that their sum \(K\) satisfies \(\|P_L-P_K\|\le1/\sqrt2\). Alternatively, ask for any bound \(c<1\) independent of \(n\).
\par\noindent\textit{Occurrences:} \hyperref[item:1981-557.8]{1981-557.8}, \hyperref[item:1981-557.9]{1981-557.9}.
\par\noindent\textit{Status (2026-11):} As of 11 November 2026: the sharp \(1/\sqrt2\) assertion is disproved here by an explicit real Lagrangian six-plane. A Fourier family forces any universal constant to be at least \(\cos(\pi/8)\). A fixed-dimension upper bound is \(\cos(\pi/(2n))\); the weaker existence of a uniform constant below one is unresolved by this work and no exact later answer was found in the documented search.

\subsection{GP-1981-41: Geometric realization of compact-group weight multiplicities}\label{prob:GP-1981-41}
Express the weight multiplicity of an integral dominant representation as a quantization index, specifically a Riemann--Roch number, of the appropriate torus reduction of its prequantized coadjoint orbit.
\par\noindent\textit{Occurrences:} \hyperref[item:1981-556.1]{1981-556.1}, \hyperref[item:1981-556.6]{1981-556.6}.
\par\noindent\textit{Status (2026-11):} As of 11 November 2026: established in the compact Hamiltonian setting, including the specified treatment of singular quotients. The notebook supplies evidence of a preprint by March 1982, before the 1982 publication.

\subsection{GP-1981-42: Local extension of a two-form to an ambient symplectic form}\label{prob:GP-1981-42}
For a submanifold germ \(N\subset M\) with \(\dim M\) even, determine when a smooth two-form \(\beta\) on \(N\) is the pullback of a symplectic form near a specified point.
\par\noindent\textit{Occurrences:} \hyperref[item:1981-557.3]{1981-557.3}.
\par\noindent\textit{Status (2026-11):} As of 11 November 2026: solved locally by closedness and \(\operatorname{corank}\beta_x\le\operatorname{codim}N\); see Arnold--Givental, Chapter 2, \S1.3.

\subsection{GP-1981-43: Eta regularity without a bounding manifold}\label{prob:GP-1981-43}
For a self-adjoint classical elliptic operator of positive order on a closed manifold, is its eta function regular at zero without assuming any geometric filling?
\par\noindent\textit{Occurrences:} 1972-296.2, 1972-296.3, 1972-296.5, \hyperref[item:1981-559.9]{1981-559.9}.
\par\noindent\textit{Status (2026-11):} As of 11 November 2026: established, in all dimensions, by Gilkey 1981 Theorem 0.1 after the earlier odd-dimensional APS theorem. The notebook’s possible reference to a more general nonself-adjoint setting is not covered by this formulation.

\subsection{GP-1981-44: Spectral zeta constructions without self-adjointness}\label{prob:GP-1981-44}
Can complex powers, spectral zeta functions and a comparison measuring spectral asymmetry be defined for a nonself-adjoint elliptic operator on a closed manifold?
\par\noindent\textit{Occurrences:} \hyperref[item:1981-559.8]{1981-559.8}.
\par\noindent\textit{Status (2026-11):} As of 11 November 2026: yes with spectral cuts admissible for both the principal symbol and the operator, together with resolvent hypotheses. This does not automatically establish the notebook’s separate normal-operator boundary index formula.

\section{Corrections to the transcription and translation}\label{corr:1981}
The evidence distinguishes a transcription correction from a mathematical emendation of wording faithfully preserved in the notebook.
\subsection*{P000981 (both; provenance): record 530/531 boundary}
\noindent\textit{Transcription reads.} Catalog range of record 530 extends to P000982\par
\noindent\textit{Critical reading.} Only the first dated heading belongs to 530; mathematical content after 21/IX belongs to 531\par
\noindent\textit{Evidence.} Actual scan P000981 has the new dated heading before the reading list and Lie-theory exposition\par
\subsection*{P000981 (both; ambiguity): trace formula for Cartan/Killing form}
\noindent\textit{Transcription reads.} \(\operatorname{tr}(\operatorname{ad}x-\operatorname{ad}y)\)\par
\noindent\textit{Critical reading.} \(\operatorname{tr}(\operatorname{ad}x\circ\operatorname{ad}y)\)\par
\noindent\textit{Evidence.} Russian and scan show a short separating stroke; standard formula in Macdonald 1981 p259 fixes the intended operation but the glyph is ambiguous\par
\subsection*{P000981-P000983 (both; ambiguity): root-vector bracket convention}
\noindent\textit{Transcription reads.} \([g_\alpha,h]=\alpha(h)g_\alpha\)\par
\noindent\textit{Critical reading.} \([h,x_\alpha]=\alpha(h)x_\alpha\), or reverse every root sign consistently\par
\noindent\textit{Evidence.} Scan/Russian order conflicts with the subsequent plus-alpha weight-raising convention\par
\subsection*{P000982 (both; faithful error): Weyl group of sl(n)}
\noindent\textit{Transcription reads.} symmetric group of order n\par
\noindent\textit{Critical reading.} symmetric group on n letters, of order n!\par
\noindent\textit{Evidence.} Russian \ru{n-го порядка} and scan; standard permutation action gives \(n!\)\par
\subsection*{P000982 (both; ambiguity): sl3 root diagram}
\noindent\textit{Transcription reads.} three roots at 120 degrees with sum zero\par
\noindent\textit{Critical reading.} three cyclic roots drawn; full A2 root system has six roots and a simple base has two\par
\noindent\textit{Evidence.} Scan agrees with the three-vector drawing; complete roots are \(\pm(e_i-e_j)\)\par
\subsection*{P000982-P000983 (both; faithful error): Cartan matrix restrictions}
\noindent\textit{Transcription reads.} \(a_{ij}\) are nonpositive; matrix positive definite\par
\noindent\textit{Critical reading.} \(a_{ii}=2\); \(a_{ij}\le0\) for \(i\ne j\); a positive diagonal symmetrization is positive definite\par
\noindent\textit{Evidence.} Source omits off-diagonal restriction; finite Cartan matrices need not be symmetric\par
\subsection*{P000983 (both; ambiguity): Serre exponents}
\noindent\textit{Transcription reads.} \((\operatorname{ad}e_i)^{1-a_{ij}}e_j=0\) without index restriction\par
\noindent\textit{Critical reading.} The Serre relations are imposed only for \(i\ne j\)\par
\noindent\textit{Evidence.} At \(i=j\) the displayed exponent would be \(-1\); Serre VI.4 Theorem 6 explicitly excludes diagonal indices\par
\subsection*{P000983 (both; faithful error): lattice of representation weights}
\noindent\textit{Transcription reads.} lattice generated by roots (coroots)\par
\noindent\textit{Critical reading.} weight lattice P; each irreducible weight set lies in a coset of root lattice Q\par
\noindent\textit{Evidence.} The sl2 two-dimensional representation has weights plus/minus alpha/2; Macdonald p260 distinguishes Q and P\par
\subsection*{P000983 (both; ambiguity): raising-operator subscript}
\noindent\textit{Transcription reads.} \(e_\alpha v_\alpha\) in a calculation about \(v_\lambda\)\par
\noindent\textit{Critical reading.} \(e_\alpha v_\lambda\)\par
\noindent\textit{Evidence.} The adjacent formula and the derivation \(h(ev)=(\lambda+\alpha)(h)ev\) determine the subscript\par
\subsection*{P000987 (both; ambiguity): terminal symbol after closure definition}
\noindent\textit{Transcription reads.} \(S^{\mathrm{cl}}\) = intersection of closed supersets = \((0)\)\par
\noindent\textit{Critical reading.} Keep the closure formula; use \(0=\operatorname{cl}(\varnothing)=\operatorname{Fix}(G)\) only for the bottom closed set\par
\noindent\textit{Evidence.} The final (0) is unexplained in the scan and cannot literally equal every closure; later intersection axioms require a bottom element\par
\subsection*{P000988 (both; faithful error): singleton closure under transitivity}
\noindent\textit{Transcription reads.} points with the same stabilizer as x\par
\noindent\textit{Critical reading.} points \(y\) whose stabilizers contain \(G_x\); equality needs additional hypotheses\par
\noindent\textit{Evidence.} Russian and scan explicitly say same; infinite transitive affine-group coset action gives strict containment\par
\subsection*{P000988 (both; faithful error): non-singleton closure obstructs 2-transitivity}
\noindent\textit{Transcription reads.} \(\operatorname{cl}\{x\}\) properly contains \(\{x\}\) implies not doubly transitive\par
\noindent\textit{Critical reading.} Require at least three points, or \(\operatorname{cl}\{x\}\) to be a proper non-singleton subset of \(X\)\par
\noindent\textit{Evidence.} The regular degree-two action is 2-transitive and \(\operatorname{cl}\{x\}=X\)\par
\subsection*{P000988 (both; ambiguity): axiom(iv) subgroup}
\noindent\textit{Transcription reads.} \(G^S\cap G^{\operatorname{cl}\{x\}}\)\par
\noindent\textit{Critical reading.} Both superscript stabilizers are setwise; do not replace the second with pointwise fixation\par
\noindent\textit{Evidence.} Actual superscripts and the preceding restriction diagram distinguish setwise \(G^S\) from pointwise \(G_S\)\par
\subsection*{P000989 (both; ambiguity): double-coset exercises U(n), GL(n)}
\noindent\textit{Transcription reads.} U(n), GL(n), with no definitions\par
\noindent\textit{Critical reading.} For the answers supplied, read \(U(n)\) as unitary and \(GL(n)\) as \(GL_n(\mathbb C)\)\par
\noindent\textit{Evidence.} The source specifies only B as triangular; U may otherwise mean a unipotent subgroup\par
\subsection*{P000985 (both; faithful error): nondegenerate A and central quotient}
\noindent\textit{Transcription reads.} if A is nondegenerate the operation is not needed\par
\noindent\textit{Critical reading.} No added derivations are needed if one starts minimally; in the displayed enlarged presentation the central quotient remains necessary\par
\noindent\textit{Evidence.} \(h_i-\sum_j a_{ij}\epsilon_j\) are central even when \(A\) is invertible; \(A=(2)\) gives \(h-2\epsilon\)\par
\subsection*{P000985 (both; ambiguity): example following (2)}
\noindent\textit{Transcription reads.} (2), matrices of order 2 with trace 0\par
\noindent\textit{Critical reading.} The resulting Lie algebra for Cartan matrix (2) is sl2, consisting of trace-zero 2 by 2 matrices\par
\noindent\textit{Evidence.} Layout in scan separates Cartan example and explanation of its Lie algebra\par
\subsection*{P000985 (both; faithful error): affine alternative II}
\noindent\textit{Transcription reads.} \(Au\ge0\) implies \(u=0\), while there exists \(u>0\) with \(Au=0\)\par
\noindent\textit{Critical reading.} \(Au\ge0\) implies \(Au=0\); the kernel is generated by a positive vector\par
\noindent\textit{Evidence.} Both contradictory lines appear in the scan; left positive null vector proves the corrected implication\par
\subsection*{P000986 (both; ambiguity): scope of Verma Hom theorem}
\noindent\textit{Transcription reads.} For G(A), Hom is zero or C without further restriction\par
\noindent\textit{Critical reading.} The cited BGG/Verma theorem is verified for finite-dimensional complex semisimple g; extra hypotheses are required before enlarging the scope\par
\noindent\textit{Evidence.} BGG 1971 explicitly fixes finite-dimensional complex semisimple g in section2\par
\subsection*{P000986 (both; faithful error): Weyl-orbit and linkage rho shifts}
\noindent\textit{Transcription reads.} \(\lambda-\rho=w(\mu-\rho)\), and the same minus shift in the chain\par
\noindent\textit{Critical reading.} \(\lambda+\rho=w(\mu+\rho)\) under \(hv_0=\lambda(h)v_0\); use dot-action chain\par
\noindent\textit{Evidence.} Both minus signs are visible in the scan; sl2 embedding M(-2) into M(0) gives a direct counterexample\par
\subsection*{P000986 (both; ambiguity): reflection-chain indexing}
\noindent\textit{Transcription reads.} \(\delta_{i-1}\) appears in a chain indexed \(\delta_1,\ldots,\delta_k\)\par
\noindent\textit{Critical reading.} Use \(\nu_j=s_{\beta_j}(\nu_{j-1}+\rho)-\rho\), with positive integral step at each \(j\)\par
\noindent\textit{Evidence.} Source has an unbound boundary index; BGG condition(A) supplies a consistent chain\par
\subsection*{P000986 (both; ambiguity): rho as half-sum after the passage to G(A)}
\noindent\textit{Transcription reads.} rho is the half-sum of all positive roots\par
\noindent\textit{Critical reading.} For infinite Kac--Moody systems use \(\rho(h_i)=1\) in a realization; a literal infinite half-sum is not defined\par
\noindent\textit{Evidence.} The earlier finite-dimensional half-sum convention is carried into a broader setting without qualification\par
\subsection*{P000989 (both; faithful error): last arrow of the \(A_m\) chain}
\noindent\textit{Transcription reads.} last arrow \(A_m\) into \(V_m\)\par
\noindent\textit{Critical reading.} \(A_{m-1}:V_{m-1}\to V_m\)\par
\noindent\textit{Evidence.} The scan really labels the last arrow \(A_m\); an \(A_m\) quiver has \(m\) vertices and \(m-1\) adjacent arrows\par
\subsection*{P000990 (both; ambiguity): Schubert-cell analogy}
\noindent\textit{Transcription reads.} orbits analogous to Schubert cells\par
\noindent\textit{Critical reading.} Keep as an analogy; exact later theorem concerns open Schubert patches of orbit closures\par
\noindent\textit{Evidence.} Rank-one orbit for a single 1 by 1 matrix is C*, not an affine Schubert cell; Zelevinsky1985 Theorem1 supplies the precise relationship\par
\subsection*{P000990 (both; ambiguity): p-adic polynomials remark}
\noindent\textit{Transcription reads.} polynomials describe representations of p-adic groups\par
\noindent\textit{Critical reading.} Interpret as the identified GL(n) multiplicity conjecture with local IC stalks, not a proved general statement about arbitrary p-adic groups\par
\noindent\textit{Evidence.} Zelevinsky1981 Introduction and section1.9 explicitly label the formula a conjecture; the notebook omits group and formula\par
\subsection*{P000990 (both; ambiguity): dimension-table question marks}
\noindent\textit{Transcription reads.} ? at \((a,b)=(3,2),(2,3)\)\par
\noindent\textit{Critical reading.} Both dimensions are infinite\par
\noindent\textit{Evidence.} ab=6 is indefinite rank-two type; scanned question marks are faithful omissions, not evidence of an open classification\par
\subsection*{P000990 (both; ambiguity): Catalan parenthesization example}
\noindent\textit{Transcription reads.} parentheses in x x ... x, number of factors omitted\par
\noindent\textit{Critical reading.} For the nth Catalan number the product has n+1 factors\par
\noindent\textit{Evidence.} \(d(n,n+1)=\binom{2n}{n}/(n+1)\), which counts full binary bracketings of \(n+1\) factors\par
\subsection*{P000991 (both; ambiguity): Hilbert-series coefficient convention}
\noindent\textit{Transcription reads.} \(p(t)=\sum_k c_k t^k/k\), \(C_k=\sum_{d\mid k}c_d\)\par
\noindent\textit{Critical reading.} Keep both formulas; \(c_k=k\dim V_k\)\par
\noindent\textit{Evidence.} The actual scan explicitly has denominator k, and the Euler-product logarithm checks this normalization\par
\subsection*{P000991 (both; faithful error): parenthetical interpretation of interlacing arrays}
\noindent\textit{Transcription reads.} number of representations of SL(4) with a given highest weight\par
\noindent\textit{Critical reading.} dimension, or number of Gelfand-Tsetlin basis vectors, of the one irreducible module with that highest weight\par
\noindent\textit{Evidence.} Russian text and scan say number of representations; fixed dominant integral highest weight determines one irreducible isomorphism class\par
\subsection*{P000992-P000993 (both; ambiguity): complex tangent dimension (0,2)}
\noindent\textit{Transcription reads.} four real equations imply complex tangent dimension (0,2)\par
\noindent\textit{Critical reading.} Assume invertibility of the even complex Jacobian \(f_x\) on the chosen chart\par
\noindent\textit{Evidence.} The written elimination of dx through dtheta requires this extra matrix-rank condition; real regularity alone is weaker\par
\subsection*{P000993 (both; ambiguity): the symbol L}
\noindent\textit{Transcription reads.} L first denotes all holomorphic transformations and later the action symmetry\par
\noindent\textit{Critical reading.} For the action use the Berezin-volume-preserving subpseudogroup \(\operatorname{Ber}J=1\)\par
\noindent\textit{Evidence.} The source explicitly narrows the group before reusing L; read with Schwarz1982 section3\par
\subsection*{P000993 (both; ambiguity): surface–superfield correspondence}
\noindent\textit{Transcription reads.} surfaces and superfields are in one-to-one correspondence\par
\noindent\textit{Critical reading.} Locally, on charts where the projection to \((\xi,\theta,\bar\theta)\) is invertible\par
\noindent\textit{Evidence.} Graph form is a local implicit-function statement; source omits the projection qualification\par
\subsection*{P000993-P000994 (both; faithful error): Levi density exponent}
\noindent\textit{Transcription reads.} \(A=|\det\Gamma|^{-1/2}\)\par
\noindent\textit{Critical reading.} In the natural odd-coordinate frame use \(|\det f_x|^{4/3}|\det\Gamma|^{-1/3}\)\par
\noindent\textit{Evidence.} Both formulas appear in scans; P994 explicitly says one must then write the second, agreeing with Schwarz1982 section3 equations7–8\par
\subsection*{P000994 (both; ambiguity): fractional power of first determinant}
\noindent\textit{Transcription reads.} \((\det(\partial f^k/\partial z^a))^{4/3}\) without absolute-value or branch convention\par
\noindent\textit{Critical reading.} Use \(|\det f_x|^{4/3}\) for the real density, or specify the equivalent oriented branch\par
\noindent\textit{Evidence.} Published Schwarz1982 equation8 uses absolute values for W and Gamma; scanned first determinant has no visible absolute-value bars\par
\subsection*{P000994 (both; ambiguity): Yadernaya Fizika 1981 N=10}
\noindent\textit{Transcription reads.} Yadernaya Fizika, 1981, N=10\par
\noindent\textit{Critical reading.} October 1981: volume34, issue4, pages1144–1149; English translation34(4),635–637\par
\noindent\textit{Evidence.} Author bibliography identifies the precise article; the numeral10 denotes the month, not volume34 issue10\par
\subsection*{P000994 (both; provenance): Nuclear Physics 1981 allusion}
\noindent\textit{Transcription reads.} Nuclear Physics, 1981, preliminary large work\par
\noindent\textit{Critical reading.} Leave the exact reference unresolved; related field-space paper is Nuclear Physics B171(1980),154–166\par
\noindent\textit{Evidence.} The only securely identified matching Schwarz title in author bibliography is dated1980; do not silently substitute it for the unidentified1981 reference\par
\subsection*{P000995 (both; ambiguity): simple-generator diagram}
\noindent\textit{Transcription reads.} \(e_1,\ldots,e_n\) bracketed as \(\mathfrak n_+\)\par
\noindent\textit{Critical reading.} Read as Lie generators of \(\mathfrak n_+\), not a vector-space basis.\par
\noindent\textit{Evidence.} Russian and scan show only the \(n\) simple generators; \(n\) is the rank. Root spaces for nonsimple positive roots are additionally needed.\par
\subsection*{P000995 (both; ambiguity): first highest-vector definition}
\noindent\textit{Transcription reads.} \(E_jx=0\) defines highest-weight vector\par
\noindent\textit{Critical reading.} Also require a common Cartan eigenvector; or call the kernel a space of singular vectors.\par
\noindent\textit{Evidence.} Scan first gives only annihilation, then later explicitly writes \(hv=\chi(h)v\).\par
\subsection*{P000995 (both; faithful error): after finite composition series}
\noindent\textit{Transcription reads.} Any subquotient has this form\par
\noindent\textit{Critical reading.} Every simple subquotient / composition factor has the form \(L(\psi)\).\par
\noindent\textit{Evidence.} Russian wording любой подфактор and scan agree; an arbitrary subquotient need not be simple.\par
\subsection*{P000995 (both; faithful error): multiplicity question}
\noindent\textit{Transcription reads.} Multiplicity of \(M(\psi)\) in a composition series\par
\noindent\textit{Critical reading.} Multiplicity of \(L(\psi)\).\par
\noindent\textit{Evidence.} Scan confirms \(M(\psi)\); the very next multiplicity formula uses \(L(\psi)\).\par
\subsection*{P000996 (both; faithful error): dominant integral condition}
\noindent\textit{Transcription reads.} Condition on \(\lambda\) after defining \(\chi\)\par
\noindent\textit{Critical reading.} Use the same weight variable \(\chi\), or explicitly rename it \(\lambda\).\par
\noindent\textit{Evidence.} Russian, English and scan contain the variable switch; this is a notational slip.\par
\subsection*{P000996 (both; ambiguity): category O marginal correction}
\noindent\textit{Transcription reads.} Here only \(\mathfrak n_+\), and not \(\mathfrak n_-\), but this makes no difference\par
\noindent\textit{Critical reading.} Changing the Borel gives an equivalent convention; the two local-nilpotence conditions are not equivalent on a fixed module.\par
\noindent\textit{Evidence.} Scan confirms both signs and margin wording. Generic Verma modules distinguish the two conditions.\par
\subsection*{P000996 (both; ambiguity): dual module construction}
\noindent\textit{Transcription reads.} The dual representation / Cartan-finite functions\par
\noindent\textit{Critical reading.} Use the restricted weight-space dual and a Chevalley anti-involution to preserve the chosen category O.\par
\noindent\textit{Evidence.} The manuscript singles out Cartan-finite functionals but does not state the action or twisting.\par
\subsection*{P000996-P000997 (both; faithful error): bilinear form orthogonality}
\noindent\textit{Transcription reads.} Different unordered sets of generator indices give zero\par
\noindent\textit{Critical reading.} Different total root weights; for simple-root words this means different multisets of indices.\par
\noindent\textit{Evidence.} Discarding repetitions loses weight data. The displayed shorthand is faithfully copied.\par
\subsection*{P000996 (both; ambiguity): bilinear form numerator}
\noindent\textit{Transcription reads.} e carrying the same index sequence as the lowering monomial\par
\noindent\textit{Critical reading.} The anti-involution reverses product order: \(e_{i_k}\cdots e_{i_1}\).\par
\noindent\textit{Evidence.} Scan has compressed multiindices, not a clear ordered product; modern contravariance fixes the interpretation.\par
\subsection*{P000997 (both; faithful error): first line}
\noindent\textit{Transcription reads.} Form invariant with respect to the Lie algebra\par
\noindent\textit{Critical reading.} Contravariant under the Chevalley anti-involution.\par
\noindent\textit{Evidence.} Russian reads инвар. отн. алг. Ли; ordinary bilinear invariance is different.\par
\subsection*{P000997 (both; ambiguity): PBW identification}
\noindent\textit{Transcription reads.} Identify all \(M(\chi)\)\par
\noindent\textit{Critical reading.} Identify underlying PBW vector spaces, not g-module structures.\par
\noindent\textit{Evidence.} The monomial map in the scan fixes negative-root vectors and changes highest weight.\par
\subsection*{P000997 (both; faithful error): definition of \(M_i\)}
\noindent\textit{Transcription reads.} Fixed \(x\) has pairing \(O(t^i)\) against every fixed y\par
\noindent\textit{Critical reading.} There exists a parameter-dependent lift of \(x\) whose pairing with the deformed lattice is divisible by \(t^i\).\par
\noindent\textit{Evidence.} Scan confirms fixed-vector shorthand. The standard Jantzen lattice definition requires moving lifts.\par
\subsection*{P000997 (both; faithful error): change \(\rho\) question}
\noindent\textit{Transcription reads.} Apparently layer multiplicities do not change\par
\noindent\textit{Critical reading.} True under appropriate dominant deformation hypotheses; false for all regular directions.\par
\noindent\textit{Evidence.} Question is explicit in Russian and scan. Williamson section 8.4 supplies a type-\(A_3\) counterexample.\par
\subsection*{P000997 (both; faithful error): polynomial multiplicity bracket}
\noindent\textit{Transcription reads.} \([L_\psi:M_i/M_{i+1}]\)\par
\noindent\textit{Critical reading.} \([M_i/M_{i+1}:L_\psi]\) in standard composition-multiplicity notation.\par
\noindent\textit{Evidence.} Scan specifically confirms the reversed order.\par
\subsection*{P000997 (both; faithful error): polynomial grading}
\noindent\textit{Transcription reads.} \(P(q)=\sum_iq^i\) times multiplicity in layer i\par
\noindent\textit{Critical reading.} Layer j corresponds to \(q^{(\ell(y)-\ell(x)-j)/2}\).\par
\noindent\textit{Evidence.} Even \(\mathfrak{sl}_2\) has \(P_{e,s}=1\) while the lower constituent lies in layer 1; this is a faithful normalization error.\par
\subsection*{P000997 (both; faithful error): Schubert-cell cohomology}
\noindent\textit{Transcription reads.} \(H^i\) of one Schubert cell near another; cell contained in other cell\par
\noindent\textit{Critical reading.} Even-degree stalk intersection cohomology of a Schubert closure at a smaller cell.\par
\noindent\textit{Evidence.} Russian and scan omit intersection and closure. Distinct cells are disjoint; Lusztig Theorem 4.1 fixes the interpretation.\par
\subsection*{P000997 (both; faithful error): weight-space annotation}
\noindent\textit{Transcription reads.} \(L(\chi)_\lambda\) is a weight space of weight \(\chi\)\par
\noindent\textit{Critical reading.} Weight \(\lambda\).\par
\noindent\textit{Evidence.} The scan confirms \(\chi\) in the annotation, while the subscript and next-page formula use \(\lambda\).\par
\subsection*{P000997-P000998 (both; faithful error): alternant character formula}
\noindent\textit{Transcription reads.} \(\chi=(a_1,\ldots,a_n)\) with denominator powers \(0,\ldots,n-1\)\par
\noindent\textit{Critical reading.} Use shifted exponents \(a_i=\lambda_{n+1-i}+i-1\) for descending dominant \(\lambda\).\par
\noindent\textit{Evidence.} The trivial representation already disproves unshifted repeated zero exponents. Source determinants and statement were visually inspected.\par
\subsection*{P000998 (both; faithful error): affine-polynomial formula}
\noindent\textit{Transcription reads.} \(m(\lambda,\chi)=P_{n_\lambda,n_\chi}(1)\), after \(m(\chi,\lambda)=\dim L(\chi)_\lambda\)\par
\noindent\textit{Critical reading.} Keep highest weight and weight argument order consistent; \(\dim L(\lambda)_\mu=P_{n_\mu,n_\lambda}(1)\).\par
\noindent\textit{Evidence.} The manuscript reverses m arguments. Lusztig 1983 Theorem 6.1 gives the exact intended formula.\par
\subsection*{P000998 (both; ambiguity): Kac--Moody conjecture}
\noindent\textit{Transcription reads.} The same for all Kac--Moody algebras; crossed-out character expression\par
\noindent\textit{Critical reading.} Retain as an unidentified conjectural extension; do not restore the erased expression as an established theorem.\par
\noindent\textit{Evidence.} Russian and scan show the full formula crossed out and leave the indices undefined.\par
\subsection*{P000998 (both; provenance): Jacobson attribution}
\noindent\textit{Transcription reads.} Last year below Jacobson\par
\noindent\textit{Critical reading.} Do not infer a first proposal in 1980.\par
\noindent\textit{Evidence.} The question occurs in supplied 1964-37.1; Elashvili 1982 cites Jacobson and explicitly dates Ooms answer 1976. The marginal phrase may refer to an unrecorded occasion.\par
\subsection*{P000998 (both; faithful error): excluded algebras}
\noindent\textit{Transcription reads.} All semisimple and nilpotent algebras unsuitable\par
\noindent\textit{Critical reading.} All nonzero such algebras; the zero algebra has enveloping algebra \(\mathbb C\).\par
\noindent\textit{Evidence.} Faithful wording needs this elementary edge case.\par
\subsection*{P000998 (both; ambiguity): affine example}
\noindent\textit{Transcription reads.} Last column zero\par
\noindent\textit{Critical reading.} Retain: a valid row/column convention, isomorphic via minus transpose to the customary last-row-zero realization.\par
\noindent\textit{Evidence.} Scan confirms column. This is an ambiguity resolution, not a transcription correction.\par
\subsection*{P000998 (both; faithful error): four-dimensional irreducible representation}
\noindent\textit{Transcription reads.} Unique four-dimensional irreducible \(\mathfrak{gl}_2\) representation\par
\noindent\textit{Critical reading.} Unique \(\mathfrak{sl}_2\) type \(\operatorname{Sym}^3(\mathbb C^2)\), with a specified nonzero scalar-center action.\par
\noindent\textit{Evidence.} The center acts by a free scalar on irreducible \(\mathfrak{gl}_2\) modules. Zero scalar makes the semidirect product non-Frobenius.\par
\subsection*{P000998 (both; ambiguity): classification theorem}
\noindent\textit{Transcription reads.} \(R\) simple\par
\noindent\textit{Critical reading.} The matching publication uses \(R\) reductive and/or \([R,R]\) simple, plus reduction and specified radical conditions.\par
\noindent\textit{Evidence.} Scan has a cancelled word before simple. The paper received 30 November 1981 has a broader ten-family theorem; exact recovery of the erased word is impossible.\par
\subsection*{P000999 (both; provenance): Bourbaki heading}
\noindent\textit{Transcription reads.} Reeb talk, November 1981; Cartier survey\par
\noindent\textit{Critical reading.} Cartier was the Bourbaki speaker reporting work of Reeb and his school.\par
\noindent\textit{Evidence.} Original exposé 580 title page reads par Pierre Cartier, Novembre 1981. Manuscript attribution is retained diplomatically.\par
\subsection*{P000999 (both; ambiguity): separated reduced equation}
\noindent\textit{Transcription reads.} Divide by \(x\) without exclusions\par
\noindent\textit{Critical reading.} Work away from \(x=0\) and \(x=\pm1\); restore the solution \(x=0\).\par
\noindent\textit{Evidence.} Direct differentiation confirms the displayed integral but division loses the equilibrium.\par
\subsection*{P000999 (both; faithful error): probability codomain}
\noindent\textit{Transcription reads.} \(p:2^\Omega\to(0,1)\)\par
\noindent\textit{Critical reading.} For the intended Cartier construction: \(p:2^\Omega\to\{0,1\}\), with \(p(\Omega)=1\).\par
\noindent\textit{Evidence.} The scan truly uses round brackets. Finite additivity forces \(p(\varnothing)=0\), equality requires 1, and Cartier section 2 assumptions (a)-(c) explicitly require zero-one values.\par
\subsection*{P000999 (both; ambiguity): equivalence probability symbol}
\noindent\textit{Transcription reads.} Capital \(P\) in \(P(x^{\prime}=x^{\prime\prime})=1\) after lowercase \(p\) definition\par
\noindent\textit{Critical reading.} Read \(P\) as the same finitely additive \(p\), absent a newly defined probability.\par
\noindent\textit{Evidence.} Scan confirms capital \(P\). This is a harmless notational switch.\par
\subsection*{P001000 (both; provenance): Kruskal numbers heading}
\noindent\textit{Transcription reads.} Kruskal numbers, simplest field extending \(R\)\par
\noindent\textit{Critical reading.} Identify as surreal sign expansions; retain Kruskal as the notebook attribution and distinguish proper-class No from set-sized ultrapowers.\par
\noindent\textit{Evidence.} Visible arrows match surreal arithmetic; Conway 1976 predates the talk, and Gonshor gives the systematic sign-sequence treatment.\par
\subsection*{P001000 (both; faithful error): finite-arrow explanatory rule}
\noindent\textit{Transcription reads.} Step halves when direction changes\par
\noindent\textit{Critical reading.} After the first change, halve the step at every subsequent sign.\par
\noindent\textit{Evidence.} Scan confirms the conditional wording. The sequence \(+\!-\!-\) equals \(\tfrac14\), not 0.\par
\subsection*{P001000 (both; ambiguity): infinite-arrow notation}
\noindent\textit{Transcription reads.} Infinity, double-hatted up, \(\varepsilon\)\par
\noindent\textit{Critical reading.} Use \(\omega\) for an \(\omega\)-long positive block, \(\omega^2\) for an \(\omega^2\)-long block, and \(\omega^{-1}\) for + followed by \(\omega\) minuses.\par
\noindent\textit{Evidence.} The hats do not themselves specify ordinal lengths. They must be fixed before the field identities have a precise meaning.\par
\subsection*{P001000 (both; ambiguity): no predecessor parenthesis}
\noindent\textit{Transcription reads.} There is no predecessor\par
\noindent\textit{Critical reading.} Specify ordinal predecessor, or clarify that the ordered field has no immediate neighbors.\par
\noindent\textit{Evidence.} The same passage correctly writes \(\omega-1\); lack of an ordinal predecessor is distinct from lack of a predecessor in field arithmetic.\par
\subsection*{P001000 (both; provenance): joint-work first surname}
\noindent\textit{Transcription reads.} Brozovskii / Брозовского\par
\noindent\textit{Critical reading.} Bibliographic identification: Brazovskii / Бразовского; preserve the written o in a diplomatic transcription.\par
\noindent\textit{Evidence.} The scan visibly has o. The exact collaborator list is confirmed by the 1982 joint paper and Krichever institutional bibliography.\par
\subsection*{P001000 (both; ambiguity): ascending union \(G=\bigcup_nG_n\)}
\noindent\textit{Transcription reads.} Editorial marker group-chain-tail may suggest groups\par
\noindent\textit{Critical reading.} The manuscript defines no group law; later work suggests a filtered associative algebra.\par
\noindent\textit{Evidence.} Only inclusions and a union are written. Krichever 1983 equation (1.91) defines an infinite-dimensional algebra, giving a partial identification.\par
\subsection*{P001001 (both; provenance): Joint-work attribution}
\noindent\textit{Transcription reads.} Schräder\par
\noindent\textit{Critical reading.} Schrader in bibliographic identification; preserve the manuscript spelling diplomatically.\par
\noindent\textit{Evidence.} Inspected scan; opened Cheeger--Müller--Schrader article and Springer chapter identify Robert Schrader.\par
\subsection*{P001001 (both; faithful error): Equation (2)}
\noindent\textit{Transcription reads.} \(i\hbar\partial/\partial t=\widehat Hf\)\par
\noindent\textit{Critical reading.} \(i\hbar\partial_t f=\widehat Hf\)\par
\noindent\textit{Evidence.} The scan omits the function on the left. Differentiating the evolution in equation (3) gives the corrected identity.\par
\subsection*{P001001 (both; faithful error): Kernel interpretation}
\noindent\textit{Transcription reads.} Transition probability\par
\noindent\textit{Critical reading.} Transition amplitude; its probabilistic interpretation requires a state and a measurement convention.\par
\noindent\textit{Evidence.} The displayed object is a matrix element of unitary Schrödinger evolution, not a nonnegative normalized density. The source wording was copied faithfully.\par
\subsection*{P001001 (both; faithful error): Equations (4)--(5)}
\noindent\textit{Transcription reads.} \(e^{-iS[x]}\), \(S=\int(\dot x^2-V)\,dt\)\par
\noindent\textit{Critical reading.} For the evolution written in (3), use formal phase \(e^{iS/\hbar}\) with \(S=\int(m\dot x^2/2-V)\,dt\), or state a consistent alternative convention.\par
\noindent\textit{Evidence.} Both signs and omitted constants are visible in the scan. Legendre transformation of the stated Hamiltonian and the free-particle kernel independently determine the conventional action and phase.\par
\subsection*{P001002 (both; faithful error): Equation (6)}
\noindent\textit{Transcription reads.} \(\hbar\partial_t f=+\widehat Hf\)\par
\noindent\textit{Critical reading.} \(\hbar\partial_t f=-\widehat Hf\) for the semigroup in (7).\par
\noindent\textit{Evidence.} Scan confirms absence of a minus sign. Differentiating the negative-exponent semigroup in the next line proves the inconsistency.\par
\subsection*{P001002 (both; faithful error): Equation (7), exponential weight}
\noindent\textit{Transcription reads.} \(\exp(-\int_0^t x(\tau)d\tau)\)\par
\noindent\textit{Critical reading.} \(\exp(-\hbar^{-1}\int_0^t V(x(\tau))d\tau)\), with Brownian generator \(\hbar\Delta/(2m)\).\par
\noindent\textit{Evidence.} Russian text and inspected scan both show x(tau), not V(x(tau)). The corrected weight follows from the specified Schrödinger Hamiltonian and heat generator.\par
\subsection*{P001002 (both; ambiguity): Equation (7), fixed endpoints}
\noindent\textit{Transcription reads.} Wiener measure on paths with both endpoints fixed\par
\noindent\textit{Critical reading.} Use Brownian-bridge expectation multiplied by the free heat kernel, or an explicitly unnormalized bridge measure.\par
\noindent\textit{Evidence.} The endpoint conditions and label Wiener measure are explicit. The source does not supply the conditioning or kernel normalization needed for equality.\par
\subsection*{P001002 (both; ambiguity): Equation (7), semigroup exponent}
\noindent\textit{Transcription reads.} \(Hx\) after \(-t/\hbar\)\par
\noindent\textit{Critical reading.} Read the operator as \(\widehat H\); keep the trailing x/index reading unresolved in the diplomatic text.\par
\noindent\textit{Evidence.} The Russian transcription flags the additional x-like stroke and the scan does not decide whether it is an operator index. The mathematics alone cannot justify deleting it as a transcription error.\par
\subsection*{P001002 (both; faithful error): Answer to Gelfand curvature-action question}
\noindent\textit{Transcription reads.} Einstein and Pontryagin classes\par
\noindent\textit{Critical reading.} For the precise natural second-order tensor problem, use the dimension-dependent Lovelock family; top-degree characteristic numbers have zero interior variation.\par
\noindent\textit{Evidence.} The manuscript answer is faithfully copied. Navarro--Navarro Theorem 2.6 supplies a bounded exact classification; dimensions at least five include nonzero Gauss--Bonnet tensors.\par
\subsection*{P001002 (both; faithful error): Equation (10)}
\noindent\textit{Transcription reads.} \((\det g)^{5/2}\prod dg_{\mu\nu}\)\par
\noindent\textit{Critical reading.} For the covariant matrix entries written, the pointwise invariant density has \((\det g)^{-5/2}\); a continuum product measure remains formal.\par
\noindent\textit{Evidence.} The positive exponent is visible in the scan and conflicts with the correctly negative general matrix exponent in (22). Congruence on symmetric matrices has Jacobian \(|\det A|^{n+1}\).\par
\subsection*{P001003 (both; ambiguity): Regge scalar-curvature normalization (14)--(16)}
\noindent\textit{Transcription reads.} \(\sum_h\delta_h\operatorname{vol}(h)=\int\operatorname{Scal}\,dV\)\par
\noindent\textit{Critical reading.} Specify the scalar-curvature convention; with scalar curvature twice Gaussian curvature on surfaces, the hinge sum needs a factor two.\par
\noindent\textit{Evidence.} The source uses unnormalized angle defect. In dimension two the sum of angle defects is 2 pi times Euler characteristic, while the scalar integral is 4 pi times Euler characteristic.\par
\subsection*{P001003 (en; translation): Refinement warning}
\noindent\textit{Transcription reads.} Barycentric subdivision is unnecessary.\par
\noindent\textit{Critical reading.} Do not use this subdivision when it destroys the required angle bound; the following note says the angles become small.\par
\noindent\textit{Evidence.} Russian says barycentrically subdividing is not to be done, and gives the geometric reason. Unnecessary weakens the practical warning.\par
\subsection*{P001003 (both; ambiguity): Convergence theorem and angle condition}
\noindent\textit{Transcription reads.} Mesh tends to zero; angles bounded away from zero.\par
\noindent\textit{Critical reading.} Use a uniformly fat admissible approximation family, with shape control on all simplex faces, not just mesh convergence or a selected angle.\par
\noindent\textit{Evidence.} The source appends a nondegeneracy warning. Cheeger--Müller--Schrader Section 1 and Theorem 5.1 supply the full approximation hypotheses.\par
\subsection*{P001003-P001005 (both; faithful error): Equations (17), (19), (26)--(29)}
\noindent\textit{Transcription reads.} \(\operatorname{tr}(\Omega^k)\) as the general limiting curvature density\par
\noindent\textit{Critical reading.} Use the corresponding normalized Lipschitz--Killing density, including the required coframe factors in lower curvature degree.\par
\noindent\textit{Evidence.} The original trace is visible and consistently copied. The full CMS article treats Lipschitz--Killing measures, not this raw trace family; degree alone rules out an unmodified trace as an n-form when 2k differs from n.\par
\subsection*{P001004 (both; faithful error): Top-degree Euler assertion}
\noindent\textit{Transcription reads.} \(\int\operatorname{tr}(\Omega^{n/2})=\chi(M)\)\par
\noindent\textit{Critical reading.} For even n, \(\int\operatorname{Pf}(\Omega/(2\pi))=\chi(M)\), with compatible orientation convention.\par
\noindent\textit{Evidence.} Scan and both transcriptions agree on trace. For an oriented surface the orthogonal curvature matrix has trace zero, while the sphere has nonzero Euler characteristic. CMS introduction identifies the Euler density.\par
\subsection*{P001004 (both; ambiguity): Solid-angle flag indices}
\noindent\textit{Transcription reads.} Index in \([\sigma^i,\sigma^{ij}]\) or overwritten adjacent superscript\par
\noindent\textit{Critical reading.} Retain the diplomatic uncertainty; the intended geometric angle is for nested faces with their dimensions specified.\par
\noindent\textit{Evidence.} The inspected stroke is not sufficiently clear to call a specific superscript a transcription error. The published flag definition supplies mathematical conventions but cannot rewrite the handwriting.\par
\subsection*{P001004 (both; provenance): Brin reference and date}
\noindent\textit{Transcription reads.} Similar Brin work around 1946\par
\noindent\textit{Critical reading.} Retain the uncertain name/date; do not silently change the decade or assign a particular Brin.\par
\noindent\textit{Evidence.} Scan inspected; searches of the exact surname with solid angles, polyhedral curvature and Euler formulas did not establish an exact bibliographic match.\par
\subsection*{P001004 (both; faithful error): Equation (23)}
\noindent\textit{Transcription reads.} \(\operatorname{vol}(\sigma,\ell)^{-(n+1)}\prod d\ell_{ij}\)\par
\noindent\textit{Critical reading.} \(C\operatorname{vol}(\sigma,\ell)^{-(n+1)}\prod\ell_{ij}\,d\ell_{ij}\), or use differentials of squared lengths.\par
\noindent\textit{Evidence.} Length variables are explicitly defined on P001003 and the scan contains no product of lengths. Gram-coordinate Jacobian calculation, independent symbolic checks and the one-edge case confirm the missing factor.\par
\subsection*{P001004 (both; faithful error): Equation (24)}
\noindent\textit{Transcription reads.} \(\deg g(x)^{5/2}\)\par
\noindent\textit{Critical reading.} Intended determinant notation; for covariant entries the invariant exponent is negative, as in (22).\par
\noindent\textit{Evidence.} The g-like final letter and positive exponent are copied faithfully. The nearby formula (10) and the matrix-density derivation determine the intended mathematical expression, not an alternative handwriting reading.\par
\subsection*{P001005 (both; ambiguity): Final convergence argument}
\noindent\textit{Transcription reads.} Eventually equal intersection Euler characteristics almost everywhere, therefore integrals converge.\par
\noindent\textit{Critical reading.} Include transversality and a valid integrable domination or equivalent uniform-integrability argument.\par
\noindent\textit{Evidence.} The manuscript records a proof idea. Almost-everywhere convergence alone does not imply convergence of integrals; the full geometric convergence theorem supplies additional control.\par
\subsection*{P001006 (both; ambiguity): Contraction exactness}
\noindent\textit{Transcription reads.} Globally (?) exact; possibly locally exact.\par
\noindent\textit{Critical reading.} A Hamiltonian Lie-algebra lift needs globally defined Hamiltonians. Preserve the disputed word as uncertain in the diplomatic reading.\par
\noindent\textit{Evidence.} Scan inspected and Russian Unclear note retained. The local Poincaré lemma would only give infinitesimal symplecticity; the later global moment map needs a stronger condition.\par
\subsection*{P001006 (both; ambiguity): Moment-map group equivariance}
\noindent\textit{Transcription reads.} Commutes with the A-action follows from the Lie-algebra homomorphism.\par
\noindent\textit{Critical reading.} Assume A connected, or impose equivariance for the remaining components separately.\par
\noindent\textit{Evidence.} Infinitesimal equivariance integrates on the identity component. The source specifies no connectedness hypothesis.\par
\subsection*{P001006 (both; provenance): Speaker initials}
\noindent\textit{Transcription reads.} V. L. Ginzburg in catalog; Ginzburg in manuscript.\par
\noindent\textit{Critical reading.} Retain catalog attribution with an explicit provenance qualifier, without assigning a full personal identity from the surname alone.\par
\noindent\textit{Evidence.} The scan heading contains only the surname; no initials are visible.\par
\subsection*{P001007 (both; faithful error): Parallel-body polynomial}
\noindent\textit{Transcription reads.} \(\sum_{i=0}^nV_k\varepsilon^k\)\par
\noindent\textit{Critical reading.} Use one summation index consistently, e.g. \(\sum_{j=0}^nV_j\varepsilon^j\).\par
\noindent\textit{Evidence.} The inspected scan genuinely mixes i and k; the transcription is faithful.\par
\subsection*{P001007 (both; faithful error): Curvature coefficient and principal curvatures}
\noindent\textit{Transcription reads.} \(V_i=C_i\int\sigma_i(k)dS\), with n principal curvatures.\par
\noindent\textit{Critical reading.} For the coefficient of \(\varepsilon^j\), j at least one, use \(j^{-1}\int\sigma_{j-1}(\kappa_1,\ldots,\kappa_{n-1})dS\); the constant coefficient is volume.\par
\noindent\textit{Evidence.} The source has no index shift. Integrating the outer normal Jacobian \(\prod_i(1+t\kappa_i)\) gives the corrected formula and the count of principal curvatures.\par
\subsection*{P001007 (both; ambiguity): Projection and affine-plane dimensions}
\noindent\textit{Transcription reads.} One coefficient index k labels both projected volume and plane-section integrals.\par
\noindent\textit{Critical reading.} For the intrinsic volume \(V_j\), average j-dimensional projection volume or Euler characteristic of affine \((n-j)\)-plane sections.\par
\noindent\textit{Evidence.} The manuscript mixes polynomial-power and intrinsic-volume conventions. McMullen pp. 491--492 and dimensional scaling distinguish them.\par
\subsection*{P001008 (both; faithful error): Partition generating-function corollary}
\noindent\textit{Transcription reads.} \(f_{\lambda,n}(t)+t^{m+n-1}f_{\lambda,m}(t)\)\par
\noindent\textit{Critical reading.} \(f_{\lambda,n}(t)+t^{m+n-1}f_{\lambda',m}(t^{-1})\), where lambda-prime is the conjugate partition.\par
\noindent\textit{Evidence.} Scan confirms neither conjugation nor inverse argument is written. Macdonald I.1 equation (1.7 prime), the preceding complementary-exponent identity, and exact finite checks independently fix both omissions.\par
\subsection*{P001008 (both; ambiguity): Symplectic cone examples and contact quotient}
\noindent\textit{Transcription reads.} Full \(\mathbb R^{2n}\) and full cotangent bundle, then smooth contact quotient.\par
\noindent\textit{Critical reading.} Remove the origin and the zero section for the stated examples, and require the scaling action to be free and proper.\par
\noindent\textit{Evidence.} The omitted zero loci are fixed points, so homogeneity alone does not give the asserted smooth contact quotient.\par
\subsection*{P001008 (both; faithful error): Weyl momentum convention}
\noindent\textit{Transcription reads.} \(\xi_k\mapsto+i\partial_{x_k}\) with \(e^{i\xi\cdot(x-y)}\)\par
\noindent\textit{Critical reading.} For the displayed phase, \(\xi_k\mapsto-i\partial_{x_k}\); changing phase changes this convention.\par
\noindent\textit{Evidence.} Both signs are unambiguous in the scan. Differentiation of the Fourier kernel independently verifies the correction.\par
\subsection*{P001008 (both; faithful error): Input function in Weyl kernel}
\noindent\textit{Transcription reads.} \((Pf)(x)\) on the left and \(u(y)\) on the right.\par
\noindent\textit{Critical reading.} Use the same input function on both sides.\par
\noindent\textit{Evidence.} The scan confirms the mismatch already flagged by the Russian transcription.\par
\subsection*{P001008-P001009 (both; ambiguity): Symbol sum and operator axioms}
\noindent\textit{Transcription reads.} An infinite ordinary sum of homogeneous symbols; unqualified compactness statements.\par
\noindent\textit{Critical reading.} Read the symbol expansion asymptotically with derivative remainder estimates; specify a global classical calculus for negative-order compactness and the unit-order principal-symbol drop.\par
\noindent\textit{Evidence.} The later asymptotic-summation axiom and the exact operator examples support this interpretation. Noncompact local calculi do not satisfy the compactness axiom without additional decay.\par
\subsection*{P001010 (both; ambiguity): MacPherson title and diagram}
\noindent\textit{Transcription reads.} Holonomic D-modules approximately equivalent to perverse sheaves.\par
\noindent\textit{Critical reading.} For ordinary constructible perverse sheaves, specify regular holonomic modules and the de Rham degree convention.\par
\noindent\textit{Evidence.} The scan omits regular. Kashiwara 1980 main theorem states the regular correspondence. The omitted hypothesis is mathematical, not a missing OCR word.\par
\subsection*{P001010 (both; provenance): Record type for MacPherson}
\noindent\textit{Transcription reads.} A completed lecture inferred from the catalog speaker entry.\par
\noindent\textit{Critical reading.} Treat as proposed subjects: the manuscript explicitly says the title of a paper he would speak about, followed by another title.\par
\noindent\textit{Evidence.} English words are part of the original scan, reproduced in the Russian diplomatic file. The record contains no completed proof narrative.\par
\subsection*{P001010 (both; ambiguity): Fuchs subgroup fibration}
\noindent\textit{Transcription reads.} G/H is used as a general model for the discrete-to-topological fiber.\par
\noindent\textit{Critical reading.} The displayed ordinary bundle model needs closed-subgroup or local-section hypotheses. For G-delta to G use the homotopy fiber, not the set quotient.\par
\noindent\textit{Evidence.} The word usually is in the source. The comparison homomorphism is the identity on underlying sets, so the naive quotient is a point and cannot encode local homology.\par
\subsection*{P001011 (both; faithful error): Superscript on diffeomorphism group in the cycle definition}
\noindent\textit{Transcription reads.} \(\operatorname{Diff}^{k}_{\omega_0}M\)\par
\noindent\textit{Critical reading.} \(\operatorname{Diff}^{c}_{\omega_0}M\) in the intended compact-support notation.\par
\noindent\textit{Evidence.} Scan really shows k; surrounding definitions use c for compact support and k for cycle degree. This is a source-level notational slip, not a transcription correction.\par
\subsection*{P001011 (both; ambiguity): Integral over K times a point}
\noindent\textit{Transcription reads.} \(x\mapsto\int_{K\times x}\Omega\) as an (n-k)-form.\par
\noindent\textit{Critical reading.} Interpret as integration along the K fiber of the appropriate mixed-degree form, with compact support in the M direction.\par
\noindent\textit{Evidence.} Literal restriction to a k-dimensional slice does not by itself describe the unintegrated M-form components. McDuff 1982 pp. 610--612 supplies the universal transverse-form framework.\par
\subsection*{P001011 (en; ambiguity): Sphere parity proof sketch}
\noindent\textit{Transcription reads.} Integral over a horizontal leaf is zero.\par
\noindent\textit{Critical reading.} Retain uncertainty between a horizontal slice and a leaf in the diagram; the sketch alone does not prove the sphere parity theorem.\par
\noindent\textit{Evidence.} Russian layer/fiber word and inspected diagram do not determine the intended cycle. Integrating a transverse volume form over an actual tangent leaf is trivially zero and would not explain parity.\par
\subsection*{P001012 (both; provenance): Name preceding rational vanishing}
\noindent\textit{Transcription reads.} Chern--Weyl\par
\noindent\textit{Critical reading.} Chern--Weil in mathematical attribution; preserve the source spelling diplomatically.\par
\noindent\textit{Evidence.} Scan and transcription agree on Weyl. Milnor Appendix and Lemma 8 proof explicitly use Chern--Weil characteristic forms.\par
\subsection*{P001012 (both; faithful error): Sah--Wagoner quotient}
\noindent\textit{Transcription reads.} \(\bigsqcup\mathbb Q\) labelled uncountable sum.\par
\noindent\textit{Critical reading.} A direct sum of copies of Q, or equivalently an uncountable rational-vector-space quotient.\par
\noindent\textit{Evidence.} The U-shaped source symbol is faithfully reproduced. A disjoint union is not an abelian group; Milnor Introduction p. 72 gives the precise vector-space quotient statement.\par
\subsection*{P001012 (both; ambiguity): Integral cohomology corollary}
\noindent\textit{Transcription reads.} A double-headed arrow between the two integral cohomology groups.\par
\noindent\textit{Critical reading.} The verified conclusion is the injection from H*(BG;Z) into H*(BG-delta;Z), for finitely many components.\par
\noindent\textit{Evidence.} Both arrowheads are visible in the scan. Milnor Corollary 1 p. 76 and its proof determine the intended mathematical conclusion; this is not evidence that a hook was omitted by OCR.\par
\subsection*{P001001-P001013 (both; provenance): Record boundaries and order}
\noindent\textit{Transcription reads.} Overlapping catalog windows treated as disjoint lecture ranges.\par
\noindent\textit{Critical reading.} Use actual dated and speaker headings: 546 ends on P001005, 547 occupies P001006, 549 precedes 548, Fuchs 551 intervenes within 553, and 555 ends on P001013.\par
\noindent\textit{Evidence.} All thirteen source scans inspected; Kirillov 556 begins at P001014. Heading record 545 has no additional mathematical content.\par
\subsection*{P001014 (both; ambiguity): Hyperbolic Cartan matrix}
\noindent\textit{Transcription reads.} Having a negative eigenvalue\par
\noindent\textit{Critical reading.} Add indecomposability, indefinite type, and finite/affine proper connected principal submatrices for the standard hyperbolic definition.\par
\noindent\textit{Evidence.} The scan contains only the negative-eigenvalue parenthesis; it does not specify a matrix.\par
\subsection*{P001015 (both; faithful error): Orbit label after the invariants}
\noindent\textit{Transcription reads.} \(\Omega_{c_1,c_2}\)\par
\noindent\textit{Critical reading.} \(\Omega_{c_2,c_3}\)\par
\noindent\textit{Evidence.} The same discussion fixes \(\operatorname{tr}X^2\) and \(\operatorname{tr}X^3\); the inconsistent subscripts are present in the scan.\par
\subsection*{P001015 (both; faithful error): Riemann–Roch multiplicity formula}
\noindent\textit{Transcription reads.} \(X=p^{-1}(y_0)\)\par
\noindent\textit{Critical reading.} \(X=\mu^{-1}(\nu)/T\), with descended prequantum data\par
\noindent\textit{Evidence.} A four-dimensional level set is not the two-dimensional reduced symplectic space used in the multiplicity theorem. The scan has no quotient.\par
\subsection*{P001016 (both; faithful error): Fractional derivative integral}
\noindent\textit{Transcription reads.} Integral without a prefactor\par
\noindent\textit{Critical reading.} Insert \(1/\Gamma(-p)\), initially \(\operatorname{Re}p<0\), then continue\par
\noindent\textit{Evidence.} Direct beta-integral evaluation and the scan confirm the missing normalization.\par
\subsection*{P001017 (both; ambiguity): Darboux coordinates for odd binary forms}
\noindent\textit{Transcription reads.} Overwritten sign before the \(p_{n-1}\) term\par
\noindent\textit{Critical reading.} Use a minus sign before \(p_{n-1}\) and alternating signs thereafter.\par
\noindent\textit{Evidence.} The full alternating formula is forced by the displayed symplectic form and checked by differentiating F(w+t). The overwritten mark in the scan alone is not decisive.\par
\subsection*{P001017 (both; ambiguity): Space of lines identified with cotangent space}
\noindent\textit{Transcription reads.} \(T^*\mathbb R^2\)\par
\noindent\textit{Critical reading.} Locally \(T^*\mathbb R^2\); globally oriented lines form \(T^*S^2\).\par
\noindent\textit{Evidence.} Givental 1983, Example 1, explicitly uses the reduced \(T^*S^{n-1}\). Orientation is not stated in the notes.\par
\subsection*{P001017 (both; faithful error): Any variational problem}
\noindent\textit{Transcription reads.} Extremals of any variational problem form a symplectic space\par
\noindent\textit{Critical reading.} Require regularity of the variational problem and a smooth reduction where used.\par
\noindent\textit{Evidence.} A degenerate Lagrangian yields a degenerate Legendre transform; the literal universal quantifier exceeds the valid construction.\par
\subsection*{P001019 (both; faithful error): Index formula with a possible boundary}
\noindent\textit{Transcription reads.} \(M\) possibly has boundary, followed by \(\operatorname{ind}D=\int\alpha_0\)\par
\noindent\textit{Critical reading.} Assume \(M\) is closed for this equality, or specify a boundary realization and correction.\par
\noindent\textit{Evidence.} Ellipticity of the differential expression alone does not give a Fredholm problem with boundary.\par
\subsection*{P001020 (both; faithful error): Hypothesis preceding the uncorrected signature formula}
\noindent\textit{Transcription reads.} \(\partial M\ne\varnothing\)\par
\noindent\textit{Critical reading.} \(\partial M=\varnothing\) for formulas (6)--(7) without a boundary term\par
\noindent\textit{Evidence.} The scan shows the nonempty-boundary sign; formula (11) later introduces the necessary eta correction.\par
\subsection*{P001020 (both; ambiguity): div–grad–curl base dimension}
\noindent\textit{Transcription reads.} In \(\mathbb R^4\)\par
\noindent\textit{Critical reading.} On \(\mathbb R^3\), acting on a scalar plus a three-component vector\par
\noindent\textit{Evidence.} The same example explicitly says \(\mathbb R^3\) on P001024; standard curl in this matrix is three-dimensional.\par
\subsection*{P001020;P001024 (both; faithful error): Square of the vector-calculus matrix}
\noindent\textit{Transcription reads.} All positive \(\operatorname{div},\operatorname{grad},\operatorname{curl}\) entries square to \(\Delta I\)\par
\noindent\textit{Critical reading.} With conventional operators use upper-right \(-\operatorname{div}\) and obtain \(-\Delta I\); the displayed all-positive matrix has vector block \(2\operatorname{grad}\operatorname{div}-\Delta I_3\).\par
\noindent\textit{Evidence.} Independent symbolic multiplication and \(\operatorname{curl}\operatorname{curl}=\operatorname{grad}\operatorname{div}-\Delta\) verify the correction.\par
\subsection*{P001021 (both; faithful error): Boundary restriction in (15)}
\noindent\textit{Transcription reads.} \(P_A(f|_A)=0\)\par
\noindent\textit{Critical reading.} \(P_A(f|_Y)=0\), \(Y=\partial M\)\par
\noindent\textit{Evidence.} The scan genuinely uses A as the restriction subscript, but A is the operator, not the boundary.\par
\subsection*{P001021-P001022 (both; ambiguity): Normal tangential operator attributed to APS}
\noindent\textit{Transcription reads.} \(A\) normal; eta defined through real parts of eigenvalues\par
\noindent\textit{Critical reading.} Classical APS assumes self-adjoint A; retain the normal extension as a separate assertion with additional hypotheses to be identified.\par
\noindent\textit{Evidence.} The scan itself says APS used self-adjoint A; APS I Theorem (3.10) confirms this scope.\par
\subsection*{P001022 (both; ambiguity): Boundary operator domain and codomain}
\noindent\textit{Transcription reads.} \(A:E|_Y\to F|_Y\) with spectral eigenspaces\par
\noindent\textit{Critical reading.} \(D=\sigma(\partial_u+A)\), \(A:E|_Y\to E|_Y\) after identifying bundles by the normal symbol\par
\noindent\textit{Evidence.} APS I Theorem (3.10) includes this isomorphism; eigenvalues are otherwise undefined.\par
\subsection*{P001022 (both; ambiguity): Positive spectral projector and zero-mode term}
\noindent\textit{Transcription reads.} \(P_+\) projects to positive eigenfunctions, with correction \(-(h+\eta)/2\)\par
\noindent\textit{Critical reading.} Use nonnegative projection for this kernel-term convention, or change the index correction.\par
\noindent\textit{Evidence.} Allowing the h-dimensional zero eigenspace changes the index by h; the scan does not settle the convention.\par
\subsection*{P001023 (both; ambiguity): The sum of the two spectral halves}
\noindent\textit{Transcription reads.} \(\zeta(A)=\zeta_+(A)+\zeta_-(A)\)\par
\noindent\textit{Critical reading.} For functions of \(s\) use \(\zeta_{|A|}\); at zero state the phase and regularity conventions.\par
\noindent\textit{Evidence.} The manuscript denotes zero values; extending the identity to the cut zeta function at arbitrary s adds an unjustified phase omission.\par
\subsection*{P001023 (both; provenance): Bibliographic page 645}
\noindent\textit{Transcription reads.} Stredder, Epstein, JDG 10 (1975), 645\par
\noindent\textit{Critical reading.} Epstein 631–645; Stredder 647–660\par
\noindent\textit{Evidence.} The titles and ranges are corroborated in Fegan–Steer 2005 references; original full-text access failed.\par
\subsection*{P001023 (both; ambiguity): Filling qualification}
\noindent\textit{Transcription reads.} Without a filling it has not been proved\par
\noindent\textit{Critical reading.} For self-adjoint elliptic eta regularity, no filling is required; retain possible reference to the unspecified normal extension.\par
\noindent\textit{Evidence.} Gilkey 1981 Theorem 0.1 establishes the self-adjoint result before this talk.\par
\subsection*{P001024 (both; ambiguity): First argument of zeta in (16)}
\noindent\textit{Transcription reads.} Overwritten \(0\) or \(\theta\)\par
\noindent\textit{Critical reading.} \(\zeta_\theta^A(0,x)\)\par
\noindent\textit{Evidence.} The local constant-value formula and the earlier P001020 expression require zero. The rewritten glyph alone is uncertain.\par
\subsection*{P001025 (both; ambiguity): Integer symbol in (19)}
\noindent\textit{Transcription reads.} \(\rho(i),\rho_\ell,b_{m(i-1)-n}/\lambda^\ell\)\par
\noindent\textit{Critical reading.} Use \(\ell\) consistently for the integer argument and symbol index\par
\noindent\textit{Evidence.} The scan and Russian transcription vary the glyphs, but the preceding statement defines ell as the integer.\par
\subsection*{P001025 (both; ambiguity): Range in the hint following (d)}
\noindent\textit{Transcription reads.} All integer points \(>0\)\par
\noindent\textit{Critical reading.} Do not silently replace: (d) states nonpositive integers; the hint’s intended sign or variable remains uncertain.\par
\noindent\textit{Evidence.} The scan clearly has \(>0\), as already noted diplomatically.\par
\subsection*{P001025 (both; faithful error): Pseudodifferential scope of (d)}
\noindent\textit{Transcription reads.} \(\rho(\ell)=0\) for every nonpositive integer after mentioning pseudodifferential operators\par
\noindent\textit{Critical reading.} Keep general zero-value independence; other integer values are governed by \(\operatorname{Res}(\Pi A^{-\ell})\). Even differential operators can have nonzero cut differences at negative integers.\par
\noindent\textit{Evidence.} Ponge Theorems 4.2–4.3 and Loya–Moroianu 2012 section 4: negative eta/absolute-zeta residues produce nonzero cut differences.\par
\subsection*{P001025 (both; faithful error): Second sum of the shift formula (20)}
\noindent\textit{Transcription reads.} \(t^k\sum_{\mu=1}^n(-1)^\mu\delta_\mu/[\mu\binom{k+\mu}{k}]\)\par
\noindent\textit{Critical reading.} Insert \(t^\mu\) in each summand, giving \(t^{k+\mu}\).\par
\noindent\textit{Evidence.} The extra factor is absent in the scan. The binomial complex-power expansion and the \(t=0\) check both show it is necessary.\par
